%DIF 1-5d1
%DIF LATEXDIFF DIFFERENCE FILE
%DIF DEL ../03_draft_revision/01_draft_original/20260706_Topology_Study_submitted (1).tex   Wed Aug 12 16:43:27 2026
%DIF ADD paper_COMPILE_flat.tex                                                             Thu Aug 27 16:30:29 2026
%DIF < %% ============================================================
%DIF < %% PAPER 1: TOPOLOGY ABSTRACTION & PHYSICS FIDELITY EFFECTS
%DIF < %% Target Journal: Applied Energy
%DIF < %% Version: 14.1 — Terminology patch: "synthetic sweep" removed
%DIF < %% ============================================================
%DIF -------

%DIF 7c2-14
%DIF < \documentclass[a4paper,fleqn]{cas-dc}
%DIF -------
\documentclass[a4paper,fleqn]{cas-sc} %DIF > 
\input glyphtounicode %DIF > 
\pdfgentounicode=1 %DIF > 
\pdfglyphtounicode{ff}{0066 0066} %DIF > 
\pdfglyphtounicode{fi}{0066 0069} %DIF > 
\pdfglyphtounicode{fl}{0066 006C} %DIF > 
\pdfglyphtounicode{ffi}{0066 0066 0069} %DIF > 
\pdfglyphtounicode{ffl}{0066 0066 006C} %DIF > 
\pdfglyphtounicode{uniFB00}{0066 0066} %DIF > 
\pdfglyphtounicode{uniFB01}{0066 0069} %DIF > 
\pdfglyphtounicode{uniFB02}{0066 006C} %DIF > 
\pdfglyphtounicode{uniFB03}{0066 0066 0069} %DIF > 
\pdfglyphtounicode{uniFB04}{0066 0066 006C} %DIF > 
%DIF -------

%DIF 9-10c16-26
%DIF < %% --- Bibliographie ---
%DIF < \usepackage[numbers,sort&compress]{natbib}
%DIF -------
\newcommand{\CP}{\ensuremath{\mathrm{CP}}}      % T0P0 copperplate, no loss   (control) %DIF > 
\newcommand{\CPL}{\ensuremath{\mathrm{CP{+}L}}} % T0P1 copperplate + aggregate loss (control) %DIF > 
\newcommand{\ZN}{\ensuremath{\mathrm{ZN}}}      % T1P1 zone-aggregated + loss %DIF > 
\newcommand{\NDz}{\ensuremath{\mathrm{ND^{0}}}} % T2P0 full nodes, no loss     (topology control) %DIF > 
\newcommand{\Lone}{\ensuremath{\mathrm{L1}}}    % T2P1 nodes + trunk loss  BASELINE %DIF > 
\newcommand{\Ltwo}{\ensuremath{\mathrm{L2}}}    % T2P2 + temperature propagation (PWL) %DIF > 
\newcommand{\Lthree}{\ensuremath{\mathrm{L3}}}  % T2P3 + trunk pressure drop & pumping %DIF > 
\newcommand{\Lfour}{\ensuremath{\mathrm{L4}}}   % T2P4 + station resolution + service laterals (flat Dp) %DIF > 
\newcommand{\Lfive}{\ensuremath{\mathrm{L5}}}   % T2P5 + dynamic flow-dependent station Dp & pumping %DIF > 
\newcommand{\Lsix}{\ensuremath{\mathrm{L6}}}    % T2P6 + transport delay %DIF > 
\newcommand{\NLref}{\ensuremath{\mathrm{NL}}}   % nonlinear reference (solved on representative 72-hour windows; global re-solve intractable) %DIF > 
%DIF -------

%DIF 12-19d28
%DIF < %% --- Mathematik ---
%DIF < \usepackage{amsmath,amssymb,amsfonts}
%DIF < 
%DIF < %% --- Tabellen ---
%DIF < \usepackage{booktabs}
%DIF < \usepackage{multirow}
%DIF < \usepackage{array}
%DIF < \usepackage{tabularx}
%DIF -------
\usepackage{threeparttable}
%DIF 21-34c29-30
%DIF < 
%DIF < %% --- Grafik & Floats ---
%DIF < \usepackage{graphicx}
%DIF < \usepackage{float}
%DIF < \usepackage{subfig}
%DIF < \usepackage[section]{placeins}
%DIF < 
%DIF < %% --- Einheiten & Chemie ---
%DIF < \usepackage{siunitx}
%DIF < \usepackage[version=4]{mhchem}
%DIF < 
%DIF < %% --- Sonstiges ---
%DIF < \usepackage{xcolor}
%DIF < \usepackage{enumitem}
%DIF -------
\usepackage{booktabs} %DIF > 
\usepackage{makecell} %DIF > 
%DIF -------
\usepackage{pifont}
%DIF 36-57c32-39
%DIF < \usepackage{adjustbox}
%DIF < 
%DIF < %% --- Hyperlinks (immer spät laden) ---
%DIF < \usepackage{hyperref}
%DIF < \usepackage{cleveref}
%DIF < \usepackage[numbers]{natbib}
%DIF < %% --- SI-Setup ---
%DIF < \sisetup{
%DIF <   per-mode=symbol,
%DIF <   range-phrase=--,
%DIF <   range-units=single
%DIF < }
%DIF < 
%DIF < %% --- Hyperref-Setup ---
%DIF < \hypersetup{
%DIF <   colorlinks=true,
%DIF <   linkcolor=blue,
%DIF <   citecolor=blue,
%DIF <   urlcolor=blue
%DIF < }
%DIF < 
%DIF < %% --- Float-Platzierung maximal lockern ---
%DIF -------
\renewcommand{\checkmark}{\ding{51}} %DIF > 
\usepackage{caption} %DIF > 
\usepackage{xcolor} %DIF > 
\usepackage[numbers,sort&compress]{natbib}   % numeric [n] citations; without this the elsarticle-num-names bst prints the verbose author-year label %DIF > 
\usepackage{float} %DIF > 
\usepackage{multicol}            % two-column nomenclature list (Applied Energy style) %DIF > 
\restylefloat{figure}            % let the float package own the figure env so [H] (place-here) is honoured under cas %DIF > 
\usepackage[section]{placeins}   % \FloatBarrier at each section so floats stay in-section %DIF > 
%DIF -------
\renewcommand{\topfraction}{0.95}
\renewcommand{\bottomfraction}{0.95}
\renewcommand{\textfraction}{0.05}
\renewcommand{\floatpagefraction}{0.85}
\setcounter{topnumber}{4}
\setcounter{bottomnumber}{4}
\setcounter{totalnumber}{8}
%DIF 65-70d47
%DIF < 
%DIF < %% --- Speziell für double-column (table*/figure*) ---
%DIF < \renewcommand{\dbltopfraction}{0.95}
%DIF < \renewcommand{\dblfloatpagefraction}{0.85}
%DIF < 
%DIF < %% --- Seiten dürfen unten unterschiedlich lang sein ---
%DIF -------
\raggedbottom
%DIF 72a48-62
\newif\ifmarkup \markupfalse          % \markuptrue for the marked-up build %DIF > 
\ifmarkup %DIF > 
  \newcommand{\new}[1]{\textcolor[HTML]{1A5FB4}{#1}}   % new in this revision %DIF > 
  \newcommand{\chg}[1]{\textcolor[HTML]{0F7A5A}{#1}}   % v1 text, materially rewritten %DIF > 
  \newcommand{\gone}[1]{\textcolor[HTML]{9A9996}{[removed: #1]}} %DIF > 
\else %DIF > 
  \newcommand{\new}[1]{#1} %DIF > 
  \newcommand{\chg}[1]{#1} %DIF > 
  \newcommand{\gone}[1]{} %DIF > 
\fi %DIF > 
\providecommand{\citet}[1]{\cite{#1}} %DIF > 
\providecommand{\citep}[1]{\cite{#1}} %DIF > 
\providecommand{\citealt}[1]{\cite{#1}} %DIF > 
\providecommand{\citeauthor}[1]{\cite{#1}} %DIF > 
\providecommand{\degree}{\ensuremath{^\circ}} %DIF > 
%DIF -------

%DIF < %% --- Custom commands ---
%DIF -------
% suppress the cas-sc "Preprint submitted to Elsevier" first-page footer %DIF > 
%DIF < \newcommand{\placeholder}[1]{\textcolor{red}{[#1]}}
\ExplSyntaxOn %DIF > 
%DIF < \newcommand{\Ltwo}{\mathrm{L2}}
\cs_set:Npn \__first_footerline: { } %DIF > 
%DIF < \newcommand{\Lthree}{\mathrm{L3}}
\ExplSyntaxOff %DIF > 
%DIF < \newcommand{\Lplus}{\mathrm{L3}^{+}}
%DIF < \newcommand{\Lnl}{\mathrm{L3}^{\mathrm{NL}}}
%DIF < 
%DIF < %% --- Einheitliche Formel-Formatierung ---
%DIF < \setlength{\mathindent}{1em}
%DIF < \setlength{\abovedisplayskip}{8pt plus 2pt minus 4pt}
%DIF < \setlength{\belowdisplayskip}{8pt plus 2pt minus 4pt}
%DIF < \setlength{\abovedisplayshortskip}{4pt plus 2pt}
%DIF < \setlength{\belowdisplayshortskip}{6pt plus 2pt minus 2pt}
%DIF < 
%DIF < %% --- Custom column types ---
%DIF < \newcolumntype{P}[1]{>{\centering\arraybackslash}m{#1}}
%DIF < \newcolumntype{Q}[1]{>{\raggedright\arraybackslash}m{#1}}
%DIF < 
%DIF < %% --- Fix caption alignment for cas-dc ---
%DIF < \makeatletter
%DIF < \long\def\@makecaption#1#2{%
%DIF <   \vskip\abovecaptionskip
%DIF <   \raggedright
%DIF <   {\bfseries #1.}\enspace #2\par
%DIF <   \vskip\belowcaptionskip}
%DIF < \makeatother
%DIF PREAMBLE EXTENSION ADDED BY LATEXDIFF
%DIF UNDERLINE PREAMBLE %DIF PREAMBLE
\RequirePackage[normalem]{ulem} %DIF PREAMBLE
\RequirePackage{color}\definecolor{RED}{rgb}{1,0,0}\definecolor{BLUE}{rgb}{0,0,1} %DIF PREAMBLE
\providecommand{\DIFadd}[1]{{\protect\color{blue}\uwave{#1}}} %DIF PREAMBLE
\providecommand{\DIFdel}[1]{{\protect\color{red}\sout{#1}}} %DIF PREAMBLE
%DIF SAFE PREAMBLE %DIF PREAMBLE
\providecommand{\DIFaddbegin}{} %DIF PREAMBLE
\providecommand{\DIFaddend}{} %DIF PREAMBLE
\providecommand{\DIFdelbegin}{} %DIF PREAMBLE
\providecommand{\DIFdelend}{} %DIF PREAMBLE
\providecommand{\DIFmodbegin}{} %DIF PREAMBLE
\providecommand{\DIFmodend}{} %DIF PREAMBLE
%DIF FLOATSAFE PREAMBLE %DIF PREAMBLE
\providecommand{\DIFaddFL}[1]{\DIFadd{#1}} %DIF PREAMBLE
\providecommand{\DIFdelFL}[1]{\DIFdel{#1}} %DIF PREAMBLE
\providecommand{\DIFaddbeginFL}{} %DIF PREAMBLE
\providecommand{\DIFaddendFL}{} %DIF PREAMBLE
\providecommand{\DIFdelbeginFL}{} %DIF PREAMBLE
\providecommand{\DIFdelendFL}{} %DIF PREAMBLE
%DIF COLORLISTINGS PREAMBLE %DIF PREAMBLE
\RequirePackage{listings} %DIF PREAMBLE
\RequirePackage{color} %DIF PREAMBLE
\lstdefinelanguage{DIFcode}{ %DIF PREAMBLE
%DIF DIFCODE_UNDERLINE %DIF PREAMBLE
  moredelim=[il][\color{red}\sout]{\%DIF\ <\ }, %DIF PREAMBLE
  moredelim=[il][\color{blue}\uwave]{\%DIF\ >\ } %DIF PREAMBLE
} %DIF PREAMBLE
\lstdefinestyle{DIFverbatimstyle}{ %DIF PREAMBLE
	language=DIFcode, %DIF PREAMBLE
	basicstyle=\ttfamily, %DIF PREAMBLE
	columns=fullflexible, %DIF PREAMBLE
	keepspaces=true %DIF PREAMBLE
} %DIF PREAMBLE
\lstnewenvironment{DIFverbatim}[1][]{\lstset{style=DIFverbatimstyle}}{} %DIF PREAMBLE
\lstnewenvironment{DIFverbatim*}[1][]{\lstset{style=DIFverbatimstyle,showspaces=true}}{} %DIF PREAMBLE
\lstset{extendedchars=\true,inputencoding=utf8}

%DIF END PREAMBLE EXTENSION ADDED BY LATEXDIFF

%DIF OLD-MACRO STUBS (so deleted original text compiles under the revised preamble)
\usepackage{siunitx}
\providecommand{\Lplus}{\ensuremath{\mathrm{L3^{+}}}}
\providecommand{\Lnl}{\ensuremath{\mathrm{L3^{\mathrm{NL}}}}}
\providecommand{\placeholder}[1]{}
\providecommand{\textpagefraction}{.001}
\providecommand{\floatpagepagefraction}{1}
\providecommand{\ce}[1]{\ensuremath{\mathrm{#1}}}
\providecommand{\Cref}[1]{\ref{#1}}
\providecommand{\cref}[1]{\ref{#1}}
\begin{document}
\DIFdelbegin %DIFDELCMD < \let\WriteBookmarks\relax
%DIFDELCMD < \def\floatpagepagefraction{1}
%DIFDELCMD < \def%%%
\DIFdel{\textpagefraction{.001}
}%DIFDELCMD < 

%DIFDELCMD < \shorttitle{Topology and Physics Fidelity in District Heating Dispatch}
%DIFDELCMD < \shortauthors{L.~Ruess and A.~Sauer}
%DIFDELCMD < 

%DIFDELCMD < \title[mode = title]{Topology Resolution Dominates Dispatch Accuracy in District Heating Networks with Industrial Demands: A Five-Level MILP Comparison}
%DIFDELCMD < %%%
\DIFdelend 

\DIFdelbegin %DIFDELCMD < \author[1,2]{Lukas Ruess}[orcid=0009-0005-5010-3520]
%DIFDELCMD < \cormark[1]
%DIFDELCMD < \ead{lukas.ruess@ipa.fraunhofer.de}
%DIFDELCMD < \credit{Conceptualization, Methodology, Software, Validation,
%DIFDELCMD <         Formal analysis, Investigation, Data curation,
%DIFDELCMD <         Writing -- original draft, Visualization, Funding acquisition}
%DIFDELCMD < 

%DIFDELCMD < \author[1,2]{Alexander Sauer}
%DIFDELCMD < \credit{Supervision, Writing -- review \& editing}
%DIFDELCMD < 

%DIFDELCMD < \affiliation[1]{organization={Institute for Energy Efficiency in
%DIFDELCMD <                 Production EEP, University of Stuttgart},
%DIFDELCMD <                 addressline={Allmandring 35},
%DIFDELCMD <                 city={Stuttgart},
%DIFDELCMD <                 postcode={70569},
%DIFDELCMD <                 country={Germany}}
%DIFDELCMD < \affiliation[2]{organization={Fraunhofer Institute for Manufacturing
%DIFDELCMD <                 Engineering and Automation IPA},
%DIFDELCMD <                 addressline={Nobelstr.~12},
%DIFDELCMD <                 city={Stuttgart},
%DIFDELCMD <                 postcode={70569},
%DIFDELCMD <                 country={Germany}}
%DIFDELCMD < 

%DIFDELCMD < \cortext[1]{Corresponding author.}
%DIFDELCMD < %%%
\DIFdelend \DIFaddbegin \title[mode=title]{%
  Loss Visibility versus Spatial Detail in District-Heating Dispatch Optimisation}
\DIFaddend 

\begin{abstract}
\DIFdelbegin \DIFdel{The widespread assumption that higher physics fidelity yields
proportional accuracy gains in district heating dispatch is
untested}\DIFdelend \DIFaddbegin \DIFadd{District-heating dispatch models are simplified, yet their decision costs are
seldom quantified}\DIFaddend . This study \DIFdelbegin \DIFdel{provides a comparison across
184 optimisation  instances on a validated industrial network and
36~synthetic topologies, pressure-drop hydraulics, temperature
propagation, and nonlinear formulations collectively contribute
less than }\DIFdelend \DIFaddbegin \DIFadd{distinguishes estimation bias, the objective-value
difference between formulations, from decision regret, defined as an
\(L_1\)-relative }\emph{\DIFadd{forward-valued}} \DIFadd{schedule gap based on each schedule's cost
when re-simulated with a high-fidelity forward model incorporating exponential
heat losses, computed-friction hydraulics, and consumer differential-pressure
constraints. Both are assessed through a one-phenomenon-at-a-time fidelity ladder
from a copperplate model to station-resolved hydraulics. Aggregate losses in the
copperplate enable exact separation of loss and topology effects. Tests use the
Memmingen network and 135 factorially designed synthetic networks. The copperplate
underestimates operating costs by 15.1\% but produces schedules whose unprovisioned
loss, valued at the marginal unit, adds 46.1\%; loss-aware schedules provision this
loss and remain hydraulically deliverable. Loss visibility, not spatial topology,
explains the discrepancy. Losses account for 95.8\% of the
copperplate-to-baseline gap, while topology alone, with physics fixed, adds 0.25\%.
Across synthetic networks, losses explain a median 100.0\% of the gap. The topology
effect remains within \(\pm 0.6\%\) for trunk lengths of at least \(5~\mathrm{km}\)
and never exceeds 2.4\%. Aggregate loss adders may conceal the deficiency in one
network but do not transfer across scenarios. Forward evaluation of all 174
transmission stations and service laterals using real component data adds less
than }\DIFaddend 1\DIFdelbegin \DIFdel{\,\% to dispatch cost while the
computationally
trivial step of resolving network topology captures
13--39\,\% of hidden cost bias.
}%DIFDELCMD < 

%DIFDELCMD < %%%
\DIFdel{A five-level taxonomy (L1 copperplate through $\Lnl$ MIQCP)
isolates three effectsthrough pairwise comparison. Results show
that topology dominates: the
copperplate to detailed gap reaches
13\,\% cost underestimation and 74\,t\,}%DIFDELCMD < \ce{CO2}%%%
\DIFdel{/year on the primary case, driven almost entirely by thermal losses made
visible through spatial resolution. Extended hydraulic physics add ${\leq}\,0.3\,\%$; the combined linearisation and transport-delay effect remains at or below solver tolerance (${\leq}\,0.5\,\%$) across two independent winter windows,
confirming that MILP linearisation is adequate for planning-accuracy dispatch. These findings yield quantitative
selection thresholds linking pipe length and demand heterogeneity
to the minimum warranted model fidelity. }%DIFDELCMD < 

%DIFDELCMD < %%%
\DIFdel{Within the tested scope (radial trees, }\DIFdelend \DIFaddbegin \DIFadd{\% to the fixed schedule and finds no hydraulic violation. Spatial routing is immaterial under central generation; distributed
generation remains open. Results assume fixed capacities, a }\DIFaddend fixed heating curve,
\DIFdelbegin \DIFdel{dispatch with given capacities), practical implications could be determined: resolve the network graph before refining the physics.
Ring topologies, variable supply temperature, and joint
dispatch--sizing problems may shift these thresholds and are
identified as priority follow-up studies}\DIFdelend \DIFaddbegin \DIFadd{and radial networks}\DIFaddend .
\end{abstract}

\begin{keywords}
District heating \sep Mixed-integer linear programming \sep \DIFdelbegin \DIFdel{Mixed-integer quadratically constrained programming }\DIFdelend \DIFaddbegin \DIFadd{Model fidelity }\DIFaddend \sep
\DIFdelbegin \DIFdel{Topology abstraction }\DIFdelend \DIFaddbegin \DIFadd{Decision regret }\DIFaddend \sep Thermal \DIFdelbegin \DIFdel{storage dispatch }\DIFdelend \DIFaddbegin \DIFadd{losses }\sep \DIFadd{Network aggregation }\DIFaddend \sep
Heat pump optimisation
\DIFdelbegin %DIFDELCMD < \sep
%DIFDELCMD < %%%
\DIFdel{Linearisation error
}\DIFdelend \end{keywords}



\DIFaddbegin \shorttitle{Loss Visibility versus Spatial Detail in District-Heating Dispatch}

\shortauthors{L.~Ruess and A.~Sauer}

\author[1,2]{Lukas Ruess}[orcid=0009-0005-5010-3520]
\cormark[1]
\ead{lukas.ruess@eep.uni-stuttgart.de}
\credit{Conceptualization, Methodology, Software, Validation,
        Formal analysis, Investigation, Data curation,
        Writing: original draft, Visualization, Funding acquisition}

\author[1,2]{\DIFadd{Alexander Sauer}}
\credit{Supervision, Writing: review \& editing}

\affiliation[1]{organization={Institute for Energy Efficiency in
                Production EEP, University of Stuttgart},
                addressline={Allmandring 35},
                city={Stuttgart},
                postcode={70569},
                country={Germany}}
\affiliation[2]{organization={Fraunhofer Institute for Manufacturing
                Engineering and Automation IPA},
                addressline={Nobelstr.~12},
                city={Stuttgart},
                postcode={70569},
                country={Germany}}

\cortext[1]{Corresponding author.}



\DIFaddend \maketitle

%DIF < % ============================================================
%DIF < % 1. INTRODUCTION
%DIF < % ============================================================
\DIFdelbegin %DIFDELCMD < \section{Introduction}
%DIFDELCMD < \label{sec:introduction}
%DIFDELCMD < %%%
\DIFdelend \DIFaddbegin \ifmarkup
\begin{center}
\footnotesize
\DIFadd{\fbox{\parbox{0.9\linewidth}{%
\textbf{Marked-up copy.}
\textcolor[HTML]{1A5FB4}{Blue} marks material added in this revision;
\textcolor[HTML]{0F7A5A}{green} marks text carried over from the original
submission and materially rewritten; and \textcolor[HTML]{9A9996}{grey}
marks deletions considered relevant to readers. Unmarked text is unchanged.
Because the manuscript was restructured and the model levels renamed,
marking is applied at paragraph rather than sentence level.}}
}\end{center}
\fi
\DIFaddend 


%DIF < \subsection{Motivation and problem statement}
%DIF < \label{sec:motivation}
\DIFaddbegin \section*{\DIFadd{Nomenclature}}
{
{\footnotesize
\DIFaddend 

\DIFdelbegin \DIFdel{District heating networks (DHN) for industrial usage }\DIFdelend \DIFaddbegin \noindent\textbf{\DIFadd{Fidelity levels and evaluation quantities}}\par\nopagebreak\smallskip
\noindent
\begin{tabular*}{\columnwidth}{@{\extracolsep{\fill}}
  >{\raggedright\arraybackslash}p{1.15cm}
  >{\raggedright\arraybackslash}p{0.95cm}
  >{\raggedright\arraybackslash}p{1.85cm}
  l @{}}
\toprule
Symbol & Code & Status & Description \\
\midrule
\CP      & T0P0 & optimised   & Copperplate, no loss (decomposition control) \\
\CPL     & T0P1 & optimised   & Copperplate $+$ exogenous aggregate loss (control) \\
\ZN      & T1P1 & optimised   & Zone-aggregated $+$ loss (aggregation sensitivity, Section~\ref{subsec:clustering}) \\
\NDz     & T2P0 & optimised   & Node-resolved, no loss (topology control) \\
\Lone    & T2P1 & optimised   & Node-resolved $+$ trunk loss (comparison baseline) \\
\Ltwo    & T2P2 & fwd.-eval.  & $+$ temperature propagation (free-variable form degenerate) \\
\Lthree  & T2P3 & optimised   & $+$ trunk pressure drop and pumping \\
\Lfour   & T2P4 & fwd.-eval.  & $+$ station resolution and service laterals \\
\Lfive   & T2P5 & fwd.-eval.  & $+$ flow-dependent station pressure drop \\
\Lsix    & T2P6 & optimised   & $+$ transport delay \\
\NLref   &      & reference   & Nonlinear forward evaluator and bounded 72-hour re-solve; not a ladder rung \\
\midrule
$b_\ell$   & & & Estimation bias of level $\ell$ (optimised cost difference) \\
$r_\ell$   & & & Decision regret of level $\ell$ (forward-evaluated difference) \\
$\lambda$  & & & Loss number, annual loss / annual demand \\
$b$        & & & Loss burden, $=\lambda/(1+\lambda)$ \\
$L_t$      & & & Exogenous aggregate loss supplied to \CPL \\
$\kappa$   & & & Loss switch in the nodal balance, $\{0,1\}$ \\
\bottomrule
\end{tabular*}

\smallskip

\newcommand{\nom}[2]{\par\noindent\hangindent=2.0cm\hangafter=1\makebox[1.9cm][l]{$#1$}#2}
\setlength{\columnsep}{16pt}
\begin{multicols}{2}
\raggedright
\setlength{\parindent}{0pt}

\textbf{\DIFadd{Sets and indices}}\par\nopagebreak\smallskip
\nom{t \in \mathcal{T}}{Time steps ($|\mathcal{T}|=8760$)}
\nom{n \in \mathcal{N}}{Network nodes}
\nom{p \in \mathcal{P}}{Pipe segments}
\nom{g \in \mathcal{G}}{Thermal generators}
\nom{h \in \mathcal{H}}{Heat pumps}
\nom{s \in \mathcal{S}}{Thermal storage units}
\nom{r \in \mathcal{R}}{Electrode boilers (power-to-heat)}
\nom{k \in \{1,\ldots,K\}}{PWL segment index}
\nom{\mathcal{G}^{\mathrm{CHP}}}{Subset of CHP generators}
\nom{\mathcal{P}_n^{\mathrm{in/out}}}{Pipes entering/leaving node~$n$}
\nom{\mathcal{G}_n,\mathcal{H}_n,\mathcal{R}_n,\mathcal{S}_n}{Generators/heat pumps/electrode boilers/storage at node~$n$}
\nom{\mathcal{T}_{\mathrm{val}}}{Validation time steps}

\medskip
\textbf{\DIFadd{Parameters}}\par\nopagebreak\smallskip
\nom{U^{s}_p, U^{r}_p}{Supply/return pipe heat-transfer coefficient [W\,m$^{-1}$\,K$^{-1}$]}
\nom{L_p}{Pipe length [m]}
\nom{D_p}{Pipe inner diameter [m]}
\nom{A_p}{Pipe cross-sectional area [m$^2$]}
\nom{R_p}{Pipe hydraulic resistance [Pa\,MW$^{-2}$]}
\nom{f_D}{Darcy friction factor}
\nom{\rho}{Fluid density [kg\,m$^{-3}$]}
\nom{c_p}{Specific heat capacity of water [J\,(kg\,K)$^{-1}$]}
\nom{\tau_p}{Transport delay [s] (Eq.~\eqref{eq:delay_time}; $k_p=\lfloor\tau_p/\Delta t\rfloor$ with $\Delta t=3600$\,s)}
\nom{k_p}{Delay in time steps}
\nom{\phi_{p,t}}{Temperature decay factor, $\in(0,1]$}
\nom{\phi_p^{\mathrm{nom}}}{Nominal decay factor}
\nom{T_{\mathrm{sup}}^{\mathrm{nom}}(t)}{Nominal supply temperature [$^\circ$C]}
\nom{T_{\mathrm{ret}}^{\mathrm{nom}}(t)}{Nominal return temperature [$^\circ$C]}
\nom{T^{\mathrm{nom}}(t)}{Nominal operating temperature [$^\circ$C]}
\nom{T_{\mathrm{ground}}(t)}{Ground temperature [$^\circ$C]}
\nom{\bar{Q}_g,\bar{Q}_h,\bar{Q}_r}{Installed thermal capacities [MW]}
\nom{\bar{Q}_p}{Rated pipe capacity [MW]}
\nom{\phi_h^{\mathrm{min}}}{Heat-pump minimum load fraction}
\nom{\eta_g}{Generator thermal efficiency}
\nom{\eta_s^{\mathrm{ch/dis}}}{Storage charge/discharge efficiency}
\nom{\eta_{\mathrm{pump}}}{Pump efficiency}
\nom{\eta_{\mathrm{Lorenz}}}{Lorenz cycle efficiency}
\nom{\alpha_s}{Storage standing-loss fraction per hour}
\nom{\mathrm{COP}_{h,t}^{\mathrm{WHR}}}{Waste-heat COP, time-varying}
\nom{\mathrm{COP}_h^{\mathrm{def}}}{Default-source COP, constant}
\nom{m_k^{(p)},\,c_k^{(p)}}{PWL segment slope and intercept (chord of $R_pQ^2$)}
\nom{\Delta T_{\mathrm{pp}}}{Heat-pump pinch-point temperature difference [K]}
\nom{\lambda_t^{\mathrm{buy/sell}}}{Electricity buy/sell price [EUR\,MWh$^{-1}$]}
\nom{\lambda_g^{\mathrm{fuel}}}{Specific fuel cost [EUR\,MWh$^{-1}$]}
\nom{\lambda^{\mathrm{dump}}}{Heat-curtailment penalty [EUR\,MWh$^{-1}$]}
\nom{\lambda^{\mathrm{CO_2}}}{Carbon price [EUR\,t$^{-1}$]}
\nom{\lambda^{\mathrm{dem}}}{Demand-charge rate [EUR\,MW$^{-1}$]}
\nom{\lambda^{\mathrm{cyc}}}{Storage-cycling cost [EUR\,MWh$^{-1}$]}
\nom{e_t^{\mathrm{grid}}}{Grid emission factor [kg\,MWh$^{-1}$]}
\nom{e_g^{\mathrm{fuel}}}{Fuel emission intensity [kg\,MWh$^{-1}$]}
\nom{f^{\mathrm{yr}}}{Annual scaling factor}
\nom{M^{\mathrm{grid}}}{Grid capacity, Big-M [MW]}
\nom{\Delta t}{Time-step length [h], $=1$}
\nom{K}{Number of PWL segments}
\nom{\mathrm{HI}}{Demand heterogeneity index (Appendix~\ref{app:hi_definition})}

\medskip
\textbf{\DIFadd{Cost terms}}\par\nopagebreak\smallskip
\nom{C^{\mathrm{en}}}{Electricity procurement net of feed-in revenue}
\nom{C^{\mathrm{fuel}}}{Fuel expenditure}
\nom{C^{\mathrm{dump}}}{Heat-curtailment penalty}
\nom{C^{\mathrm{CO_2}}}{Carbon cost}
\nom{C^{\mathrm{dem}}}{Demand charge on peak import}
\nom{C^{\mathrm{cyc}}}{Storage-cycling penalty}
\nom{Z}{Solver objective}
\nom{z^{\mathrm{econ}}}{Economic operating cost (reported basis)}

\medskip
\textbf{\DIFadd{Variables}}\par\nopagebreak\smallskip
\nom{Q_{g,t}^{\mathrm{th}}}{Generator thermal output [MW]}
\nom{Q_{h,t}}{Heat-pump thermal output [MW]}
\nom{Q_{h,t}^{\mathrm{WHR/def}}}{Heat-pump waste-heat/default output [MW]}
\nom{Q_{r,t}}{Electrode-boiler output [MW]}
\nom{Q_{p,t}}{Pipe heat flow [MW]}
\nom{Q_{p,t}^{\mathrm{loss,pipe}}}{Pipe thermal loss [MW]}
\nom{Q_{s,t}^{\mathrm{ch/dis}}}{Storage charge/discharge rate [MW]}
\nom{Q_{n,t}^{\mathrm{dem}}}{Nodal heat demand [MW]}
\nom{Q_{n,t}^{\mathrm{dump}}}{Curtailed heat at node~$n$ [MW]}
\nom{P_t^{\mathrm{buy/sell}}}{Grid import/export [MW]}
\nom{P_{g,t}^{\mathrm{el}}}{CHP electrical output [MW]}
\nom{P_{h,t}^{\mathrm{el}}}{Heat-pump electrical input [MW]}
\nom{P_{r,t}^{\mathrm{el}}}{Electrode-boiler electrical input [MW]}
\nom{P_t^{\mathrm{pump}}}{Network pumping power [MW]}
\nom{P^{\mathrm{peak}}}{Annual peak grid import [MW]}
\nom{\Delta p_{p,t}}{Pipe pressure drop [Pa]}
\nom{\dot{m}_{p,t}}{Pipe mass-flow rate [kg\,s$^{-1}$]}
\nom{T_{n,t}^{\mathrm{sup}}}{Node supply temperature [$^\circ$C]}
\nom{T_{p,t}^{\mathrm{out}}}{Pipe outlet temperature [$^\circ$C]}
\nom{E_{s,t}}{Storage state of charge [MWh]}
\nom{F_{g,t}}{Fuel input, heating-value basis [MW]}
\nom{E_t^{\mathrm{CO_2,grid}}}{Grid CO$_2$ emissions [kg]}
\nom{E_{g,t}^{\mathrm{CO_2,fuel}}}{Fuel CO$_2$ emissions [kg]}
\nom{W_{p,t}}{Bilinear auxiliary, $=Q_{p,t}^2$ [MW$^2$]}
\nom{u_{h,t}}{Heat-pump on/off, binary}
\nom{z_t}{Grid import mode, binary}
\nom{z_{p,t,k}}{Binary PWL segment selector, $\{0,1\}$ ($Q_{p,t,k}$ the flow on segment $k$ [MW])}

\end{multicols}

}%DIF >  end footnotesize}

\section{Introduction}
\label{sec:introduction}
{
\DIFadd{District-heating networks serving industrial demand }\DIFaddend are central to European
decarbonisation\DIFdelbegin \DIFdel{strategies}\DIFdelend , integrating heat pumps\DIFdelbegin \DIFdel{(HP)}\DIFdelend , combined heat and power \DIFdelbegin \DIFdel{(CHP), }\DIFdelend and thermal storage \DIFdelbegin \DIFdel{(TES) within
}\DIFdelend \DIFaddbegin \DIFadd{in
}\DIFaddend multi-source systems~\DIFdelbegin \DIFdel{\mbox{%DIFAUXCMD
\cite{Lund2014_4GDH,Werner2017,Lund2025_clean}}\hskip0pt%DIFAUXCMD
. A recent bibliometric
analysis confirms that DH research output has grown steadily over
the past decade, with optimisation  and modelling now constituting
one of the field's dominant thematic clusters~\mbox{%DIFAUXCMD
\cite{PereaMoreno2026}}\hskip0pt%DIFAUXCMD
.
Operational optimisation of such networks increasingly relies
}\DIFdelend \DIFaddbegin \DIFadd{\mbox{%DIFAUXCMD
\cite{Lund2014_4GDH}}\hskip0pt%DIFAUXCMD
, and their operational optimisation relies
increasingly }\DIFaddend on mathematical programming\DIFdelbegin \DIFdel{. The level of detail in
}\DIFdelend \DIFaddbegin \DIFadd{~\mbox{%DIFAUXCMD
\cite{PereaMoreno2026}}\hskip0pt%DIFAUXCMD
. District heating already
meets a substantial share of building heat demand in mature markets~\mbox{%DIFAUXCMD
\citep{Werner2017} }\hskip0pt%DIFAUXCMD
and
remains integral to clean-energy-system designs~\mbox{%DIFAUXCMD
\citep{Lund2025_clean}}\hskip0pt%DIFAUXCMD
. Large heat pumps in
particular are becoming central to its decarbonisation~\mbox{%DIFAUXCMD
\citep{David2017}}\hskip0pt%DIFAUXCMD
, which raises the
economic stakes of }\DIFaddend the \DIFdelbegin \DIFdel{modelling of heating networks depends heavily on whether the topology and thermo-hydraulics are simplified: }%DIFDELCMD < 

%DIFDELCMD < %%%
%DIF < 
%DIFDELCMD < \begin{enumerate}[nosep]
%DIFDELCMD <   \item %%%
\textbf{\DIFdel{Topology abstraction:}} %DIFAUXCMD
\DIFdel{the level of spatial
detail
        }\DIFdelend \DIFaddbegin \DIFadd{operational models that schedule them. How much of the network
such a model represents is a design choice, and it is made along two axes: the spatial
resolution }\DIFaddend at which the network \DIFdelbegin \DIFdel{is represented }\DIFdelend \DIFaddbegin \DIFadd{appears, }\DIFaddend from a single-bus copperplate to a pipe-by-pipe
graph\DIFaddbegin \DIFadd{; and the physical phenomena admitted, from a lossless energy balance to a model with
pressure-coupled hydraulics, temperature propagation and transport delay}\DIFaddend .
\DIFdelbegin %DIFDELCMD < \item %%%
\textbf{\DIFdel{Thermo-hydraulic accuracy:}} %DIFAUXCMD
\DIFdel{the set of physical phenomena modelled from zero loss energy balances to pressure drop coupled, temperature propagating, delayaware
        thermo-hydraulic models.
}%DIFDELCMD < \end{enumerate}
%DIFDELCMD < %%%
\DIFdelend 

\DIFdelbegin \DIFdel{At one extreme, }\emph{\DIFdel{copperplate}} %DIFAUXCMD
\DIFdelend \DIFaddbegin \DIFadd{The two extremes are well populated. Copperplate }\DIFaddend models aggregate all components and
demands \DIFaddbegin \DIFadd{on}\DIFaddend to a single bus, ignoring spatial structure and thermal losses entirely\DIFdelbegin \DIFdel{. \mbox{%DIFAUXCMD
\citet{Wirtz2020} }\hskip0pt%DIFAUXCMD
explicitly acknowledge
this limitation : }\DIFdelend \DIFaddbegin \DIFadd{, a
limitation their authors acknowledge: \mbox{%DIFAUXCMD
\citet{Wirtz2020} }\hskip0pt%DIFAUXCMD
note that }\DIFaddend ``the
validity of the global energy balance assumption depends on the network topology, size and
the operation strategy\DIFdelbegin \DIFdel{,''yet no quantification of }\DIFdelend \DIFaddbegin \DIFadd{'', without quantifying }\DIFaddend the resulting error\DIFdelbegin \DIFdel{has been
published. At the other extreme, }\emph{\DIFdel{full thermo-hydraulic}} %DIFAUXCMD
\DIFdelend \DIFaddbegin \DIFadd{. Full thermo-hydraulic
}\DIFaddend models resolve the \DIFdelbegin \DIFdel{physical pipe-by-pipe network , capturing thermal losses }\DIFdelend \DIFaddbegin \DIFadd{network pipe by pipe and capture the thermal losses that amount to
}\DIFaddend 5--15\,\% of annual heat supply~\DIFdelbegin \DIFdel{\mbox{%DIFAUXCMD
\cite{Frederiksen2013,Talebi2016}}\hskip0pt%DIFAUXCMD
, }\DIFdelend \DIFaddbegin \DIFadd{\mbox{%DIFAUXCMD
\cite{Frederiksen2013}}\hskip0pt%DIFAUXCMD
, along with }\DIFaddend pressure-dependent
pumping\DIFdelbegin \DIFdel{costs}\DIFdelend , temperature propagation \DIFdelbegin \DIFdel{, and transport delays arising
}\DIFdelend \DIFaddbegin \DIFadd{and the delays that follow }\DIFaddend from finite fluid velocity. A
\DIFdelbegin \DIFdel{further complication arises from the
}\emph{\DIFdel{mathematical formulation}}%DIFAUXCMD
\DIFdel{: extended thermo-hydraulic physics introduce }\DIFdelend \DIFaddbegin \DIFadd{third complication sits across both axes: extended physics introduces }\DIFaddend nonlinear and
bilinear terms incompatible with standard \DIFdelbegin \DIFdel{Mixed integer linear programming(MILP) solvers}\DIFdelend \DIFaddbegin \DIFadd{mixed-integer linear programming}\DIFaddend , forcing a
choice between linearisation\DIFdelbegin \DIFdel{(preserving
MILP tractability but introducing approximation
error) and nonlinear/quadratic formulations(preserving physical accuracy but
potentially sacrificing solve
time)}\DIFdelend \DIFaddbegin \DIFadd{, which preserves tractability at the cost of approximation
error, and native nonlinear formulations, which preserve the physics at the cost of solve
time}\DIFaddend .

\DIFdelbegin \DIFdel{A common assumption is that higher thermo-hydraulic accuracy and }\DIFdelend \DIFaddbegin \DIFadd{The working assumption in practice is that }\DIFaddend finer spatial resolution \DIFdelbegin \DIFdel{both yield proportional accuracygains, which can justify the associated computational cost }\DIFdelend \DIFaddbegin \DIFadd{and richer physics
both buy proportionate accuracy, and that the computational cost is therefore justified in
proportion. No quantitative evidence establishes that proportionality, or identifies which
axis delivers the return}\DIFaddend . \DIFdelbegin \DIFdel{Yet no quantitative evidence exists to confirm or refute this
proportionality, nor to identify which
modelling dimension delivers the dominant returnon complexity. A closely related effect has
been quantified for electricity transmission systems: low
spatial network }\DIFdelend \DIFaddbegin \DIFadd{The closest evidence comes from a neighbouring field: in
electricity systems, coarse spatial }\DIFaddend resolution causes capacity-expansion models to \DIFdelbegin \DIFdel{ignore }\DIFdelend \DIFaddbegin \DIFadd{miss
}\DIFaddend intra-regional congestion and \DIFdelbegin \DIFdel{underestimate system costs, whereas spatial aggregation choices dominate }\DIFdelend \DIFaddbegin \DIFadd{understate system cost, and the spatial aggregation choice
dominates }\DIFaddend other modelling simplifications in scenario
comparisons~\DIFdelbegin \DIFdel{\mbox{%DIFAUXCMD
\cite{Frysztacki2021,
FrewJacobson2016}}\hskip0pt%DIFAUXCMD
. Whether an analogous dominance of topology over
thermo-hydraulic detail holds for DH dispatch optimisation  has not been established.
}%DIFDELCMD < 

%DIFDELCMD < \subsection{Research gap}
%DIFDELCMD < \label{sec:gap}
%DIFDELCMD < %%%
\DIFdelend \DIFaddbegin \DIFadd{\mbox{%DIFAUXCMD
\cite{Frysztacki2021}}\hskip0pt%DIFAUXCMD
. A parallel power-systems analysis shows spatial and
temporal resolution trading off against one another under a fixed computational
budget~\mbox{%DIFAUXCMD
\citep{FrewJacobson2016}}\hskip0pt%DIFAUXCMD
, the same class of resolution-versus-physics trade this
paper prices for district heating. It is natural to expect the analogue in district
heating, and that expectation is worth stating precisely because this paper finds it
holds only in part. Once the losses are supplied by other means, the spatial structure itself
contributes little to dispatch cost, at least where all generation is co-located at a
single source; whether that holds under distributed generation is left as a
controlled open question. The distinction is
invisible to any comparison that changes both at once, which is the first of the two
limitations this study addresses.
}}
\DIFaddend 

\DIFdelbegin \DIFdel{The existing literature focuses on studies involving multidimensional adjustments: published comparisons change
multiple modelling dimensions simultaneously. For instance, comparing the full nonlinear pipe model with temperature-dependent
Coefficient of performance (COP) used by \mbox{%DIFAUXCMD
\citet{Vesterlund2017} }\hskip0pt%DIFAUXCMD
against the copperplate formulation of \mbox{%DIFAUXCMD
\citet{Gabrielli2018}
}\hskip0pt%DIFAUXCMD
conflates topology effects with (i)~None Linear Programming (NLP) vs.\ constant-loss
formulations, (ii)~temperature-dependent vs. \ constant COP, (iii)~the presence or absence of hydraulic coupling, and
(iv)~differing temporal resolutions.
The same pattern appears
systematically across the literature
(}%DIFDELCMD < \Cref{sec:related_work}%%%
\DIFdel{). Four specific gaps persist: %DIF < 
}%DIFDELCMD < \begin{enumerate}[nosep]
%DIFDELCMD <   \item %%%
\DIFdel{The isolated topology abstraction as the sole
        experimental variable while holding physics constant.
  }%DIFDELCMD < \item %%%
\DIFdel{The isolated thermo-hydraulic effect (adding
        pressure drop, temperature propagation, transport delay)
        with a similar topology}\DIFdelend \DIFaddbegin \DIFadd{Two features of the existing fidelity literature limit what its comparisons can
conclude. First, the comparisons are confounded: moving from a copperplate or a
zone-aggregated model to a node-resolved one changes the spatial representation of the
network and the representation of thermal losses }\emph{\DIFadd{together}}\DIFadd{, so a resulting cost
difference cannot be attributed to either: the very objection that a coarse model
``misses losses'' is indistinguishable from ``misses topology''. Second, and more
fundamentally, fidelity is judged by the wrong quantity. Studies report how far the
}\emph{\DIFadd{objective value}} \DIFadd{of a simplified formulation departs from a detailed one, which
measures the models against each other; but a planning or scheduling model is used to
}\emph{\DIFadd{make a decision}}\DIFadd{, and what matters is the cost of having decided with the simpler
model and then operated the network}\DIFaddend . \DIFdelbegin %DIFDELCMD < \item %%%
\DIFdel{The linearisation error introduced by MILP-compatible
        approximations of thermo-hydraulic physics has not been
        bounded against a quadratic reference in a DH dispatch
        context.
  }%DIFDELCMD < \item %%%
\DIFdel{A quantitative guide to
selecting piping topology and thermo-hydraulic design based on observable networkcharacteristics. }%DIFDELCMD < \end{enumerate}
%DIFDELCMD < 

%DIFDELCMD < \subsection{Contributions and research questions}
%DIFDELCMD < \label{sec:contributions}
%DIFDELCMD < 

%DIFDELCMD < %%%
\DIFdel{This paper addresses these four gaps through a controlled
experimental design with three orthogonal comparison dimensions.
The controlled comparison simultaneously holds topology
constant while varying physics, and holds thermo-hydraulics constant while
varying topology}\DIFdelend \DIFaddbegin \DIFadd{A formulation can misstate the objective yet
produce a schedule that executes well, or state the objective plausibly yet schedule
poorly; objective-difference comparisons cannot tell these apart}\DIFaddend . This paper \DIFdelbegin \DIFdel{makes four contributions, each answering one
research question: %DIF < 
}%DIFDELCMD < \begin{description}[style=nextline, nosep, leftmargin=1.5em]
%DIFDELCMD <   \item[Contribution 1 / RQ1 Topology effect:]
%DIFDELCMD <     %%%
\DIFdel{A five-level taxonomy and topology comparison within a unified MILP
    framework where scenarios share thermo-hydraulics. This
    isolates topology as the sole experimental variable and is
    validated against measured industrial operational data.
    }\emph{\DIFdel{How much does topology abstraction in the dispatch optimisation  change annual
    operational cost and }%DIFDELCMD < \ce{CO2} %%%
\DIFdel{emissions?}}
%DIFAUXCMD
%DIFDELCMD < 

%DIFDELCMD <   \item[Contribution 2 / RQ2 thermo-hydraulic effect:]
%DIFDELCMD <     %%%
\DIFdel{A controlled thermo-hydraulic accuracy comparison
    that isolates the effect of adding pressure drop and
    temperature-dependent loss propagation while holding topologyand formulation class constant.
}\emph{\DIFdel{Do advanced hydraulic and thermal calculations
    have a significant impact on load distribution in the dispatch optimisation ?}}
%DIFAUXCMD
%DIFDELCMD < 

%DIFDELCMD <   \item[Contribution 3 / RQ3 Linearisation error bounds:]
%DIFDELCMD <     %%%
\DIFdel{An upper-bound characterisation of the combined linearisation
    and transport-delay effect against an mixed integer quadratic constraint programming (MIQCP) reference
    ($\Lplus$ vs.\ $\Lnl$).
    }\emph{\DIFdel{What is the maximum cost deviation attributable to MILP
    linearisation and transport-delay omission?}}
%DIFAUXCMD
%DIFDELCMD < 

%DIFDELCMD <   \item[Contribution 4 / RQ4 Generalisability:]
%DIFDELCMD <     %%%
\DIFdel{Parametric analysis across 36~synthetic configurations
    yielding joint selection rules based on observable networkproperties.
    }\emph{\DIFdel{Which network properties determine whether detailed
    topology and/or extended physics representation is warranted in dispatch optimisation ?}}
%DIFAUXCMD
%DIFDELCMD < \end{description}
%DIFDELCMD < 

%DIFDELCMD < %%%
%DIF < % ============================================================
%DIF < % 2. RELATED WORK
%DIF < % ============================================================
%DIFDELCMD < \section{Related work}
%DIFDELCMD < \label{sec:related_work}
%DIFDELCMD < 

%DIFDELCMD < %%%
\DIFdel{In the field of dispatch strategies for DHN numerous studies have been published.
Given the wide range of applications, MILP modelling strategies for the topologies of DHN are a frequent topic of discussion (}%DIFDELCMD < \Cref{sec:rw_milp_topo}%%%
\DIFdel{).
The linearisation of thermo-hydraulic networks varies as well across these studies (}%DIFDELCMD < \Cref{sec:rw_thermo-hydraulic}%%%
\DIFdel{). }\DIFdelend \DIFaddbegin \DIFadd{addresses
both: a control condition that separates loss from topology, and a decision-regret
metric that separates estimation error from decision quality.
}\DIFaddend 

\DIFdelbegin %DIFDELCMD < \subsection{MILP optimisation  and spatial scale in DH systems}
%DIFDELCMD < \label{sec:rw_milp_topo}
%DIFDELCMD < %%%
\DIFdelend \DIFaddbegin \subsection{Literature review and research gap}
{
\DIFadd{Two strands of the district-heating literature bear on this paper: how finely a
dispatch model resolves the network, and how finely it resolves the network's physics.
Neither strand has been evaluated by the quantity that matters for a scheduling model:
the cost of the decisions it produces.
}\DIFaddend 

\DIFdelbegin \DIFdel{MILP has become the dominant methodology for operational optimisation  of DH systems~\mbox{%DIFAUXCMD
\cite{Fang2016,Guelpa2019review}}\hskip0pt%DIFAUXCMD
.
Open-source frameworks like }\DIFdelend \DIFaddbegin \DIFadd{Reviews of district-heating modelling and optimisation document a broad methodological
spectrum~\mbox{%DIFAUXCMD
\citep{Talebi2016}}\hskip0pt%DIFAUXCMD
. Mixed-integer linear programming dominates operational district-heating
optimisation~\mbox{%DIFAUXCMD
\cite{Fang2016}}\hskip0pt%DIFAUXCMD
, and the general-purpose frameworks in which such studies are
built (}\DIFaddend oemof~\cite{Hilpert2018_oemof}, PyPSA~\cite{Brown2018_pypsa},
FINE~\cite{Welder2018_FINE}, Calliope~\cite{Pfenninger2018_calliope},
urbs~\cite{Dorfner2016_urbs}\DIFdelbegin \DIFdel{support multi-carrier energy
system optimisation  but share common limitations for detailed DH
analysis: thermal losses are typically }\DIFdelend \DIFaddbegin \DIFadd{) are typically applied with aggregated or component-based
thermal networks: losses }\DIFaddend time-invariant \DIFdelbegin \DIFdel{, heat pump
COP is constant or seasonal, hydraulic coupling is absent, and
network topology is either fully aggregated or fixed at
a single
detail level . All five frameworks operate in the default setting at L1 in the taxonomy proposed here; the cost of this simplification
has not been
}\DIFdelend \DIFaddbegin \DIFadd{or absent, no hydraulic coupling, topology fixed at
one level of detail, at what this paper calls the copperplate level $\CP$; detailed
thermo-hydraulic coupling requires additional formulations. The cost of that simplification
is not }\DIFaddend quantified.

\DIFdelbegin \DIFdel{The DH modelling literature exhibits }\DIFdelend \DIFaddbegin \DIFadd{Between the copperplate and the full graph lies }\DIFaddend a continuous spectrum of spatial
aggregation. \DIFdelbegin \DIFdel{The reduction of detailed pipe networks to a
small number of aggregated branches has a long tradition in DH
simulation: }\DIFdelend \citet{Larsen2004} compare aggregation methods on a 23-substation network
and show that pipe count can be reduced by roughly an order of magnitude before
production and return-temperature errors become significant\DIFdelbegin \DIFdel{, while \mbox{%DIFAUXCMD
\citet{Falay2020} }\hskip0pt%DIFAUXCMD
extend this line
of work to large-scale }\DIFdelend \DIFaddbegin \DIFadd{; \mbox{%DIFAUXCMD
\citet{Falay2020} }\hskip0pt%DIFAUXCMD
carry the
approach to }\DIFaddend dynamic simulation of several thousand substations. \DIFdelbegin \DIFdel{\mbox{%DIFAUXCMD
\citet{Schweiger2018_paradigms} }\hskip0pt%DIFAUXCMD
survey the resulting
landscape of DH modelling paradigms and general-purpose tools. These
studies }\DIFdelend \DIFaddbegin \DIFadd{Both }\DIFaddend establish
\emph{which} \DIFdelbegin \DIFdel{aggregation preserves }\DIFdelend \DIFaddbegin \DIFadd{aggregations preserve }\DIFaddend simulation fidelity, \DIFdelbegin \DIFdel{but none }\DIFdelend \DIFaddbegin \DIFadd{and neither }\DIFaddend quantifies the
resulting bias in \DIFdelbegin \emph{\DIFdel{dispatch-optimisation  cost}}%DIFAUXCMD
\DIFdel{. Single-node models are standard in national-scale analyses~\mbox{%DIFAUXCMD
\cite{Connolly2014_EnergyPLAN,
Brown2018_pypsa,Lund2017_SmartEnergy}}\hskip0pt%DIFAUXCMD
.
\mbox{%DIFAUXCMD
\citet{Gabrielli2018} }\hskip0pt%DIFAUXCMD
formulated a MILP for distributed
}\DIFdelend \DIFaddbegin \DIFadd{dispatch-optimisation cost. The same emphasis persists in the most
recent work: \mbox{%DIFAUXCMD
\citet{Vrain2025} }\hskip0pt%DIFAUXCMD
formalise an aggregation method with an explicit measure of
the error it induces in large-scale }\DIFaddend multi-energy \DIFdelbegin \DIFdel{systems including thermal and electrical storage but
treat the thermal network as a single aggregated efficiency
factor}\DIFdelend \DIFaddbegin \emph{\DIFadd{simulation}}\DIFadd{, and \mbox{%DIFAUXCMD
\citet{Cassetti2025}
}\hskip0pt%DIFAUXCMD
show how spatial resolution shapes the outcome of regional heat-network }\emph{\DIFadd{planning}}\DIFadd{.
Both confirm that aggregation error is resolution-dependent, yet both judge simulated or
planned quantities rather than the bias a coarse network imparts to the }\emph{\DIFadd{operational
dispatch decision}}\DIFadd{, and neither isolates thermal loss from spatial topology}\DIFaddend . Zone-based models \DIFdelbegin \DIFdel{(3--15~nodes) }\DIFdelend aggregate demand into \DIFaddbegin \DIFadd{a
handful of }\DIFaddend geographic clusters~\DIFdelbegin \DIFdel{\mbox{%DIFAUXCMD
\cite{Wirtz2020,Mertz2016,Buffa2019}}\hskip0pt%DIFAUXCMD
;
\mbox{%DIFAUXCMD
\citet{Buenning2020} }\hskip0pt%DIFAUXCMD
demonstrate that zone-level aggregation
captures 85--95\,\% of dynamic temperature variation. Full-graph
models resolve
node-by-node }\DIFdelend \DIFaddbegin \DIFadd{\mbox{%DIFAUXCMD
\cite{Wirtz2020}}\hskip0pt%DIFAUXCMD
, while full-graph models resolve
the network node by node }\DIFaddend from GIS data~\DIFdelbegin \DIFdel{\mbox{%DIFAUXCMD
\cite{Vesterlund2017,Schweiger2017,vanderHeijde2017,
Blommaert2020}}\hskip0pt%DIFAUXCMD
; \mbox{%DIFAUXCMD
\citet{Giraud2015} }\hskip0pt%DIFAUXCMD
show tractability up to
${\approx}\,100$~nodes for hourly annual horizons. }\DIFdelend \DIFaddbegin \DIFadd{\mbox{%DIFAUXCMD
\cite{vanderHeijde2017}}\hskip0pt%DIFAUXCMD
, an approach demonstrated in
detailed dynamic-simulation libraries~\mbox{%DIFAUXCMD
\cite{Giraud2015}}\hskip0pt%DIFAUXCMD
. \mbox{%DIFAUXCMD
\citet{Vesterlund2017} }\hskip0pt%DIFAUXCMD
apply
full-graph optimisation to a real multi-source network, again to operate it rather than to
price the cost of aggregating it. Single-node
representations remain standard in national-scale
analysis~\mbox{%DIFAUXCMD
\cite{Connolly2014_EnergyPLAN}}\hskip0pt%DIFAUXCMD
, and \mbox{%DIFAUXCMD
\citet{Gabrielli2018} }\hskip0pt%DIFAUXCMD
likewise reduce the
thermal network of a distributed multi-energy system to a single aggregate efficiency.
Bidirectional, low-temperature fifth-generation designs~\mbox{%DIFAUXCMD
\citep{Buffa2019} }\hskip0pt%DIFAUXCMD
add flow-direction
choices beyond the radial networks considered here.
}\DIFaddend 

\DIFdelbegin \DIFdel{Across this spectrum, the level of spatial aggregation has been
shown to materially affect dispatch
results~\mbox{%DIFAUXCMD
\cite{Wirtz2020,Buffa2019}}\hskip0pt%DIFAUXCMD
, but no scalar metric links
observable networkproperties to
the required resolution level. \mbox{%DIFAUXCMD
\citet{Leitner2019} }\hskip0pt%DIFAUXCMD
compared copperplate and full-graph
representations for an Austrian DH network and found cost
differences of 3--8\,\%, but simultaneously changed the loss (more instances in }%DIFDELCMD < \Cref{sec:rw_summary}%%%
\DIFdel{).
Three very recent studies extend this line without isolating
topology
as an independent factor. \mbox{%DIFAUXCMD
\citet{Trentmann2026} }\hskip0pt%DIFAUXCMD
integrate
}\DIFdelend \DIFaddbegin \DIFadd{Where the spectrum has been traversed, the comparison has typically confounded topology
with losses: moving from a lumped to a resolved network introduces pipe losses at the same
time as spatial structure, so a reported cost difference cannot be attributed to either.
This confound is structural rather than incidental: the objection that a coarse model
``misses losses'' is empirically indistinguishable from ``misses topology''. Recent }\DIFaddend full-graph \DIFdelbegin \DIFdel{DHN expansion into }\DIFdelend \DIFaddbegin \DIFadd{studies do not resolve it,
because they do not vary resolution at all: \mbox{%DIFAUXCMD
\citet{Trentmann2026} }\hskip0pt%DIFAUXCMD
embed full-graph
network expansion in }\DIFaddend an energy-hub MILP with \DIFdelbegin \DIFdel{temporally-resolved }\DIFdelend \DIFaddbegin \DIFadd{temporally resolved }\DIFaddend heat-pump \DIFdelbegin \DIFdel{COPs, showing that network supply
temperature strongly shifts the optimal centralized-vs-building-level
supply split; their model always resolves the network at
full-graph detail,
so the copperplate/aggregated comparison studied
here is outside their scope. \mbox{%DIFAUXCMD
\citet{Ceruti2025} }\hskip0pt%DIFAUXCMD
likewise }\DIFdelend \DIFaddbegin \DIFadd{performance,
\mbox{%DIFAUXCMD
\citet{Ceruti2025} }\hskip0pt%DIFAUXCMD
}\DIFaddend formulate a full-graph \DIFdelbegin \DIFdel{MILP for DHN design with TES and find that
temporal (not spatial ) }\DIFdelend \DIFaddbegin \DIFadd{design MILP with storage and conclude that
temporal rather than spatial }\DIFaddend resolution dominates design quality\DIFdelbegin \DIFdel{. \mbox{%DIFAUXCMD
\citet{Vieth2026} }\hskip0pt%DIFAUXCMD
address }\DIFdelend \DIFaddbegin \DIFadd{, and
\mbox{%DIFAUXCMD
\citet{Vieth2026} }\hskip0pt%DIFAUXCMD
treat }\DIFaddend expansion planning of existing full-graph networks under
resilience constraints\DIFdelbegin \DIFdel{, showing that considering generation-unit outages
lowers the connectable peak demand by up to 97\,\%. 
}%DIFDELCMD < 

%DIFDELCMD < %%%
\DIFdel{A parallel and largely separate literature quantifies the dispatch
value of the generation technologies studied in the present case
(HP, electrode boiler, CHP, TES) without addressing spatial
resolution. \mbox{%DIFAUXCMD
\citet{Nielsen2016} }\hskip0pt%DIFAUXCMD
and \mbox{%DIFAUXCMD
\citet{Boettger2015} }\hskip0pt%DIFAUXCMD
show, for
the Danish and German electricity markets respectively, that
deterministic dispatch models systematically overestimate the economic value of HP and electric-boiler flexibility relative to
stochastic or control-power-aware formulations. \mbox{%DIFAUXCMD
\citet{Galteland2026} }\hskip0pt%DIFAUXCMD
confirm this technology-value
question is still actively studied : for a hybrid industrial
electrode-boiler/steam-accumulator/battery system, they show optimal
sizing is highly sensitive to reserve-market participation and
storage cost assumptions. \mbox{%DIFAUXCMD
\citet{Mollenhauer2018} }\hskip0pt%DIFAUXCMD
formulate an MILP
unit-commitment and dispatch problem for German CHP plants
retrofitted with HP and TES at hourly resolution over a full year,
methodologically close to the present dispatch model but, like
\mbox{%DIFAUXCMD
\citet{Schweiger2017_p2h,Averfalk2017,David2017}}\hskip0pt%DIFAUXCMD
, representing the DH network itself as a single aggregated demand
node. \mbox{%DIFAUXCMD
\citet{Walden2023} }\hskip0pt%DIFAUXCMD
extend this line to a nonlinear industrial
power-to-heat system coupling a high-temperature heat pump, sensible
TES, and wind generation, optimizing operation with an accurate
nonlinear component model rather than a network representation; their setting shares this paper's technology mix (HP-based
electrification, TES buffering of variable renewable supply) but ,
like \mbox{%DIFAUXCMD
\citet{Mollenhauer2018}}\hskip0pt%DIFAUXCMD
, treats the distribution network as a single lumped point rather than a resolved topology.
\mbox{%DIFAUXCMD
\citet{Riedl2026} }\hskip0pt%DIFAUXCMD
propose, most recently, a general flexibility
assessment framework for DHC systems that explicitly incorporates
network-side constraints alongside generation-side flexibility ,
though their published
case studies do not yet vary spatial resolution as an experimental
factor. \mbox{%DIFAUXCMD
\citet{Buhler2017} }\hskip0pt%DIFAUXCMD
instead quantify the industrial excess-heat
resource available to DH networks spatially, but not within a
dispatch-optimisation  framework. None of these studies varies
network topologyas an experimental factor, reinforcing that the
topology/technology interaction remains unexplored}\DIFdelend . \DIFdelbegin %DIFDELCMD < 

%DIFDELCMD < \subsection{Thermo-hydraulic network models and linearisation}
%DIFDELCMD < \label{sec:rw_thermo-hydraulic}
%DIFDELCMD < %%%
\DIFdelend \DIFaddbegin \DIFadd{Each resolves the network fully throughout, so the aggregated
comparison studied here lies outside their scope. Closest to the decision-oriented view
taken here, \mbox{%DIFAUXCMD
\citet{Friedrich2025} }\hskip0pt%DIFAUXCMD
trace how the fidelity of the network model propagates
into operational efficiency; but they too assess the simulated operation of a single model
rather than the regret of executing a lower-fidelity schedule under a higher-fidelity one,
and do not separate loss visibility from spatial topology.
}\DIFaddend 

Beyond topology, \DIFdelbegin \DIFdel{the }\DIFdelend \DIFaddbegin \DIFadd{pipe models differ substantially in }\DIFaddend thermo-hydraulic \DIFdelbegin \DIFdel{fidelity of pipe models varies
substantially }\DIFdelend \DIFaddbegin \DIFadd{fidelity}\DIFaddend .
Steady-state thermal losses follow from the fluid--ground temperature difference and \DIFdelbegin \DIFdel{pipe insulation
properties }\DIFdelend \DIFaddbegin \DIFadd{the
pipe insulation, }\DIFaddend and are MILP-compatible \DIFdelbegin \DIFdel{when
temperatures are }\DIFdelend \DIFaddbegin \DIFadd{whenever temperatures enter as }\DIFaddend fixed
parameters~\cite{Frederiksen2013}. \DIFdelbegin \DIFdel{The }\DIFdelend \DIFaddbegin \DIFadd{Pressure drop is not: the }\DIFaddend Darcy--Weisbach \DIFdelbegin \DIFdel{equation,
in contrast, relates pressure drop to the square of flow velocity ,
introducing a quadratic nonlinearity; pumping power is proportional
to $\dot{m} \cdot \Delta p$, creating }\DIFdelend \DIFaddbegin \DIFadd{relation is
quadratic in flow velocity and pumping power adds }\DIFaddend a cubic term\DIFdelbegin \DIFdel{. Authors have
addressed this through }\DIFdelend \DIFaddbegin \DIFadd{, handled in the literature
either by }\DIFaddend piecewise linearisation~\DIFdelbegin \DIFdel{\mbox{%DIFAUXCMD
\cite{Bordin2016,Mertz2016} }\hskip0pt%DIFAUXCMD
or }\DIFdelend \DIFaddbegin \DIFadd{\mbox{%DIFAUXCMD
\cite{Bordin2016} }\hskip0pt%DIFAUXCMD
or by }\DIFaddend second-order cone
\DIFdelbegin \DIFdel{relaxations~\mbox{%DIFAUXCMD
\cite{Blommaert2020}}\hskip0pt%DIFAUXCMD
. \mbox{%DIFAUXCMD
\citet{Hari2024} }\hskip0pt%DIFAUXCMD
provide a
comprehensive thermal network flow optimisation  framework
integrating }\DIFdelend \DIFaddbegin \DIFadd{relaxation~\mbox{%DIFAUXCMD
\cite{Blommaert2020}}\hskip0pt%DIFAUXCMD
; \mbox{%DIFAUXCMD
\citet{Hari2024} }\hskip0pt%DIFAUXCMD
integrate }\DIFaddend Darcy--Weisbach pressure drop
with explicit pump placement \DIFaddbegin \DIFadd{in a general thermal network flow framework. \mbox{%DIFAUXCMD
\citet{Mertz2016}
}\hskip0pt%DIFAUXCMD
instead pose the joint network configuration and design as a mixed-integer nonlinear program}\DIFaddend .
Reported pumping costs \DIFdelbegin \DIFdel{range from }\DIFdelend \DIFaddbegin \DIFadd{span }\DIFaddend 1--5\,\% of thermal energy cost\DIFdelbegin \DIFdel{s}\DIFdelend  ~\cite{Frederiksen2013}. \DIFdelbegin %DIFDELCMD < 

%DIFDELCMD < %%%
\DIFdelend \DIFaddbegin \DIFadd{This paper returns to
that range, since the network studied here falls two orders of magnitude below it for
reasons that are physical rather than numerical. }\DIFaddend Temperature propagation along \DIFdelbegin \DIFdel{pipes }\DIFdelend \DIFaddbegin \DIFadd{a pipe
}\DIFaddend introduces exponential and bilinear terms~\cite{vanderHeijde2017}\DIFdelbegin \DIFdel{; \mbox{%DIFAUXCMD
\citet{Maldonado2024} }\hskip0pt%DIFAUXCMD
validate this model }\DIFdelend \DIFaddbegin \DIFadd{, validated by
\mbox{%DIFAUXCMD
\citet{Maldonado2024} }\hskip0pt%DIFAUXCMD
}\DIFaddend against industrial data \DIFdelbegin \DIFdel{, achieving supply
temperature Mean absolute error
(MAE) below }%DIFDELCMD < \SI{0.5}{\kelvin}%%%
\DIFdel{. Transport delays of 1--3}\DIFdelend \DIFaddbegin \DIFadd{to a supply-temperature mean absolute error
below 0.5}\DIFaddend \,\DIFdelbegin \DIFdel{h arise for
}\DIFdelend \DIFaddbegin \DIFadd{K, and transport delays of one to three hours arise in }\DIFaddend long
pipes~\DIFdelbegin \DIFdel{\mbox{%DIFAUXCMD
\cite{vanderHeijde2017,Giraud2015}}\hskip0pt%DIFAUXCMD
; the
fixed-delay approximation is most accurate when flow variations
are moderate ($\pm 30\,\%$ of nominal). }%DIFDELCMD < 

%DIFDELCMD < %%%
\DIFdel{MIQCP formulationsoffer a middle ground between MILP and full
NLP~\mbox{%DIFAUXCMD
\cite{Ommen2016}}\hskip0pt%DIFAUXCMD
: }\DIFdelend \DIFaddbegin \DIFadd{\mbox{%DIFAUXCMD
\cite{Giraud2015}}\hskip0pt%DIFAUXCMD
. Equation-based Modelica libraries capture this thermo-hydraulic
coupling for dynamic simulation~\mbox{%DIFAUXCMD
\citep{Schweiger2017}}\hskip0pt%DIFAUXCMD
. Across dispatch formulations, methods span linear, mixed-integer
and nonlinear programming~\mbox{%DIFAUXCMD
\cite{Ommen2016}}\hskip0pt%DIFAUXCMD
; MIQCP occupies a middle ground in which
}\DIFaddend quadratic constraints are handled natively \DIFdelbegin \DIFdel{by modern solvers }\DIFdelend while integer variables retain
unit-commitment logic. \citet{Hering2021} \DIFdelbegin \DIFdel{and \mbox{%DIFAUXCMD
\citet{Hering2022} }\hskip0pt%DIFAUXCMD
demonstrate this
approach for DH }\DIFdelend \DIFaddbegin \DIFadd{demonstrate the approach for }\DIFaddend design and
dispatch with multiple heat pumps \DIFdelbegin \DIFdel{,
formulating temperature-dependent constraints as MIQCP for a 4GDH
network; their case studies show the MIQCP remains tractable }\DIFdelend at the scale of tens of consumers, but \DIFdelbegin \DIFdel{do not report a reference-model
}\DIFdelend \DIFaddbegin \DIFadd{report no
reference }\DIFaddend comparison isolating the cost of \DIFdelbegin \DIFdel{PWL-MILP approximation relative to
}\DIFdelend \DIFaddbegin \DIFadd{the piecewise-linear approximation against }\DIFaddend the
native quadratic formulation. \DIFdelbegin \DIFdel{However, MIQCP problems lack
LP-relaxation tightness, potentially increasing solve times. }\DIFdelend The magnitude of \DIFdelbegin \DIFdel{the }\DIFdelend \DIFaddbegin \DIFadd{that }\DIFaddend linearisation error in a
\DIFdelbegin \DIFdel{DH dispatch context i.e.,
the cost gap between Piece wise linear (PWL)-MILP and native MIQCP
formulations }\DIFdelend \DIFaddbegin \DIFadd{district-heating dispatch setting }\DIFaddend has not been published.

\DIFdelbegin %DIFDELCMD < \subsection{Positioning and research gap}
%DIFDELCMD < \label{sec:rw_summary}
%DIFDELCMD < \Cref{tab:prior_work} %%%
\DIFdel{reveals the systematic confounding of eleven representative studies spanning copperplate to full-graph
resolution, }\emph{\DIFdel{none}} %DIFAUXCMD
\DIFdel{isolates topology from physics as an
independent experimental variable.
}%DIFDELCMD < \Cref{tab:prior_work} %%%
\DIFdel{positions representative studies. No prior work isolates topology from physics, nor bounds linearisation error against a quadratic
reference in a
controlled DH dispatch setting. }%DIFDELCMD < 

%DIFDELCMD < %%%
\DIFdel{Two studies come closest to the present design. \mbox{%DIFAUXCMD
\citet{Quaggiotto2021}
}\hskip0pt%DIFAUXCMD
report, }\DIFdelend \DIFaddbegin \DIFadd{Two features of this literature limit what its comparisons can conclude, and both are
visible in Table~\ref{tab:priorwork}. The first is the confound above: no prior study
separates the effect of representing thermal losses from the effect of representing
spatial structure, so the two are never independently priced. The second is more
fundamental. Fidelity is judged by comparing }\emph{\DIFadd{objective values}} \DIFadd{across formulations,
which measures the models against one another; but a scheduling model is used to make a
decision, and what matters is the cost of having decided with the simpler model and then
operated the network. A formulation can misstate its objective yet produce a schedule
that executes well, or state its objective plausibly yet schedule badly, and an
objective-difference comparison cannot distinguish the two. The closest prior work
approaches this without closing it: \mbox{%DIFAUXCMD
\citet{Quaggiotto2021} }\hskip0pt%DIFAUXCMD
report }\DIFaddend for the Verona \DIFdelbegin \DIFdel{DH network
, that simplifying the network's
}\DIFdelend \DIFaddbegin \DIFadd{network
that simplifying }\DIFaddend thermal behaviour in a dispatch optimisation underestimates operational
cost relative to a model-predictive controller using a detailed network model, \DIFdelbegin \DIFdel{but without a taxonomy that separates topology from
thermo-hydraulic fidelity or quantifies the gap as a function of
network properties. \mbox{%DIFAUXCMD
\citet{Jansen2024} }\hskip0pt%DIFAUXCMD
independently confirm this
}\DIFdelend \DIFaddbegin \DIFadd{and
\mbox{%DIFAUXCMD
\citet{Jansen2024} }\hskip0pt%DIFAUXCMD
confirm the }\DIFaddend direction with a mixed-integer \DIFdelbegin \DIFdel{non-linear MPC for DH networks,
reporting a similar fidelity/performance trade-off for a real
}\DIFdelend \DIFaddbegin \DIFadd{nonlinear MPC on a }\DIFaddend Belgian
network, but \DIFdelbegin \DIFdel{again without isolating topology as a factor. \mbox{%DIFAUXCMD
\citet{Wirtz2021_complexity} }\hskip0pt%DIFAUXCMD
systematically }\DIFdelend \DIFaddbegin \DIFadd{neither separates topology from thermo-hydraulic fidelity nor expresses the
gap as a function of observable network properties. \mbox{%DIFAUXCMD
\citet{Wirtz2021_complexity} }\hskip0pt%DIFAUXCMD
}\DIFaddend vary the
level of detail across 24 MILP formulations of a multi-energy system and identify which
simplifications are cost-relevant\DIFdelbegin \DIFdel{; their
experimental design closely parallels }\DIFdelend \DIFaddbegin \DIFadd{, }\DIFaddend the controlled-comparison logic \DIFdelbegin \DIFdel{adopted here}\DIFdelend \DIFaddbegin \DIFadd{this paper adopts}\DIFaddend ,
but their \DIFdelbegin \DIFdel{model }\DIFdelend complexity axis spans technology and part-load representation rather than
network topology and thermo-hydraulics, and the building-scale \DIFdelbegin \DIFdel{multi-energy }\DIFdelend system they study does not
resolve pipe-by-pipe \DIFdelbegin \DIFdel{network }\DIFdelend structure. \citet{Kotzur2021} review \DIFdelbegin \DIFdel{complexity-management methods for
energy system }\DIFdelend \DIFaddbegin \DIFadd{complexity management for
energy-system }\DIFaddend optimisation more broadly and \DIFdelbegin \DIFdel{note }\DIFdelend \DIFaddbegin \DIFadd{observe }\DIFaddend that, while temporal aggregation is
\DIFdelbegin \DIFdel{well-studied}\DIFdelend \DIFaddbegin \DIFadd{well studied}\DIFaddend , spatial aggregation \DIFdelbegin \DIFdel{"}\DIFdelend \DIFaddbegin \DIFadd{``}\DIFaddend has only a small community\DIFdelbegin \DIFdel{" }\DIFdelend \DIFaddbegin \DIFadd{'' }\DIFaddend and lacks systematic
guidance.
\DIFdelbegin \DIFdel{The present
paper closes the gap by comparing
systematic, axis-isolated logic specifically to the topology/physics
trade-off in DH dispatch , and within a single framework spanning all three dimensions (topology, thermo-hydraulics,
linearisation).
}%DIFDELCMD < 

%DIFDELCMD < %%%
%DIF < % ============================================================
%DIF < % 3. METHODOLOGY
%DIF < % ============================================================
%DIF < % ============================================================
%DIF < % 3. METHODOLOGY
%DIF < % ============================================================
%DIFDELCMD < \section{Methodology}
%DIFDELCMD < \label{sec:methodology}
%DIFDELCMD < %%%
\DIFdelend 

\DIFdelbegin \DIFdel{This section defines the controlled experimental design
(}%DIFDELCMD < \Cref{sec:exp_design}%%%
\DIFdel{), presents the mathematical formulationsin
order of increasing fidelity (}%DIFDELCMD < \Cref{sec:milp,sec:extended}%%%
\DIFdel{),
describes the validation protocol (}%DIFDELCMD < \Cref{sec:val_protocol}%%%
\DIFdel{), and closes with implementation details (}%DIFDELCMD < \Cref{sec:model_impl}%%%
\DIFdel{). }\DIFdelend \DIFaddbegin \DIFadd{A separate literature quantifies the dispatch value of the generation technologies
studied here while leaving the network lumped. The operational role of thermal storage in
these systems is surveyed by \mbox{%DIFAUXCMD
\citet{Guelpa2019review}}\hskip0pt%DIFAUXCMD
. \mbox{%DIFAUXCMD
\citet{Mollenhauer2018} }\hskip0pt%DIFAUXCMD
formulate hourly
annual unit commitment for German CHP plants retrofitted with heat pumps and thermal
storage (methodologically close to the dispatch model used here, but representing the
network as a single aggregated demand node), and \mbox{%DIFAUXCMD
\citet{Walden2023} }\hskip0pt%DIFAUXCMD
optimise an
industrial power-to-heat system with an accurate nonlinear component model and no network
representation at all. \mbox{%DIFAUXCMD
\citet{Nielsen2016} }\hskip0pt%DIFAUXCMD
show for the Danish market that deterministic
dispatch systematically overestimates the economic value of heat-pump and electric-boiler
flexibility relative to stochastic formulations. Large heat pumps are increasingly deployed
in Swedish district-heating systems~\mbox{%DIFAUXCMD
\citep{Averfalk2017}}\hskip0pt%DIFAUXCMD
. Power-to-heat conversion more
broadly couples the heat and electricity sectors~\mbox{%DIFAUXCMD
\citep{Schweiger2017_p2h}}\hskip0pt%DIFAUXCMD
. \mbox{%DIFAUXCMD
\citet{Riedl2026} }\hskip0pt%DIFAUXCMD
propose a flexibility
assessment framework for district heating and cooling that incorporates network-side
constraints alongside generation-side flexibility, though the published case studies do
not yet vary spatial resolution as an experimental factor. None of this work treats
network representation as a factor, so the interaction between topology and technology
value remains open.
}\DIFaddend 

\DIFdelbegin %DIFDELCMD < \begin{table*}[!tbp]
%DIFDELCMD < %%%
\DIFdelendFL \DIFaddbeginFL \DIFaddFL{This paper addresses both gaps directly: a control condition that separates loss
visibility from spatial resolution, and a decision-regret metric that separates
estimation error from decision quality. The second is not a new idea in optimisation:
it is analogous to the value of the stochastic solution~\mbox{%DIFAUXCMD
\citep{BirgeLouveaux2011}}\hskip0pt%DIFAUXCMD
, transposed
from the uncertainty axis to the resolution axis. The value of the stochastic solution itself
extends to multistage decision problems~\mbox{%DIFAUXCMD
\citep{Escudero2007VSS}}\hskip0pt%DIFAUXCMD
. The analogy is structural rather than exact:
the forward valuation is an overlay under a common high-fidelity model, so, unlike the classical
value of the stochastic solution, it need not be sign-definite (a loss-blind schedule can
forward-evaluate as cheaper on a mispriced term). The claimed contribution lies in the
transposition and the empirical finding rather than in the concept itself.
}}
\begin{table}[t]
\DIFaddendFL \centering
\DIFdelbeginFL %DIFDELCMD < \caption{DH optimisation  studies with
%DIFDELCMD < isolated topology and thermo-hydraulics as experimental
%DIFDELCMD < variables.}
%DIFDELCMD < \label{tab:prior_work}
%DIFDELCMD < \begin{threeparttable}
%DIFDELCMD < \footnotesize
%DIFDELCMD < \setlength{\tabcolsep}{4pt}
%DIFDELCMD < \begin{tabular}{@{}l c c c c c c c@{}}
%DIFDELCMD < \toprule
%DIFDELCMD < Reference & Topo. & Physics & Form. & COP &
%DIFDELCMD <   Scale & Temp. & Comparison \\
%DIFDELCMD < \midrule
%DIFDELCMD < \citet{Connolly2014_EnergyPLAN}
%DIFDELCMD <   & Cop. & Fix.~$\eta$ & LP & Const. & Nat. & Ann. & No \\
%DIFDELCMD < \citet{Gabrielli2018}
%DIFDELCMD <   & Cop. & Fix.~$\eta$ & MILP & Const. & 10\,MW & Ann. & No \\
%DIFDELCMD < \citet{Wirtz2020}
%DIFDELCMD <   & Zone(5) & Lin.\ loss & MILP & Seas. & 25\,MW & Ann. & No \\
%DIFDELCMD < \citet{Mertz2016}
%DIFDELCMD <   & Zone(8) & Lin.+$\Delta p$ & MINLP & Const.
%DIFDELCMD <   & 20\,MW & Ann. & No \\
%DIFDELCMD < \citet{Hari2024}
%DIFDELCMD <   & Full & DW+pumps & NLP & -- & 50\,MW & Steady & No \\
%DIFDELCMD < \citet{Leitner2019}
%DIFDELCMD <   & L1+L3 & Diff.\ phys. & MILP & Const. & 12\,MW
%DIFDELCMD <   & Ann. & No \\
%DIFDELCMD < \citet{Bordin2016}
%DIFDELCMD <   & Full & PWL~$\Delta p$ & MILP & -- & 40\,MW & Ann. & No \\
%DIFDELCMD < \citet{Giraud2015}
%DIFDELCMD <   & Full & Steady-st. & NLP & $T$-dep. & 30\,MW & Ann. & No \\
%DIFDELCMD < \citet{Vesterlund2017}
%DIFDELCMD <   & Full & NLP therm. & NLP & $T$-dep. & 50\,MW & Ann. & No \\
%DIFDELCMD < \citet{vanderHeijde2017}
%DIFDELCMD <   & Full & Pipe comp. & Var. & -- & Review & --
%DIFDELCMD <   & Part.\tnote{a} \\
%DIFDELCMD < \citet{Blommaert2020}
%DIFDELCMD <   & Full & Adj.\ NLP & NLP & $T$-dep. & 30\,MW & Des. & No \\
%DIFDELCMD < \citet{Quaggiotto2021}
%DIFDELCMD <   & Full & FD therm. & MPC & -- & Dist. & Days & No \\
%DIFDELCMD < \citet{Jansen2024}
%DIFDELCMD <   & Full & MINLP & MPC & -- & Dist. & Days & No \\
%DIFDELCMD < \citet{Hering2021}
%DIFDELCMD <   & Fix. & MIQCP & MIQCP & $T$-dep. & 7 cons. & Des. & No \\
%DIFDELCMD < \citet{Wirtz2021_complexity}
%DIFDELCMD <   & Cop. & 24 var. & MILP & Var. & Bldg. & Ann. & Part.\tnote{b} \\
%DIFDELCMD < \midrule
%DIFDELCMD < \textbf{This work}
%DIFDELCMD <   & \textbf{$L1$--$\Lnl$}
%DIFDELCMD <   & Bas.+ext.
%DIFDELCMD <   & MILP+QCP
%DIFDELCMD <   & Pre-c.
%DIFDELCMD <   & 12\,MW
%DIFDELCMD <   & 8760\,h
%DIFDELCMD <   & Yes \\
%DIFDELCMD < \bottomrule
%DIFDELCMD < \end{tabular}
%DIFDELCMD < \begin{tablenotes}\footnotesize
%DIFDELCMD <   \item[a] %%%
\DIFdelFL{Compared pipe-loss formulations but not spatial
    aggregation levels. }%DIFDELCMD < \item[b] %%%
\DIFdelFL{Systematically varied model complexity (24 variants)
    but along technology}\DIFdelendFL \DIFaddbeginFL \small
\caption{Positioning against prior model-fidelity and network-aggregation studies. No prior
study isolates the loss effect from the topology effect, and none evaluates fidelity by the
cost of the resulting \emph{decisions} rather than by objective values. This study closes both
gaps.}
\label{tab:priorwork}
\setlength{\tabcolsep}{5pt}
\begin{tabular}{@{}p{0.52\linewidth}ccc@{}}
\toprule
Study & \makecell{Varies\\fidelity} & \makecell{Isolates\\loss/topology} & \makecell{Evaluates\\decisions} \\
\midrule
Wirtz et al.~\cite{Wirtz2021_complexity}, 24 MILP variants, multi-energy & \checkmark & $\times$ & $\times$ \\
Kotzur et al.~\cite{Kotzur2021}, complexity-management review & \checkmark & $\times$ & $\times$ \\
Larsen~\cite{Larsen2004}, Falay~\cite{Falay2020}, network aggregation & \checkmark & $\times$ & $\times$ \\
van der Heijde et al.~\cite{vanderHeijde2017}, district-heating physics model & $\times$ & $\times$ & $\times$ \\
Hering et al.~\cite{Hering2021}, district-heating MIQCP dispatch & $\times$ & $\times$ & $\times$ \\
\textbf{This study} & \checkmark & \checkmark & \checkmark \\
\bottomrule
\end{tabular}
\end{table}


\subsection{Contributions and research questions}
\DIFaddFL{The individual formulations compared here are established, and comparative fidelity studies
already exist: Wirtz et al.\ compared MILP variants of a multi-energy system, Larsen and
Falay studied network aggregation, and Kotzur et al.\ reviewed complexity management. The
contribution lies not in the formulations themselves but in the experimental structure and
the quantity evaluated, and it rests on two pillars. The first is a decision-oriented
evaluation of fidelity. Each model's dispatch schedule is re-costed under a common
high-fidelity forward model to quantify decision regret and tested for physical
deliverability. This evaluation is paired with a control condition, a copperplate supplied
with the aggregate loss it cannot itself compute, that turns the comparison into an exact
loss}\DIFaddendFL /\DIFdelbeginFL \DIFdelFL{part-load axes, not network topology}\DIFdelendFL \DIFaddbeginFL \DIFaddFL{topology decomposition. Together, these show that loss visibility, not spatial routing,
sets the fidelity requirement; that a model can be both a biased estimator and an inadequate
controller, with its estimation bias and signed forward-cost deviation pointing in opposite
directions; and that a per-network loss adder that makes a copperplate appear sufficient does
not transfer across networks or operating regimes, making the resolved model the transferable
one. The regret idea is positioned as }\emph{\DIFaddFL{analogous}} \DIFaddFL{to the value of the stochastic
solution~\mbox{%DIFAUXCMD
\citep{BirgeLouveaux2011}}\hskip0pt%DIFAUXCMD
, transposed from the uncertainty axis to the resolution
axis; the contribution is the transfer and the empirical finding, not the concept itself.
The second pillar establishes where the fidelity threshold lies by pushing physical detail
to its limit: resolving the demand side to all 174 individual transmission stations and
their service laterals using real component data, with the resulting model cross-checked
against a nonlinear pipe-flow solver, changes dispatch cost by less than one percent and
produces essentially no change in decisions. This null result at maximal fidelity directly
addresses the hydraulic-validation concern. Two further findings are reported as secondary:
the spatial-routing effect is null under central generation and remains a controlled open
question for distributed generation; and an }\emph{\DIFaddFL{a priori}} \DIFaddFL{estimator of bias from observable
network properties predicts held-out synthetic networks without refitting. The paper is
organised around four research questions: the loss/topology decomposition, decision regret
and deliverability, the transferability of a lumped-loss correction, and the location of the
hydraulic threshold, with the two secondary questions addressed alongside}\DIFaddendFL .
\DIFdelbeginFL %DIFDELCMD < \end{tablenotes}
%DIFDELCMD < \end{threeparttable}
%DIFDELCMD < \end{table*}
%DIFDELCMD < %%%
\DIFdelend 

%DIF < % ============================================================
%DIF < % 3.1 EXPERIMENTAL DESIGN
%DIF < % ============================================================
\DIFdelbegin %DIFDELCMD < \subsection{Controlled experimental design}
%DIFDELCMD < \label{sec:exp_design}
%DIFDELCMD < %%%
\DIFdelend \DIFaddbegin \DIFadd{The remainder of the paper is organised as follows.
Section~\ref{sec:methodology} sets out the experimental design (the fidelity ladder, the
decomposition controls and the contrasts each isolates), then the two case studies, the
formulations, the forward evaluator, the cost basis, the validation protocol and the
computational setup. Section~\ref{sec:results} reports and discusses the results in the
order of the research questions, beginning with what the measurements can and cannot
validate and ending with the limitations of the study.
Section~\ref{sec:conclusion} concludes.
}\DIFaddend 

\DIFdelbegin \DIFdel{The experimental design isolates three effects through pairwise
comparisons (}%DIFDELCMD < \Cref{fig:comparison_design}%%%
\DIFdel{):
%DIF < 
}%DIFDELCMD < \begin{enumerate}[nosep]
%DIFDELCMD <   \item %%%
\textbf{\DIFdel{Topology effect}} %DIFAUXCMD
\DIFdel{(L1\,$\to$\,L2\,$\to$\,L3): spatial resolution varies; physics held constant. }%DIFDELCMD < \item %%%
\textbf{\DIFdel{Physics fidelity effect}} %DIFAUXCMD
\DIFdelend \DIFaddbegin \section{Experimental design and methodology}
\label{sec:methodology}

\subsection{Experimental design: fidelity ladder and decomposition controls}
\DIFadd{This study compares a family of formulations that differ from one another in exactly one
modelled phenomenon at a time, so that each pairwise contrast isolates a single effect and no
comparison is confounded. The family has two parts. Four }\emph{\DIFadd{decomposition controls}} \DIFadd{form a
$2\times2$ factorial of spatial topology against thermal loss: a copperplate with no loss
($\CP$), the copperplate supplied with the aggregate loss ($\CPL$), the node-resolved network
with no loss ($\NDz$), and the node-resolved network with trunk loss ($\Lone$), which is also
the comparison baseline. A monotone }\emph{\DIFadd{fidelity ladder}} \DIFadd{then adds, one at a time, on the
same node-resolved network: temperature propagation }\DIFaddend (\DIFdelbegin \DIFdel{L3\,$\to$\,$\Lplus$):
physics scope varies; topology and formulation class (MILP)
        held constant. }%DIFDELCMD < \item %%%
\textbf{\DIFdel{Linearisation error}} %DIFAUXCMD
\DIFdelend \DIFaddbegin \DIFadd{$\Ltwo$), trunk pressure and pumping
($\Lthree$), station resolution with service laterals ($\Lfour$), a dynamic flow-dependent
station pressure requirement ($\Lfive$), and transport delay }\DIFaddend (\DIFdelbegin \DIFdel{$\Lplus$\,$\to$\,$\Lnl$):
mathematical formulation varies (MILP\,$\to$}\DIFdelend \DIFaddbegin \DIFadd{$\Lsix$), with the fully
nonlinear formulation ($\NLref$) as a reference that is not solved globally (evaluated
forward and on bounded 72-hour re-solves). Table~\ref{tab:contrasts} states
what each arrow isolates. Two design principles hold. First, within the four decomposition controls, each contrast changes exactly one
phenomenon, so its cost difference is attributable to that phenomenon. The
extended-physics levels (\(\Ltwo\) to \(\Lsix\), \(\NLref\)) form a
physics-evaluation framework rather than an exact one-factor-at-a-time solved
ladder. Where a step also changes the representation or evaluation mode, as
\(\Ltwo\), \(\Lfour\), \(\Lfive\), and \(\NLref\) do, this is stated explicitly. Second, the 174 transmission stations
are present at every level; only their spatial aggregation changes, from all lumped at one
point through node and zone aggregation to individually resolved, so the taxonomy is consistent
across the whole family. The four controls satisfy an exact additive identity,
$\text{cost}(\Lone)-\text{cost}(\CP)=\text{loss}+\text{topology}+\text{interaction}$, where the
loss main effect is $\text{cost}(\CPL)-\text{cost}(\CP)$, the topology main effect is
$\text{cost}(\NDz)-\text{cost}(\CP)$, and the interaction is the residual; it is an identity,
not an approximation, and closes to machine precision when the four models are solved to a
0.01}\DIFaddend \,\DIFdelbegin \DIFdel{MIQCP);
        physics and topology held identical.
}%DIFDELCMD < \end{enumerate}
%DIFDELCMD < 

%DIFDELCMD < \begin{figure}[ht]
%DIFDELCMD <   \centering
%DIFDELCMD <   \includegraphics[width=\linewidth]{paper1_dh_fidelity/Figure_1.png}
%DIFDELCMD <   \caption{Experimental design with three orthogonal comparison
%DIFDELCMD <   dimensions. Each arrow isolates exactly one modelling dimension.}
%DIFDELCMD <   \label{fig:comparison_design}
%DIFDELCMD < \end{figure}
%DIFDELCMD < 

%DIFDELCMD < \Cref{tab:model_scope} %%%
\DIFdel{summarises which physical phenomena enter at
each level. The five variants are built up as follows, inspired by \mbox{%DIFAUXCMD
\citet{Blommaert2020}}\hskip0pt%DIFAUXCMD
.}%DIFDELCMD < 

%DIFDELCMD < \begin{table}[htbp]
%DIFDELCMD < %%%
\DIFdelendFL \DIFaddbeginFL \DIFaddFL{\% optimality gap.
Two reference roles must be kept distinct. The baseline $\Lone$ is the detailed
mixed-integer model a practitioner would actually build; every estimation-bias percentage is
normalised to it. The }\emph{\DIFaddFL{physics}} \DIFaddFL{reference, against which regret and physical feasibility
are judged, is the forward evaluator of Section~\ref{subsec:evaluator}, }\emph{\DIFaddFL{not}} \DIFaddFL{any solved
level: even the finest solved level linearises the temperature-decay factor, so using it as
ground truth would carry the same linearisation into the reference. Model-to-model differences
are therefore estimation bias; only the forward evaluation speaks to decision quality.
}\begin{table}[t]
\DIFaddendFL \centering
\DIFdelbeginFL %DIFDELCMD < \caption{Model scope across all five variants. ``PWL'': piecewise-linear;
%DIFDELCMD < ``Quad./Bilin.'': native quadratic/bilinear constraint.}
%DIFDELCMD < \label{tab:model_scope}
%DIFDELCMD < \begin{threeparttable}
%DIFDELCMD < \scriptsize
%DIFDELCMD < \setlength{\tabcolsep}{2pt}
%DIFDELCMD < \begin{tabular*}{\columnwidth}{@{\extracolsep{\fill}}
%DIFDELCMD <                     l c c c c c @{}}
%DIFDELCMD < \toprule
%DIFDELCMD < Phenomenon & L1 & L2 & L3 & $\Lplus$ & $\Lnl$ \\
%DIFDELCMD < \midrule
%DIFDELCMD < \multicolumn{6}{@{}l@{}}{\textit{Thermal}} \\[1pt]
%DIFDELCMD < \quad Pipe losses ($UL\Delta T$)
%DIFDELCMD <   & -- & \checkmark & \checkmark & \checkmark & \checkmark \\
%DIFDELCMD < \quad Temp.\ propagation
%DIFDELCMD <   & -- & -- & -- & PWL & Bilin. \\
%DIFDELCMD < \quad Time-varying COP
%DIFDELCMD <   & \checkmark & \checkmark & \checkmark & \checkmark & \checkmark \\
%DIFDELCMD < \quad Storage losses \& $\eta$
%DIFDELCMD <   & \checkmark & \checkmark & \checkmark & \checkmark & \checkmark \\
%DIFDELCMD < \midrule
%DIFDELCMD < \multicolumn{6}{@{}l@{}}{\textit{Hydraulic}} \\[1pt]
%DIFDELCMD < \quad Pressure drop (DW)
%DIFDELCMD <   & -- & -- & -- & PWL & Quad. \\
%DIFDELCMD < \quad Pumping power
%DIFDELCMD <   & -- & -- & -- & \checkmark & \checkmark \\
%DIFDELCMD < \midrule
%DIFDELCMD < \multicolumn{6}{@{}l@{}}{\textit{Temporal}} \\[1pt]
%DIFDELCMD < \quad Transport delay
%DIFDELCMD <   & -- & -- & -- & -- & \checkmark \\
%DIFDELCMD < \quad Hourly grid $e$-factors
%DIFDELCMD <   & \checkmark & \checkmark & \checkmark & \checkmark & \checkmark \\
%DIFDELCMD < \midrule
%DIFDELCMD < \multicolumn{6}{@{}l@{}}{\textit{Outside the scope of the research}} \\[1pt]
%DIFDELCMD < \quad Sub-hourly transients
%DIFDELCMD <   & \multicolumn{5}{c}{excluded} \\
%DIFDELCMD < \quad Capacity sizing
%DIFDELCMD <   & \multicolumn{5}{c}{excluded} \\
%DIFDELCMD < \quad $T_{\mathrm{sup}}$ optim.
%DIFDELCMD <   & \multicolumn{5}{c}{excluded} \\
%DIFDELCMD < \quad Ring / bidir.\ flow
%DIFDELCMD <   & \multicolumn{5}{c}{excluded} \\
%DIFDELCMD < \bottomrule
%DIFDELCMD < \end{tabular*}
%DIFDELCMD < \begin{tablenotes}\scriptsize
%DIFDELCMD <   \item[] %%%
\DIFdelFL{Physics scope as shown applies to the
    }\emph{\DIFdelFL{primary case}}%DIFAUXCMD
\DIFdelFL{. In the synthetic parametric study, level labels are reassigned to enable the additive
    decomposition; see }%DIFDELCMD < \Cref{tab:synth_level_mapping}%%%
\DIFdelFL{.
}%DIFDELCMD < \end{tablenotes}
%DIFDELCMD < \end{threeparttable}
%DIFDELCMD < %%%
\DIFdelendFL \DIFaddbeginFL \caption{The experimental contrasts. Within the four decomposition controls (topology
$\times$ loss) each arrow changes exactly one phenomenon, so those comparisons are unconfounded;
the remaining steps form an extended-physics evaluation framework in which a step may change
more than one thing (evaluation mode, laterals, pressure, delay), noted where it occurs.}
\label{tab:contrasts}
\begin{tabular}{ll}
\hline
Contrast & Isolates \\
\hline
$\CP\!\to\!\CPL$ & loss visibility (no topology) \\
$\CP\!\to\!\NDz$ & spatial topology (no loss) \\
$\Lone$ vs $\CPL,\NDz$ & loss $\times$ topology interaction \\
$\Lone\!\to\!\Ltwo$ & temperature propagation \\
$\Ltwo\!\to\!\Lthree$ & trunk pressure drop \& pumping \\
$\Lthree\!\to\!\Lfour$ & station resolution + service laterals \\
$\Lfour\!\to\!\Lfive$ & dynamic (flow-dependent) station $\Delta p$ \\
$\Lfive\!\to\!\Lsix\!\to\!\NLref$ & transport delay; linearisation error \\
\hline
\end{tabular}
\DIFaddendFL \end{table}


\DIFdelbegin \paragraph{\DIFdel{Level~1 (Copperplate).}}
%DIFAUXCMD
\addtocounter{paragraph}{-1}%DIFAUXCMD
\DIFdel{Single bus, no spatial structure, no pipe losses. All assets are
co-located at }\DIFdelend \DIFaddbegin \DIFadd{Where the contrasts in Table~\ref{tab:contrasts} name what each step }\emph{\DIFadd{changes}}\DIFadd{,
Table~\ref{tab:designgrid} gives the complete picture: every fidelity level against every
phenomenon. In that grid }\checkmark\DIFadd{\ marks a represented phenomenon and a dash its absence;
}\texttt{\DIFadd{A}} \DIFadd{an aggregated form and }\texttt{\DIFadd{exo}} \DIFadd{an exogenous input; }\texttt{\DIFadd{PWL}}\DIFadd{, }\texttt{\DIFadd{Bil}}
\DIFadd{and }\texttt{\DIFadd{Quad}} \DIFaddend a \DIFdelbegin \DIFdel{single node; the network is represented by a single scalar heat balance.
}%DIFDELCMD < 

%DIFDELCMD < %%%
\paragraph{\DIFdel{Level~2 (Simplified multi-node).}}
%DIFAUXCMD
\addtocounter{paragraph}{-1}%DIFAUXCMD
\DIFdel{Aggregated demand zones (7 in the primary case)connected by
simplified pipe segments that preserve the tree topology. The total
$U \cdot L$ product matches L3, ensuring identical aggregate annual
loss energy.
Asset locations correspond to their physical positions
(}%DIFDELCMD < \Cref{sec:case_studies}%%%
\DIFdel{).
}%DIFDELCMD < 

%DIFDELCMD < %%%
\paragraph{\DIFdel{Level~3 (Detailed multi-node, basic MILP).}}
%DIFAUXCMD
\addtocounter{paragraph}{-1}%DIFAUXCMD
\DIFdel{Full spatial resolution (15 junction nodes). Individual pipe
parameters with steady-state losses, fixed supply temperature
(heating curve)}\DIFdelend \DIFaddbegin \DIFadd{piecewise-linear}\DIFaddend , \DIFaddbegin \DIFadd{native-bilinear or native-quadratic treatment; and
}\texttt{\DIFadd{flat}} \DIFadd{or }\texttt{\DIFadd{dyn}} \DIFadd{the static or flow-dependent station pressure requirement. The
four decomposition controls stand to the left of the monotone $\Lone$ through $\Lsix$ ladder,
with the nonlinear reference $\NLref$ closing it on the right.
The four decomposition controls, together with $\Lthree$ }\DIFaddend and \DIFdelbegin \DIFdel{pre-computed COP (}%DIFDELCMD < \Cref{app:cop}%%%
\DIFdel{). No pressure
drop, no temperature propagation , no transport delay.
}%DIFDELCMD < 

%DIFDELCMD < %%%
\paragraph{\DIFdel{$\Lplus$ (Detailed multi-node, extended MILP).}}
%DIFAUXCMD
\addtocounter{paragraph}{-1}%DIFAUXCMD
\DIFdel{Same topology as L3. Adds piecewise-linear pressure drop
(Darcy--Weisbach), linearised temperature propagation, and the
resulting in the pumping power. Transport delay is configured but bypassed:
the linearised pipe model cannot represent the source-to-consumer
lag, so this effect is characterised exclusively through the
$\Lplus\!\to\!\Lnl$ comparison. }%DIFDELCMD < 

%DIFDELCMD < %%%
\paragraph{\DIFdel{$\Lnl$ (Detailed multi-node, extended MIQCP).}}
%DIFAUXCMD
\addtocounter{paragraph}{-1}%DIFAUXCMD
\DIFdel{Same topology and physics scope as $\Lplus$, but pressure drop
enters as a
native quadratic constraint, temperature propagation as
bilinear products, and transport delay is dispatched natively. $\Lnl$ is the only variant in which transport delay is active.
}%DIFDELCMD < 

%DIFDELCMD < %%%
\paragraph{\DIFdel{Auxiliary variant: L1$_{\mathrm{topo}}$.}}
%DIFAUXCMD
\addtocounter{paragraph}{-1}%DIFAUXCMD
\DIFdel{For the additive gap decomposition (}%DIFDELCMD < \Cref{sec:case_synthetic}%%%
\DIFdel{),
an additionally L1$_{\mathrm{topo}}$ was defined: full network routing
(nodal balance as in L3) but with zero pipe losses. This isolates
the
pure routing effect from the thermal-loss effect within the L1$\to$L3 gap.
}%DIFDELCMD < 

%DIFDELCMD < %%%
\paragraph{\DIFdel{Demand concentration metrics.}}
%DIFAUXCMD
\addtocounter{paragraph}{-1}%DIFAUXCMD
\DIFdel{The topology effect depends on how unevenly demand is distributed
across nodes. The study quantifies this via (i)~the }\textbf{\DIFdel{Gini
coefficient}} %DIFAUXCMD
\DIFdel{of nodal annual demand shares and (ii)~the
}\textbf{\DIFdel{top-3 demand share}}%DIFAUXCMD
\DIFdel{. For the primary
case,
Gini\,$= 0.41$ and top-3 share}\DIFdelend \DIFaddbegin \DIFadd{$\Lsix$, are solved as
mixed-integer linear programs to a $0.01$}\DIFaddend \,\DIFdelbegin \DIFdel{$= 51.7\,\%$, indicating
moderate spatial concentration. The
synthetic parametric study uses a
normalised variance index HI as
a discrete parameter
(Low/Med/High) ; its definition and mapping to Gini are given in }%DIFDELCMD < \Cref{app:hi_definition}%%%
\DIFdel{.
}\DIFdelend \DIFaddbegin \DIFadd{\% gap; the non-convex levels
($\Ltwo$, $\Lfour$, $\Lfive$, $\NLref$) are quantified with the validated forward model
rather than solved. 
One further formulation sits outside both parts of the family. The zone-aggregated level
$\ZN$ resolves demand into geographic clusters rather than junctions and computes trunk
loss on the aggregated graph; it is not a decomposition control, since it occupies no
corner of the $2\times2$, and not a ladder rung, since the ladder is defined on the
node-resolved network. It is used solely to test whether the choice of zone boundary is a
hidden degree of freedom in the decomposition (Section~\ref{subsec:clustering}), and is
reported there rather than among the bias and regret results.
}\begin{table}[htbp]
\centering
\caption{Model scope grid: which phenomenon each fidelity level represents, and in what
form: \emph{fwd} = forward-evaluated (not in the solved dispatch), \emph{A} = aggregated
station, PWL/Bil/Quad the linearisation form. The solved levels $\Lone$, $\Lthree$, $\Lsix$
fix node temperatures (the free-variable form is degenerate), so $\Lsix=\Lthree+$ transport
delay and does \emph{not} inherit the station-resolved rows; temperature propagation and the
station levels ($\Ltwo$, $\Lfour$, $\Lfive$) are forward-evaluated. The four decomposition
controls ($\CP$, $\CPL$, $\NDz$, $\Lone$) are separated from the zone-aggregated variant
$\ZN$, used only for the aggregation-sensitivity study of Section~\ref{subsec:zoneclustering}.}
\label{tab:designgrid}
\footnotesize
\setlength{\tabcolsep}{3pt}
\begin{tabular}{@{}lcc|c|cccccccc@{}}
\hline
Phenomenon & $\CP$ & $\CPL$ & $\ZN$ & $\NDz$ & $\Lone$ & $\Ltwo$ & $\Lthree$ & $\Lfour$ & $\Lfive$ & $\Lsix$ & $\NLref$ \\
\hline
Demand aggregation & 1pt & 1pt & zones & nodes & nodes & nodes & nodes & stat. & stat. & nodes & stat. \\
Station representation & A & A & A & A & A & A & A & \checkmark & \checkmark & A & \checkmark \\
Trunk losses ($U\!L\Delta T$, defensible $U$) & -- & exo & \checkmark & -- & \checkmark & \checkmark & \checkmark & \checkmark & \checkmark & \checkmark & \checkmark \\
Service-lateral losses & -- & -- & -- & -- & -- & -- & -- & \checkmark & \checkmark & -- & \checkmark \\
Temperature propagation & -- & -- & -- & -- & -- & fwd & -- & fwd & fwd & -- & Bil \\
Time-varying COP & \checkmark & \checkmark & \checkmark & \checkmark & \checkmark & \checkmark & \checkmark & \checkmark & \checkmark & \checkmark & \checkmark \\
Storage losses \& $\eta$ & \checkmark & \checkmark & \checkmark & \checkmark & \checkmark & \checkmark & \checkmark & \checkmark & \checkmark & \checkmark & \checkmark \\
Trunk pressure drop (DW) & -- & -- & -- & -- & -- & -- & PWL & PWL & PWL & PWL & Quad \\
Trunk pumping power & -- & -- & -- & -- & -- & -- & \checkmark & \checkmark & \checkmark & \checkmark & \checkmark \\
Station $\Delta p$ requirement & -- & -- & -- & -- & -- & -- & -- & flat & dyn. & -- & dyn. \\
Station \& lateral pumping & -- & -- & -- & -- & -- & -- & -- & \checkmark & \checkmark & -- & \checkmark \\
Transport delay & -- & -- & -- & -- & -- & -- & -- & -- & -- & \checkmark & \checkmark \\
Hourly grid emission factors & \checkmark & \checkmark & \checkmark & \checkmark & \checkmark & \checkmark & \checkmark & \checkmark & \checkmark & \checkmark & \checkmark \\
\hline
\end{tabular}
\end{table}
\DIFaddend 


%DIF < % ============================================================
%DIF < % 3.2 BASIC MILP
%DIF < % ============================================================
\DIFdelbegin %DIFDELCMD < \subsection{Basic MILP formulation (L1/L2/L3)}
%DIFDELCMD < \label{sec:milp}
%DIFDELCMD < 

%DIFDELCMD < %%%
\DIFdel{All five variants }\DIFdelend \DIFaddbegin \subsection{Base formulation}
{
\DIFadd{All levels }\DIFaddend share the objective function and component models \DIFdelbegin \DIFdel{below; L1/L2/L3 }\DIFdelend \DIFaddbegin \DIFadd{set out below. They }\DIFaddend differ
only in \DIFdelbegin \DIFdel{the energy-balance topology and
loss treatment }\DIFdelend \DIFaddbegin \DIFadd{two places (the topology of the energy balance and the treatment of thermal
losses), and it is precisely because these two enter as separable terms that the
decomposition of Section~\ref{sec:results} is exact rather than fitted}\DIFaddend .

\paragraph{Objective function.}
The optimiser minimises total annual \DIFdelbegin \DIFdel{operational cost, which
consists of electricity procurement and }\DIFdelend \DIFaddbegin \DIFadd{operating cost: electricity procurement net of
}\DIFaddend feed-in \DIFdelbegin \DIFdel{revenues, fuelexpenditure}\DIFdelend \DIFaddbegin \DIFadd{revenue, fuel}\DIFaddend , heat-curtailment penalties, carbon charges, \DIFdelbegin \DIFdel{and }\DIFdelend a demand charge on
peak grid import\DIFdelbegin \DIFdel{:
%DIF < 
}\DIFdelend \DIFaddbegin \DIFadd{, and a thermal-storage cycling penalty,
}\DIFaddend \begin{equation}
  \min \; Z = C^{\mathrm{en}} + C^{\mathrm{fuel}}
    + C^{\mathrm{dump}} + C\DIFdelbegin \DIFdel{^{\ce{CO2}} }\DIFdelend \DIFaddbegin \DIFadd{^{\mathrm{CO_2}} }\DIFaddend + C^{\mathrm{dem}}
    + C\DIFdelbegin \DIFdel{^{\mathrm{pump}}
  }\DIFdelend \DIFaddbegin \DIFadd{^{\mathrm{cyc}} .
  }\DIFaddend \label{eq:objective}
\end{equation}
%DIF < 
\DIFdelbegin \DIFdel{The pumping cost $C^{\mathrm{pump}} = 0$ for L1/L2/L3 and enters
only in $\Lplus$/$\Lnl$ (}%DIFDELCMD < \Cref{sec:extended}%%%
\DIFdel{). The individual
terms are defined as follows, where $\Delta t = 1\,\mathrm{h}$ is
the time-step length }\DIFdelend \DIFaddbegin \DIFadd{Network pumping does not appear as a separate cost term. The pump electrical power
$P_t^{\mathrm{pump}}$ enters the electricity balance, Eq.~\eqref{eq:balance_el}, and is
therefore paid for through grid procurement inside $C^{\mathrm{en}}$; it is identically
zero at every level below $\Lthree$ and is computed from Eq.~\eqref{eq:pump_power} from
$\Lthree$ upward (Section~\ref{sec:extended}). With $\Delta t = 1$\,h }\DIFaddend and $t$ \DIFdelbegin \DIFdel{indexes }\DIFdelend \DIFaddbegin \DIFadd{indexing }\DIFaddend all
8\DIFdelbegin %DIFDELCMD < {%%%
\DIFdel{,}%DIFDELCMD < }%%%
\DIFdelend \DIFaddbegin \DIFadd{\,}\DIFaddend 760 hours of the year\DIFdelbegin \DIFdel{:
%DIF < 
}\DIFdelend \DIFaddbegin \DIFadd{,
}\DIFaddend \begin{align*}
C^{\mathrm{en}}   &= \Delta t \sum_t
  \bigl(P_t^{\mathrm{buy}} \lambda_t^{\mathrm{buy}}
        - P_t^{\mathrm{sell}} \lambda_t^{\mathrm{sell}}\bigr)\DIFdelbegin %DIFDELCMD < \\
%DIFDELCMD < %%%
\DIFdelend \DIFaddbegin \DIFadd{, }&
\DIFaddend C^{\mathrm{fuel}} &= \Delta t \sum_{g,t}
  F_{g,t}\, \lambda_g^{\mathrm{fuel}}\DIFaddbegin \DIFadd{, }\DIFaddend \\
C^{\mathrm{dump}} &= \Delta t \sum_{n,t}
  Q_{n,t}^{\mathrm{dump}}\, \lambda^{\mathrm{dump}}\DIFaddbegin \DIFadd{, }&
\DIFadd{C^{\mathrm{dem}}  }&\DIFadd{= \lambda^{\mathrm{dem}} f^{\mathrm{yr}}
  P^{\mathrm{peak}}, }\DIFaddend \\
C\DIFdelbegin \DIFdel{^{\ce{CO2}}      }\DIFdelend \DIFaddbegin \DIFadd{^{\mathrm{CO_2}} }\DIFaddend &= \DIFdelbegin \DIFdel{\frac{\lambda^{\ce{CO2}}}{1000} }\DIFdelend \DIFaddbegin \DIFadd{\frac{\lambda^{\mathrm{CO_2}}}{1000} }\DIFaddend \sum_t
  \Bigl(E_t\DIFdelbegin \DIFdel{^{\ce{CO2},\mathrm{grid}} }\DIFdelend \DIFaddbegin \DIFadd{^{\mathrm{CO_2,grid}} }\DIFaddend +
        \DIFaddbegin \DIFadd{\textstyle}\DIFaddend \sum_g E_{g,t}\DIFdelbegin \DIFdel{^{\ce{CO2},\mathrm{fuel}}}\DIFdelend \DIFaddbegin \DIFadd{^{\mathrm{CO_2,fuel}}}\DIFaddend \Bigr)\DIFdelbegin %DIFDELCMD < \\
%DIFDELCMD < %%%
\DIFdelend \DIFaddbegin \DIFadd{, }&
\DIFaddend C\DIFdelbegin \DIFdel{^{\mathrm{dem}}  }\DIFdelend \DIFaddbegin \DIFadd{^{\mathrm{cyc}} }\DIFaddend &= \lambda\DIFdelbegin \DIFdel{^{\mathrm{dem}} f^{\mathrm{yr}}
  P^{\mathrm{peak}}
}\DIFdelend \DIFaddbegin \DIFadd{^{\mathrm{cyc}}\,\Delta t \sum_{s,t}
  }\bigl(\DIFadd{Q_{s,t}^{\mathrm{ch}} + Q_{s,t}^{\mathrm{dis}}}\bigr)\DIFadd{.
}\DIFaddend \end{align*}
%DIF < 
\DIFdelbegin \DIFdel{Here , }\DIFdelend \DIFaddbegin \DIFadd{Here }\DIFaddend $P_t^{\mathrm{buy}}$ and $P_t^{\mathrm{sell}}$ [MW] are \DIFdelbegin \DIFdel{the }\DIFdelend grid import and export at
time \DIFdelbegin \DIFdel{~}\DIFdelend $t$ \DIFdelbegin \DIFdel{;
$\lambda_t^{\mathrm{buy}}$ and }\DIFdelend \DIFaddbegin \DIFadd{and $\lambda_t^{\mathrm{buy}}$, }\DIFaddend $\lambda_t^{\mathrm{sell}}$ [EUR/MWh] \DIFdelbegin \DIFdel{are }\DIFdelend the
corresponding prices; $F_{g,t}$ [MW] is the fuel input to generator \DIFdelbegin \DIFdel{~}\DIFdelend $g$ \DIFdelbegin \DIFdel{at heating value basis ;
}\DIFdelend \DIFaddbegin \DIFadd{on a
heating-value basis at specific cost }\DIFaddend $\lambda_g^{\mathrm{fuel}}$\DIFdelbegin %DIFDELCMD < [%%%
\DIFdel{EUR/MWh}%DIFDELCMD < ] %%%
\DIFdel{is its specific fuel cost }\DIFdelend ;
$Q_{n,t}^{\mathrm{dump}}$ [MW] is \DIFdelbegin \DIFdel{curtailed heat at node~}\DIFdelend \DIFaddbegin \DIFadd{heat curtailed at node }\DIFaddend $n$, penalised at
$\lambda^{\mathrm{dump}}$\DIFdelbegin %DIFDELCMD < [%%%
\DIFdel{EUR/MWh}%DIFDELCMD < ]%%%
\DIFdel{; $\lambda^{\ce{CO2}}$ }\DIFdelend \DIFaddbegin \DIFadd{; $\lambda^{\mathrm{CO_2}}$ }\DIFaddend [EUR/t\DIFdelbegin \DIFdel{$_{\mathrm{CO_2}}$}\DIFdelend ] is the carbon price;
$f^{\mathrm{yr}}$ \DIFdelbegin \DIFdel{a fractional }\DIFdelend \DIFaddbegin \DIFadd{an }\DIFaddend annualisation factor; \DIFdelbegin \DIFdel{and
}\DIFdelend $P^{\mathrm{peak}}$ [MW] the annual peak
grid import entering the demand charge $\lambda^{\mathrm{dem}}$\DIFaddbegin \DIFadd{; and
$\lambda^{\mathrm{cyc}}$ }\DIFaddend [\DIFdelbegin \DIFdel{MWh$^{-1}$}\DIFdelend \DIFaddbegin \DIFadd{EUR/MWh}\DIFaddend ] \DIFaddbegin \DIFadd{the storage-cycling cost, applied to the combined
charge-and-discharge throughput $Q_{s,t}^{\mathrm{ch}}+Q_{s,t}^{\mathrm{dis}}$ of each
storage $s$}\DIFaddend .

\DIFaddbegin \DIFadd{Two elements of $Z$ depart from operator cash flow: the carbon charge $C^{\mathrm{CO_2}}$
enters on a gross basis, and the cycling penalty $C^{\mathrm{cyc}}$ is a wear-and-arbitrage
deterrent rather than a cash outlay. The reported }\emph{\DIFadd{economic}} \DIFadd{cost applies the standard
cogeneration self-use carbon credit and excludes $C^{\mathrm{cyc}}$, so it differs
systematically from $Z$; the economic cost is used for every reported comparison. The
distinction and its magnitude are set out in Section~\ref{subsec:costacct}.
}

\DIFaddend \paragraph{Energy balance\DIFaddbegin \DIFadd{, and how the controls are constructed}\DIFaddend .}
\DIFdelbegin \DIFdel{L1 enforces }\DIFdelend \DIFaddbegin \DIFadd{The copperplate levels $\CP$ and $\CPL$ enforce }\DIFaddend a single bus-level balance \DIFdelbegin \DIFdel{across all
generatorsand
storage units, ignoring network losses and }\DIFdelend \DIFaddbegin \DIFadd{over all
generators, heat pumps, electrode boilers and storage, with no }\DIFaddend spatial routing:
%DIF < 
\begin{equation}
  \sum_{g} Q_{g,t}^{\mathrm{th}}
    + \sum\DIFaddbegin \DIFadd{_{h} Q_{h,t} + \sum_{r} Q_{r,t}
    + \sum}\DIFaddend _{s} Q_{s,t}^{\mathrm{dis}}
  = Q_t^{\mathrm{dem}} + Q_t^{\mathrm{dump}}
    + \sum_{s} Q_{s,t}^{\mathrm{ch}}
  \DIFaddbegin \DIFadd{+ }\underbrace{L_t}\DIFadd{_{\CPL\ \mathrm{only}}
  }\DIFaddend \;\; \forall\, t\DIFdelbegin %DIFDELCMD < \label{eq:balance_L1}
%DIFDELCMD < %%%
\DIFdelend \DIFaddbegin \DIFadd{,
  }\label{eq:balance_cp}
\DIFaddend \end{equation}
%DIF < 
where $Q_{g,t}^{\mathrm{th}}$ [MW] is the thermal output of generator \DIFdelbegin \DIFdel{~}\DIFdelend $g$,
\DIFaddbegin \DIFadd{$Q_{h,t}$ and $Q_{r,t}$ }[\DIFadd{MW}] \DIFadd{the heat-pump and electrode-boiler thermal outputs,
}\DIFaddend $Q_{s,t}^{\mathrm{dis}}$ and $Q_{s,t}^{\mathrm{ch}}$ [MW] \DIFdelbegin \DIFdel{are }\DIFdelend the discharge and charge of
storage \DIFdelbegin \DIFdel{unit~}\DIFdelend $s$, and $Q_t^{\mathrm{dem}}$ [MW] \DIFdelbegin \DIFdel{is }\DIFdelend the total network \DIFdelbegin \DIFdel{heat demand. }\DIFdelend \DIFaddbegin \DIFadd{demand. The term $L_t$ is
the exogenous aggregate loss supplied to the control condition $\CPL$ and is zero for
$\CP$; its calibration is given in Section~\ref{subsec:lumped}.
}\DIFaddend 

\DIFdelbegin \DIFdel{L2/L3/$\Lplus$/$\Lnl$ }\DIFdelend \DIFaddbegin \DIFadd{All node-resolved levels ($\ZN$, $\NDz$ and $\Lone$ through $\Lsix$) }\DIFaddend instead enforce
a \DIFdelbegin \emph{\DIFdel{nodal}} %DIFAUXCMD
\DIFdelend \DIFaddbegin \DIFadd{nodal }\DIFaddend balance at each junction \DIFdelbegin \DIFdel{~}\DIFdelend $n$, accounting for thermal losses along incoming pipes
and \DIFaddbegin \DIFadd{for }\DIFaddend the spatial distribution of assets and demands:
%DIF < 
\begin{align}
  &\sum_{p \in \mathcal{P}_n^{\mathrm{in}}}
    \bigl(Q_{p,t} - \DIFaddbegin \DIFadd{\kappa\, }\DIFaddend Q_{p,t}^{\mathrm{loss,pipe}}\bigr)
  + \sum_{g \in \mathcal{G}_n} Q_{g,t}^{\mathrm{th}}
  + \sum\DIFaddbegin \DIFadd{_{h \in \mathcal{H}_n} Q_{h,t}
  + \sum_{r \in \mathcal{R}_n} Q_{r,t}
  + \sum}\DIFaddend _{s \in \mathcal{S}_n} Q_{s,t}^{\mathrm{dis}} \notag \\
  &= Q_{n,t}^{\mathrm{dem}} + Q_{n,t}^{\mathrm{dump}}
  + \sum_{s \in \mathcal{S}_n} Q_{s,t}^{\mathrm{ch}}
  + \sum_{p \in \mathcal{P}_n^{\mathrm{out}}} Q_{p,t}
  \;\; \forall\, n,\, t\DIFaddbegin \DIFadd{,
  }\DIFaddend \label{eq:balance_nodal}
\end{align}
%DIF < 
\DIFdelbegin \DIFdel{Here }\DIFdelend \DIFaddbegin \DIFadd{with }\DIFaddend $\mathcal{P}_n^{\mathrm{in}}$\DIFdelbegin \DIFdel{and }\DIFdelend \DIFaddbegin \DIFadd{, }\DIFaddend $\mathcal{P}_n^{\mathrm{out}}$ \DIFdelbegin \DIFdel{are the sets of }\DIFdelend \DIFaddbegin \DIFadd{the }\DIFaddend pipes entering and
leaving node \DIFdelbegin \DIFdel{~}\DIFdelend $n$\DIFdelbegin \DIFdel{;
}\DIFdelend \DIFaddbegin \DIFadd{, }\DIFaddend $Q_{p,t}$ [MW] \DIFdelbegin \DIFdel{is }\DIFdelend the heat flow in pipe \DIFdelbegin \DIFdel{~}\DIFdelend $p$\DIFdelbegin \DIFdel{; }\DIFdelend \DIFaddbegin \DIFadd{, }\DIFaddend and
$Q_{p,t}^{\mathrm{loss,pipe}}$ [MW] \DIFdelbegin \DIFdel{is the corresponding thermal
loss (defined below}\DIFdelend \DIFaddbegin \DIFadd{its thermal loss.
}

\DIFadd{The switch $\kappa \in \{0,1\}$ is the point of the design. Equations
\eqref{eq:balance_cp} and \eqref{eq:balance_nodal} differ in two respects at once
(spatial resolution and loss visibility), and prior fidelity comparisons have moved both
together, which is why their cost differences cannot be attributed to either. Setting the
two independently generates the four corners of a $2\times2$ design: $\CP$ is
\eqref{eq:balance_cp} with $L_t = 0$ (neither), $\CPL$ is \eqref{eq:balance_cp} with
$L_t$ supplied (loss, no topology), $\NDz$ is \eqref{eq:balance_nodal} with $\kappa = 0$
(topology, no loss), and $\Lone$ is \eqref{eq:balance_nodal} with $\kappa = 1$ (both}\DIFaddend ).
The \DIFdelbegin \DIFdel{contrast between
}%DIFDELCMD < \cref{eq:balance_L1,eq:balance_nodal} %%%
\DIFdel{encodes the
central
experimental variable: the L1 model's missing pipe-loss and spatial-flow terms are the mechanism driving the L1\,$\to$\,L2
}\DIFdelend cost gap \DIFaddbegin \DIFadd{$\CP \to \Lone$ therefore decomposes additively into a loss main effect, a
topology main effect and their interaction, with no residual by construction. This is the
control condition the confound requires, and it is the reason the decomposition in
Section~\ref{subsec:decomp} closes to machine precision rather than to a fitted
tolerance}\DIFaddend .

The global electricity balance and grid mutual-exclusivity constraints close the system.
All electrical generators, \DIFdelbegin \DIFdel{heat
pumps (}\DIFdelend power-to-heat units \DIFdelbegin \DIFdel{~}\DIFdelend $r$\DIFdelbegin \DIFdel{), and pumps must be balanced
}\DIFdelend \DIFaddbegin \DIFadd{, heat pumps and pumps balance }\DIFaddend against
grid import, export \DIFdelbegin \DIFdel{, }\DIFdelend and in-network consumption:
%DIF < 
\begin{align}
  P_t^{\mathrm{buy}}
    &+ \sum_{g \in \mathcal{G}^{\mathrm{CHP}}} P_{g,t}^{\mathrm{el}}
  = \sum_{h \in \mathcal{H}} P_{h,t}^{\mathrm{el}}
    + \sum_{r \in \mathcal{R}} P_{r,t}^{\mathrm{el}} \notag \\
    &+ P_t^{\mathrm{pump}} + P_t^{\mathrm{sell}}
  \quad \forall\, t\DIFaddbegin \DIFadd{,
  }\DIFaddend \label{eq:balance_el} \\[4pt]
  P_t^{\mathrm{buy}} &\leq M^{\mathrm{grid}}\, z_t,
  \;\;
  P_t^{\mathrm{sell}} \leq M^{\mathrm{grid}}\, (1 - z_t)
  \;\; \forall\, t\DIFaddbegin \DIFadd{,
  }\DIFaddend \label{eq:grid_bigm}
\end{align}
%DIF < 
where $P_{g,t}^{\mathrm{el}}$ [MW] is \DIFdelbegin \DIFdel{the }\DIFdelend CHP electrical output, $P_{h,t}^{\mathrm{el}}$ \DIFdelbegin %DIFDELCMD < [%%%
\DIFdel{MW}%DIFDELCMD < ] %%%
\DIFdel{the
heat pump }\DIFdelend \DIFaddbegin \DIFadd{the
heat-pump }\DIFaddend electrical input, $P_{r,t}^{\mathrm{el}}$ \DIFdelbegin %DIFDELCMD < [%%%
\DIFdel{MW}%DIFDELCMD < ]%%%
\DIFdel{, the electrode boiler input , }\DIFdelend \DIFaddbegin \DIFadd{the electrode-boiler input }\DIFaddend and
$P_t^{\mathrm{pump}}$ \DIFdelbegin %DIFDELCMD < [%%%
\DIFdel{MW}%DIFDELCMD < ] %%%
\DIFdelend the network pumping power. The binary \DIFdelbegin \DIFdel{variable }\DIFdelend $z_t \in \{0,1\}$ enforces
that the grid \DIFdelbegin \DIFdel{operates in either import or export mode at any hour;
}\DIFdelend \DIFaddbegin \DIFadd{connection either imports or exports in any hour, with }\DIFaddend $M^{\mathrm{grid}}$
\DIFdelbegin \DIFdel{is }\DIFdelend a Big-M \DIFdelbegin \DIFdel{constant }\DIFdelend set to the \DIFdelbegin \DIFdel{maximum
}\DIFdelend grid capacity.

\paragraph{Thermal loss model\DIFdelbegin \DIFdel{(L2/L3)}\DIFdelend .}
\DIFdelbegin %DIFDELCMD < \label{sec:model_losses}
%DIFDELCMD < %%%
\DIFdelend For each pipe \DIFdelbegin \DIFdel{~}\DIFdelend $p$, steady-state conductive \DIFdelbegin \DIFdel{heat }\DIFdelend losses are computed separately for the supply
and return circuit. Because the supply temperature \DIFdelbegin \DIFdel{is determined by }\DIFdelend \DIFaddbegin \DIFadd{follows }\DIFaddend the heating curve rather than
\DIFaddbegin \DIFadd{being }\DIFaddend optimised, the losses are \DIFdelbegin \emph{\DIFdel{fixed parameters}}%DIFAUXCMD
\DIFdel{. This is a key
simplification that preserves MILP linearity:
%DIF < 
}\DIFdelend \DIFaddbegin \DIFadd{fixed parameters, the simplification that keeps the
formulation linear:
}\DIFaddend \begin{align}
  Q_{p,t}^{\mathrm{loss,sup}} &= \DIFdelbegin \DIFdel{\frac{U_p\, L_p\,
    \bigl(T_{\mathrm{sup}}^{\mathrm{nom}}(t)
    - T_{\mathrm{ground}}(t)\bigr)}{10^6}
  }\DIFdelend \DIFaddbegin \DIFadd{\frac{U^{s}_p\, L_p\,
    \bigl(T_{\mathrm{sup}}^{\mathrm{nom}}(t)
    - T_{\mathrm{ground}}(t)\bigr)}{10^6},
  }\DIFaddend \label{eq:loss_supply} \\[4pt]
  Q_{p,t}^{\mathrm{loss,ret}} &= \DIFdelbegin \DIFdel{\frac{U_p\, L_p\,
    \bigl(T_{\mathrm{ret}}^{\mathrm{nom}}(t)
    - T_{\mathrm{ground}}(t)\bigr)}{10^6}
  }\DIFdelend \DIFaddbegin \DIFadd{\frac{U^{r}_p\, L_p\,
    \bigl(T_{\mathrm{ret}}^{\mathrm{nom}}(t)
    - T_{\mathrm{ground}}(t)\bigr)}{10^6},
  }\DIFaddend \label{eq:loss_return}
\end{align}
%DIF < 
\DIFdelbegin \DIFdel{where $U_p$ }\DIFdelend \DIFaddbegin \DIFadd{with $U^{s}_p$, $U^{r}_p$ }\DIFaddend [\DIFdelbegin \DIFdel{MW\,}\DIFdelend \DIFaddbegin \DIFadd{W}\DIFaddend /(m\,K)] \DIFdelbegin \DIFdel{is the linear heat transfer
coefficient of pipe~$p$ (determined by pipe }\DIFdelend \DIFaddbegin \DIFadd{the linear heat-transfer coefficients of the supply
and return pipe, set by }\DIFaddend diameter and insulation class\DIFdelbegin \DIFdel{)}\DIFdelend , $L_p$ [m] \DIFdelbegin \DIFdel{its length,
$T_{\mathrm{sup}}^{\mathrm{nom}}(t)$
and $T_{\mathrm{ret}}^{\mathrm{nom}}(t)$ }\DIFdelend \DIFaddbegin \DIFadd{the length,
$T_{\mathrm{sup}}^{\mathrm{nom}}$, $T_{\mathrm{ret}}^{\mathrm{nom}}$ }\DIFaddend [\DIFdelbegin %DIFDELCMD < \si{\celsius}%%%
\DIFdelend \DIFaddbegin \DIFadd{°C}\DIFaddend ] the nominal
\DIFdelbegin \DIFdel{supply and return }\DIFdelend temperatures prescribed by the heating curve \DIFdelbegin \DIFdel{, and $T_{\mathrm{ground}}(t)$ }\DIFdelend \DIFaddbegin \DIFadd{and $T_{\mathrm{ground}}$ }\DIFaddend [\DIFdelbegin %DIFDELCMD < \si{\celsius}%%%
\DIFdelend \DIFaddbegin \DIFadd{°C}\DIFaddend ] the \DIFdelbegin \DIFdel{ambient ground
temperature. The factor }\DIFdelend \DIFaddbegin \DIFadd{ground
temperature; }\DIFaddend $10^6$ converts \DIFdelbegin \DIFdel{from }\DIFdelend W to MW. \DIFdelbegin \DIFdel{For L1,
all pipe losses }\DIFdelend \DIFaddbegin \DIFadd{Losses }\DIFaddend are identically zero \DIFaddbegin \DIFadd{for $\CP$ and $\NDz$,
enter exogenously as $L_t$ for $\CPL$, and are computed by \eqref{eq:loss_supply} and
\eqref{eq:loss_return} for $\ZN$ and $\Lone$ upward}\DIFaddend .

\DIFaddbegin \DIFadd{Two points about the calibration matter for what follows. First, the $U$ values are
constrained to a defensible range for the pipe class rather than fitted freely: forcing the
trunk-only loss to match the total inferred from the delivered-energy record by inflating
terminal-pipe coefficients cannot be attributed to insulation ageing. Second, what such a multiplier would absorb is the
}\emph{\DIFadd{last mile}}\DIFadd{, the service laterals and substations a trunk-only model does not represent,
and it is therefore placed where it physically belongs, entering only at the station-resolved
level $\Lfour$. The node-resolved levels consequently }\emph{\DIFadd{undercount}} \DIFadd{the network
loss by construction, and the modelled total matches the measured }\emph{\DIFadd{delivered energy}} \DIFadd{to
about one percent only once the laterals are added (Section~\ref{subsec:validation}); the
network loss this implies is model-derived, not independently metered. This is a result rather
than a defect: a quantitative statement of what a trunk-resolved model cannot see.
}

\DIFaddend \paragraph{Heat pump\DIFdelbegin \DIFdel{model}\DIFdelend .}
Each heat pump \DIFdelbegin \DIFdel{~}\DIFdelend $h$ \DIFdelbegin \DIFdel{can recover }\DIFdelend \DIFaddbegin \DIFadd{either recovers }\DIFaddend industrial waste heat \DIFdelbegin \DIFdel{(WHR source
,
temperature-varying) or fall }\DIFdelend \DIFaddbegin \DIFadd{at a time-varying source
temperature or falls }\DIFaddend back on a \DIFdelbegin \DIFdel{default }\DIFdelend low-temperature source \DIFdelbegin \DIFdel{(constant COP). The total thermal output is }\DIFdelend \DIFaddbegin \DIFadd{at constant performance, the
thermal output being }\DIFaddend the sum of \DIFdelbegin \DIFdel{both
contributions:
%DIF < 
}\DIFdelend \DIFaddbegin \DIFadd{the two contributions:
}\DIFaddend \begin{align}
  Q_{h,t} &= Q_{h,t}^{\mathrm{WHR}} + Q_{h,t}^{\mathrm{def}}\DIFaddbegin \DIFadd{,
  }\DIFaddend \label{eq:hp_split} \\[4pt]
  P_{h,t}^{\mathrm{el}} &= \frac{Q_{h,t}^{\mathrm{WHR}}}{
    \mathrm{COP}_{h,t}^{\mathrm{WHR}}}
    + \frac{Q_{h,t}^{\mathrm{def}}}{\mathrm{COP}_h^{\mathrm{def}}}\DIFaddbegin \DIFadd{.
  }\DIFaddend \label{eq:hp_pel}
\end{align}
%DIF < 
\DIFdelbegin %DIFDELCMD < \Cref{eq:hp_split} %%%
\DIFdel{splits the thermal output $Q_{h,t}$ }%DIFDELCMD < [%%%
\DIFdel{MW}%DIFDELCMD < ]
%DIFDELCMD < %%%
\DIFdel{between the two source circuits; }%DIFDELCMD < \Cref{eq:hp_pel} %%%
\DIFdel{computes the
corresponding electrical demand $P_{h,t}^{\mathrm{el}}$ }%DIFDELCMD < [%%%
\DIFdel{MW}%DIFDELCMD < ] %%%
\DIFdel{using
the time-varying waste-heat $\mathrm{COP}_{h,t}^{\mathrm{WHR}}$
and the constant default $\mathrm{COP}_h^{\mathrm{def}}$.
}\DIFdelend The waste-heat \DIFdelbegin \DIFdel{COP }\DIFdelend \DIFaddbegin \DIFadd{coefficient of performance $\mathrm{COP}_{h,t}^{\mathrm{WHR}}$ }\DIFaddend is
pre-computed offline from an analytical Lorenz-cycle model
(\DIFdelbegin %DIFDELCMD < \Cref{app:cop}%%%
\DIFdel{); this avoids introducing
}\DIFdelend \DIFaddbegin \DIFadd{Appendix~\ref{app:cop}), which avoids }\DIFaddend bilinear terms in the optimisation\DIFdelbegin \DIFdel{model}\DIFdelend \DIFaddbegin \DIFadd{. That
pre-computation is a deliberate scope restriction rather than an approximation of
convenience, and it is revisited in the limitations: a model with endogenous COP would
give the heat pump more room to respond to the network state than it has here}\DIFaddend .

Operation is bounded by a minimum load fraction $\phi_h^{\mathrm{min}}$ and the installed
capacity $\bar{Q}_h$, \DIFdelbegin \DIFdel{controlled }\DIFdelend \DIFaddbegin \DIFadd{switched }\DIFaddend by the on/off binary $u_{h,t} \in \{0,1\}$:
%DIF < 
\begin{equation}
  \phi_h^{\mathrm{min}}\, \bar{Q}_h\, u_{h,t}
  \leq Q_{h,t} \leq \bar{Q}_h\, u_{h,t}
  \quad \forall\, h,\DIFdelbegin \DIFdel{\; \forall}\DIFdelend \, t\DIFaddbegin \DIFadd{,
  }\DIFaddend \label{eq:hp_cap}
\end{equation}
%DIF < 
\DIFdelbegin \DIFdel{When $u_{h,t} = 0$, the heat pump is off, and both outputs are
forced to zero; when $u_{h,t} = 1$ the output must reach the
minimum loadfraction $\phi_h^{\mathrm{min}}$}\DIFdelend \DIFaddbegin \DIFadd{so that the unit is either off or above its minimum load}\DIFaddend , preventing uneconomic part-load
cycling.

\paragraph{Thermal storage.}
The state of charge \DIFdelbegin \DIFdel{~}\DIFdelend $E_{s,t}$ [MWh] of storage unit \DIFdelbegin \DIFdel{~}\DIFdelend $s$ \DIFdelbegin \DIFdel{evolves
according to }\DIFdelend \DIFaddbegin \DIFadd{follows }\DIFaddend a discrete-time \DIFdelbegin \DIFdel{energy balance
that accounts for
}\DIFdelend \DIFaddbegin \DIFadd{balance
with }\DIFaddend self-discharge \DIFdelbegin \DIFdel{, charging efficiency $\eta_s^{\mathrm{ch}}$, and discharging efficiency $\eta_s^{\mathrm{dis}}$:
%DIF < 
}\DIFdelend \DIFaddbegin \DIFadd{and separate charging and discharging efficiencies:
}\DIFaddend \begin{equation}
  E_{s,t} = (1 - \alpha_s)\DIFdelbegin \DIFdel{^{\Delta t}}\DIFdelend \DIFaddbegin \DIFadd{^{\Delta t/\mathrm{1\,h}}}\DIFaddend \, E_{s,t-1}
    + \eta_s^{\mathrm{ch}}\, Q_{s,t}^{\mathrm{ch}}\, \Delta t
    - \frac{Q_{s,t}^{\mathrm{dis}}\, \Delta t}{\eta_s^{\mathrm{dis}}}\DIFaddbegin \DIFadd{.
  }\DIFaddend \label{eq:storage_soc}
\end{equation}
%DIF < 
\DIFdelbegin \DIFdel{The self-discharge }\DIFdelend \DIFaddbegin \DIFadd{The }\DIFaddend coefficient $\alpha_s$\DIFdelbegin %DIFDELCMD < [%%%
\DIFdel{$h^{-1}$}%DIFDELCMD < ] %%%
\DIFdel{models standing heat }\DIFdelend \DIFaddbegin \DIFadd{, the fractional standing loss per hour, represents }\DIFaddend losses of
the \DIFdelbegin \DIFdel{storage vessel;
}\DIFdelend \DIFaddbegin \DIFadd{vessel, and }\DIFaddend the exponential form \DIFdelbegin \DIFdel{$(1-\alpha_s)^{\Delta t}$ provides a
consistent }\DIFdelend \DIFaddbegin \DIFadd{makes the }\DIFaddend discretisation independent of the time step.
A cyclic \DIFdelbegin \DIFdel{boundary }\DIFdelend condition $E_{s,|\mathcal{T}|} = E_{s,0}$ enforces energy-neutral operation over
the \DIFdelbegin \DIFdel{full simulation }\DIFdelend year, preventing artificial depletion or accumulation.

\paragraph{Remaining \DIFdelbegin \DIFdel{generation }\DIFdelend assets and emissions.}
Gas boiler, biomass boiler, electrode boiler \DIFdelbegin \DIFdel{, }\DIFdelend and CHP follow standard linear formulations\DIFdelbegin \DIFdel{(capacity-bounded,
constant
efficiency; detailed in
}%DIFDELCMD < \Cref{app:component_models}%%%
\DIFdel{). }%DIFDELCMD < 

%DIFDELCMD < %%%
\DIFdelend \DIFaddbegin \DIFadd{,
capacity-bounded with constant efficiency, detailed in
Appendix~\ref{app:component_models}. }\DIFaddend Grid and fuel emissions are \DIFdelbegin \DIFdel{computed as }\DIFdelend the product of energy
consumed and the \DIFdelbegin \DIFdel{respective emission intensity. This allows the
optimiser to trade off electricity procurement against fuel
combustion in a time-varying carbon context:
%DIF < 
}\DIFdelend \DIFaddbegin \DIFadd{corresponding emission intensity,
}\DIFaddend \begin{equation}
  E_t\DIFdelbegin \DIFdel{^{\ce{CO2},\mathrm{grid}} }\DIFdelend \DIFaddbegin \DIFadd{^{\mathrm{CO_2,grid}} }\DIFaddend =
    P_t^{\mathrm{buy}}\, e_t^{\mathrm{grid}}\, \Delta t,
  \qquad
  E_{g,t}\DIFdelbegin \DIFdel{^{\ce{CO2},\mathrm{fuel}} }\DIFdelend \DIFaddbegin \DIFadd{^{\mathrm{CO_2,fuel}} }\DIFaddend =
    F_{g,t}\, e_g^{\mathrm{fuel}}\, \Delta t\DIFaddbegin \DIFadd{,
  }\DIFaddend \label{eq:emissions}
\end{equation}
%DIF < 
\DIFdelbegin \DIFdel{where }\DIFdelend \DIFaddbegin \DIFadd{with }\DIFaddend $e_t^{\mathrm{grid}}$ [kg\DIFdelbegin \DIFdel{\,MWh$^{-1}$}\DIFdelend \DIFaddbegin \DIFadd{/MWh}\DIFaddend ] \DIFdelbegin \DIFdel{is }\DIFdelend the hourly marginal grid emission factor and
$e_g^{\mathrm{fuel}}$ [kg\DIFdelbegin \DIFdel{\,MWh$^{-1}$}\DIFdelend \DIFaddbegin \DIFadd{/MWh}\DIFaddend ] the fuel-specific \DIFdelbegin \DIFdel{emission }\DIFdelend intensity. Both \DIFdelbegin \DIFdel{values are fixed input parameters; }\DIFdelend \DIFaddbegin \DIFadd{are fixed inputs: }\DIFaddend the
optimiser is \DIFdelbegin \DIFdel{thus
}\DIFdelend carbon-aware but does not optimise the \DIFdelbegin \DIFdel{grid emission factor .
}\DIFdelend \DIFaddbegin \DIFadd{emission factor itself.
}}
\DIFaddend 

%DIF < % ============================================================
%DIF < % 3.3 EXTENDED PHYSICS (L+ and Lnl)
%DIF < % ============================================================
\DIFdelbegin %DIFDELCMD < \subsection{Extended physics formulation ($\Lplus$ / $\Lnl$)}
%DIFDELCMD < \label{sec:extended}
%DIFDELCMD < %%%
\DIFdelend \DIFaddbegin \subsection{Copperplate with aggregate losses}\label{subsec:lumped}
\DIFadd{The control condition $\CPL$ is the copperplate supplied with the aggregate network loss it
cannot itself compute, so that the loss/topology decomposition is exact. The loss enters
exogenously, by inflating the single-node demand series; the formulation stays a linear
program with the same variable count as the copperplate, and no network is added. Two calibrations of the adder are used. The first is a constant adder
\(L_a = E_\text{loss}/8760\), which spreads the annual network loss evenly over
the year; the second is a heating-curve adder
}\[
\DIFadd{L_b(t)
= \sum_p U_p\,\ell_p\bigl(T_\text{sup}(t)-T_\text{g}\bigr)
+ \sum_p U^r_p\,\ell_p\bigl(T_\text{ret}(t)-T_\text{g}\bigr),
}\]
\DIFadd{which tracks the seasonal supply and return temperatures. The sums run over
pipes \(p\) with length \(\ell_p\) and supply and return heat-loss coefficients
\(U_p\) and \(U^r_p\), respectively. Both are frozen before
the copperplate is solved, so the copperplate never sees the reference model's dispatch, only
its aggregate loss; the two calibrations bracket the effect of the temporal loss profile and,
as reported below, give the same bias and regret to within a few tenths of a percentage point.
The adder is deliberately the network's }\emph{\DIFadd{own}} \DIFadd{realised loss: this is what makes $\CPL$ a
best case for the lumped-loss shortcut, and the finding that even this best case does not
transfer across networks (Section~\ref{subsec:generalis}) is therefore conservative.
}\DIFaddend 

\DIFdelbegin \DIFdel{$\Lplus$ and $\Lnl$ augment L3 with three additional physical
phenomena: pressure-dependent pumping losses, temperature-dependent
heat }\DIFdelend \DIFaddbegin \subsection{Forward evaluator and decision regret}
\label{subsec:evaluator}
\DIFadd{The evaluator is a forward simulator, not an optimisation model: it takes a }\emph{\DIFadd{fixed}}
\DIFadd{dispatch schedule (the generation, storage and temperature trajectory a level produced)
and computes the forward-evaluated (policy-dependent) cost of executing it under high-fidelity
physics, reporting any physical
constraint the schedule violates. Because it never optimises it has no tractability limit and
contains no piecewise linearisation anywhere. Heat transport along each pipe uses the native
exponential temperature decay $T_\text{out}=T_\text{g}+(T_\text{in}-T_\text{g})\exp(-U L/\dot m
c_p)$ on both the supply and the return circuit; hydraulics use Darcy--Weisbach with a
computed (Swamee--Jain) friction factor on supply and return, plus the differential-pressure
requirement at the critical consumer's substation. This is strictly a higher-fidelity
reference than the finest solved level, whose temperature-decay factor is still piecewise
linear, the caveat that made a solved level an unsound ``ground truth'' in the original
submission. Cost is accounted with the model's own conventions (energy purchase less sales,
fuel, CO$_2$, dump and demand charge), so that a schedule executed under the physics it was
optimised for reproduces its own economic cost to within the pump and residual-loss terms
the physics adds; regret then measures physics, not an accounting mismatch. Schedules from the
copperplate and zone levels, which carry no node-resolved dispatch, are disaggregated to the
network by the per-node demand share (with an equal-share alternative reported as a
robustness check); this rule is part of the method and stated as such. Under this common evaluation, for each level \(\ell\) relative to the comparison baseline,
estimation }\emph{\DIFadd{bias}} \DIFadd{is defined as the difference in optimised economic cost and decision
}\emph{\DIFadd{regret}} \DIFadd{as the difference between the forward-evaluated costs of the two schedules.
The hours and energy associated with any velocity, differential-pressure, or unmet-demand
}\emph{\DIFadd{violations}} \DIFadd{are also reported. This regret is a }\emph{\DIFadd{forward valuation}} \DIFadd{of each
schedule under the common high-fidelity model, measured against the baseline schedule's
forward-evaluated cost, not against the (uncomputed) high-fidelity optimum: it reports what a
level's own schedule is worth, valued forward, relative to the baseline, not the distance to
the best schedule the finer physics would admit. It is accordingly a valuation overlay rather
than a constrained recourse cost: the loss shortfall is priced at the marginal unit, not
re-dispatched under commitment and ramp constraints. The evaluator
is itself validated: it reproduces the baseline schedule's economic cost as a self-consistency
check, its forward losses match the model's exported network loss to about two percent (the
exponential lying slightly below the linear approximation, as expected), and its trunk
pressures agree with an independent nonlinear pipe-flow solver to better than $0.007$\,bar.
}

\paragraph{\DIFadd{Fixed schedule, recomputed state, and shortfall valuation.}}
\DIFadd{The mapping $M_\ell$ transfers a level's decision vector to 
the evaluator while holding all first stage }\emph{\DIFadd{dispatch decisions}} 
\DIFadd{fixed exactly as produced. These decisions include generation set points, 
generator and heat pump commitment, storage charging and discharging, grid imports and exports, 
and dumped heat. The evaluator then recomputes the }\emph{\DIFadd{physical state}} \DIFadd{using the high fidelity 
model, including network flows, node temperatures, pumping power, pipe losses, and the resulting 
storage state of charge. Storage bounds and the cyclic condition are checked using this recomputed 
state. Recomputed pumping power is priced at the electricity tariff and added to the reported cost 
without adjusting the fixed grid import. If recomputed native losses are lower than scheduled losses,
 the surplus receives no value and no additional dumping is introduced. The resulting forward cost 
 therefore values the fixed schedule without representing a cheaper redispatch. A shortfall can 
 remain only when the source model does not provision the full network losses. This study does 
 not permit redispatch to cover such a shortfall. Instead, three }\emph{\DIFadd{shortfall valuation policies}}
  \DIFadd{are applied as cost overlays rather than execution simulations (Table~\ref{tab:regret}). 
  The }\emph{\DIFadd{marginal}} \DIFadd{policy prices each MWh of shortfall at the cost of the marginal heat unit. 
  The }\emph{\DIFadd{peak}} \DIFadd{policy uses the cost of the most expensive unit dispatched during that hour. 
  The }\emph{\DIFadd{unmet}} \DIFadd{policy charges the shortfall as unserved energy. 
  Because these policies are valuation overlays rather than executable recourse, 
  they do not impose commitment, capacity, ramping, grid, or CHP coupling constraints. 
  Any electricity or fuel required for a physical top up is represented only through the 
  applicable cost per MWh and is not rebalanced within the dispatch. No other adjustment 
  is permitted. The schedule cannot exploit the finer physics through redispatch, storage 
  operation is not reoptimised, and any residual unmet demand is represented by penalised 
  slack and reported in the }\emph{\DIFadd{violation}} \DIFadd{columns. By fixing decisions, recomputing 
  physical states, and applying valuation only, this procedure measures schedule quality 
  rather than the outcome of reoptimisation.
}

{
\subsection{Cost accounting: economic cost versus the solver objective}
\label{subsec:costacct}

\DIFadd{Every bias and regret figure in this paper is reported on the }\emph{\DIFadd{economic}} \DIFadd{operating
cost, not on the raw solver objective, and the two differ systematically. On Memmingen the
objective exceeds the economic cost by $38.7$--$41.0$\,\% depending on the level
(Table~S1, Supplementary Material). Because that difference is large enough to change the
apparent size of every comparison, this section states what it is and why the economic cost is the
right quantity.
}

\DIFadd{Two mechanisms account for it, and neither is a distortion of the dispatch. First, carbon
enters the objective on a }\emph{\DIFadd{gross}} \DIFadd{basis (every tonne from fuel combustion and grid
import priced at the carbon price), whereas the reported cost applies the standard
cogeneration self-use allocation, crediting back the emissions of combined-heat-and-power
electricity that is exported rather than self-consumed. This is the avoided-emissions
convention conventional in district-heating studies, and on this network it amounts to
$55$--$58$\,kEUR. Second, the objective carries a thermal-storage cycling penalty of
$2$\,EUR per MWh of throughput, worth $24$--$26$\,kEUR, which shapes the solution away from
physically implausible rapid cycling but corresponds to no operator cash flow. Together
these reproduce the residual to within about $1.5$\,kEUR per run, under one percent of the
objective.
}

\DIFadd{The economic cost is reported because it represents what an operator pays and is the primary
operator-cost basis for the reported comparisons. Dispatch also changes emissions, storage
cycling, and exports, which are accounted for separately and not folded into the bias and regret
figures. Every bias and regret figure is computed }\emph{\DIFadd{directly}} \DIFadd{from this economic cost, so
these separate accounting terms do not enter any reported comparison. The loss/topology
decomposition (Table~\ref{tab:decomposition}) is formed from the four levels' economic costs and
closes with zero residual by construction. The objective residual, including its roughly
\(5.6\,\mathrm{kEUR}\) variation across levels, lies entirely outside the decomposition. The
residual is not assumed to cancel; it does not enter the comparison in the first place. The
copperplate's estimation bias is \(-11.8\,\%\) on the raw objective and
\(\mathbf{-15.1}\,\%\) on the economic cost: the same finding, measured against a base that
either includes or excludes terms outside the operator's economic cost.
}

\DIFadd{The carbon convention is itself a modelling decision and is stated explicitly: exported
cogeneration electricity is credited with the emissions it displaces at the grid margin. The
alternative gross convention shifts the magnitudes by two to three percentage points: the
copperplate bias is \(-12.2\,\%\) rather than \(-15.1\,\%\), and the loss share of the gap is
\(97.0\,\%\) rather than \(95.8\,\%\). It changes neither the sign of any bias or regret figure
nor the loss-dominance result. Indeed, the gross basis strengthens that result because the
credit is nearly identical at every level. The per-level split is released as
}\texttt{\DIFadd{objective\_decomposition.csv}}\DIFadd{, allowing either basis to be reconstructed.
}

\DIFadd{One remaining concern is that the schedules are optimised using the solver objective, so
aligning the objective with the reported economic cost could, in principle, shift the
decomposition. This possibility is tested directly. The storage-cycling penalty is the only
objective term that both differs from the economic cost and constitutes a well-posed
optimisation term. The carbon self-use credit is a reporting allocation that, if placed in the
objective, would reward selling electricity solely to shed attributed carbon. It is therefore
excluded, and the gross-basis figures above provide the objective-consistent carbon results.
Removing the cycling penalty and re-solving the four levels to the same \(0.01\,\%\) gap lowers
every economic cost by a near-uniform \(3.5\,\%\): the loss share changes only from \(95.8\,\%\)
to \(95.6\,\%\). The decomposition is therefore stable to the distinction between the objective
and economic cost, rather than merely being arithmetically separated from it
(}\texttt{\DIFadd{loss\_share\_econobj\_sensitivity.csv}}\DIFadd{).
}}

\subsection{Extended thermo-hydraulic formulation}\label{sec:extended}
{
\DIFadd{Three phenomena are added above the node-resolved baseline $\Lone$, each on its own rung so
that no comparison changes more than one thing: temperature }\DIFaddend propagation along pipes
\DIFaddbegin \DIFadd{($\Ltwo$), trunk pressure drop and the pumping power it implies ($\Lthree$)}\DIFaddend , and transport
delay \DIFdelbegin \DIFdel{. They share
identical physics scope but differ in mathematical treatment:
$\Lplus$ linearises all nonlinear terms for MILP solvability, while
$\Lnl$ retains native quadratic and bilinear constraints (MIQCP). All PWL approximations in $\Lplus$ use SOS2 machinery 
with $K{=}3$ segments
and $K{+}1=4$ support points. A complete side-by-side
comparison is provided in }%DIFDELCMD < \Cref{app:extended_formulations}%%%
\DIFdelend \DIFaddbegin \DIFadd{($\Lsix$). A fourth object, the nonlinear reference $\NLref$, retains all three in
native form and is used for forward evaluation and bounded re-solves rather than as a rung
of the ladder. Piecewise-linear approximations use a binary multiple-choice representation with $K = 3$ segments
throughout. Both the supply and the return circuit are resolved at every level from
$\Lthree$ upward}\DIFaddend .

\paragraph{Pressure drop and pumping power \DIFaddbegin \DIFadd{($\Lthree$)}\DIFaddend .}
\DIFdelbegin %DIFDELCMD < \label{sec:model_pressure}
%DIFDELCMD < %%%
\DIFdel{Fluid friction in pipe ~}\DIFdelend \DIFaddbegin \DIFadd{Friction in pipe }\DIFaddend $p$ \DIFdelbegin \DIFdel{is governed by the }\DIFdelend \DIFaddbegin \DIFadd{follows }\DIFaddend Darcy--Weisbach\DIFdelbegin \DIFdel{equation \mbox{%DIFAUXCMD
\cite{vanderHeijde2017}}\hskip0pt%DIFAUXCMD
, which relates }\DIFdelend \DIFaddbegin \DIFadd{, relating the }\DIFaddend pressure drop $\Delta p_p$ [Pa]
to the square of the mass flow \DIFdelbegin \DIFdel{rate }\DIFdelend $\dot{m}_p$ [kg\DIFdelbegin \DIFdel{\,s$^{-1}$}\DIFdelend \DIFaddbegin \DIFadd{/s}\DIFaddend ]:
%DIF < 
\begin{equation}
  \Delta p_{p} = f_D\, \frac{L_p}{D_p}\,
    \frac{8\,\dot{m}_p^2}{\rho \pi^2 D_p^4}\DIFaddbegin \DIFadd{,
  }\DIFaddend \label{eq:darcy}
\end{equation}
%DIF < 
\DIFdelbegin \DIFdel{where }\DIFdelend \DIFaddbegin \DIFadd{with }\DIFaddend $f_D$ \DIFdelbegin %DIFDELCMD < [%%%
\DIFdel{-}%DIFDELCMD < ] %%%
\DIFdel{is }\DIFdelend the Darcy friction factor, $D_p$ [m] the inner \DIFdelbegin \DIFdel{pipe diameter , }\DIFdelend \DIFaddbegin \DIFadd{diameter }\DIFaddend and $\rho$ [kg\DIFdelbegin \DIFdel{\,m$^{-3}$}\DIFdelend \DIFaddbegin \DIFadd{/m$^3$}\DIFaddend ]
the fluid density. \DIFdelbegin \DIFdel{Since }\DIFdelend \DIFaddbegin \DIFadd{Because the }\DIFaddend supply and return temperatures are prescribed by the
heating curve, the mass flow \DIFdelbegin \DIFdel{rate is a linear function of }\DIFdelend \DIFaddbegin \DIFadd{is linear in }\DIFaddend the pipe heat flow $Q_{p,t}$ [MW]:
%DIF < 
\begin{equation}
  \dot{m}_{p,t} = \DIFdelbegin \DIFdel{\frac{Q_{p,t}}{
    c_p \bigl(T_{\mathrm{sup}}^{\mathrm{nom}}(t)
    - T_{\mathrm{ret}}^{\mathrm{nom}}(t)\bigr)}
  }\DIFdelend \DIFaddbegin \DIFadd{\frac{10^{6}\, Q_{p,t}}{
    c_p \bigl(T_{\mathrm{sup}}^{\mathrm{nom}}(t)
    - T_{\mathrm{ret}}^{\mathrm{nom}}(t)\bigr)},
  }\DIFaddend \label{eq:massflow}
\end{equation}
%DIF < 
\DIFdelbegin \DIFdel{where $c_p$ }%DIFDELCMD < [%%%
\DIFdel{J\,(kg\,K)$^{-1}$}%DIFDELCMD < ] %%%
\DIFdel{is the specific heat capacity of
water. Substituting }%DIFDELCMD < \Cref{eq:massflow} %%%
\DIFdel{into }%DIFDELCMD < \Cref{eq:darcy} %%%
\DIFdel{reveals
that pressure drop is }\emph{\DIFdel{quadratic}} %DIFAUXCMD
\DIFdel{in $Q_{p,t}$; it can be
written compactly as $\Delta p_{p,t} = R_p \cdot Q_{p,t}^2$, where
}\DIFdelend \DIFaddbegin \DIFadd{so that $\Delta p_{p,t} = R_p\, Q_{p,t}^2$ with }\DIFaddend $R_p$ [Pa\DIFdelbegin \DIFdel{\,MW$^{-2}$}\DIFdelend \DIFaddbegin \DIFadd{/MW$^2$}\DIFaddend ] \DIFdelbegin \DIFdel{is the pipe }\DIFdelend \DIFaddbegin \DIFadd{the }\DIFaddend hydraulic resistance.
\DIFdelbegin %DIFDELCMD < 

%DIFDELCMD < %%%
\DIFdel{In $\Lnl$ this quadratic relation enters directly as a native
quadratic constraint. In $\Lplus$, it }\DIFdelend \DIFaddbegin \DIFadd{In $\Lthree$ this quadratic }\DIFaddend is approximated by a \DIFaddbegin \DIFadd{multiple-choice }\DIFaddend piecewise-linear \DIFdelbegin \DIFdel{(PWL) function }\DIFdelend \DIFaddbegin \DIFadd{model }\DIFaddend over $K$
equal \DIFdelbegin \DIFdel{segments using SOS2
weights $w_{p,t,k} \in [0,1]$:
%DIF < 
}\DIFdelend \DIFaddbegin \DIFadd{flow segments with breakpoints $0 = Q^{(0)} < Q^{(1)} < \dots < Q^{(K)} = \bar{Q}_p$. A
binary $z_{p,t,k}$ selects the active segment and the flow is disaggregated onto it,
}\DIFaddend \begin{equation}
  \DIFdelbegin \DIFdel{\Delta p}\DIFdelend \DIFaddbegin \DIFadd{Q}\DIFaddend _{p,t} \DIFdelbegin \DIFdel{^{\mathrm{PWL}} }\DIFdelend = \sum_{k=1}^{K} \DIFdelbegin \DIFdel{m_k^{(p)} \cdot w}\DIFdelend \DIFaddbegin \DIFadd{Q}\DIFaddend _{p,t,k}\DIFdelbegin %DIFDELCMD < \label{eq:pressure_pwl}
%DIFDELCMD < %%%
\DIFdelend \DIFaddbegin \DIFadd{, \quad
  \sum_{k=1}^{K} z_{p,t,k} = 1, \quad
  Q^{(k-1)} z_{p,t,k} }\le \DIFadd{Q_{p,t,k} }\le \DIFadd{Q^{(k)} z_{p,t,k}, \quad z_{p,t,k} \in \{0,1\},
  }\label{eq:pressure_pwl_sel}
\DIFaddend \end{equation}
%DIF < 
\DIFdelbegin \DIFdel{where }\DIFdelend \DIFaddbegin \DIFadd{so exactly one segment carries the flow (it is these segment binaries the solver branches on,
not SOS2 sets); the pressure drop is the chord of $R_p Q^2$ on the active segment,
}\begin{equation}
  \DIFadd{\Delta p_{p,t}^{\mathrm{PWL}} = \sum_{k=1}^{K}
    \bigl(m_k^{(p)}\, Q_{p,t,k} + c_k^{(p)}\, z_{p,t,k}\bigr),
  \label{eq:pressure_pwl}
}\end{equation}
\DIFadd{with }\DIFaddend $m_k^{(p)}$ \DIFdelbegin \DIFdel{are the slopes oneach segment. The }\DIFdelend \DIFaddbegin \DIFadd{and $c_k^{(p)}$ the slope and intercept of the $k$-th chord of $R_p Q^2$. Its
}\DIFaddend maximum approximation error \DIFdelbegin \DIFdel{of this PWL is (derivation in }%DIFDELCMD < \Cref{app:pwl}%%%
\DIFdel{):
%DIF < 
}\begin{displaymath}
  \DIFdel{\varepsilon_{\Delta p}^{\max}
    = \frac{R_p\, \bar{Q}_p^2}{4K^2}
    \approx 2.8\,\% \text{ of peak pressure drop for } K}{\DIFdel{=}}\DIFdel{3
}\end{displaymath}%DIFAUXCMD
%DIF < 
\DIFdel{where $\bar{Q}_p$ }%DIFDELCMD < [%%%
\DIFdel{MW}%DIFDELCMD < ] %%%
\DIFdel{is the rated pipe capacity}\DIFdelend \DIFaddbegin \DIFadd{is $\varepsilon^{\max}_{\Delta p} = R_p \bar{Q}_p^2/(4K^2)$,
about $2.8$\,\% of the peak pressure drop at $K = 3$ (Appendix~\ref{app:pwl}). In $\NLref$
the quadratic enters natively}\DIFaddend .

\DIFdelbegin \DIFdel{The network pumping power $P_t^{\mathrm{pump}}$ }%DIFDELCMD < [%%%
\DIFdel{MW}%DIFDELCMD < ] %%%
\DIFdel{is the sum
of hydraulic work across all pipes , normalised }\DIFdelend \DIFaddbegin \DIFadd{Pumping power is the hydraulic work summed over pipes and divided }\DIFaddend by the pump efficiency\DIFdelbegin \DIFdel{$\eta_{\mathrm{pump}}$:
%DIF < 
}\begin{displaymath}
  \DIFdel{P_t^{\mathrm{pump}} = \frac{1}{\eta_{\mathrm{pump}}}
    \sum_{p \in \mathcal{P}}
    \frac{\dot{m}_{p,t} \cdot \Delta p_{p,t}}{\rho}
  \label{eq:pump_power}
}\end{displaymath}%DIFAUXCMD
\DIFdelend \DIFaddbegin \DIFadd{,
}\begin{equation}
  \DIFadd{P_t^{\mathrm{pump}} = \frac{10^{-6}}{\eta_{\mathrm{pump}}}
    \sum_{p \in \mathcal{P}}
    \frac{\dot{m}_{p,t}\, \Delta p_{p,t}}{\rho},
  \label{eq:pump_power}
}\end{equation}\DIFaddend 
%DIF < 
\DIFdelbegin \DIFdel{Because this involves $\dot{m}_{p,t} \propto Q_{p,t}$ and
$\Delta p_{p,t} \propto Q_{p,t}^2$,
pumping power }\DIFdelend \DIFaddbegin \DIFadd{which }\DIFaddend is cubic in $Q_{p,t}$ \DIFdelbegin \DIFdel{. In $\Lplus$, the cubic relation is directly
}\DIFdelend \DIFaddbegin \DIFadd{and is }\DIFaddend piecewise-linearised \DIFdelbegin \DIFdel{. In $\Lnl$, it is modelled via the bilinear
auxiliary variable }\DIFdelend \DIFaddbegin \DIFadd{directly in $\Lthree$; in $\NLref$
it is carried by the auxiliary }\DIFaddend $W_{p,t} = Q_{p,t}^2$ \DIFdelbegin \DIFdel{, entering the objective
as $Q_{p,t} \cdot W_{p,t}$. }\DIFdelend \DIFaddbegin \DIFadd{entering as $Q_{p,t} W_{p,t}$; the
auxiliary keeps the formulation quadratically constrained (a bilinear $Q_{p,t}W_{p,t}$ with the
quadratic definition $W_{p,t}=Q_{p,t}^2$) rather than cubic, solved globally with
}\texttt{\DIFadd{NonConvex=2}}\DIFadd{. The $10^{6}$ in
Eq.~\eqref{eq:massflow} and the $10^{-6}$ here reconcile the MW heat and power units with the
SI $c_p$ }[\DIFadd{J\,kg$^{-1}$\,K$^{-1}$}]\DIFadd{, $\Delta p$ }[\DIFadd{Pa}] \DIFadd{and $\dot m$ }[\DIFadd{kg\,s$^{-1}$}] \DIFadd{of the physical
relations, matching the implementation.
}\DIFaddend 

\DIFdelbegin \paragraph{\DIFdel{Temperature-dependent loss propagation.}}
%DIFAUXCMD
\addtocounter{paragraph}{-1}%DIFAUXCMD
%DIFDELCMD < \label{sec:model_temp_prop}
%DIFDELCMD < %%%
\DIFdel{In L3, pipe losses are computed }\DIFdelend \DIFaddbegin \paragraph{\DIFadd{Temperature propagation ($\Ltwo$).}}
\DIFadd{Where $\Lone$ computes pipe losses }\DIFaddend from fixed nominal temperatures\DIFdelbegin \DIFdel{. In
$\Lplus$ and $\Lnl$, the actual outlet temperature is tracked }\DIFdelend \DIFaddbegin \DIFadd{, $\Ltwo$ tracks the
outlet temperature }\DIFaddend as the fluid travels\DIFdelbegin \DIFdel{along pipe~$p$. Under steady-state conditions,
the outlet temperature $T_{p,t}^{\mathrm{out}}$ }%DIFDELCMD < [\si{\celsius}]
%DIFDELCMD < %%%
\DIFdelend \DIFaddbegin \DIFadd{. In steady state it }\DIFaddend decays exponentially from the
inlet temperature \DIFdelbegin \DIFdel{$T_{n_1,t}^{\mathrm{sup}}$ }%DIFDELCMD < [\si{\celsius}] %%%
\DIFdelend towards the ground temperature \DIFdelbegin \DIFdel{$T_{\mathrm{gr}}(t)$ }%DIFDELCMD < [\si{\celsius}]%%%
\DIFdel{:
%DIF < 
}\DIFdelend \DIFaddbegin \DIFadd{$T_{\mathrm{gr}}$:
}\DIFaddend \begin{equation}
  T_{p,t}^{\mathrm{out}} = T_{\mathrm{gr}}(t)
    + \bigl(T_{n_1,t}^{\mathrm{sup}} - T_{\mathrm{gr}}(t)\bigr)
    \cdot \underbrace{\exp\!\left(
      \frac{-U_p L_p}{\dot{m}_{p,t}\, c_p}
    \right)}_{\displaystyle\phi_p(\dot{m}_{p,t})}\DIFaddbegin \DIFadd{.
  }\DIFaddend \label{eq:temp_prop}
\end{equation}
%DIF < 
The decay factor \DIFdelbegin \DIFdel{$\phi_p \in (0,\,1]$ quantifies what }\DIFdelend \DIFaddbegin \DIFadd{$\phi_p \in (0,1]$ is the }\DIFaddend fraction of the inlet temperature excess
\DIFdelbegin \DIFdel{is preserved at the outlet; it
approaches unity for high flowrates (little cooling) and zero
for near-zero flow (complete equalisation to ground temperature}\DIFdelend \DIFaddbegin \DIFadd{surviving to the outlet: near unity at high flow, near zero as flow approaches zero; the
forward evaluator floors $\dot{m}_p$ at a small positive value so the exponent stays finite,
and at zero flow $\phi_p$ is taken as its limit, $0$ (full decay}\DIFaddend ). Two
nonlinearities \DIFdelbegin \DIFdel{arise from this expression}\DIFdelend \DIFaddbegin \DIFadd{follow}\DIFaddend : the exponential \DIFdelbegin \DIFdel{$\phi_p(\dot{m}_{p,t})$ }\DIFdelend \DIFaddbegin \DIFadd{itself }\DIFaddend and the bilinear product
\DIFdelbegin \DIFdel{$T_{n_1,t}^{\mathrm{sup}} \cdot \phi_p$. }%DIFDELCMD < 

%DIFDELCMD < %%%
\DIFdel{Both $\Lplus$ and $\Lnl$ approximate $\phi_p$ with a shared
PWL over $K$ breakpoints }\DIFdelend \DIFaddbegin \DIFadd{$T_{n_1,t}^{\mathrm{sup}} \phi_p$. The exponential is approximated by a PWL }\DIFaddend in $\dot{m}$ \DIFdelbegin \DIFdel{(this shared error is
excluded from the
measured linearisation gap )}\DIFdelend \DIFaddbegin \DIFadd{in the
PWL-linearised level $\Ltwo$, and enters natively in $\NLref$ and the nonlinear 72-hour re-solves; the
reported PWL-versus-native gap therefore combines this exponential approximation with the
product approximation below, rather than cancelling it}\DIFaddend . The bilinear product \DIFdelbegin \DIFdel{is handled differently in each level. In $\Lnl$,
it enters directly as $T_{n_1,t}^{\mathrm{sup}} \cdot \phi_{p,t}^{\mathrm{PWL}}$.
In $\Lplus$, }\DIFdelend \DIFaddbegin \DIFadd{likewise enters
natively in $\NLref$; in $\Ltwo$ it is linearised by }\DIFaddend a first-order Taylor expansion \DIFdelbegin \DIFdel{around }\DIFdelend \DIFaddbegin \DIFadd{about }\DIFaddend the
nominal operating point\DIFdelbegin \DIFdel{$(T^{\mathrm{nom}}(t),\, \phi_p^{\mathrm{nom}})$
linearises the product:
%DIF < 
}\DIFdelend \DIFaddbegin \DIFadd{,
}\DIFaddend \begin{align}
  T_{p,t}^{\mathrm{out}} &\approx T_{\mathrm{gr}}(t)
    + \bigl(T^{\mathrm{nom}}(t) - T_{\mathrm{gr}}(t)\bigr)
      \DIFdelbegin \DIFdel{\cdot }\DIFdelend \phi_{p,t}^{\mathrm{PWL}} \notag \\
    &\DIFaddbegin \DIFadd{\quad }\DIFaddend + \phi_p^{\mathrm{nom}}
      \DIFdelbegin \DIFdel{\cdot }\DIFdelend \bigl(T_{n_1,t}^{\mathrm{sup}} - T^{\mathrm{nom}}(t)\bigr)\DIFaddbegin \DIFadd{,
  }\DIFaddend \label{eq:temp_prop_lin}
\end{align}
%DIF < 
\DIFdelbegin \DIFdel{The }\DIFdelend \DIFaddbegin \DIFadd{with the }\DIFaddend neglected second-order cross-term \DIFdelbegin \DIFdel{is bounded by
$|\Delta T_{\max}| \cdot |\Delta\phi_{\max}| < 2.5\,\text{K}$
(derivation in }%DIFDELCMD < \Cref{app:taylor_derivation}%%%
\DIFdelend \DIFaddbegin \DIFadd{bounded by
$|\Delta T_{\max}||\Delta\phi_{\max}| < 2.5$\,K (Appendix~\ref{app:taylor_derivation}}\DIFaddend ).

\DIFaddbegin \DIFadd{One property of this formulation determines how it is reported. Freeing the node temperatures
makes the programme non-convex and, absent a hydraulic penalty on the resulting flow
increase, degenerate: the optimiser reduces the supply temperature uniformly to its lower
bound, collapsing the loss term rather than reproducing the physical decay. In the solve it
does so for about $91$\,\% of hours, against a mean supply temperature of $85.9$\,°C at the
neighbouring levels. The solved objective of $\Ltwo$ is therefore not a meaningful fidelity
measurement and is not used; the effect of temperature propagation is quantified
instead through the forward evaluator of Section~\ref{subsec:evaluator}, which applies the
native exponential to a fixed schedule and cannot exploit the degeneracy. $\Ltwo$ remains in
the ladder as the formulation whose linearisation error is measured, not as a source of
dispatch cost.
}

\DIFaddend \paragraph{Transport delay \DIFaddbegin \DIFadd{($\Lsix$)}\DIFaddend .\DIFdelbegin %DIFDELCMD < \MBLOCKRIGHTBRACE
%DIFDELCMD < \label{sec:model_delay}
%DIFDELCMD < %%%
\DIFdel{At finite flow }\DIFdelend \DIFaddbegin }
\DIFadd{At finite }\DIFaddend velocity, heat injected at the source at time \DIFdelbegin \DIFdel{~}\DIFdelend $t$ \DIFdelbegin \DIFdel{does not arrive at the consumer until }\DIFdelend \DIFaddbegin \DIFadd{arrives downstream }\DIFaddend later. The
nominal travel time \DIFdelbegin \DIFdel{$\tau_p$ }%DIFDELCMD < [%%%
\DIFdel{h}%DIFDELCMD < ] %%%
\DIFdel{through pipe ~}\DIFdelend \DIFaddbegin \DIFadd{through pipe }\DIFaddend $p$\DIFdelbegin \DIFdel{is determined by the
pipe cross-sectional area }\DIFdelend \DIFaddbegin \DIFadd{, with cross-section }\DIFaddend $A_p = \pi D_p^2/4$\DIFdelbegin %DIFDELCMD < [%%%
\DIFdel{m$^2$}%DIFDELCMD < ] %%%
\DIFdel{and
the nominal mass flow $\dot{m}_p^{\mathrm{nom}}$ }%DIFDELCMD < [%%%
\DIFdel{kg\,s$^{-1}$}%DIFDELCMD < ]%%%
\DIFdel{:
%DIF < 
}\DIFdelend \DIFaddbegin \DIFadd{, is
}\DIFaddend \begin{equation}
  \tau_p = \frac{L_p\, \rho\, A_p}{\dot{m}_p^{\mathrm{nom}}}\DIFaddbegin \DIFadd{,
  }\DIFaddend \label{eq:delay_time}
\end{equation}
%DIF < 
\DIFdelbegin \DIFdel{Discretising to the nearest integer number of time steps $k_p = \mathrm{round}(\tau_p / \Delta t)$, }\DIFdelend \DIFaddbegin \DIFadd{discretised to $k_p = \lfloor \tau_p/\Delta t \rfloor$ time steps (with $\tau_p$ from
Eq.~\eqref{eq:delay_time} in seconds and the hourly step taken as $\Delta t = 3600$\,s, so
$k_p$ counts whole hours), so }\DIFaddend the heat delivered \DIFdelbegin \DIFdel{at the downstream end }\DIFdelend \DIFaddbegin \DIFadd{downstream }\DIFaddend is the heat injected $k_p$ steps
earlier \DIFdelbegin \DIFdel{,
reduced by the corresponding pipe losses:
%DIF < 
}\DIFdelend \DIFaddbegin \DIFadd{less the corresponding losses:
}\DIFaddend \begin{equation}
  Q_{p,t}^{\mathrm{delivered}} =
    Q_{p,\,t-k_p} - Q_{p,\,t-k_p}^{\mathrm{loss,pipe}}
  \quad \forall\, p,\; t \geq k_p{+}\DIFdelbegin \DIFdel{1
  }\DIFdelend \DIFaddbegin \DIFadd{1.
  }\DIFaddend \label{eq:delayed_balance}
\end{equation}
%DIF < 
\DIFdelbegin \DIFdel{Transport delay is active only in $\Lnl$. Under $\Lplus$ the linearised pipemodel bypasses the lag because the Taylor
linearisation of $\phi_p$ is performed at fixed nominal flow, which implicitly assumes instantaneous transport
(cf. \ }%DIFDELCMD < \Cref{tab:model_scope}%%%
\DIFdel{)}\DIFdelend \DIFaddbegin \DIFadd{Isolating delay on its own rung allows it to be separated from the linearisation error rather 
than confounded with it. It also produces a clean negative result, whose reason is visible 
in \eqref{eq:delay_time}: every trunk travel time in this network is under an hour, with 
the longest being about 35 minutes, so \(k_p=0\) for every pipe and \eqref{eq:delayed_balance} 
collapses to the undelayed balance. \(\Lsix\) solves \(\Lthree\) plus this integer-step delay; 
it does }\emph{\DIFadd{not}} \DIFadd{inherit the forward-only station levels \(\Lfour\) and \(\Lfive\), which 
lie outside the solved ladder. With \(k_p=0\) at hourly resolution, \(\Lsix\) and \(\Lthree\) 
are therefore the same programme, with identical variable, binary, and constraint counts and 
an identical optimum to the last digit. The represented integer-step delay effect is zero at 
hourly resolution rather than merely small, and this is a statement about temporal resolution: 
at sub-hourly resolution or in a network with several additional kilometres of trunk, \(k_p\) 
would not vanish. This is checked at the }\emph{\DIFadd{path}} \DIFadd{level, not only per pipe: the longest 
source-to-consumer path, the approximately \(2.8\,\mathrm{km}\) trunk to the far end, has a 
cumulative nominal travel time below one hour, so flooring each pipe separately does not 
erase a cumulative delay that would otherwise exceed the time step.}}
\DIFadd{The extended formulation resolves both the supply and the return circuit. Trunk pressure drop
(from $\Lthree$) follows Darcy--Weisbach, linearised as a piecewise function of flow, and the
associated pumping power enters the objective at the pump efficiency taken from the installed
units' datasheets (Wilo IL-E). Every consumer additionally requires a differential pressure
across its substation, set to the operator's specified $0.6$\,bar; this is the term coarser
trunk-only treatments omit. The resulting pumping-energy share is small and, rather than being
calibrated away, is reconciled physically by the compact network and its heavily oversized pumps;
the magnitudes, and the agreement with an independent nonlinear pipe-flow solver, are reported
with the results in Section~\ref{subsec:hydraulics}.
}

\subsubsection{Station-resolved hydraulics (L4, L5): the test that locates the threshold}\label{subsec:stationhyd}
\DIFadd{The station levels are evaluated with the forward model rather than solved as additional
mixed-integer dispatch levels, and this is the rigorous choice, not a compromise, for four
reasons. First, the forward module uses the native nonlinear station pressure relation,
validated against real component data and the nonlinear solver (exactly the validated
nonlinear reference), whereas a mixed-integer level would re-linearise
that relation and reintroduce the very linearisation this forward treatment avoids. Second,
re-costing the baseline schedule under the full station physics measures the cost of executing
that fixed schedule: it rises by under one percent, with no feasibility violation. This is a
measurement of the fixed schedule under the finer physics, not a bound over all schedules
(re-optimising under that physics could in principle yield a different schedule), but it shows
the baseline schedule stays cheap and deliverable when the station detail is added. Third, it is consistent with the paper's position that the physics
reference is the forward evaluator and not further piecewise-linear levels. Fourth, station
pumping is load-following (it tracks flow and hence the fixed demand), so there is little scope
to re-optimise around it, and bias, regret and the forward adder nearly coincide in this case}\DIFaddend .

\DIFdelbegin %DIFDELCMD < \Cref{tab:lplus_vs_lnl} %%%
\DIFdel{in }%DIFDELCMD < \Cref{app:extended_formulations}
%DIFDELCMD < %%%
\DIFdel{provides a complete side-by-side comparison of all extended-physics
treatments in $\Lplus$ and $\Lnl$. }\DIFdelend \DIFaddbegin \DIFadd{Temperature propagation ($\Ltwo$) is treated in the same manner for a related reason.
For a prescribed plant supply temperature, downstream nodes become cooler as heat
dissipates along the pipes. Representing this process requires products of temperature
and mass flow, making the physics bilinear and therefore nonconvex. Treating nodal
temperatures as free variables within the mixed integer formulation would be both
computationally intractable and }\emph{\DIFadd{degenerate}}\DIFadd{. The optimiser would gain the freedom
to lower the overall temperature level. Without a hydraulic penalty, it would drive
every node to a common lower bound rather than reproduce the physical temperature
decay. This outcome represents temperature }\emph{\DIFadd{optimisation}}\DIFadd{, which is a separate
question addressed in Section~\ref{subsec:tsup}, rather than temperature propagation.
This study therefore isolates the propagation effect using the validated nonlinear
forward reference. The evaluator fixes the plant supply temperature at the heating
curve value prescribed by the schedule and propagates the native exponential
temperature decay downstream. The difference relative to the fixed temperature
baseline defines the temperature propagation effect. This effect is reported in
Section~\ref{subsec:linearisation} and isolated through the exact decomposition.
Temperature propagation reduces total network losses by about 2\%, corresponding to
only a few tenths of a percent of total cost. Like station hydraulics, its effect
therefore remains below the decision relevance threshold.
Transport delay ($\Lsix$) is a linear time shift of delivered heat and is isolated as its own
level. The fully nonlinear reference (bilinear temperature propagation and quadratic
pressure together) is not solved }\emph{\DIFadd{globally}}\DIFadd{: the full-year mixed-integer
}\emph{\DIFadd{nonlinear}} \DIFadd{program is intractable (no incumbent within a day). Its distance from the
linearised levels is instead characterised by the measured linearisation error and by the
forward evaluator, and by bounded fixed-integer re-solves on representative 72-hour windows
(Section~\ref{subsec:linearisation}), without asserting an unattainable full-year global solve.
}\DIFaddend 

%DIF < % ============================================================
%DIF < % 3.4 VALIDATION
%DIF < % ============================================================
\subsection{Validation protocol}
\DIFdelbegin %DIFDELCMD < \label{sec:val_protocol}
%DIFDELCMD < 

%DIFDELCMD < %%%
\DIFdel{The network underwent a technology upgrade (HP}\DIFdelend \DIFaddbegin {
\DIFadd{The network was upgraded with the heat pump}\DIFaddend , electrode boiler \DIFdelbegin \DIFdel{,
TES) }\DIFdelend \DIFaddbegin \DIFadd{and thermal storage
}\DIFaddend \emph{after} the measurement record was collected\DIFdelbegin \DIFdel{. Following
\mbox{%DIFAUXCMD
\citet{Kus2025} }\hskip0pt%DIFAUXCMD
and \mbox{%DIFAUXCMD
\citet{Maldonado2024}}\hskip0pt%DIFAUXCMD
, this work adopts }\DIFdelend \DIFaddbegin \DIFadd{, so validation follows }\DIFaddend a two-stage
strategy \DIFaddbegin \DIFadd{after \mbox{%DIFAUXCMD
\citet{Kus2025} }\hskip0pt%DIFAUXCMD
and \mbox{%DIFAUXCMD
\citet{Maldonado2024}}\hskip0pt%DIFAUXCMD
}\DIFaddend .

\paragraph{Stage \DIFdelbegin \DIFdel{~}\DIFdelend 1: \DIFdelbegin \DIFdel{Direct }\DIFdelend \DIFaddbegin \DIFadd{direct }\DIFaddend network validation.}
The pipe model is validated against pre-upgrade monitoring data \DIFdelbegin \DIFdel{(}\DIFdelend \DIFaddbegin \DIFadd{with the }\DIFaddend new assets set to
zero capacity\DIFdelbegin \DIFdel{). Calibration follows the branch-sequential $U_p$-tuning procedure of \mbox{%DIFAUXCMD
\citet{Maldonado2024}}\hskip0pt%DIFAUXCMD
}\DIFdelend . \DIFdelbegin %DIFDELCMD < \Cref{tab:val_targets} %%%
\DIFdel{lists
the KPIs with
gatesdefined }\DIFdelend \DIFaddbegin \DIFadd{Following the calibration methodology of \mbox{%DIFAUXCMD
\citet{Maldonado2024}}\hskip0pt%DIFAUXCMD
, the trunk
heat-transfer coefficients $U_p$ are constrained to a defensible range for the pipe class;
this avoids the terminal-pipe inflation that an unconstrained fit to delivered energy would
require, as discussed in Section~\ref{subsec:validation}. Table~\ref{tab:valtargets} lists
the key performance indicators and their gates, fixed }\DIFaddend \emph{ex ante} from published practice
\DIFdelbegin \DIFdel{. The
$T_{\mathrm{sup}}$ gates are assigned to $\Lnl$ because only this
variant propagates supply temperature along pipes; L3 uses fixed
nominal temperatures and therefore cannot predict downstream
temperature deviations. }%DIFDELCMD < 

%DIFDELCMD < \begin{table}[ht]
%DIFDELCMD < \centering
%DIFDELCMD < \begin{threeparttable}
%DIFDELCMD < \caption{Stage~1 validation KPIs and gates (defined ex ante).}
%DIFDELCMD < \label{tab:val_targets}
%DIFDELCMD < \footnotesize
%DIFDELCMD < \setlength{\tabcolsep}{2pt}
%DIFDELCMD < \begin{tabular*}{\columnwidth}{@{\extracolsep{\fill}} l l l l @{}}
%DIFDELCMD < \toprule
%DIFDELCMD < KPI & Model & Gate & Source \\
%DIFDELCMD < \midrule
%DIFDELCMD < $T_{\mathrm{sup}}$, mean
%DIFDELCMD <   & $\Lnl$ & ${<}\SI{1.5}{\kelvin}$
%DIFDELCMD <   & \citet{Kus2025} \\
%DIFDELCMD < $T_{\mathrm{sup}}$, worst
%DIFDELCMD <   & $\Lnl$ & ${<}\SI{2.5}{\kelvin}$
%DIFDELCMD <   & \citet{Maldonado2024} \\
%DIFDELCMD < $T_{\mathrm{ret}}$ at source
%DIFDELCMD <   & L3 & ${<}\SI{2.5}{\kelvin}$
%DIFDELCMD <   & \citet{Kus2025} \\
%DIFDELCMD < Flow (MAPE)
%DIFDELCMD <   & L3 & n/a\tnote{a} & -- \\
%DIFDELCMD < Energy balance
%DIFDELCMD <   & L3 & ${<}2\,\%$ & Conserv. \\
%DIFDELCMD < \bottomrule
%DIFDELCMD < \end{tabular*}
%DIFDELCMD < \begin{tablenotes}\footnotesize
%DIFDELCMD <   \item[a] %%%
\DIFdelFL{${\pm}15\,\%$ instrument error; no gate
.
}%DIFDELCMD < \end{tablenotes}
%DIFDELCMD < \end{threeparttable}
%DIFDELCMD < \end{table}
%DIFDELCMD < %%%
\DIFdelend \DIFaddbegin \DIFadd{before the calibration was run.
}

\DIFadd{The use of these gates is deliberate. All are reported in
Section~\ref{subsec:validation}, each against the formulation to which it applies. The
supply-temperature MAE gates are met by the propagating formulation on the validation nodes
(Table~\ref{tab:valtargets}); what the instrumentation cannot support is the }\emph{\DIFadd{annual,
point}} \DIFadd{temperature field, whose far-end and corridor comparisons are dominated by a
mixing-valve sensor offset and are therefore reported as diagnostics rather than as gated
metrics. Retaining and reporting these diagnostics is the point rather than an oversight:
they support the claim that the resolution of a model cannot be verified beyond that of its
metering. The conclusions are instead anchored to annual delivered energy, for which the gate
is met; the annual network loss is then }\emph{\DIFadd{model-implied}} \DIFadd{(computed from the calibrated
pipe model rather than measured), with delivered-energy agreement fixing the demand scale to
which that loss is referenced, not the loss itself.
}\DIFaddend 

\paragraph{Stage \DIFdelbegin \DIFdel{~}\DIFdelend 2: \DIFdelbegin \DIFdel{Necessary-condition }\DIFdelend \DIFaddbegin \DIFadd{necessary-condition }\DIFaddend verification.}
Dispatch of \DIFaddbegin \DIFadd{the }\DIFaddend new assets is verified \DIFdelbegin \DIFdel{through plausibility bounds :
HP COP within $[2.5,\,5.5]$; electrode boiler price-response
correlation $r > 0.3$; }\DIFdelend \DIFaddbegin \DIFadd{against plausibility bounds rather than against
measurements: the heat-pump coefficient of performance must remain within \([2.5,5.5]\),
the thermal-storage state of charge within \([0.05,0.95]\bar{E}_s\), simultaneous charging
and discharging are prohibited, and }\DIFaddend per-step energy closure \DIFdelbegin \DIFdel{${<}\,1\,\%$; TES
SOC within $[0.05,\,0.95] \cdot \bar{E}_s$. These are necessary,
not sufficient ; }\DIFdelend \DIFaddbegin \DIFadd{is required. These checks are
necessary but not sufficient conditions. A }\DIFaddend post-upgrade \DIFdelbegin \DIFdel{data are recommended for }\DIFdelend \DIFaddbegin \DIFadd{measurement campaign is therefore
the natural }\DIFaddend follow-up \DIFaddbegin \DIFadd{and is identified as a limitation}\DIFaddend .

\paragraph{Structural sanity checks.}
Two \DIFdelbegin \DIFdel{additional checks ensure model consistency: (i)~$\Lnl$ with all }\DIFdelend \DIFaddbegin \DIFadd{checks confirm internal consistency. The nonlinear reference with its }\DIFaddend quadratic
constraints deactivated reproduces \DIFdelbegin \DIFdel{L3's objective within 0.01}\DIFdelend \DIFaddbegin \DIFadd{the baseline objective to within $0.01$}\DIFaddend \,\%; \DIFdelbegin \DIFdel{(ii)~the
relaxation property
$\mathrm{cost}(\mathrm{L1}) \leq \mathrm{cost}(\mathrm{L3})$ holds for every synthetic instance}\DIFdelend \DIFaddbegin \DIFadd{and the
ordering $\mathrm{cost}(\CP) \le \mathrm{cost}(\Lone)$ holds on every one of the
135 synthetic instances, the expected direction when the copperplate is a loss-free relaxation
of the resolved model on the same network, though not a guaranteed inequality once features
such as CHP export and unit commitment differ between the two. A departure from either check
on this test set would flag a formulation error rather than a finding.
}}
\begin{table}[t]
\centering
\caption{Validation targets. Gates were fixed \emph{ex ante} from the values reported for
comparable validations by \citet{Kus2025} and \citet{Maldonado2024}, before calibration.
The supply-temperature gates apply to the propagating formulation only and are
metering-limited; the return-temperature check tests the assumed differential-temperature
parameterisation of the dispatch model rather than a resolved temperature field
(Section~\ref{subsec:validation}).}
\label{tab:valtargets}
\begin{tabular}{lccc}
\toprule
Quantity & Gate & Formulation & Status \\
\midrule
\multicolumn{4}{l}{\emph{Gated}}\\
Annual delivered-energy error       & $\le 2$\,\%   & MILP & met \\
Energy-balance closure              & $\le 2$\,\%   & MILP & met \\
Supply-temperature MAE, mean (val.\ nodes, 744\,h)  & $\le 1.5$\,K  & NLP  & met (1.32\,K) \\
Supply-temperature MAE, worst node (744\,h)         & $\le 2.5$\,K  & NLP  & met (2.27\,K) \\
\addlinespace
\multicolumn{4}{l}{\emph{Calibration / diagnostic (not gated validation)}}\\
Return-temperature MAE at source    & (calib.) & MILP ($\Delta T$ param.) & calibration check \\
Far-end supply-temp MAE (annual)    & (diag.)  & NLP  & 9.21\,K (metering-limited) \\
\addlinespace
\multicolumn{4}{l}{\emph{Reported, not gated}}\\
Flow MAPE, source (reported, not gated)  & 33.8\,\%  & MILP \\
Demand MAPE, total (reported, not gated) & 29.1\,\%  & MILP \\
\bottomrule
\end{tabular}
\end{table}
\DIFadd{Validation resolution bounds what fidelity can be claimed, and the Memmingen metering makes the
point concretely. Its consumer sensors are billing instruments downstream of mixing valves, so
they read several degrees below the primary junction temperature; intermediate node temperatures
are therefore not directly testable, and even the source-to-far-end corridor match is
flow-sensitive under billing-grade metering. The defensible thermal validation is therefore the
}\emph{\DIFadd{annual delivered energy}} \DIFadd{(matched to about one percent), from which the network loss the
conclusions price is model-implied, rather than the point temperature field. The station hydraulics are validated separately against real component data and an
independent nonlinear solver. The general principle is that a model resolved finer than its
metering cannot have that extra detail verified: its additional resolution is asserted, not
tested. This has a practical
consequence and a converse. The practical consequence is that, for an operator metering at the
substation or zone, the node or zone thermal level is the finest }\emph{\DIFadd{defensible}} \DIFadd{level, not
merely the cheapest. The converse, established in the results, is that forward-evaluating
hydraulics all the way to the individual substation on the fixed schedule adds under one percent
to its cost and reveals no violation, so the fidelity requirement sits below the level at which
extra detail could even be validated}\DIFaddend .

%DIF < % ============================================================
%DIF < % 3.5 IMPLEMENTATION
%DIF < % ============================================================
\subsection{Implementation}
\DIFdelbegin %DIFDELCMD < \label{sec:model_impl}
%DIFDELCMD < 

%DIFDELCMD < %%%
\DIFdel{All five variants }\DIFdelend \DIFaddbegin {
\DIFadd{All levels }\DIFaddend are implemented in Python \DIFdelbegin \DIFdel{using PYOMO 
}\DIFdelend \DIFaddbegin \DIFadd{with Pyomo~}\DIFaddend \cite{pyomo}, instantiated from a single
parameterised codebase \DIFdelbegin \DIFdel{via }\DIFdelend \DIFaddbegin \DIFadd{through }\DIFaddend YAML configuration files\DIFdelbegin \DIFdel{. Solver: }\DIFdelend \DIFaddbegin \DIFadd{, so that a fidelity level is a
configuration of one model rather than a separate model. This is what makes the
one-phenomenon-per-step contrasts exact: two levels differ in the terms their configuration
activates and in nothing else, and a level with a phenomenon disabled reproduces the level
below it to the cent. Solves use }\DIFaddend Gurobi~13.0 with \DIFdelbegin \DIFdel{MIP gap
${\leq}\,0.5\,\%$ for both MILP and MIQCP
(}\texttt{\DIFdel{NonConvex=2}}%DIFAUXCMD
\DIFdel{); }\DIFdelend identical random seeds \DIFdelbegin \DIFdel{. $\Lnl$ is warm-started from the $\Lplus$ solution}\DIFdelend \DIFaddbegin \DIFadd{on a 66-core,
180\,GB machine, with }\texttt{\DIFadd{NonConvex=2}} \DIFadd{where a quadratic formulation is involved}\DIFaddend .

\DIFdelbegin \DIFdel{The $\Lplus$ formulation adds ${\approx}\,753{,}000$ variables
relative to L3, of which ${\approx}\,368{,}000$ are binary
(SOS2 mode selection for piecewise linearisation) and ${\approx}\,385{,}000$ continuous; $\Lnl$ adds
${\approx}\,245{,}000$ quadratic constraints. A single
744\,h MIQCP solve requires ${>}\,4$}\DIFdelend \DIFaddbegin \DIFadd{Every comparison that involves a nonlinear or mixed-integer quadratic formulation is reported
with a rigorous optimality bound alongside the point estimate, using the incumbent of one model
and the best bound of the other so that the reported effect is conservative, together with both
solver gaps. No effect smaller than the attained tolerance is claimed; where the bound does not
support the point estimate this is stated, and treated as a result rather than a failure. The
linear levels are solved to a $0.01$}\DIFaddend \,\DIFdelbegin \DIFdel{h at 1.99}\DIFdelend \DIFaddbegin \DIFadd{\% gap. Temperature propagation is }\emph{\DIFadd{not}} \DIFadd{reported as
a solved dispatch: its free-variable form is bilinear and degenerate, so it is forward-evaluated
(Section~\ref{sec:extended}) rather than optimised. Where its diagnostic MILP is solved at all it
terminates only at a $0.1$}\DIFaddend \,\% gap\DIFdelbegin \DIFdel{; the full-year formulation is intractable on single-workstation hardware from the
University of Stuttgart. The synthetic parametric study is therefore restricted
to L1/L2/L3/$\Lplus$, while the
linearisation error is characterised on the primary case. }\DIFdelend \DIFaddbegin \DIFadd{, and that solve is retained as a diagnostic, not a reported
cost.
}\DIFaddend 

%DIF < % ============================================================
%DIF < % 4. CASE STUDIES
%DIF < % ============================================================
\DIFdelbegin %DIFDELCMD < \section{Case study}
%DIFDELCMD < \label{sec:case_studies}
%DIFDELCMD < %%%
\DIFdel{The case study comprises data from a real industrial heating network and several synthetic networks to generate generalisable results.
}\DIFdelend \DIFaddbegin \DIFadd{Not every level is solved as a dispatch problem, and the distinction is stated wherever the
levels are used. The station-resolved levels $\Lfour$ and $\Lfive$ are quantified with the
validated forward pressure module for the reasons given in
Section~\ref{subsec:stationhyd}, and temperature propagation ($\Ltwo$) likewise, for the
degeneracy reason given in Section~\ref{sec:extended}. In both cases the forward evaluation
reports the forward-evaluated cost of a fixed schedule under the finer physics; it is a
measurement of that schedule, not a bound over all schedules the finer physics would admit. The instance
count is 135 synthetic networks plus the Memmingen case, each solved through the four
decomposition configurations, together with the ladder levels solved on Memmingen.
}}
\DIFaddend 

\DIFdelbegin %DIFDELCMD < \subsection{Primary case: industrial district heating network}
%DIFDELCMD < \label{sec:case_stadtbach}
%DIFDELCMD < %%%
\DIFdelend \DIFaddbegin \subsection{Case studies and data}
\label{sec:cases}
\DIFaddend 

\DIFdelbegin \DIFdel{The primary case is an industrial DH network operated by a
medium-sized utility in southern Germany. The networkhas a radial
tree structure with a central production node~(j$_1$) and two
distribution arms (15~junctions, 14~pipe segments, 27~consumers;
}%DIFDELCMD < \Cref{fig:topology}%%%
\DIFdel{). All generation assets (CHP, gas boiler, biomass boiler, }\DIFdelend \DIFaddbegin \subsubsection{Memmingen}
\DIFadd{The real case is the Memmingen district-heating network, resolved as fifteen junction nodes
carrying twenty-seven metered consumer substations, supplied from a single plant that holds
the whole generation portfolio: combined heat and power, gas and biomass boilers, a }\DIFaddend heat
pump, \DIFdelbegin \DIFdel{and electrode boiler ) are modelled at the central production node j1. The heat pump recovers industrial waste heat via a time-varying Lorenz-cycle COP (Appendix A.1), independent of its position in the network graph. }%DIFDELCMD < 

%DIFDELCMD < %%%
\DIFdel{Supply temperature follows a piecewise-linear heating curve from
}%DIFDELCMD < \SI{70}{\celsius} %%%
\DIFdel{(low load) to }%DIFDELCMD < \SI{100}{\celsius} %%%
\DIFdel{(peak load), with corresponding return temperatures of }%DIFDELCMD < \SI{55}{\celsius} %%%
\DIFdel{to
}%DIFDELCMD < \SI{40}{\celsius} %%%
\DIFdel{(}%DIFDELCMD < \Cref{app:heating_curve}%%%
\DIFdel{). Simulated annual
thermal pipe losses are approximately 5\,\% of network heat
demand. Absolute demand values are anonymised per non-disclosure
agreement, and the key results are reported as relative differences.
}%DIFDELCMD < \Cref{tab:case_summary} %%%
\DIFdel{consolidates the key parameters.
}\DIFdelend \DIFaddbegin \DIFadd{an electrode boiler and thermal storage. Because all generation is co-located, Memmingen
is a purely central-generation case: spatial resolution can reveal loss but not a source-routing
decision. Whether that distinction matters under distributed generation is left as a controlled
open question, not a factor in the synthetic factorial, which is likewise central-generation
throughout. The metering is
billing-grade, which is itself the point of the resolution argument: it bounds what any model
of this network can be validated against. What that metering can and cannot support, and how the
thermal and hydraulic models are anchored (annual delivered energy, the mixing-valve offset, and the
independent nonlinear hydraulic cross-check), is set out in Section~\ref{subsec:validation}. A separate
limitation applies to the portfolio: the measurements predate the heat pump, electrode boiler
and storage, so the network transport physics they validate is asset-independent, but the
dispatch of those units is modelled with standard component relations and plausibility-checked
rather than validated against post-upgrade metering (Section~\ref{sec:conclusion} limitations).
}\begin{table}[t]
\centering
\caption{The Memmingen case at a glance. All generation is co-located at the plant (central
generation); the demand side resolves to 15 junction nodes, 27 metered consumer substations and
174 transmission stations.}
\label{tab:casesummary}
\begin{tabular}{lr}
\hline
Property & Value \\
\hline
Network nodes / consumer substations / stations & 15 / 27 / 174 \\
Annual heat demand (measured / simulated) & $9.9$ / $9.8$\,GWh \\
Model-implied annual network loss (defensible-$U$) & $\approx 1.2$\,GWh \\
CHP (thermal / electric) & $0.2$ / $0.1$\,MW \\
Gas boiler / biomass boiler & $13$ / $3.3$\,MW \\
Heat pump / electrode boiler & $5$ / $5$\,MW \\
Thermal storage (energy / power) & $500$\,MWh / $30$\,MW \\
\hline
\end{tabular}
\end{table}
\DIFaddend 


\DIFdelbegin \DIFdel{L2 aggregates the
15~L3 junctions into seven~zones by merging
adjacent nodes sharing a branch segment
(}%DIFDELCMD < \Cref{app:l2_zone_mapping}%%%
\DIFdel{). The total $\sum U_p L_p$ is preserved exactly (}%DIFDELCMD < \SI{1,011}{\watt\per\kelvin}%%%
\DIFdel{).
}\DIFdelend \DIFaddbegin \begin{figure}[H]
  \centering
  \includegraphics[width=\linewidth,height=0.5\textheight,keepaspectratio]{figure1}
  \caption{The Memmingen network: fifteen junctions (j$_1$--j$_{15}$) serving twenty-seven metered consumer substations (V$_1$--V$_{27}$) on a radial tree, with all generation (CHP, thermal storage, gas and biomass boilers, heat pump and electrode boiler) co-located at the production hub j$_1$. Pipe diameters are labelled DN; the trunk and the two branch arms are distinguished by colour.}
  \label{fig:network}
\end{figure}
\DIFaddend 

\DIFdelbegin %DIFDELCMD < \begin{figure}[ht]
%DIFDELCMD <   %%%
\DIFdelendFL \DIFaddbeginFL \begin{figure}[H]
  \DIFaddendFL \centering
  \DIFdelbeginFL %DIFDELCMD < \includegraphics[width=\linewidth]{paper1_dh_fidelity/Figure_2.png}
%DIFDELCMD <   \caption{Network topology. Assets at j$_1$ (CHP, TES, boilers, HP, electrode boiler). Transport delays annotated.}
%DIFDELCMD <   \label{fig:topology}
%DIFDELCMD < %%%
\DIFdelendFL \DIFaddbeginFL \includegraphics[width=\linewidth,height=0.32\textheight,keepaspectratio]{figure2}
  \caption{Normalised total heat demand of the Memmingen network over 2025 (cleaned consumption): the summed load of the twenty-seven consumer substations, shown as its daily mean and scaled to the annual peak. The seasonal heating profile is the exogenous demand the dispatch must serve.}
  \label{fig:demand}
\DIFaddendFL \end{figure}

\DIFdelbegin %DIFDELCMD < \begin{table}[ht]
%DIFDELCMD < %%%
\DIFdelendFL \DIFaddbeginFL \subsubsection{Synthetic factorial}
\DIFaddFL{The synthetic factorial is a balanced $3\times5\times3\times3$ family of radial networks that
varies node count, trunk pipe length, demand heterogeneity and storage size. Each of the 135
networks is solved through the four decomposition controls ($\CP$, $\CPL$, $\NDz$, $\Lone$)
only, all with central generation; the extended Memmingen fidelity ladder, distributed
generation and parameterised station-level hydraulics are }\emph{\DIFaddFL{not}} \DIFaddFL{applied across the
factorial. Its purpose is twofold: to test whether the loss-dominance result generalises beyond
one network, and to expose the loss burden across the order-of-magnitude range of pipe lengths
that drives the frozen-adder non-transfer. A stated capacity-sizing convention (rather than
ad-hoc feasibility filtering) sizes the assets, and its influence on the results is reported.
Full-network solves use
the same 0.01\,\% optimality gap as the real case, so the per-network decomposition terms sit
well above solver noise.
}\begin{table}[t]
\DIFaddendFL \centering
\DIFdelbeginFL %DIFDELCMD < \caption{Primary case study: consolidated parameters.}
%DIFDELCMD < \label{tab:case_summary}
%DIFDELCMD < \begin{threeparttable}
%DIFDELCMD < \scriptsize
%DIFDELCMD < \setlength{\tabcolsep}{2pt}
%DIFDELCMD < \begin{tabular*}{\columnwidth}{@{\extracolsep{\fill}} l l l @{}}
%DIFDELCMD < \toprule
%DIFDELCMD < Parameter & Value & Notes \\
%DIFDELCMD < \midrule
%DIFDELCMD < \multicolumn{3}{@{}l}{\textit{Network}} \\
%DIFDELCMD < Nodes/consumers/pipes
%DIFDELCMD <   & 15/27/14 & Radial tree \\
%DIFDELCMD < Total pipe length
%DIFDELCMD <   & \SI{3255}{\metre} & Longest: \SI{2125}{\metre} \\
%DIFDELCMD < Max.\ $\tau$
%DIFDELCMD <   & ${\approx}\SI{35}{\minute}$ & $k_p{=}1$ \\
%DIFDELCMD < \midrule
%DIFDELCMD < \multicolumn{3}{@{}l}{\textit{Thermal}} \\
%DIFDELCMD < $T_{\mathrm{sup}}$/$T_{\mathrm{ret}}$
%DIFDELCMD <   & 70--100/40--\SI{55}{\celsius}
%DIFDELCMD <   & PWL heating curve \\
%DIFDELCMD < $T_{\mathrm{ground}}$
%DIFDELCMD <   & \SI{10}{\celsius} & Fixed \\
%DIFDELCMD < $U$-value (sup/ret)
%DIFDELCMD <   & 0.28--0.34\,\si{\watt\per\metre\per\kelvin}
%DIFDELCMD <   & Per-pipe \\
%DIFDELCMD < \midrule
%DIFDELCMD < \multicolumn{3}{@{}l}{\textit{Generation at~j$_1$}} \\
%DIFDELCMD < Gas boiler/biomass
%DIFDELCMD <   & 13/\SI{3.3}{\mega\watt_{\mathrm{th}}}
%DIFDELCMD <   & $\eta{=}0.90$/$0.85$ \\
%DIFDELCMD < CHP (th/el)
%DIFDELCMD <   & 0.2/\SI{0.1}{\mega\watt}
%DIFDELCMD <   & Backpressure\tnote{a} \\
%DIFDELCMD < Thermal storage
%DIFDELCMD <   & \SI{500}{\mega\watt\hour}/\SI{30}{\mega\watt}
%DIFDELCMD <   & Cyclic BC \\
%DIFDELCMD < HP (WHR)/electrode
%DIFDELCMD <   & 5/\SI{5}{\mega\watt_{\mathrm{th}}}
%DIFDELCMD <   & Lorenz $\eta{=}0.6$/$0.95$ \\
%DIFDELCMD < \midrule
%DIFDELCMD < \multicolumn{3}{@{}l}{\textit{Demand}} \\
%DIFDELCMD < Peak/Gini/top-3
%DIFDELCMD <   & ${\approx}\SI{12}{\mega\watt}$/0.41/51.7\,\%
%DIFDELCMD <   & Moderate concentration \\
%DIFDELCMD < \bottomrule
%DIFDELCMD < \end{tabular*}
%DIFDELCMD < \begin{tablenotes}\scriptsize
%DIFDELCMD <   \item[a] %%%
\DIFdelFL{CHP supplies ${<}2\,\%$ of annual output. }%DIFDELCMD < \end{tablenotes}
%DIFDELCMD < \end{threeparttable}
%DIFDELCMD < %%%
\DIFdelendFL \DIFaddbeginFL \caption{The synthetic factorial: a balanced $3\times5\times3\times3$ grid of 135 radial
networks, each solved through the same four decomposition controls
($\CP$, $\CPL$, $\NDz$, $\Lone$) as Memmingen. All 135 networks use central generation;
distributed generation is out of scope and left to future work.}
\label{tab:synthetic}
\begin{tabular}{ll}
\hline
Factor & Levels \\
\hline
Node count & 5, 15, 30 \\
Trunk pipe length (km) & 1, 5, 15, 30, 50 \\
Demand heterogeneity index & 0.1, 0.4, 0.8 \\
Storage horizon (h) & 2, 6, 12 \\
\hline
\end{tabular}
\DIFaddendFL \end{table}


\DIFdelbegin %DIFDELCMD < \subsection{Synthetic parametric network}
%DIFDELCMD < \label{sec:case_synthetic}
%DIFDELCMD < %%%
\DIFdelend \DIFaddbegin \section{Results and discussion}
\label{sec:results}
\DIFaddend 

\DIFdelbegin \DIFdel{To generalise beyond the primary case, a synthetic tree network is parameterised by four properties (}%DIFDELCMD < \Cref{tab:synthetic}%%%
\DIFdel{)
. Of the $3^4 = 81$ parameter combinations, 36 are retained  after feasibility filtering. Configurations are classified 
as infeasible when the
LP relaxation of the L1 copperplate 
model has no solution, which occurs primarily at extreme 
parameter combinations: high demand heterogeneity 
(HI}\DIFdelend \DIFaddbegin \subsection{Validation}\label{subsec:validation}
{
\DIFadd{Table~\ref{tab:validation} reports the validation outcome against the gates fixed
}\emph{\DIFadd{ex ante}}\DIFadd{, with each metric assessed against the formulation to which it applies. That
attribution is essential when reading the table: the baseline prescribes supply temperature
from the heating curve and does not propagate it along pipes, so downstream temperature
metrics are structurally outside its scope and are marked as such rather than as failures.
}

\paragraph{\DIFadd{What the baseline is validated on.}}
\DIFadd{Three quantities test the formulation that produces every cost result in this paper, and all
three pass. Annual delivered energy is reproduced to }\textbf{\DIFadd{1.23}}\DIFaddend \,\DIFdelbegin \DIFdel{$=\,0.8$) combined with long pipe networks 
($L\,{=}\,15$}\DIFdelend \DIFaddbegin \DIFadd{\% against a \(2\,\%\)
gate. This anchors the energy scale priced by the decomposition. Network loss is the model's
computed difference between generation and delivered demand and is therefore
}\emph{\DIFadd{model-implied}}\DIFadd{, not measured: an independent value would require metered supply-side
energy that the site does not record. Delivered-energy agreement fixes the demand scale to
which that loss is referenced, rather than the loss itself.
The source return temperature matches to }\textbf{\DIFadd{2.08}}\DIFaddend \,\DIFdelbegin \DIFdel{km) and low storage capacity 
($\tau_s\,{=}\,2$}\DIFdelend \DIFaddbegin \DIFadd{°C within a \(2.5\)}\DIFaddend \,\DIFdelbegin \DIFdel{h) creates supply demand imbalances 
that no asset configuration can resolve. All 36 retained 
configurations are feasible across all five model levels 
(L1 through $\Lplus$) }\DIFdelend \DIFaddbegin \DIFadd{°C calibration
tolerance and,
more informatively, has a bias of only \(-0.31\)\,°C. The source return sensor is the only
temperature instrument in the network that does }\emph{\DIFadd{not}} \DIFadd{sit downstream of a mixing valve,
and the model is essentially unbiased against it. Heat-pump performance remains within its
expected band, with a mean coefficient of performance of }\textbf{\DIFadd{2.99}}\DIFaddend .

\DIFdelbegin %DIFDELCMD < \begin{table}[ht]
%DIFDELCMD < \centering
%DIFDELCMD < \caption{Synthetic network parameter space (36~feasible).}
%DIFDELCMD < \label{tab:synthetic}
%DIFDELCMD < \footnotesize
%DIFDELCMD < \begin{tabular*}{\columnwidth}{@{\extracolsep{\fill}} l c c c l @{}}
%DIFDELCMD < \toprule
%DIFDELCMD < Parameter & Low & Med & High & Rationale \\
%DIFDELCMD < \midrule
%DIFDELCMD < Pipe [km]    & 1   & 5   & 15  & Campus--urban \\
%DIFDELCMD < Demand HI    & 0.1 & 0.4 & 0.8 & Uniform--conc. \\
%DIFDELCMD < Node count   & 5   & 15  & 30  & Sparse--dense \\
%DIFDELCMD < Storage [h]  & 2   & 6   & 12  & Buffer--shift \\
%DIFDELCMD < \bottomrule
%DIFDELCMD < \end{tabular*}
%DIFDELCMD < \end{table}
%DIFDELCMD < 

%DIFDELCMD < %%%
\DIFdel{The parametric study thus comprises $36 \times 5$ instances
(L1/L1$_{\mathrm{topo}}$/L2/L3/$\Lplus$ per configuration)
plus 4~primary-case levels, totalling 184~optimisation  runs.
}%DIFDELCMD < 

%DIFDELCMD < %%%
\paragraph{\DIFdel{Additive gap decomposition.}}
%DIFAUXCMD
\addtocounter{paragraph}{-1}%DIFAUXCMD
\DIFdel{The parametric study includes a zero-physics topology stage
L1$_{\mathrm{topo}}$ (full network routing, zero losses) to decompose the cumulative L1$\to$L3 gap: %DIF < 
}\begin{displaymath}
\DIFdel{\label{eq:cost_decomposition}
}\begin{split}
  \DIFdel{\underbrace{\mathrm{cost}(L3) - \mathrm{cost}(L1)}_{\text{total gap}}
  \;=\;}&
  \DIFdel{\underbrace{\mathrm{cost}(L1_{\mathrm{topo}}) - \mathrm{cost}(L1)}_{\text{routing}} }\\[6pt]
  &\DIFdel{+\;
  \underbrace{\mathrm{cost}(L2) - \mathrm{cost}(L1_{\mathrm{topo}})}_{\text{heat loss}} }\\[6pt]
  &\DIFdel{+\;
  \underbrace{\mathrm{cost}(L3) - \mathrm{cost}(L2)}_{\text{pressure drop}}
}\end{split}
\DIFdel{}\end{displaymath}%DIFAUXCMD
%DIF < 
\DIFdel{To enable this three-way decomposition, the synthetic study
redefines the physics scope of L2 and L3 relative to the
primary-case taxonomy (}%DIFDELCMD < \Cref{tab:model_scope}%%%
\DIFdel{): synthetic~L2
includes steady-state thermal losses but no pressure drop, while synthetic~L3 adds Darcy--Weisbach pressure drop on top.
}%DIFDELCMD < \Cref{tab:synth_level_mapping} %%%
\DIFdel{clarifies this mapping.
}\DIFdelend \DIFaddbegin \paragraph{\DIFadd{Temperature field: what cannot be validated, and why.}}
\DIFaddend The \DIFdelbegin \DIFdel{L3$\to\Lplus$ step in the synthetic study therefore
isolates only transport-delay activation (which produces zero
cost change at hourly resolution when
$\tau_{\max} < 1$}\DIFdelend \DIFaddbegin \DIFadd{supply-temperature field belongs to the temperature-propagating formulation, and two
obstacles prevent quantitative validation here. The first is formulational: the baseline does
not propagate supply temperature, so its trunk drop is identically zero compared with a
measured \(8.2\)}\DIFaddend \,\DIFdelbegin \DIFdel{h); the pressure-drop effect itself is
captured within the L2$\to$L3 decompositionterm. For the
primary case, pressure drop is excluded from L3 and enters
exclusively at $\Lplus$, consistent with
}%DIFDELCMD < \Cref{tab:model_scope}%%%
\DIFdel{. }%DIFDELCMD < 

%DIFDELCMD < \begin{table}[ht]
%DIFDELCMD < \centering
%DIFDELCMD < \caption{Physics-scope mapping between primary case and
%DIFDELCMD < synthetic parametric study. The reassignment enables the
%DIFDELCMD < three-term additive decomposition (\Cref{eq:cost_decomposition}).}
%DIFDELCMD < \label{tab:synth_level_mapping}
%DIFDELCMD < \begin{threeparttable}
%DIFDELCMD < \footnotesize
%DIFDELCMD < \begin{tabular*}{\columnwidth}{@{\extracolsep{\fill}}
%DIFDELCMD <   >{\raggedright\arraybackslash}p{0.14\columnwidth}
%DIFDELCMD <   >{\raggedright\arraybackslash}p{0.34\columnwidth}
%DIFDELCMD <   >{\raggedright\arraybackslash}p{0.14\columnwidth}
%DIFDELCMD <   >{\raggedright\arraybackslash}p{0.20\columnwidth} @{}}
%DIFDELCMD < \toprule
%DIFDELCMD < Synth.\ label & Physics content & Primary equiv.
%DIFDELCMD <   & Decomp.\ role \\
%DIFDELCMD < \midrule
%DIFDELCMD < L1
%DIFDELCMD <   & No losses, no $\Delta p$
%DIFDELCMD <   & L1
%DIFDELCMD <   & Baseline \\
%DIFDELCMD < L1$_{\mathrm{topo}}$
%DIFDELCMD <   & Routing only, no losses
%DIFDELCMD <   & --
%DIFDELCMD <   & Routing ref. \\
%DIFDELCMD < L2
%DIFDELCMD <   & Thermal losses, no $\Delta p$
%DIFDELCMD <   & $\approx$\,L3\tnote{a}
%DIFDELCMD <   & Heat-loss ref. \\
%DIFDELCMD < L3
%DIFDELCMD <   & Thermal losses + $\Delta p$
%DIFDELCMD <   & $\approx$\,$\Lplus$\tnote{b}
%DIFDELCMD <   & $\Delta p$ target \\
%DIFDELCMD < $\Lplus$
%DIFDELCMD <   & Losses + $\Delta p$ + delay
%DIFDELCMD <   & --
%DIFDELCMD <   & Delay check \\
%DIFDELCMD < \bottomrule
%DIFDELCMD < \end{tabular*}
%DIFDELCMD < \begin{tablenotes}\scriptsize
%DIFDELCMD <   \item[a] %%%
\DIFdelFL{Equivalent in loss treatment; differs in node
    aggregation (synthetic L2 uses full node resolution, not
    the 7-zone aggregation of the primary-case L2). }%DIFDELCMD < \item[b] %%%
\DIFdelFL{Without Taylor-linearised temperature propagation;
    pressure drop enters as PWL (Darcy--Weisbach) in both.
}%DIFDELCMD < \end{tablenotes}
%DIFDELCMD < \end{threeparttable}
%DIFDELCMD < \end{table}
%DIFDELCMD < 

%DIFDELCMD < %%%
%DIF < % ============================================================
%DIF < % 5. RESULTS
%DIF < % ============================================================
%DIFDELCMD < \section{Results}
%DIFDELCMD < \label{sec:results}
%DIFDELCMD < %%%
\DIFdel{In the results section,
}\DIFdelend \DIFaddbegin \DIFadd{°C, and }\DIFaddend the \DIFdelbegin \DIFdel{validation results are presented first, followed by the topology and thermo-hydraulic influencing factors. Finally, the generalisability, sensitivity, and computational performance are examined. 
}%DIFDELCMD < 

%DIFDELCMD < %%%
%DIF < % ============================================================
%DIF < % 5.1 VALIDATION RESULTS
%DIF < % ============================================================
%DIFDELCMD < \subsection{Validation results}
%DIFDELCMD < \label{sec:results_validation}
%DIFDELCMD < 

%DIFDELCMD < \Cref{tab:validation_stage1} %%%
\DIFdel{reports the Stage~1 results; all gates defined ex ante (}%DIFDELCMD < \Cref{tab:val_targets}%%%
\DIFdel{) are met. }%DIFDELCMD < \Cref{fig:val_timeseries} %%%
\DIFdel{shows a representative winter week
(Dec. ~22--28) with dispatch composition by asset class (top panel)
and thermal storage state-of-charge (bottom panel); measured
boundary conditions are reproduced within the prescribed gates. The elevated flow MAPE (36}\DIFdelend \DIFaddbegin \DIFadd{far-end and corridor gates do not apply to it. The second is
instrumental: for the propagating formulation, the far-end error is \(9.21\)}\DIFaddend \,\DIFdelbegin \DIFdel{\%) is attributable to
pressure-differential metering at partial load; the 1.23}\DIFdelend \DIFaddbegin \DIFadd{°C, with a bias
of the same magnitude. Mean absolute error and bias coincide to four decimal places, so every
residual has the same sign, and the magnitude lies squarely within the \(5\) to \(15\)}\DIFaddend \,\DIFdelbegin \DIFdel{\%
annual energy errorconfirms that energy-weighted flow is correctly
reproduced}\DIFdelend \DIFaddbegin \DIFadd{°C
offset associated with a three-way mixing valve. The corridor comparison shows the same
signature, with a bias of \(-7.73\)\,°C. This is }\emph{\DIFadd{consistent with}} \DIFadd{a fixed sensor-location
offset rather than model error. However, the primary and secondary temperatures were not
measured simultaneously, so the available instrumentation cannot distinguish sensor-location
bias from model error, and the temperature field is therefore not independently validated.
Neither obstacle can be removed through a better pipe model, and the first is a deliberate
property of the formulation used for the reported cost results}\DIFaddend .

\DIFdelbegin %DIFDELCMD < \begin{figure}[!htbp]
%DIFDELCMD <   \centering
%DIFDELCMD <   \includegraphics[width=\linewidth]{paper1_dh_fidelity/Figure_3.png}
%DIFDELCMD <   \caption{Stage~1 validation, representative winter week
%DIFDELCMD <   (Dec.~22--28): dispatch composition (top) and TES
%DIFDELCMD <   state-of-charge (bottom).}
%DIFDELCMD <   \label{fig:val_timeseries}
%DIFDELCMD < \end{figure}
%DIFDELCMD < %%%
\DIFdelend \DIFaddbegin \DIFadd{This is the concrete form of an argument that runs throughout the paper. A model resolved more
finely than its metering cannot have that additional resolution verified, and a model that
does not represent a quantity cannot be tested on it. The conclusions are therefore anchored
to annual delivered energy, which the measurements resolve, and hence to the model-implied
network loss on which both formulations agree, rather than to a point temperature field that
neither the instrumentation nor the baseline can support.
}\DIFaddend 

\DIFdelbegin \DIFdel{Node j$_1$ is the production source and serves as boundary
condition (measured inlet temperature); itis excluded from the
temperature-error evaluation. Of the 14~remaining network nodes, 6 enter the gate evaluation (``validation nodes'', cf.
\ }%DIFDELCMD < \Cref{tab:val_targets}%%%
\DIFdel{); the remaining 8 are excluded for two
distinct reasons (per-node MAE breakdown in
}%DIFDELCMD < \Cref{app:per_node_mae}%%%
\DIFdel{). Four nodes (j$_6$, j$_7$, j$_8$, j$_{14}$) served as }\emph{\DIFdel{calibration nodes}} %DIFAUXCMD
\DIFdel{during the branch-sequential $U_p$-tuning procedure and are therefore
in-sample residuals; their group mean is }%DIFDELCMD < \SI{2.32}{\celsius}%%%
\DIFdel{, and
node j$_7$ at }%DIFDELCMD < \SI{3.07}{\celsius} %%%
\DIFdel{would exceed }\DIFdelend \DIFaddbegin \paragraph{\DIFadd{Scope of the temperature evaluation.}}
\DIFadd{Two views of the temperature error are consistent. On the ex-ante validation set, comprising
the held-out nodes over the 744-hour winter window for which the gates in
Table~\ref{tab:valtargets} were fixed, the mean supply-temperature error is \(1.32\)\,°C and
}\DIFaddend the worst-node \DIFdelbegin \DIFdel{gate
of }%DIFDELCMD < \SI{2.5}{\celsius} %%%
\DIFdel{if held out for validation. This in-sample
deterioration is expected and
motivates the strict
calibration/validation split adopted here. The remaining four nodes
(j$_2$, j$_3$, j$_4$,
j$_5$) are excluded for measurement reasons
(mixing-valve offset above }%DIFDELCMD < \SI{2}{\celsius} %%%
\DIFdel{or absence of a single
junction reference among multiple downstream consumers); two of
them (j$_3$: }%DIFDELCMD < \SI{0.58}{\celsius}%%%
\DIFdel{; j$_4$:
}%DIFDELCMD < \SI{0.37}{\celsius}%%%
\DIFdel{) are
in fact the two best-performing nodes in the network}\DIFdelend \DIFaddbegin \DIFadd{error is \(2.27\)\,°C, both }\emph{\DIFadd{within}} \DIFadd{their respective \(1.5\)\,K and
\(2.5\)\,K gates. The \(9.21\)\,°C figure in Table~\ref{tab:validation} instead comes from a
separate, more conservative diagnostic based on the far-end node over the full annual record,
without selecting a favourable window. The far end operates the mixing valve closest to fully
open. Appendix~\ref{app:per_node_mae} confirms that the earlier selection was not favourable:
the all-node mean of \(1.68\)\,°C lies between the validation-node }\DIFaddend and \DIFdelbegin \DIFdel{would pass
either gate comfortably. Mean MAE across
all 14~nodes is }%DIFDELCMD < \SI{1.68}{\celsius}%%%
\DIFdel{.
}\DIFdelend \DIFaddbegin \DIFadd{calibration-node group
means, and two excluded nodes are the network's best-performing.
}\DIFaddend 

\DIFdelbegin %DIFDELCMD < \begin{table}[!htbp]
%DIFDELCMD < \centering
%DIFDELCMD < \caption{Stage~1 validation results.}
%DIFDELCMD < \label{tab:validation_stage1}
%DIFDELCMD < \begin{threeparttable}
%DIFDELCMD < \footnotesize
%DIFDELCMD < \begin{tabular*}{\columnwidth}{@{\extracolsep{\fill}} l c c c @{}}
%DIFDELCMD < \toprule
%DIFDELCMD < KPI & Model & Result & Gate \\
%DIFDELCMD < \midrule
%DIFDELCMD < $T_{\mathrm{sup}}$, mean\tnote{a}
%DIFDELCMD <   & $\Lnl$ & \SI{1.32}{\celsius} & ${<}\SI{1.5}{\celsius}$ \\
%DIFDELCMD < $T_{\mathrm{sup}}$, worst\tnote{a}
%DIFDELCMD <   & $\Lnl$ & \SI{2.27}{\celsius} & ${<}\SI{2.5}{\celsius}$ \\
%DIFDELCMD < $T_{\mathrm{ret}}$ at source
%DIFDELCMD <   & L3     & \SI{2.08}{\celsius} & ${<}\SI{2.5}{\celsius}$ \\
%DIFDELCMD < Flow (MAPE)
%DIFDELCMD <   & L3     & $36\,\%$            & n/a \\
%DIFDELCMD < Annual energy
%DIFDELCMD <   & L3     & $1.23\,\%$          & ${<}2\,\%$ \\
%DIFDELCMD < \bottomrule
%DIFDELCMD < \end{tabular*}
%DIFDELCMD < \begin{tablenotes}\footnotesize
%DIFDELCMD <   \item[a] %%%
\DIFdelFL{744}\DIFdelendFL \DIFaddbeginFL \paragraph{\DIFaddFL{Flow, and the per-step balance.}}
\DIFaddFL{The source flow has a mean absolute percentage error of }\textbf{\DIFaddFL{33.8}}\DIFaddendFL \,\DIFdelbeginFL \DIFdelFL{h winter window; 6~of 14~nodes used as
    validation nodes (4~calibration, 4~excluded for measurement
    reasons; see }%DIFDELCMD < \Cref{app:per_node_mae}%%%
\DIFdelFL{). }%DIFDELCMD < \end{tablenotes}
%DIFDELCMD < \end{threeparttable}
%DIFDELCMD < \end{table}
%DIFDELCMD < 

%DIFDELCMD < %%%
\DIFdel{Stage~2 plausibility checks all pass: observed HP COP within
$[2.5,\,3.0]$ (bounded by the moderate waste-heat source
temperature of 25--}%DIFDELCMD < \SI{35}{\celsius} %%%
\DIFdel{and the high sink temperature
of 70--}%DIFDELCMD < \SI{100}{\celsius}%%%
\DIFdel{); electrode-boiler price-response
correlation $r = 0.47$; }\DIFdelend \DIFaddbegin \DIFadd{\%, and disaggregation
by load band explains rather than excuses it. The absolute error is nearly constant across
bands at roughly \(12\) to \(13\,\mathrm{m^3\,h^{-1}}\), so the percentage is governed by
the denominator: the same physical error is \(46\,\%\) of a
\(30.9\,\mathrm{m^3\,h^{-1}}\) mean flow below one-quarter of peak load and \(13\,\%\) of a
\(102.6\,\mathrm{m^3\,h^{-1}}\) mean flow between one-half and three-quarters of peak load.
Low-load hours account for \(60\,\%\) of the year but carry little energy, which is why a
large relative error in those hours is compatible with an annual energy match of
\(1.23\,\%\). The top band is the exception in both directions, with the largest absolute
error of \(25.5\,\mathrm{m^3\,h^{-1}}\) and a relative error of \(23\,\%\), but comprises
only 56 hours; it is reported rather than absorbed into an average. The }\DIFaddend per-step energy
\DIFdelbegin \DIFdel{closure below 0.3\,\%; TES
SOC in $[0.18,\,0.94] \cdot \bar{E}_s$. Solver consistency:
$\Lnl$ with quadratic constraints deactivated reproduces L3 within $0.004\,\%$}\DIFdelend \DIFaddbegin \DIFadd{balance, comprising generation plus net storage against delivered demand and network loss,
closes to machine precision, with an annual residual below \(10^{-3}\,\mathrm{MWh}\) and
every hour within tolerance. This confirms that the model's internal conservation holds
exactly. The annual demand-energy agreement of \(1.23\,\%\) is the separate measured
comparison on which the cost conclusions rest}\DIFaddend .

\DIFdelbegin \paragraph{\DIFdel{Independent forward-physics cross-check.}}
%DIFAUXCMD
\addtocounter{paragraph}{-1}%DIFAUXCMD
\DIFdel{As complementary validation, a boundary-condition method (BCM)
propagates the
measured source temperature at j$_1$ through the steady-state pipe-loss model without dispatch optimisation, isolating pipe physics from the optimiser. At }\DIFdelend \DIFaddbegin \DIFadd{A percentage error can conceal either of two different failures: a model that mistakes the
level or one that fails to follow the variation. Comparing first differences separates them
because a fixed offset cancels. The model tracks the flow level with a correlation of
}\textbf{\DIFadd{0.91}} \DIFadd{and its day-to-day change with }\textbf{\DIFadd{0.80}}\DIFadd{; for demand, }\DIFaddend the \DIFdelbegin \DIFdel{far-end node
j$_{15}$ (path }%DIFDELCMD < \SI{2125}{\metre}%%%
\DIFdel{), a global trunk }\DIFdelend \DIFaddbegin \DIFadd{corresponding
values are }\textbf{\DIFadd{0.93}} \DIFadd{and }\textbf{\DIFadd{0.89}}\DIFadd{. At hourly resolution, the correlation of
differences falls to \(0.24\), which reflects the metering rather than the model: the record
is sampled at fifteen-minute intervals and resampled to hourly resolution, so hour-to-hour
differences are dominated by sampling noise. The aggregate flow error is therefore a
level-and-denominator effect at low load, not a failure to reproduce how the network varies.
That variation is the property on which the dispatch results depend, since cost accumulates
over days and seasons rather than individual hours.
}

\paragraph{\DIFadd{Post-upgrade assets and structural checks.}}
\DIFadd{The heat pump, electrode boiler, and storage were installed after the measurement record, so
their dispatch is verified against necessary conditions rather than measurements: the
coefficient of performance remains within its plausible band, the storage state of charge
remains within its operating limits, and simultaneous charging and discharging are
prohibited. Two structural checks close the section. The nonlinear reference with its
quadratic constraints deactivated reproduces the baseline objective to better than
\(0.01\,\%\). On the economic-cost basis, the ordering
\(\mathrm{cost}(\CP) \leq \mathrm{cost}(\Lone)\) also holds for every one of the 135 synthetic
instances, the expected direction for a loss-free relaxation, though not an inequality
guaranteed in general. A departure from either check on this test set would indicate a
formulation error rather than a finding.}}
\begin{table*}[t]
\centering
\caption{Validation against the Memmingen measurements. \emph{Model} is the formulation each
metric is scored against; bias is the signed mean error and scatter the RMSE, both in the unit of
the quantity (K for temperature differences). Two rows are gated outputs of the MILP formulation
that produces the cost results, namely the annual delivered-energy error and the per-step
energy-balance closure; the COP band is a consistency check, the return temperature a
parameter-calibration check. The supply-temperature gates of Table~\ref{tab:valtargets} are
scored against the propagating (NLP) formulation on the held-out validation nodes over the
744-hour window on which they were fixed; both the mean ($1.32$\,K) and worst-node ($2.27$\,K)
errors pass. The far-end and corridor rows are separate, more conservative \emph{annual}
diagnostics (\emph{diag.}), not the gated metrics: the $9.21$\,K far-end value is a near-pure
bias (mean absolute error $\approx$ bias) \emph{consistent with} a fixed metering offset at the
three-way mixing valve, though the instrumentation cannot rule out model error
(Section~\ref{subsec:validation}). The trunk-drop row is \emph{n/a}, because the baseline that
produces the cost results does not propagate temperature and so does not produce that quantity. The flow and demand MAPE are denominator-driven at low load, where a
near-constant absolute error is a large share of small flows; representative bands are shown
inline, and the annual delivered-energy match above (1.23\,\%) is the measure the cost
conclusions rest on.}
\label{tab:validation}
\footnotesize
\setlength{\tabcolsep}{4pt}
\begin{tabular}{@{}p{4.9cm}lrrrlc@{}}
\hline
Quantity & Model & Value & Bias & Scatter & Gate & Status \\
\hline
Annual delivered-energy error & MILP & 1.23\,\% & -- & -- & $\le$2\,\% & \checkmark \\
Heat-pump COP (mean) & MILP & 2.99 & -0.00 & -- & $[2.5,5.5]$ & \checkmark \\
Source return-temp calibration (const.\ param.) & param. & 2.08\,K & $-0.31$\,K & 2.87\,K & -- & report \\
Energy-balance closure (annual residual) & MILP & $<10^{-3}$\,MWh ($\approx$0\,\%) & -- & -- & $\le$2\,\% & \checkmark \\
Flow MAPE, source & MILP & \makecell[r]{33.8\,\%\\[1pt]{\scriptsize 46\,\% low-load, 13\,\% mid--high}} & -- & -- & -- & report \\
Demand MAPE, total & MILP & \makecell[r]{29.1\,\%\\[1pt]{\scriptsize 39\,\% low-load, 12\,\% mid--high}} & -- & -- & -- & report \\
Flow tracking ($r_\text{lvl}$; $r_{\Delta,\text{daily}}$) & MILP & 0.91 / 0.80 & -- & h.\ $\Delta$ 0.24 (noise) & -- & \checkmark \\
Supply-temp MAE, mean (val.\ nodes, 744\,h) & NLP & 1.32\,K & -- & -- & $\le$1.5\,K & \checkmark \\
Supply-temp MAE, worst node (744\,h) & NLP & 2.27\,K & -- & -- & $\le$2.5\,K & \checkmark \\
Far-end supply-temp MAE (annual diagnostic) & NLP & 9.21\,K & $+9.21$\,K & 11.31\,K & -- & diag. \\
Trunk supply-temp drop MAE (diagnostic) & MILP & -- (not propagated) & -- & -- & -- & n/a \\
Corridor supply-temp MAE (annual diagnostic) & NLP & 8.32\,K & $-7.73$\,K & -- & -- & diag. \\
\hline
\end{tabular}
\end{table*}
\begin{figure}[H]
  \centering
  \includegraphics[width=\linewidth,height=0.42\textheight,keepaspectratio]{figure3}
  \caption{The mixing-valve offset that bounds what the temperature field can be validated against. Consumer sensors are billing instruments installed downstream of the valve, so the metered temperature sits systematically below the primary junction temperature.}
  \label{fig:mixingvalve}
\end{figure}
\DIFadd{Pipe losses admit two calibration approaches that require reconciliation. The first
uses a global boundary-condition match and applies a single scaling factor to all
heat-transfer coefficients. The second uses a branch-sequential fit and increases
the coefficients of the }\emph{\DIFadd{terminal}} \DIFadd{pipes to several times their nominal values
until the trunk-only loss matches the total loss implied by delivered energy.
}

\DIFadd{The second approach is not physically defensible. Such increases in nominal
}\DIFaddend $U$\DIFdelbegin \DIFdel{-multiplier
(${\times}1.330$) is calibrated on the Oct--Feb winter subset
across two partial heating seasons ($n = 3173$\,h).
The resulting
error metrics are MAE\,$=$\,}%DIFDELCMD < \SI{1.56}{\celsius}%%%
\DIFdel{, RMSE\,$=$\,}%DIFDELCMD < \SI{2.15}{\celsius}%%%
\DIFdel{, and a residual positive bias of }%DIFDELCMD < \SI{+0.72}{\celsius}%%%
\DIFdel{, indicating mild systematic overestimation
consistent with unmodelled flow-rate variability. Sixty percent of hourly predictions fall within ${\pm}$}%DIFDELCMD < \SI{1.5}{\celsius}%%%
\DIFdel{; the remaining deviations correlate with rapid flow transients
(partial-load switching) that the steady-state pipe }\DIFdelend \DIFaddbegin \DIFadd{-values cannot reasonably be attributed to insulation ageing. Instead, they
absorb the last-mile losses from service laterals and substations that the 15-node
trunk }\DIFaddend model does not \DIFdelbegin \DIFdel{capture. The mean MAE (}%DIFDELCMD < \SI{1.56}{\celsius}%%%
\DIFdel{) nevertheless remains
within the Stage~1 gate.
}%DIFDELCMD < \Cref{fig:val_scatter} %%%
\DIFdel{presents the scatter plot of simulated versus measured supply temperature;
}%DIFDELCMD < \Cref{fig:val_timeseries_jan} %%%
\DIFdel{shows the
corresponding January
time-series with the ${\pm}$}%DIFDELCMD < \SI{1.5}{\celsius} %%%
\DIFdel{confidence band. The calibration approach differs methodologically from the $\Lnl$
MIQCP-output validation reported in }%DIFDELCMD < \Cref{tab:validation_stage1}%%%
\DIFdel{;
details are given in
}%DIFDELCMD < \Cref{app:bcm_calibration}%%%
\DIFdel{.
}%DIFDELCMD < 

%DIFDELCMD < \begin{figure}[!htbp]
%DIFDELCMD <   \centering
%DIFDELCMD <   \includegraphics[width=\linewidth]{paper1_dh_fidelity/Figure_4.png}
%DIFDELCMD <   \caption{BCM cross-check: simulated vs.\ measured
%DIFDELCMD <   $T_{\mathrm{sup}}$ at j$_{15}$. Dotted lines mark the
%DIFDELCMD <   ${\pm}$\SI{1.5}{\celsius} gate.}
%DIFDELCMD <   \label{fig:val_scatter}
%DIFDELCMD < \end{figure}
%DIFDELCMD < %%%
\DIFdelend \DIFaddbegin \DIFadd{represent. Because the model contains no explicit representation
of these components, the calibration assigns their losses artificially to the trunk
pipes.
}\DIFaddend 

\DIFdelbegin %DIFDELCMD < \begin{figure}[!htbp]
%DIFDELCMD <   \centering
%DIFDELCMD <   \includegraphics[width=\linewidth]{paper1_dh_fidelity/Figure_5.png}
%DIFDELCMD <   \caption{BCM cross-check (January): $T_{\mathrm{sup}}$
%DIFDELCMD <   time-series at j$_{15}$. The shaded band marks the
%DIFDELCMD <   ${\pm}$\SI{1.5}{\celsius} gate. Underneath is the ambient
%DIFDELCMD <   temperature for the week.}
%DIFDELCMD <   \label{fig:val_timeseries_jan}
%DIFDELCMD < \end{figure}
%DIFDELCMD < %%%
\DIFdelend \DIFaddbegin \DIFadd{This study therefore retains trunk heat-transfer coefficients within a defensible
range and represents the residual loss through an explicit service-lateral term.
This term enters only at the station-resolved level (L4). Consequently, the
node-resolved levels intentionally undercount total network losses because they
exclude the last mile. The model matches the total loss implied by delivered energy
only after the service laterals are included. This finding is reported in
Section~\ref{subsec:linearisation}.
}\DIFaddend 

%DIF < % ============================================================
%DIF < % 5.2 TOPOLOGY EFFECT
%DIF < % ============================================================
\DIFdelbegin %DIFDELCMD < \subsection{Topology effect: L1 vs.\ L2 vs.\ L3 (RQ1)}
%DIFDELCMD < \label{sec:results_cost}
%DIFDELCMD < %%%
\DIFdelend \DIFaddbegin \DIFadd{Importantly, the revised calibration leaves the decomposition unchanged. The
dominance of losses is therefore not an artefact of the global multiplier. This
study reports in-sample and out-of-sample metrics separately because an in-sample
fit alone does not constitute validation. Flow accuracy is reported as mean
absolute percentage error by load band, making the fit at part load, where most
operating hours occur, visible rather than obscuring it within an annual aggregate.
}\DIFaddend 

\DIFdelbegin %DIFDELCMD < \Cref{tab:cost_co2} %%%
\DIFdel{reports the topology comparison under identical
basic physics. The
cumulative L1$\to$L3 gap reaches 13.0\,\% in
operational cost and 74\,t\,}%DIFDELCMD < \ce{CO2}%%%
\DIFdel{/year. The L1$\to$L2 step
accounts for 10.4}\DIFdelend \DIFaddbegin \subsection{Estimation bias: loss visibility versus spatial resolution}\label{subsec:decomp}
\DIFadd{On Memmingen the copperplate-to-baseline cost gap of }\textbf{\DIFadd{15.1}}\DIFaddend \,\%
(\DIFdelbegin \DIFdel{63\,t\,}%DIFDELCMD < \ce{CO2}%%%
\DIFdel{/yr), while
L2}\DIFdelend \DIFaddbegin \CP\DIFadd{\ }\DIFaddend $\to$ \DIFdelbegin \DIFdel{L3 contributes the remaining 2.6}\DIFdelend \DIFaddbegin \Lone\DIFadd{) decomposes to machine precision into a loss main effect of
}\textbf{\DIFadd{95.8}}\DIFadd{\,\% of the gap and an interaction of }\textbf{\DIFadd{-0.5}}\DIFaddend \,\% (\DIFdelbegin \DIFdel{11}\DIFdelend \DIFaddbegin \DIFadd{defensible trunk $U$; the single
indefensible terminal-pipe multiplier removed). Loss visibility carries the
result: a copperplate reproduces the cost level once handed the aggregate loss it
cannot itself compute (}\CPL\DIFadd{), whereas resolving spatial topology without loss
(}\NDz\DIFadd{) recovers almost none of the gap. The as-run $\NDz-\CP$ difference of
$\approx$961}\DIFaddend \,\DIFdelbegin \DIFdel{t}\DIFdelend \DIFaddbegin \DIFadd{EUR (4.7}\DIFaddend \,\DIFdelbegin %DIFDELCMD < \ce{CO2}%%%
\DIFdel{/yr). }%DIFDELCMD < 

%DIFDELCMD < %%%
\DIFdel{The L1$\to$L2 gap arises from two mechanisms: (1)~L1 ignores
network losses entirely, the additive decomposition (}%DIFDELCMD < \Cref{sec:case_synthetic}%%%
\DIFdel{) confirms that thermal losses, not routing or locational siting, drive the
cumulative gap.
}%DIFDELCMD < 

%DIFDELCMD < \begin{table}[!htbp]
%DIFDELCMD < \centering
%DIFDELCMD < \caption{Topology comparison (basic physics, L3 = reference).}
%DIFDELCMD < \label{tab:cost_co2}
%DIFDELCMD < \begin{threeparttable}
%DIFDELCMD < \footnotesize
%DIFDELCMD < \begin{tabular*}{\columnwidth}{@{\extracolsep{\fill}} l c c c c c @{}}
%DIFDELCMD < \toprule
%DIFDELCMD < Level & $\Delta$Cost & $\Delta$Fuel &
%DIFDELCMD <   $\Delta$Elec. & $\Delta$\ce{CO2} & t/yr \\
%DIFDELCMD < \midrule
%DIFDELCMD < L1 & $-13.0\%$ & $-12.9\%$ & $-14.0\%$ & $-8.7\%$ & $772$ \\
%DIFDELCMD < L2 & $-2.6\%$  & $-2.4\%$  & $-3.2\%$  & $-1.3\%$ & $835$ \\
%DIFDELCMD < L3 & $(ref.)$  & $(ref.)$  & $(ref.)$  & $(ref.)$ & $846$ \\
%DIFDELCMD < \bottomrule
%DIFDELCMD < \end{tabular*}
%DIFDELCMD < \end{threeparttable}
%DIFDELCMD < \end{table}
%DIFDELCMD < 

%DIFDELCMD < %%%
%DIF < % ============================================================
%DIF < % 5.3 PHYSICS + LINEARISATION
%DIF < % ============================================================
%DIFDELCMD < \subsection{thermo-hydraulic effect and linearisation error (RQ2 + RQ3)}
%DIFDELCMD < \label{sec:results_extended}
%DIFDELCMD < 

%DIFDELCMD < %%%
\paragraph{\DIFdel{Extended physics: L3\,$\to$\,$\Lplus$.}}
%DIFAUXCMD
\addtocounter{paragraph}{-1}%DIFAUXCMD
\DIFdel{Adding linearised pressure drop and temperature propagation
increases annual cost by $+0.11\,\%$, attributable entirely to
pumping power (+255}\DIFdelend \DIFaddbegin \DIFadd{\% of the gap) is, however, }\emph{\DIFadd{not}} \DIFadd{spatial resolution alone: the
copperplate and node-resolved configurations differ also in three physics parameters
(return temperature, heating curve, and heat-pump coefficient of performance). Re-solving }\NDz\DIFadd{\
under the copperplate's physics isolates the topology main effect proper at only
$\approx$}\textbf{\DIFadd{52}}\DIFadd{\,EUR (}\textbf{\DIFadd{0.25}}\DIFadd{\,\% of the gap), with the remaining $\approx$909}\DIFaddend \,EUR
\DIFdelbegin \DIFdel{/yr on 2.4}\DIFdelend \DIFaddbegin \DIFadd{(4.4}\DIFaddend \,\DIFdelbegin \DIFdel{MWh additional pump
electricity) . The thermo-hydraulic gap }\DIFdelend \DIFaddbegin \DIFadd{\%) a physics-parameter difference between the two configurations, reported on its own line
in Table~\ref{tab:decomposition} so that it is not attributed to spatial resolution. Topology
proper }\DIFaddend is thus two orders of magnitude \DIFdelbegin \DIFdel{smaller than the topology gap. This answers RQ2 with a
clear negative: extended hydraulic and thermal physics do not
materially change dispatch when topology is already resolved at the scale of }\DIFdelend \DIFaddbegin \DIFadd{below the loss effect: spatial resolution matters as the
}\emph{\DIFadd{mechanism}} \DIFadd{that makes loss endogenous, not through routing. Across the 135
synthetic networks, which are generated with identical physics across the copperplate and
node-resolved variants and so carry no such confound, the same holds emphatically: loss explains
a median }\textbf{\DIFadd{100.0}}\DIFadd{\,\% of the gap and }\DIFaddend the \DIFdelbegin \DIFdel{primary case.
}%DIFDELCMD < 

%DIFDELCMD < \begin{table}[!htbp]
%DIFDELCMD < \centering
%DIFDELCMD < \caption{Thermo-hydraulics and linearisation comparison.}
%DIFDELCMD < \label{tab:cost_extended}
%DIFDELCMD < \footnotesize
%DIFDELCMD < \begin{tabular*}{\columnwidth}{@{\extracolsep{\fill}} l c c c @{}}
%DIFDELCMD < \toprule
%DIFDELCMD < Comparison & $\Delta$Cost & Window & MIP gap \\
%DIFDELCMD < \midrule
%DIFDELCMD < L3$\to\Lplus$
%DIFDELCMD <   & $+0.11\%$ & $Full year$ & ${<}\,0.5\%$ \\
%DIFDELCMD < $\Lplus\to\Lnl$, $Jan.\ 2025$
%DIFDELCMD <   & $+0.35\%$ & $744\,h$    & $1.99\%$ \\
%DIFDELCMD < $\Lplus\to\Lnl$, $Feb.\ 2025$
%DIFDELCMD <   & $+0.50\%$ & $672\,h$    & $0.46\%$ \\
%DIFDELCMD < \bottomrule
%DIFDELCMD < \end{tabular*}
%DIFDELCMD < \end{table}
%DIFDELCMD < 

%DIFDELCMD < %%%
\paragraph{\DIFdel{Linearisation error bounds: $\Lplus$\,$\to$\,$\Lnl$.}}
%DIFAUXCMD
\addtocounter{paragraph}{-1}%DIFAUXCMD
%DIFDELCMD < \Cref{tab:cost_extended} %%%
\DIFdel{reports the combined $\Lplus\to\Lnl$ gap on two independent winter windows: $+0.35\,\%$ (January, 744}\DIFdelend \DIFaddbegin \DIFadd{spatial-topology main effect stays
within $\pm$0.6}\DIFaddend \,\DIFdelbegin \DIFdel{h)and $+0.50\,\%$ (February, 672}\DIFdelend \DIFaddbegin \DIFadd{\% on every network of 5}\DIFaddend \,\DIFdelbegin \DIFdel{h) . Three observations bound the
interpretation:
%DIF < 
}%DIFDELCMD < \begin{enumerate}[nosep]
%DIFDELCMD <   \item %%%
\DIFdel{The gap combines PWL-vs.-quadratic error
        }\emph{\DIFdel{and}} %DIFAUXCMD
\DIFdel{transport-delay activation; the two cannot be
        separated without an intermediate MIQCP run with delay
        disabled.
  }%DIFDELCMD < \item %%%
\DIFdel{The February window has a tight MIP gap (0.46}\DIFdelend \DIFaddbegin \DIFadd{km trunk length or more, never exceeding 2.4\,\% even on the shortest, where the entire gap is below 6\,\% of cost (all solves at
$\le0.1$\,\% optimality gap), over node counts, pipe lengths, demand heterogeneity and
storage. The corrected Memmingen topology (0.25\,\%) sits comfortably inside this synthetic range.
The $U$-calibration choice leaves the loss share essentially unchanged (inflated-$U$: 96\,\%;
defensible-$U$: }\textbf{\DIFadd{95.8}}\DIFaddend \,\%), \DIFdelbegin \DIFdel{giving high confidence that the true linearisation effect is close to $+0.5\,\%$. The January window has a wider
        gap
(1.99}\DIFdelend \DIFaddbegin \DIFadd{so the conclusion is not an
artefact of the loss calibration. Supplying the copperplate control with the hourly,
heating-curve-consistent loss profile that the resolved model realises, in place of the flat
annual adder, likewise moves the loss main effect only from $95.8$ to $95.5$}\DIFaddend \,\DIFdelbegin \DIFdel{\%), so its $+0.35\,\%$ is less precisely
        determined.
  }%DIFDELCMD < \item %%%
\DIFdel{A spring window (May 2025) was attempted but produced no
        feasible MIQCP incumbent within the time limit because near-zero
        flow in shoulder months causes degenerate bilinear
        constraints. The bound is therefore validated for winter conditions (peak-flow regime). }%DIFDELCMD < \end{enumerate}
%DIFDELCMD < %%%
%DIF < 
\DIFdel{The January gap ($+0.35\,\%$)falls below solver tolerance
(MIP gap 1.99}\DIFdelend \DIFaddbegin \DIFadd{\% of the gap
(a $0.4$-percentage-point shift, and a $0.06$}\DIFaddend \,\% \DIFdelbegin \DIFdel{) and is therefore not independently
significant; it does, however, confirm that $\Lnl$ cannot
materially improve on the $\Lplus$ solution even when granted
substantially more optimality slack. The February gap
($+0.50\,\%$, MIP gap 0.46\,\%) is the tighter and statistically meaningful bound. Together, both windows support
the conclusion that the combined linearisation and transport-delay effect is ${\leq}\,0.5\,\%$ under winter
operating conditions. The effect is supposedly negligible for planning-accuracy dispatch. }\DIFdelend \DIFaddbegin \DIFadd{change in the $\CPL$ cost); the constant
adder is retained as the primary oracle because the decomposition is insensitive to the
adder's temporal shape.
}\begin{table}[t]
\centering
\caption{Loss$\times$topology decomposition of the copperplate-to-baseline cost gap on Memmingen
(economic cost; the Memmingen solves reach $\le$0.01\,\% gap, the synthetic factorial $\le$0.1\,\%). Main effects use baseline--treatment coding
against $\CP$: loss $=\text{cost}(\CPL)-\text{cost}(\CP)$; the topology main effect is measured
with physics held identical to $\CP$, $\text{cost}(\NDz^{\ast})-\text{cost}(\CP)$, where
$\NDz^{\ast}$ is the node-resolved no-loss network re-solved under $\CP$'s return temperature,
heating curve and heat-pump coefficient of performance. The as-run $\NDz$ differs from $\CP$ in
those three parameters as well as in topology; the resulting
$\text{cost}(\NDz)-\text{cost}(\NDz^{\ast})=$909\,EUR is a physics-parameter effect carried by the
configuration rather than by spatial resolution, and is listed separately so it is not attributed
to topology. Interaction is the residual and the closure is an exact \emph{algebraic identity} on
the optimised costs; a symmetric (Shapley) split of the small interaction changes the shares
negligibly. The synthetic column, whose copperplate and node-resolved variants share physics and
so carry no such confound, reports the 135-network distribution.}
\label{tab:decomposition}
\begin{tabular}{lrrr}
\hline
Term & EUR/yr & \% of gap & Synthetic (135 nets) \\
\hline
Total gap ($\CP\!\to\!\Lone$) & 20,591 & 100.0 & burden 3.2--81.2\,\% of cost \\
Loss main effect & 19,737 & 95.8 & median 100.0\,\% \\
Topology main effect (physics-matched) & 52 & 0.25 & $\pm$0.6\,\% (L$\geq$5\,km); $\leq$2.4\,\% at 1\,km \\
Physics-parameter difference (config) & 909 & 4.4 & --- \\
Interaction & -107 & -0.5 & --- \\
\hline
Closure residual & 0.00e+00 & 0.0e+00 & (exact identity) \\
\hline
\end{tabular}
\end{table}
\DIFaddend 

%DIF < % ============================================================
%DIF < % 5.4 SYNTHETIC PARAMETRIC STUDY
%DIF < % ============================================================
\DIFdelbegin %DIFDELCMD < \subsection{Generalisability: synthetic parametric study (RQ4)}
%DIFDELCMD < \label{sec:results_generalizability}
%DIFDELCMD < 

%DIFDELCMD < %%%
\DIFdel{The synthetic parametric study covers 36~feasible configurations
(}%DIFDELCMD < \Cref{tab:synthetic}%%%
\DIFdel{) , with the full
L1/L1$_{\mathrm{topo}}$/L2/L3/$\Lplus$ hierarchy solved for each
(180~runs + 4~primary-case = 184 total). Three findings emerge.
}\DIFdelend \DIFaddbegin \begin{figure}[H]
  \centering
  \includegraphics[width=\linewidth,height=0.42\textheight,keepaspectratio]{figure4}
  \caption{Exact decomposition of the copperplate-to-baseline cost gap for Memmingen (left) and as a distribution across the 135 synthetic networks (right). On Memmingen the topology main effect is measured with physics held identical to the copperplate; the physics-parameter difference between the copperplate and node-resolved configurations (return temperature, heating curve, heat-pump coefficient of performance) is shown separately so it is not attributed to spatial resolution. The additive identity closes to machine precision, so the terms are measured rather than fitted. The synthetic networks share physics across variants by construction and so carry no such term.}
  \label{fig:decomp}
\end{figure}
\DIFaddend 

\DIFdelbegin \paragraph{\DIFdel{Pipe-length scaling.}}
%DIFAUXCMD
\addtocounter{paragraph}{-1}%DIFAUXCMD
\DIFdel{The cumulative L1$\to$L3 gap is strictly positive in 36/36
configurations and scales monotonically with pipe length
(}%DIFDELCMD < \Cref{fig:synth_sweep}%%%
\DIFdel{c, }%DIFDELCMD < \Cref{tab:synth_by_length}%%%
\DIFdel{):
median 3.2}\DIFdelend \DIFaddbegin \subsection{Decision regret and physical feasibility}
\DIFadd{Estimation bias and decision regret diverge sharply, and for the copperplate they
carry }\emph{\DIFadd{opposite signs}}\DIFadd{. Its schedule appears }\textbf{\DIFadd{-15.1}}\DIFaddend \,\DIFdelbegin \DIFdel{\% at 1}\DIFdelend \DIFaddbegin \DIFadd{\% cheaper
than the baseline (bias), yet re-costing that loss-ignoring schedule under the common
forward physics values its unprovisioned loss at }\textbf{\DIFadd{+46.1}}\DIFaddend \,\DIFdelbegin %DIFDELCMD < \si{\kilo\metre}%%%
\DIFdel{, 15.8}\DIFdelend \DIFaddbegin \DIFadd{\% of the baseline
(its forward-valued schedule gap): a copperplate is a biased estimator }\emph{\DIFadd{and}}
\DIFadd{under-provisions the network losses it never saw. The asymmetry
($|\text{regret}|>|\text{bias}|$) is mechanistic: the loss the copperplate omits must
be valued, under the forward model, at the marginal or peak unit's cost in the winter
hours it never planned for, rather than avoided by optimal pre-planning. The schedules show this directly: the copperplate runs the heat pump 2}\DIFaddend \,\DIFdelbegin \DIFdel{\% at 5}\DIFdelend \DIFaddbegin \DIFadd{063 hours against the baseline's 2}\DIFaddend \,\DIFdelbegin %DIFDELCMD < \si{\kilo\metre}%%%
\DIFdel{, 39.0}\DIFdelend \DIFaddbegin \DIFadd{309 and gives it 77.9}\DIFaddend \,\DIFdelbegin \DIFdel{\% }\DIFdelend \DIFaddbegin \DIFadd{\% of production against 87.8\,\%: 246 fewer operating hours and ten percentage points less share, the 46\,\% regret being that substitution priced. Supplying the same
copperplate the aggregate loss it cannot compute (}\CPL\DIFadd{) collapses both bias and regret
to $\approx\textbf{-0.54}\,\%$, but that adder is calibrated ex post on this
network and does not transfer (S\ref{subsec:generalis}). The loss-aware node-resolved levels
(}\Lone\DIFadd{, }\Lthree\DIFadd{, }\Lsix\DIFadd{) reach regret $\approx$ bias with no calibration, as expected, since with
no ignored loss to cover the forward-minus-economic residual is a near-constant pump offset, so
the decision-relevant content is the copperplate's asymmetry rather than this near-equality.
Temperature propagation (}\Ltwo\DIFadd{) is forward-evaluated rather than solved (Section~\ref{subsec:evaluator}), its
free-variable form being degenerate. The sign of the
copperplate regret is invariant across three shortfall-pricing schemes (valuing the shortfall at
the marginal unit, }\DIFaddend at \DIFdelbegin \DIFdel{15}\DIFdelend \DIFaddbegin \DIFadd{the peak unit, and penalising unmet demand at the dump price, giving
$+46$/$+61$/$+832$}\DIFaddend \,\DIFdelbegin %DIFDELCMD < \si{\kilo\metre}%%%
\DIFdel{. }%DIFDELCMD < 

%DIFDELCMD < \begin{figure*}[!ht]
%DIFDELCMD <   %%%
\DIFdelendFL \DIFaddbeginFL \DIFaddFL{\% regret) and across the demand-disaggregation rule
($<0.3$\,pp), so it is not an artefact of either assumption. Physical deliverability is
checked directly: loss-aware schedules incur zero velocity, differential-pressure or
unmet-demand violations under the forward model. This is the decision-relevant reading of
fidelity that model-to-model bias alone cannot give: use the simplest model for
scheduling only if it is loss-aware, not a loss-blind copperplate.
}\begin{table*}[t]
\DIFaddendFL \centering
\DIFdelbeginFL %DIFDELCMD < \includegraphics[width=\textwidth]{paper1_dh_fidelity/Figure_6.png}
%DIFDELCMD <   \caption{Synthetic parametric study across 36~configurations.
%DIFDELCMD <   (a)~L1$\to$L3 gap distribution (median 20.0\,\%, range
%DIFDELCMD <   3.0--44.8\,\%, 36/36 positive).
%DIFDELCMD <   (b)~Additive decomposition sorted by net gap: heat losses
%DIFDELCMD <   (red) dominate; topology routing (blue) and pressure drop
%DIFDELCMD <   (orange) are negligible.
%DIFDELCMD <   (c)~Monotone gap scaling with pipe length.}
%DIFDELCMD <   \label{fig:synth_sweep}
%DIFDELCMD < \end{figure*}
%DIFDELCMD < %%%
\DIFdelend \DIFaddbegin \caption{Estimation bias versus decision regret (\% of the baseline $\Lone$ operating cost).
Decision regret is a forward valuation of the unprovisioned network loss at the marginal
replacement heat (an L1-relative forward-valued schedule gap, not a re-optimised operating
cost), shown for three pricing schemes. \emph{Short.} is the raw unprovisioned network loss
(MWh\,yr$^{-1}$) the schedule leaves uncovered, a physical under-provision, distinct from the
hydraulic \emph{Viol.} count (velocity, differential pressure); only a level with \emph{both}
zero is physically deliverable. The last two columns give the heat-pump energy share and
running hours.}
\label{tab:regret}
\begin{tabular}{lrrrrrrrr}
\hline
Level & Bias & \multicolumn{3}{c}{Regret (marginal / peak / unmet)} & Viol. & Short.\ & HP share & HP hrs \\
 & & & & & (hyd.) & (MWh) & & \\
\hline
$\CP$ & -15.1 & +46.1 & +61.2 & +831.7 & 0 & 1\,159 & 77.9 & 2\,063 \\
$\CPL$ (const.) & -0.6 & -0.5 & -0.5 & -0.5 & 0 & 0 & 78.2 & 2\,309 \\
$\CPL$ (curve) & -0.7 & -0.6 & -0.6 & -0.6 & 0 & 0 & 78.2 & 2\,310 \\
$\NDz$ & -14.4 & +46.8 & +61.9 & +832.5 & 0 & 1\,160 & 78.0 & 2\,061 \\
$\Lone$ & +0.0 & +0.0 & +0.0 & +0.0 & 0 & 0 & 87.8 & 2\,309 \\
$\Lthree$ & +1.4 & +1.4 & +1.4 & +1.4 & 0 & 0 & 87.8 & 2\,311 \\
$\Lsix$ & +1.4 & +1.4 & +1.4 & +1.4 & 0 & 0 & 87.8 & 2\,311 \\
\hline
\end{tabular}
\end{table*}
\DIFaddend 

\DIFdelbegin %DIFDELCMD < \begin{table}[ht]
%DIFDELCMD < %%%
\DIFdelendFL \DIFaddbeginFL \begin{figure}[H]
  \DIFaddendFL \centering
  \DIFdelbeginFL %DIFDELCMD < \caption{Cumulative L1$\to$L3 cost gap by pipe length.}
%DIFDELCMD < \label{tab:synth_by_length}
%DIFDELCMD < \footnotesize
%DIFDELCMD < \begin{tabular*}{\columnwidth}{@{\extracolsep{\fill}} l c c c @{}}
%DIFDELCMD < \toprule
%DIFDELCMD < Pipe length & $n$ & Median & p10--p90 \\
%DIFDELCMD < \midrule
%DIFDELCMD < 1\,km  &  6 & $+3.2\%$  & $+3.0$--$4.4\%$ \\
%DIFDELCMD < 5\,km  & 14 & $+15.8\%$ & $+14.2$--$20.1\%$ \\
%DIFDELCMD < 15\,km & 16 & $+39.0\%$ & $+34.8$--$44.2\%$ \\
%DIFDELCMD < \bottomrule
%DIFDELCMD < \end{tabular*}
%DIFDELCMD < \end{table}
%DIFDELCMD < %%%
\DIFdelend \DIFaddbegin \includegraphics[width=\linewidth,height=0.42\textheight,keepaspectratio]{figure5}
  \caption{Estimation bias against decision regret for every level, as a percentage of the baseline operating cost. For the copperplate and the topology-only control the two carry \emph{opposite signs}: a schedule that appears cheaper on paper leaves a large forward-valued loss shortfall. Loss-aware levels show regret approximately equal to bias.}
  \label{fig:regret}
\end{figure}
\DIFaddend 

\DIFdelbegin \paragraph{\DIFdel{Three-component decomposition.}}
%DIFAUXCMD
\addtocounter{paragraph}{-1}%DIFAUXCMD
\DIFdel{Decomposing the gap into topology routing, heat losses , and pressure drop (}%DIFDELCMD < \Cref{eq:cost_decomposition}%%%
\DIFdel{) reveals an
unambiguous driver (}%DIFDELCMD < \Cref{fig:synth_sweep}%%%
\DIFdel{b): the }\DIFdelend \DIFaddbegin {
\subsection{Sensitivity to the zone aggregation}
\label{subsec:clustering}

\DIFadd{The zone-aggregated level rests on one hand-made seven-zone partition of the fifteen-node
network. A partition is a modelling choice rather than a physical fact, and a different one
could in principle redistribute losses and routing differently. The aggregated model was
therefore re-solved across a family of alternative partitions, holding everything else fixed.
}

\DIFadd{The family comprises the original partition re-run through the same aggregator, three
deliberate alternatives spanning granularity (four, seven and ten zones) and a null
distribution of twenty random tree-contiguous partitions, all solved to a $0.01$\,\%
optimality gap. Every member conserves the total }\DIFaddend heat-loss \DIFdelbegin \DIFdel{component accounts for nearly the entire gap (median
$+20.1\,\%$) , while routing ($-0.7$ to $-8\,\%$ of the total gap}\DIFdelend \DIFaddbegin \DIFadd{capacity
$\sum_p U_p L_p = 1140.2$}\DIFaddend \,\DIFdelbegin \footnote{\DIFdel{Negative routing contributions indicate that
spatial flow-routing constraints marginally reduce cost by
forcing locally efficient asset utilisation; the heat-loss
component correspondingly exceeds 100\,\% of the net gap.}}%DIFAUXCMD
\addtocounter{footnote}{-1}%DIFAUXCMD
\DIFdel{) and pressure drop ($+0.02\,\%$)remain at or below the MIP-gap
noise floor. The $+0.02\,\%$ pressure-drop component is measured within the synthetic decomposition framework
(}%DIFDELCMD < \Cref{tab:synth_level_mapping}%%%
\DIFdel{), where L3 includes
Darcy--Weisbach physics; it is consistent with the $+0.11\,\%$
L3$\to\Lplus$ increment for the primary case
(}%DIFDELCMD < \Cref{tab:cost_extended}%%%
\DIFdel{), where pressure drop enters one
level later}\DIFdelend \DIFaddbegin \DIFadd{W/K of the node-resolved network, which is verified per
configuration and confirmed by an identical realised annual loss across the family
($1257.8$--$1257.9$\,MWh). Conservation requires resolving the producer as its own zone: a
partition that merges the plant with downstream nodes has no incoming pipe to carry the
producer zone's internal trunk losses and silently discards them. The producer is therefore held
isolated in every member, as in the original partition, and the consequence for the null is
stated explicitly: each draw cuts the producer edge and then
five further random edges, partitioning the remaining fourteen consumer nodes into six
zones. The null therefore tests consumer-zone boundary choice at fixed producer resolution
and fixed total loss, which is the comparison the question requires}\DIFaddend .

\DIFdelbegin \paragraph{\DIFdel{Parameter ranking.}}
%DIFAUXCMD
\addtocounter{paragraph}{-1}%DIFAUXCMD
%DIFDELCMD < \Cref{tab:synth_marginal} %%%
\DIFdel{(}%DIFDELCMD < \Cref{app:synth_detail}%%%
\DIFdel{) ranks the four design parameters by their marginal effect on the L1$\to$L3
gap. Pipe length dominates ($+35.8$}\DIFdelend \DIFaddbegin \DIFadd{Because the delivered loss is equal across the family, the only quantity free to move is
where the boundaries fall, a pure routing effect. It barely moves. The three deliberate
alternatives differ from the original by at most $0.01$}\DIFaddend \,\DIFdelbegin \DIFdel{pp across its range);
storage capacity ranks second ($+29.6$}\DIFdelend \DIFaddbegin \DIFadd{\% of annual cost, the twenty null
draws have a standard deviation of $4$}\DIFaddend \,\DIFdelbegin \DIFdel{pp) but is partially
confounded with pipe length due to uneven feasibility filtering
(}%DIFDELCMD < \Cref{app:synth_detail}%%%
\DIFdel{). Node count ($+7.1$}\DIFdelend \DIFaddbegin \DIFadd{EUR, and the full spread across all twenty-four
partitions is }\textbf{\DIFadd{11}}\DIFaddend \,\DIFdelbegin \DIFdel{pp) and demand
heterogeneity ($-4.8$}\DIFdelend \DIFaddbegin \DIFadd{EUR, or }\textbf{\DIFadd{0.008}}\DIFaddend \,\DIFdelbegin \DIFdel{pp) are an order of magnitude
weaker and
non-monotone (}%DIFDELCMD < \Cref{tab:synth_hi_effect}%%%
\DIFdelend \DIFaddbegin \DIFadd{\% of operating cost. Granularity between
four and ten zones changes essentially nothing.
}

\DIFadd{Set against the decomposition, this is more than a robustness check. With losses held
equal, the residual is directly comparable to the spatial-topology main effect measured on
the same network, and the three quantities span three orders of magnitude
(Table~\ref{tab:r16}}\DIFaddend ).

\DIFdelbegin %DIFDELCMD < \Cref{fig:synth_surface} %%%
\DIFdel{confirms this dominance across all
36~configurations: at 1}\DIFdelend \DIFaddbegin \begin{table}[htbp]
\centering
\footnotesize
\caption{Zone-boundary choice against the effects it must be compared with, on the
copperplate-to-baseline ($\CP\!\to\!\Lone$) cost gap. With the total heat-loss capacity held
equal across all twenty-four partitions, the residual routing effect is negligible.}
\label{tab:r16}
\begin{tabular}{lrr}
\hline
Effect & Cost & \% of the $\CP\!\to\!\Lone$ gap \\
\hline
Loss visibility & 19\,737\,EUR & 95.8 \\
Spatial topology (resolving routing, physics-matched) & 52\,EUR & 0.25 \\
Zone-boundary choice (this subsection) & $\lesssim$11\,EUR & at solver tolerance \\
\hline
\end{tabular}
\end{table}

\DIFadd{Choosing where to draw a zone boundary is negligible: at $\lesssim$11}\DIFaddend \,\DIFdelbegin \DIFdel{km the gap remains below 5}\DIFdelend \DIFaddbegin \DIFadd{EUR the residual is
three orders of magnitude below representing the loss and, like the physics-matched topology
effect itself, sits at the $0.01$}\DIFaddend \,\DIFdelbegin \DIFdel{\%, rising
to 14--20}\DIFdelend \DIFaddbegin \DIFadd{\% solver tolerance
($\approx$13}\DIFaddend \,\DIFdelbegin \DIFdel{\% at 5}\DIFdelend \DIFaddbegin \DIFadd{EUR on the Memmingen cost), so no zone-partition effect is resolved
above it. The aggregation is therefore not a hidden degree
of freedom in the results, and the finding corroborates the decomposition from an
independent direction: loss, not geometry, carries the cost.
}

\DIFadd{The corollary is a practical warning worth stating. An aggregation that does }\emph{\DIFadd{not}}
\DIFadd{conserve the total heat-loss capacity (for instance one that merges the producer with
downstream nodes and so discards the trunk pipes internal to that zone, here $10$--$18$}\DIFaddend \,\DIFdelbegin \DIFdel{km and exceeding 35}\DIFdelend \DIFaddbegin \DIFadd{\%
of $\sum_p U_p L_p$) shifts reported cost by up to about $6$}\DIFaddend \,\DIFdelbegin \DIFdel{\% at 15}\DIFdelend \DIFaddbegin \DIFadd{\%, against the $0.01$}\DIFaddend \,\DIFdelbegin \DIFdel{km. Storage
capacity has a secondary effect: larger reservoirs allow more
temporal dispatch flexibility that the copperplate model cannot
capture.
The nearly vertical gradient confirms that network
topology (via thermal losses) rather than operational scheduling
drives the cost underestimation. }\DIFdelend \DIFaddbegin \DIFadd{\%
agreement among conserving partitions. The binding requirement for a valid reduced network
model is preservation of the delivered loss, not fidelity of the zone geometry. That is the
same conclusion the decomposition reaches by a different route, and it is a straightforward
error to make.
}\DIFaddend 

\DIFdelbegin \DIFdel{Notably, }\DIFdelend \DIFaddbegin \DIFadd{One parameterisation note, since two values for }\DIFaddend the \DIFdelbegin \DIFdel{L3$\to\Lplus$ gap is identically zero across all
36~synthetic configurations (median and max, for each of the three pipe-length groups: 1}\DIFdelend \DIFaddbegin \DIFadd{same pipe appear in the released data.
The aggregator uses the network's per-pipe design heat-transfer coefficients, under which
the far-east lateral carries a markedly higher value than the trunk, reflecting its older
uninsulated construction. The Stage-1 validation instead calibrates a single uniform
coefficient across all pipes, one free parameter, to reproduce the annual loss implied by the delivered-energy measurement.
The two serve different purposes and are not intended to agree pipe by pipe: the first is
the design network, the second a minimal calibration of its aggregate behaviour. All
partitions compared here use the same per-pipe basis, so the comparison is unaffected.
}}\label{subsec:zoneclustering}

\subsection{Scope condition: the routing null is specific to central generation}
\DIFadd{One scope condition attaches to the routing null. All generation in Memmingen sits at a
single node, and the synthetic factorial is built the same way (also central generation), so
the null is established for centrally-supplied networks rather than for district-heating
networks in general. Whether resolution matters more when generation is distributed, making
which source serves which demand a live choice, cannot be answered by a radial generator
that injects all heat at the primary producer: a non-root source is structurally stranded,
and admitting one requires the signed-flow formulation discussed in the limitations.
Distributed generation, with meshed and bidirectional topologies, is therefore left to the
companion study.
}

\subsection{Thermo-hydraulic effect, to the transmission station}\label{subsec:hydraulics}
\DIFadd{Station-resolved hydraulics leave the fixed schedule almost unchanged in forward-evaluated cost. Resolving the demand side
to all 174 transmission stations and their service laterals (L4), and replacing the
flat station differential-pressure requirement with a flow-dependent one (L5), the
total electrical pump draw (producer friction, per-station $\Delta p$ and
service-lateral terms) peaks at $\approx 3$}\DIFaddend \,\DIFdelbegin \DIFdel{km, $n{=}6$; 5}\DIFdelend \DIFaddbegin \DIFadd{kW against the $110.8$}\DIFaddend \,\DIFdelbegin \DIFdel{km, $n{=}14$;
15}\DIFdelend \DIFaddbegin \DIFadd{kW of installed
Wilo pump capacity (five parallel units), a margin of $\approx 108$}\DIFaddend \,\DIFdelbegin \DIFdel{km, $n{=}16$).
This
arises because the synthetic L3 implementation includes
Darcy--Weisbach pressure drop (}%DIFDELCMD < \Cref{tab:synth_level_mapping}%%%
\DIFdel{);
the L3$\to\Lplus$ step therefore isolates only transport-delay
activation, which produces no cost change at hourly resolution
when $\tau_{\max} < 1$}\DIFdelend \DIFaddbegin \DIFadd{kW. The worst
real node carries $0.10$}\DIFaddend \,\DIFdelbegin \DIFdel{h.
The pressure-drop effect itself
($+0.11\,\%$)is quantified exclusively on the
primary case
(}%DIFDELCMD < \Cref{tab:cost_extended}%%%
\DIFdel{), where L3 excludes hydraulic physics
by design.
}%DIFDELCMD < 

%DIFDELCMD < \begin{figure}[!ht]
%DIFDELCMD <   \centering
%DIFDELCMD <   \includegraphics[width=\columnwidth]{paper1_dh_fidelity/Figure_7.png}
%DIFDELCMD <   \caption{L1$\to$L3 cost gap as a function of total pipe length
%DIFDELCMD <   and thermal storage capacity across all 36~synthetic
%DIFDELCMD <   configurations (averaged over node count and demand heterogeneity
%DIFDELCMD <   index).}
%DIFDELCMD <   \label{fig:synth_surface}
%DIFDELCMD < \end{figure}
%DIFDELCMD < %%%
%DIF < % ============================================================
%DIF < % 5.5 SENSITIVITY
%DIF < % ============================================================
%DIFDELCMD < \subsection{Sensitivity analysis}
%DIFDELCMD < \label{sec:results_sensitivity}
%DIFDELCMD < 

%DIFDELCMD < %%%
\DIFdel{Gap stability is assessed across 11~price and
demand scenarios
(}%DIFDELCMD < \Cref{tab:gap_stability}%%%
\DIFdel{). The L1--L3 gap remains in the
12--14}\DIFdelend \DIFaddbegin \DIFadd{bar trunk $+\,0.60$}\DIFaddend \,\DIFdelbegin \DIFdel{\% band regardless of fuel (${\pm}25\,\%$) , electricity
(${\pm}30\,\%$), or }%DIFDELCMD < \ce{CO2} %%%
\DIFdel{price assumptions
(${\pm}50\,\%$), confirming that the
topology effect is
}\emph{\DIFdel{structural}} %DIFAUXCMD
\DIFdel{(driven by pipe losses)rather than }\emph{\DIFdel{economic}} %DIFAUXCMD
\DIFdel{(driven by merit-order shifts). The L3--$\Lplus$
gap stays below 0.3}\DIFdelend \DIFaddbegin \DIFadd{bar station $+\,0.62$\,bar lateral
$=1.32$\,bar, leaving $2.7$\,bar over its floor. These figures rest on real component
data (pump datasheets, DXF-traced laterals) and agree with an independent
}\texttt{\DIFadd{pandapipes}} \DIFadd{nonlinear solve to $<0.007$\,bar on trunk pipes, a check of the trunk
pressure }\emph{\DIFadd{calculation}}\DIFadd{, not of the flows themselves nor of the station and lateral
components, which }\texttt{\DIFadd{pandapipes}} \DIFadd{does not resolve (Appendix~\ref{app:pandapipes}).
Applying a $\pm33.8$\,\% deterministic flow perturbation, matched to the aggregate flow MAPE (pump power scales as flow$^{3}$), widens the
pump-energy estimate to roughly $2$--$19$\,MWh\,yr$^{-1}$ and the peak draw to at most
$\sim$7\,kW; even at the upper end the $\approx$108\,kW capacity margin is untouched and the
$\Lone\!\to\!\Lthree$ cost effect stays near $1.4$\,\%, so the conclusion is unchanged. The
pump-energy share is therefore $\sim\!0.080$\,\%
of thermal demand ($7.9\,\mathrm{MWh}$ against $\approx$9.9\,GWh), more than an order of
magnitude below the $1$--$5$}\DIFaddend \,\% \DIFdelbegin \DIFdel{in all 11~scenarios, reinforcing the practical irrelevance of extended hydraulic physics}\DIFdelend \DIFaddbegin \DIFadd{transmission-network literature
range, reconciled by the compact network and heavily oversized pumps, not by a
modelling error. Consequently adding station-level hydraulics to the fixed schedule moves its forward-valued
cost by $<1$\,\% and reveals no evaluated hydraulic violation: even resolved to the individual
substation with flow-dependent pressure, hydraulic fidelity is below the decision-relevance
threshold, which bounds the value of re-optimising rather than the re-optimised schedule itself}\DIFaddend .

\DIFdelbegin \DIFdel{One notable exception: the warm-year scenario shows a sign
reversal for L2--L3 ($-2.2\,\%$). At reducedthermal demand, pipe losses shrink as a fraction of
load; simultaneously, L2's zone-averaged $UL$-product overestimates losses in the
warm season, because the spatial demand distribution under partial load deviates most from the calibration baseline. The sign reversal is therefore a loss-aggregation artefact of L2's zone approximation, not a locational asset effect. This is the only scenario in which L2 overestimates cost relative to L3 suggesting that L2 may be
sufficient for low-load or shoulder-season planning even in
heterogeneous networks. 
}\DIFdelend \DIFaddbegin \subsection{Robustness to a flexible supply temperature}
\label{subsec:tsup}
\DIFaddend 

%DIF < % ============================================================
%DIF < % 5.6 COMPUTATIONAL PERFORMANCE
%DIF < % ============================================================
\DIFdelbegin %DIFDELCMD < \subsection{Computational performance}
%DIFDELCMD < \label{sec:results_computation}
%DIFDELCMD < %%%
\DIFdelend \DIFaddbegin \DIFadd{The station-hydraulics null result and the loss-dominance finding are established
with the plant supply temperature fixed by an outdoor-temperature heating curve.
Because modern operators may lower the supply temperature to reduce losses, this
study stress-tests both findings against that control option.
}\DIFaddend 

\DIFdelbegin %DIFDELCMD < \Cref{tab:computation} %%%
\DIFdel{summarises problem size and solve time. All
MILP variants reach the 0.5\,\% gap within 20\,min.
$\Lnl$ is
intractable for the full year, confirming $\Lplus$ as the
highest thermo-hydraulic practical formulation.
}\DIFdelend \DIFaddbegin \DIFadd{Treating the supply temperature as a free dispatch variable introduces the
bilinear relation
}\[
    \DIFadd{Q = \dot{m}\left(T_{\mathrm{sup}} - T_{\mathrm{ret}}\right),
}\]
\DIFadd{which makes the resulting mixed-integer formulation nonconvex and computationally
intractable at the required scale. Consistent with the treatment of the nonlinear
reference, the supply-temperature variation is therefore evaluated forward rather
than through a new dispatch optimisation. Demand and return temperature remain
fixed while the plant supply temperature is reduced. The validated forward model
then recomputes network losses, mass flows, pipe velocities, and pumping power.
}\DIFaddend 

\DIFdelbegin \DIFdel{Storage dispatch is indistinguishable across abstraction levels
on the primary case: cycle counts (8~equivalent winter cycles)
and SOC timing shifts (${<}\,0.2$\,h) are identical across
L1--$\Lplus$. This equivalence is expected for the compact
primary case ($\tau_{\max} \approx 35$\,min); sensitivity to topology likely emerges in networks with longer transport delays
(}%DIFDELCMD < \Cref{sec:conclusion}%%%
\DIFdel{). }\DIFdelend \DIFaddbegin \DIFadd{Lowering the supply temperature reduces trunk losses but also narrows the
temperature difference. Consequently, mass flow, frictional pumping demand, and
pipe velocity increase. The lowest-cost tested reduction is
\(\mathbf{17.5}\,\mathrm{K}\), yielding an annual saving of
\(\mathbf{6\,220}\,\mathrm{EUR}\,\mathrm{yr}^{-1}\). This corresponds to
approximately \(\mathbf{4.6}\,\%\) of annual operating cost when the recovered
loss is valued at the marginal unit, which is its highest-priced valuation.
}\DIFaddend 

\DIFdelbegin %DIFDELCMD < \begin{table}[ht]
%DIFDELCMD < %%%
\DIFdelendFL \DIFaddbeginFL \DIFaddFL{Two findings follow. First, the saving remains a }\emph{\DIFaddFL{loss effect}}\DIFaddFL{. Its gross
benefit arises from lower evaluated thermal losses and is partly offset by higher
pumping costs. Loss visibility, rather than spatial routing, therefore remains the
operative mechanism when the supply temperature is flexible.
}

\DIFaddFL{Second, hydraulics determine the lowest-cost tested point. Pump energy increases
from \(\mathbf{7.9}\) to \(\mathbf{32.3}\,\mathrm{MWh}\), approximately a factor
of four, as the supply temperature falls. The pipe-velocity limit is first reached
at the tested reduction of \(\mathbf{20.0}\,\mathrm{K}\), making this the first
infeasible tested point rather than the exact continuous feasibility boundary.
Hydraulics therefore shift from a negligible }\emph{\DIFaddFL{cost}} \DIFaddFL{contribution under the
fixed heating curve to a binding }\emph{\DIFaddFL{constraint}} \DIFaddFL{under a flexible supply
temperature. This boundary qualifies, but does not overturn, the
station-hydraulics null result.
}\begin{table}[t]
\DIFaddendFL \centering
\DIFdelbeginFL %DIFDELCMD < \caption{Computational performance (single workstation).}
%DIFDELCMD < \label{tab:computation}
%DIFDELCMD < \begin{threeparttable}
%DIFDELCMD < \scriptsize
%DIFDELCMD < \begin{tabular*}{\columnwidth}{@{\extracolsep{\fill}} l r r r r l @{}}
%DIFDELCMD < \toprule
%DIFDELCMD < Level & Vars & Bin. & Quad.\
%DIFDELCMD <   & Solve & Gap \\
%DIFDELCMD < \midrule
%DIFDELCMD < L1       & 210k   & 53k  & 0   & 268\,s  & ${<}0.5\%$ \\
%DIFDELCMD < L2       & 1.28M  & 53k  & 0   & 450\,s  & ${<}0.5\%$ \\
%DIFDELCMD < L3       & 2.27M  & 53k  & 0   & 650\,s  & ${<}0.5\%$ \\
%DIFDELCMD < $\Lplus$ & 3.02M  & 420k & 0   & 1165\,s & ${<}0.5\%$ \\
%DIFDELCMD < $\Lnl$ (744h)
%DIFDELCMD <          & --     & --   & 65k & ${>}4$h & 1.99\% \\
%DIFDELCMD < \bottomrule
%DIFDELCMD < \end{tabular*}
%DIFDELCMD < \begin{tablenotes}\scriptsize
%DIFDELCMD <   \item %%%
\DIFdelFL{$\Lnl$: Jan. \ window, warm-started from $\Lplus$.
}%DIFDELCMD < \end{tablenotes}
%DIFDELCMD < \end{threeparttable}
%DIFDELCMD < %%%
\DIFdelendFL \DIFaddbeginFL \caption{Forward supply-temperature sensitivity on $\Lone$: lowering the plant supply
temperature by $\Delta T_\text{sup}$ cuts loss cost but raises pumping. Lowest-cost tested
reduction (over the swept grid) marked $^{\star}$; the velocity limit binds at 20\,K. Loss valued at the
marginal unit (its highest-priced valuation). Costs in EUR/yr.}
\label{tab:tsup}
\begin{tabular}{rrrrrc}
\hline
$\Delta T_\text{sup}$ (K) & Loss cost & Pump cost & Total & Vel.\ viol. & Feasible \\
\hline
0.0 & 83,721 & 784 & 84,505 & 0 & yes \\
5.0 & 81,060 & 1,012 & 82,073 & 0 & yes \\
10.0 & 78,381 & 1,466 & 79,847 & 0 & yes \\
15.0 & 76,183 & 2,224 & 78,407 & 0 & yes \\
17.5$^{\star}$ & 75,076 & 3,209 & 78,285 & 0 & yes \\
20.0 & 74,459 & 4,185 & 78,644 & 3 & no \\
25.0 & 74,059 & 5,854 & 79,913 & 33 & no \\
\hline
\end{tabular}
\DIFaddendFL \end{table}


%DIF < % ============================================================
%DIF < % 6. DISCUSSION
%DIF < % ============================================================
%DIF < % ============================================================
%DIF < % 6. DISCUSSION
%DIF < % ============================================================
\DIFdelbegin %DIFDELCMD < \section{Discussion}
%DIFDELCMD < \label{sec:discussion}
%DIFDELCMD < %%%
\DIFdelend \DIFaddbegin \subsection{Linearisation error and transport delay, separated}
\label{subsec:linearisation}
\DIFaddend 

\DIFdelbegin \DIFdel{The results highlight the role of the
topology dominance principle in modelling of DH distribution: in radial-tree topologies with fixedsupply
temperature and pipe lengths up to 15\,km, the modelling step
that delivers nearly all of the quantified accuracy gain
(spatial resolution with steady-state losses) requires zero
additional solver capability, zero nonlinear constraints, and increases solve time by merely 70\,\%. Conversely, the modelling
steps that consume the vast majority of development effort and
solver resources (thermo-hydraulic physics with PWL or MIQCP
formulations) deliver less than 1\,\% marginal accuracy. The
practical implication is direct: }\emph{\DIFdel{resolve the graph before
refining the physics}}%DIFAUXCMD
\DIFdelend \DIFaddbegin \DIFadd{Each extended phenomenon is isolated individually. Under the operating regime
considered here, each has only a limited effect relative to loss visibility}\DIFaddend .

\DIFdelbegin \DIFdel{The qualifier matters. The 1\,\% upper bound on extended-physics
effects has been established for radial-tree topologies, fixed
heating curves, and dispatch with given capacities. Ring
topologies with bidirectional flow may amplify hydraulic coupling
through loop-flow degrees of freedom; variable supply temperature
introduces $T \cdot \dot{m}$ bilinear terms that scale with
network throughput; joint dispatch--sizing problems may select
different capacities under L1 vs.
\ L3 representations, which in
turn affects pumping fractions. Each of these is identified as a priority follow-up in }%DIFDELCMD < \Cref{sec:conclusion}%%%
\DIFdelend \DIFaddbegin \paragraph{\DIFadd{Trunk pressure and pumping.}}
\DIFadd{Adding trunk pressure and pumping to the node-resolved baseline and solving the
resulting model to an optimality gap of \(0.01\,\%\) increases annual cost by about
\(1.4\,\%\), equivalent to approximately
\(1{,}900\,\mathrm{EUR}\,\mathrm{yr}^{-1}\) relative to the \(\Lone\) cost.
This increment does not represent pump energy alone. Pumping consumes
\(7.9\,\mathrm{MWh}\), costing approximately
\(780\,\mathrm{EUR}\,\mathrm{yr}^{-1}\) (Table~S2, Supplementary Material).
The remaining \(1{,}100\,\mathrm{EUR}\) arises from the demand charge and the
dispatch response induced by including pumping in the objective. Direct pump
electricity therefore accounts for less than half of the increment, while the
dispatch ordering remains unchanged}\DIFaddend .

\DIFdelbegin %DIFDELCMD < \subsection{Effect hierarchy and interpretation}
%DIFDELCMD < \label{sec:disc_hierarchy}
%DIFDELCMD < %%%
\DIFdelend \DIFaddbegin \paragraph{\DIFadd{Temperature propagation and transport delay.}}
\DIFadd{Temperature propagation is evaluated forward as described in
Section~\ref{subsec:evaluator}. Its free-variable formulation is nonconvex and
degenerate because it allows the optimiser to alter the temperature level rather
than isolate physical propagation. Native temperature propagation reduces total
network losses by approximately \(2\,\%\), corresponding to only a few tenths of
a percent of total cost. Transport delay, represented as a linear temporal shift,
has an even smaller effect.
}\DIFaddend 

\DIFdelbegin \DIFdel{The results establish the following magnitude ordering:
%DIF < 
}\begin{align*}
  \DIFdel{\underbrace{\text{L1} \to \text{L2}}_{\substack{\text{dominant} \\
    10.4\%}}
  \;>\;
  \underbrace{\text{L2} \to \text{L3}}_{\substack{\text{resolution}
    \\ 2.6\%}}
  \;\gg\;
  \underbrace{\text{L3} \to \Lplus}_{\substack{\Delta p \\
    0.11\%}}
  \;>\;
  \underbrace{\Lplus \to \Lnl}_{\substack{\text{bound} \\
    {\leq}0.5\%}}
}\end{align*}%DIFAUXCMD
%DIFDELCMD < \Cref{fig:effect_hierarchy} %%%
\DIFdel{visualises the paradox: the first
step (topology + losses) delivers nearly all accuracy gain at minimal solve-time cost, while subsequent steps invert this
ratio}\DIFdelend \DIFaddbegin \paragraph{\DIFadd{Solved nonlinear reference.}}
\DIFadd{This study quantifies the linearisation error by solving a nonlinear reference
rather than relying only on a bound. The mixed-integer schedule obtained from the
piecewise-linear model is held fixed, including unit commitment and
piecewise-linear segment selectors. Only the continuous variables, comprising
flows, nodal temperatures, pressures, and pumping power, are reoptimised under
the native physics.
}

\DIFadd{The resulting model restores exponential temperature decay and the heat-loss
bilinearities. It is therefore a nonconvex nonlinear programme rather than a
quadratic programme. For representative winter and autumn windows of 72 hours,
the nonlinear re-solves change cost by \(\mathbf{-0.15}\,\%\) and
\(\mathbf{-0.33}\,\%\), respectively. The global solver
(Gurobi, }\texttt{\DIFadd{NonConvex=2}}\DIFadd{) reports optimality gaps of \(0.009\,\%\) for the
winter window and \(0.025\,\%\) for the autumn window, both below
\(0.03\,\%\) and measured against valid global lower bounds}\DIFaddend .

\DIFdelbegin %DIFDELCMD < \begin{figure}[ht]
%DIFDELCMD <   %%%
\DIFdelendFL \DIFaddbeginFL \DIFaddFL{The negative signs are informative. They show that the piecewise linearisation
slightly overstates cost and that the native nonlinear physics is marginally
cheaper. Both effects remain below half a percent, consistent with the exact
decomposition. The friction and station-pressure terms contribute only a few
hundredths of a percent and partly offset one another.
}

\paragraph{\DIFaddFL{Limits of the nonlinear reference.}}
\DIFaddFL{The corresponding full-year re-solve produces no incumbent within the available
solver budget. No full-year value is therefore reported as a solved result.
In addition, a low-load summer window is infeasible under the native physics
when the linearised schedule is held fixed. This solver outcome is reported
without assigning it a physical interpretation. It has not been established
whether the result reflects genuine physical non-deliverability or a formulation
or numerical issue. An irreducible infeasible subsystem analysis or a
feasibility-relaxation diagnosis is left for future work.
}

\DIFaddFL{Each effect is reported relative to its rigorous optimality bound. No point
estimate is claimed when the bound cannot resolve the effect. In such cases,
the conclusion is that no effect can be demonstrated at the attained tolerance.
Across all extensions, the result remains consistent with the decomposition:
every effect beyond loss visibility is only a fraction of the loss effect.
}

\begin{table}[t]
\DIFaddendFL \centering
\DIFdelbeginFL %DIFDELCMD < \includegraphics[width=\linewidth]{paper1_dh_fidelity/Figure_8.png}
%DIFDELCMD <   \caption{Effect hierarchy and computational trade-off.
%DIFDELCMD <   (a)~Cost gap per modelling step (\% of cost(L3)); topology
%DIFDELCMD <   with losses dominates.
%DIFDELCMD <   (b)~Additional solve time per step (log scale); extended
%DIFDELCMD <   physics and MIQCP consume disproportionate resources for
%DIFDELCMD <   marginal accuracy.}
%DIFDELCMD <   \label{fig:effect_hierarchy}
%DIFDELCMD < \end{figure}
%DIFDELCMD < %%%
\DIFdelend \DIFaddbegin \caption{Extended-physics effects on dispatch cost, each isolated with the method used to
quantify it. Entries are shares of total cost except the temperature-propagation row, which is a
share of network loss and corresponds to a few tenths of a percent of cost. All lie far below the
loss main effect, and the linearisation error is small and negative (the piecewise-linear model
slightly over-states cost).}
\label{tab:linearisation}
\begin{tabular}{lrl}
\hline
Phenomenon (step) & Effect & How quantified \\
\hline
Trunk pressure \& pumping ($\Lone\!\to\!\Lthree$) & $+1.4$\,\% & solved MILP, $\le$0.01\,\% gap \\
Temperature propagation ($\Lone\!\to\!\Ltwo$) & $\approx-2$\,\% of loss & forward (native decay) \\
Transport delay ($\Lthree\!\to\!\Lsix$) & $<0.01$\,\% & solved MILP, $\le$0.01\,\% gap ($k_p{=}0$) \\
Linearisation error (PWL vs native) & $-0.15/-0.33$\,\% & solved native, $\le$0.03\,\% gap (winter/autumn 72-h) \\
\hline
\end{tabular}
\end{table}
\DIFaddend 

\DIFdelbegin \DIFdel{The synthetic decomposition refines this picture: pure topology routing without losses contributes effectively zero; the entire dominant step is attributable to thermal losses that become visible
once the network is spatially resolved. An important finding is not that topology matters (this follows from thermodynamics) but that }\emph{\DIFdel{pressure-drop physics, transport delay, and nonlinear formulations are collectively
unnecessary}} %DIFAUXCMD
\DIFdel{for planning-accuracy dispatch in }\DIFdelend \DIFaddbegin \subsection{Secondary finding: generalisability and out-of-sample prediction}
\label{subsec:generalis}
\DIFadd{The decomposition generalises. Across the 135 synthetic networks the loss main effect
accounts for a median }\textbf{\DIFadd{100.0}}\DIFadd{\,\% of the cost gap and the topology main
effect stays within $\pm$0.6\,\% on every network of 5\,km trunk length or more, never exceeding 2.4\,\% even on the shortest, where the entire gap is below 6\,\% of cost, over the full range of
node counts, pipe lengths, demand heterogeneity and storage sizes; a variance decomposition
attributes essentially all of the between-network variation in the gap to loss and almost none
to topology. The same data expose why the lumped-loss shortcut does not transfer. The loss
}\emph{\DIFadd{burden}}\DIFadd{, the share of cost the adder must reproduce, ranges from
}\textbf{\DIFadd{3.2}}\DIFadd{\,\% to }\textbf{\DIFadd{81.2}}\DIFadd{\,\% across the factorial, driven
chiefly by pipe length; a single frozen adder is by construction correct for only one point in
that range, and even the most transferable choice mis-estimates the burden by a mean of
}\textbf{\DIFadd{23.5}} \DIFadd{and up to }\textbf{\DIFadd{40.1}} \DIFadd{percentage points of cost
when carried to another network. This is the quantitative form of the non-transfer that makes
the node-resolved model, which computes the loss endogenously, the transferable one. Finally,
an a-priori estimator of the cost gap from observable network properties (pipe length, node
count, demand heterogeneity and storage) fitted on the synthetic }\DIFaddend networks up to 15\,km \DIFdelbegin \DIFdel{pipe length . This ``negative result'' for thermo-hydraulics has
direct resource-allocation implications: modelling effort should be
invested in spatial resolution (which is cheap computationally) rather
than physics refinement (which is expensive and yields $<\!1\,\%$) .
}%DIFDELCMD < 

%DIFDELCMD < %%%
\DIFdel{Four implications follow for practitioners:
%DIF < 
}%DIFDELCMD < \begin{enumerate}[nosep]
%DIFDELCMD <   \item %%%
\textbf{\DIFdel{Thermal-loss accounting is the first modelling
        investment.}} %DIFAUXCMD
\DIFdel{Losses scale with pipe length (median
        $+3\,\%$ at 1}\DIFdelend \DIFaddbegin \DIFadd{trunk
length (train $R^2\approx0.999$) predicts the held-out longer networks without refitting, an
extrapolation beyond the fitted range. It predicts the 30}\DIFaddend \,km \DIFdelbegin \DIFdel{to $+39\,\%$ at 15}\DIFdelend \DIFaddbegin \DIFadd{networks to within about seven
percentage points and under-predicts at the extreme 50}\DIFaddend \,km \DIFdelbegin \DIFdel{).
  }%DIFDELCMD < \item %%%
\textbf{\DIFdel{Pressure-drop physics are marginal.}} %DIFAUXCMD
\DIFdel{Even at
        15}\DIFdelend \DIFaddbegin \DIFadd{extrapolation (held-out mean absolute
percentage error $\approx19$}\DIFaddend \,\DIFdelbegin \DIFdel{km, }\DIFdelend \DIFaddbegin \DIFadd{\%; Table~S5, Supplementary Material). This study reports the outcome as obtained. Extrapolation performance degrades
at the extreme, as expected. A prediction specified before observing the outcome,
whether successful or not, is more informative than a curve fitted.
}

\DIFadd{The estimator above is a fitted regression on observable properties; the same physics yields a
sharper, parameter-free tool. If }\DIFaddend the \DIFdelbegin \DIFdel{pressure-drop component remains an order of
        magnitude below the topology}\DIFdelend \DIFaddbegin \DIFadd{node-resolved model differs from a copperplate only in that
it additionally generates the network loss, and cost scales with heat produced, then the burden a
copperplate misses is $b=\lambda/(1+\lambda)$, where the dimensionless }\emph{\DIFadd{loss number}}
\DIFadd{$\lambda=\text{annual loss}/\text{annual demand}$ is computable a priori from the pipe geometry, heat-transfer coefficients, supply}\DIFaddend /\DIFdelbegin \DIFdel{loss gap. }%DIFDELCMD < \item %%%
\textbf{\DIFdel{Transport delay requires MIQCP to characterise.}}
        %DIFAUXCMD
\DIFdel{Under MILP linearisation, transport delay is bypassed. If $\tau > 1$\,h, a nonlinear validation run is warranted. }%DIFDELCMD < \item %%%
\DIFdelend \DIFaddbegin \DIFadd{return/ground temperatures and annual demand
alone, $\lambda=\big(\sum_p U^{s}_p L_p (T_{\text{sup}}-T_g)+\sum_p U^{r}_p L_p
(T_{\text{ret}}-T_g)\big)\cdot 8760/(10^{6}\,D)$, with the $10^{6}$ converting W to MW,
representative annual temperatures, and $D$ the annual demand in MWh, without solving any
optimisation. This
zero-parameter curve reproduces the computed burden across all }\DIFaddend \textbf{\DIFdelbegin \DIFdel{MILP is adequate for planning.}\DIFdelend \DIFaddbegin \DIFadd{136}\DIFaddend } \DIFdelbegin \DIFdel{The combined
        $\Lplus\!\to\!\Lnl$ bound (${\leq}\,0.5\,\%$, February
        window) is itself comparable to solver tolerance; }\DIFdelend \DIFaddbegin \DIFadd{points (}\DIFaddend the
\DIFdelbegin \DIFdel{January window confirms no measurable improvement even at
        relaxed optimality (1.99}\DIFdelend \DIFaddbegin \DIFadd{135 synthetic networks and Memmingen, $\lambda$ spanning
}\textbf{\DIFadd{0.03}}\DIFadd{--}\textbf{\DIFadd{1.7}}\DIFadd{) with $R^2=\textbf{0.87}$ and a mean
absolute error of }\textbf{\DIFadd{6.7}} \DIFadd{percentage points; a one-parameter calibration that
absorbs the residual constant topology term raises this to $R^2=\textbf{0.96}$. Memmingen
($\lambda=\textbf{0.12}$) sits on the curve, its }\textbf{\DIFadd{11}}\DIFaddend \,\% \DIFdelbegin \DIFdel{MIP gap). }%DIFDELCMD < \end{enumerate}
%DIFDELCMD < %%%
\DIFdelend \DIFaddbegin \DIFadd{predicted
burden close to the }\textbf{\DIFadd{15}}\DIFadd{\,\% measured. The residual few points are
informative rather than error: they are consistent with the omitted topology main effect and the
accounting difference of Section~\ref{subsec:costacct}, so the rule under-predicts by about the
size of the terms it omits. It is a }\emph{\DIFadd{first-order screening heuristic}}\DIFadd{, not a general cost
law: it assumes heat-production cost roughly proportional to output, no commitment
discontinuities, no demand-charge or CHP-export interaction, no storage-induced timing effect,
and a consistent loss denominator. Under those assumptions it gives a one-line screen, plotted
as a nomogram in Figure~\ref{fig:rule}: a planner computes $\lambda$ from the network geometry
and reads off whether a copperplate plausibly suffices ($\lambda$ small, $b\lesssim10$\,\%),
whether a calibrated loss adder is likely adequate, or whether node resolution is worth
checking, a screening indication rather than a guarantee, since a known transferable loss
adder can itself suffice, as the decomposition shows.
}\begin{table}[t]
\centering
\caption{Balanced-design analysis of variance of the loss burden across the 135 synthetic
networks (full $3\times5\times3\times3$ factorial). Trunk pipe length dominates the
between-network variance.}
\label{tab:synthanova}
\begin{tabular}{lr}
\hline
Factor & Share of loss-burden variance \\
\hline
Trunk pipe length & 95.8\,\% \\
Storage horizon & 3.2\,\% \\
Demand heterogeneity & 0.0\,\% \\
Node count & 0.0\,\% \\
Interactions + residual & 0.8\,\% \\
\hline
\end{tabular}
\end{table}
\DIFaddend 


\DIFdelbegin %DIFDELCMD < \subsection{Integration pathway for multi-energy system frameworks}
%DIFDELCMD < \label{sec:disc_tools}
%DIFDELCMD < %%%
\DIFdelend \DIFaddbegin \begin{figure}[H]
  \centering
  \includegraphics[width=\linewidth,height=0.42\textheight,keepaspectratio]{figure6}
  \caption{Fidelity design rule. The cost gap a copperplate misses (the loss burden) against the
  a-priori loss number $\lambda=\text{annual loss}/\text{annual demand}$. The parameter-free
  physics curve $b=\lambda/(1+\lambda)$ (navy) tracks the computed burden of the 135
  synthetic networks (grey) and Memmingen (star) with $R^2=\textbf{0.87}$. Shaded bands are
  \emph{illustrative} screening indications (copperplate/adder, calibrate an adder, or check
  node resolution) read off from $\lambda$ alone; the thresholds are judgmental, not validated
  cutoffs.}
  \label{fig:rule}
\end{figure}
\begin{figure}[H]
  \centering
  \includegraphics[width=\linewidth,height=0.42\textheight,keepaspectratio]{figure7}
  \caption{Drift of a frozen loss adder against trunk pipe length across the 135 synthetic networks. An adder calibrated on any one network mis-estimates the loss burden on others by a mean of 23.5 and up to 40.1 percentage points of cost, which is why the node-resolved model, computing the loss endogenously, is the transferable one.}
  \label{fig:drift}
\end{figure}
\DIFaddend 

\DIFdelbegin \DIFdel{Open-source frameworks for multi-energy system optimisation 
including oemof~\mbox{%DIFAUXCMD
\cite{Hilpert2018_oemof}}\hskip0pt%DIFAUXCMD
, PyPSA~\mbox{%DIFAUXCMD
\cite{Brown2018_pypsa}}\hskip0pt%DIFAUXCMD
, FINE~\mbox{%DIFAUXCMD
\cite{Welder2018_FINE}}\hskip0pt%DIFAUXCMD
, Calliope~\mbox{%DIFAUXCMD
\cite{Pfenninger2018_calliope}}\hskip0pt%DIFAUXCMD
,
and urbs~\mbox{%DIFAUXCMD
\cite{Dorfner2016_urbs}}\hskip0pt%DIFAUXCMD
, typically represent thermal
networks at the copperplate level (L1)in their default
DH formulations as used in the MES literature cited above.
Several frameworks offer extensions with per-link loss representations (e.g., \ oemof-thermal, PyPSA-Heat, Calliope with loss
factors); however, in
cross-framework MES studies, the copperplate representation
remains the de facto default.
Our results quantify the cost of this default: for networks with total pipe length exceeding
}%DIFDELCMD < \SI{1}{\kilo\metre} %%%
\DIFdel{(i.e., virtually all operational DH systems), the L1 default introduces a systematic cost underestimation of 3--45}\DIFdelend \DIFaddbegin \subsection{Sensitivity and robustness}
\DIFadd{The decomposition is not an artefact of a single modelling choice: loss dominance survives every
perturbation in Table~S4 (Supplementary Material), from removing the indefensible terminal-pipe
multiplier (the defensible-$U$ loss share of }\textbf{\DIFadd{95.8}}\DIFaddend \,\% \DIFdelbegin \DIFdel{depending on network extent.
}\DIFdelend \DIFaddbegin \DIFadd{is
essentially the inflated-$U$ value) through the two lumped-loss calibrations, the
demand-disaggregation rule and the marginal-cost range to a flexible supply temperature. It also
survives tightening the optimality gap to $0.01$\,\%, at which the small topology and interaction
terms resolve out of solver noise to stable values (Table~S3, Supplementary Material). Where a result }\emph{\DIFadd{is}} \DIFadd{sensitive to an assumption, this study states this
explicitly. The magnitude of last-mile losses depends on the reconstruction of
the service laterals, while the value of supply-temperature flexibility depends
on how the recovered loss is priced. Neither sensitivity changes the ordering
on which the conclusions rest.
}\begin{figure}[H]
  \centering
  \includegraphics[width=\linewidth,height=0.42\textheight,keepaspectratio]{figure8}
  \caption{Supply-temperature flexibility. Lowering the plant supply temperature reduces thermal loss but shrinks the temperature difference, so flow and pumping rise. The lowest-cost tested reduction is 17.5\,K; beyond 20\,K the pipe velocity limit binds, and hydraulics change from a negligible cost into the binding constraint.}
  \label{fig:tsup}
\end{figure}
\DIFaddend 

\DIFdelbegin \DIFdel{The contribution of this work is
therefore to identify a }\emph{\DIFdel{pragmatic integration pathway}} %DIFAUXCMD
\DIFdel{for frameworks that have not
yet committed to a particular optimisation level:
%DIF < 
}%DIFDELCMD < \begin{enumerate}[nosep]
%DIFDELCMD <   \item %%%
\textbf{\DIFdel{L2 as a low-cost minimum.}} %DIFAUXCMD
\DIFdel{Adding per-segment
        steady-state losses (}%DIFDELCMD < \Cref{eq:loss_supply} %%%
\DIFdel{and }%DIFDELCMD < \Cref{eq:loss_return}%%%
\DIFdel{)requires only a linear loss term per
        pipe and a nodal balance reformulation
        (}%DIFDELCMD < \Cref{eq:balance_nodal}%%%
\DIFdel{). No solver change, no
        additional binaries, and minimal solve-time increase
        (${+}70\,\%$ in our implementation, }%DIFDELCMD < \Cref{tab:computation}%%%
\DIFdel{).
This step alone closes most of the 3--45}\DIFdelend \DIFaddbegin {
\subsection{Fidelity gain versus computational cost}
\label{subsec:computation}

\DIFadd{Table~\ref{tab:computation} reports model size and solve time for the levels that are
solved. Read alongside the effects each level buys, it gives the trade-off a modeller
actually faces, and the trade is unusually lopsided.
}

\DIFadd{Moving from the copperplate to the node-resolved baseline multiplies the model by roughly
a factor of eleven in variables and by about a factor of two in solve time, and buys the
entire $15.1$}\DIFaddend \,\% \DIFdelbegin \DIFdel{L1-bias gap.
    }%DIFDELCMD < \item %%%
\textbf{\DIFdel{L3 covers the main impact: topology}} %DIFAUXCMD
\DIFdel{Resolving
        individual junctions (L3) adds modest cost
        (1.44$\times$ L2 solve
time) and is warranted when demand
        is spatially heterogeneous (Gini~${>}\,0.3$) .
  }%DIFDELCMD < \item %%%
\textbf{\DIFdel{Diminishing returns beyond L3.}} %DIFAUXCMD
\DIFdel{Pressure-drop
        physics ($\Lplus$)
use 420}\DIFdelend \DIFaddbegin \DIFadd{cost gap: the whole of the estimation bias and the whole of the
$46.1$}\DIFaddend \,\DIFdelbegin \DIFdel{k binary variables in total and double
        solve time while contributing ${\leq}\,0.3\,\%$ cost
accuracy across all tested scenarios (radial trees, fixed
        heating curve; maximum $+0.2\,\%$, cf. \ }%DIFDELCMD < \Cref{tab:gap_stability}%%%
\DIFdel{). For planning studies, L3
        represents the complexity--accuracy Pareto front. }%DIFDELCMD < \item %%%
\textbf{\DIFdel{Validation discipline.}} %DIFAUXCMD
\DIFdel{Framework users
        integrating network losses should validate against
        measured data using the
two-stage protocol described in
        }%DIFDELCMD < \Cref{sec:val_protocol}%%%
\DIFdel{; uncalibrated $U$-values can
        introduce systematic errors that exceed the topology gap
        itself.
}%DIFDELCMD < \end{enumerate}
%DIFDELCMD < 

%DIFDELCMD < \Cref{fig:pareto} %%%
\DIFdel{confirms this hierarchy. In the
figure, the solving time is plotted against the accuracy. L3 achieves
${\leq}\,0.5\,\%$ cost deviation from the MIQCP reference at less
than 11}\DIFdelend \DIFaddbegin \DIFadd{\% decision regret. Every increment above that is expensive and nearly inert.
Adding trunk pressure adds $368$}\DIFaddend \,\DIFdelbegin \DIFdel{min solve time , while the full nonlinear reference
($\Lnl$) requires ${>}\,4$}\DIFdelend \DIFaddbegin \DIFadd{thousand binary variables, roughly doubling the solve
time, to move dispatch cost by $1.4$}\DIFaddend \,\DIFdelbegin \DIFdel{h for a marginal accuracy gain.
}\DIFdelend \DIFaddbegin \DIFadd{\%. Adding transport delay adds no structure at all
at hourly resolution and moves cost by exactly nothing. The station-resolved levels are
not solved as dispatch problems, and the forward evaluation that replaces them shows why
that is no loss: their effect on decisions is indistinguishable from zero.
}\DIFaddend 

\DIFdelbegin %DIFDELCMD < \begin{figure}[ht]
%DIFDELCMD <   %%%
\DIFdelendFL \DIFaddbeginFL \DIFaddFL{The practical reading is that fidelity in this setting is not a smooth trade of accuracy
against computation. It has one step, and the step is loss visibility. A modeller who
resolves the network and computes its losses has bought essentially all of the available
correction at about twice the solve time of a copperplate; a modeller who then adds
hydraulics pays again and buys almost nothing. That the expensive levels are also the ones
whose extra detail cannot be validated at billing-grade metering
(Section~\ref{subsec:validation}) makes the case against them twice over.
}

\DIFaddFL{Storage dispatch is indistinguishable across levels on this case: equivalent winter cycle
counts and state-of-charge timing shifts under $0.2$\,h from the copperplate through the
pressure-resolved level. This is expected for a compact network (the longest trunk travel
time is about 35 minutes, which is also why transport delay vanishes at hourly resolution)
and it is the mechanism behind the identity between the delay level and the level below
it. On a network with substantially longer transport times, both results would be expected
to separate.
}}
\begin{table}[t]
\DIFaddendFL \centering
\DIFdelbeginFL %DIFDELCMD < \includegraphics[width=\linewidth]{paper1_dh_fidelity/Figure_9.png}
%DIFDELCMD <   \caption{Accuracy solve-time comparison. The dashed vertical line marks a 1\,h solver run time.}
%DIFDELCMD <   \label{fig:pareto}
%DIFDELCMD < \end{figure}
%DIFDELCMD < %%%
\DIFdelend \DIFaddbegin \caption{Model size and solve time for the solved levels (single-year, hourly). Temperature
propagation (L2) and the station levels (L4/L5) are forward-evaluated, not solved -- L2's
free-variable temperature model is a non-convex mixed-integer bilinear program and degenerate. L3 adds the pressure PWL binaries.}
\label{tab:computation}
\begin{tabular}{lrrrr}
\hline
Level & Variables & Binaries & Constraints & Solve time \\
\hline
$\CP$ & 210\,246 & 52\,562 & 271\,567 & 3 min \\
$\CPL$ & 210\,246 & 52\,562 & 271\,567 & 3 min \\
$\NDz$ & 2\,268\,846 & 52\,562 & 1\,804\,567 & 4 min \\
$\Lone$ & 2\,268\,846 & 52\,562 & 1\,804\,567 & 5 min \\
$\Lthree$ & 3\,153\,606 & 420\,482 & 3\,687\,967 & 9 min \\
$\Lsix$ & 3\,153\,606 & 420\,482 & 3\,687\,967 & 9 min \\
\hline
\end{tabular}
\end{table}
\DIFaddend 

\DIFdelbegin \DIFdel{These findings do not invalidate L1-based studies for their
stated purpose (national-scale technology screening); they
characterise the magnitude of network-level dispatch error in default L1 configurations and offer framework developers a
quantitative basis for prioritising thermal-network extensions. }\DIFdelend \DIFaddbegin \begin{figure}[H]
  \centering
  \includegraphics[width=\linewidth,height=0.42\textheight,keepaspectratio]{figure9}
  \caption{Solve time and model size across the solved levels. Bar labels give the variable count. Solve time rises from three to nine minutes along the ladder; the decision-relevant step (CP$\to$L1, loss visibility) is reached at five minutes, whereas the further physics (L3, L6) nearly doubles solve time and model size while changing dispatch by under $1.4$\,\%.}
  \label{fig:solvetime}
\end{figure}
\DIFaddend 

\DIFdelbegin %DIFDELCMD < \subsection{Model selection criteria}
%DIFDELCMD < \label{sec:disc_criteria}
%DIFDELCMD < 

%DIFDELCMD < %%%
\DIFdel{The quantified gaps translate into a decision tree for practitioners (}%DIFDELCMD < \Cref{tab:criteria}%%%
\DIFdel{). The dominant cost driver is topology because it is the sole mechanism exposing thermal losses
to the optimiser. }\DIFdelend \DIFaddbegin {
\subsection{Why the extended physics moves so little here}
\label{subsec:whynull}
\DIFaddend 

\DIFdelbegin %DIFDELCMD < \begin{table*}[!ht]
%DIFDELCMD < \raggedright
%DIFDELCMD < \caption{Recommended model selection criteria.}
%DIFDELCMD < \label{tab:criteria}
%DIFDELCMD < \footnotesize
%DIFDELCMD < \begin{tabularx}{\textwidth}{@{} l X X l @{}}
%DIFDELCMD < \toprule
%DIFDELCMD < Level & Conditions & Rationale & Solve \\
%DIFDELCMD < \midrule
%DIFDELCMD < L1
%DIFDELCMD <   & Pipe ${<}\SI{1}{\kilo\metre}$ AND no storage
%DIFDELCMD <   & Losses ${<}1\,\%$
%DIFDELCMD <   & Seconds \\[4pt]
%DIFDELCMD < L2
%DIFDELCMD <   & Pipe 1--\SI{10}{\kilo\metre} AND uniform demand
%DIFDELCMD <     (Gini ${\leq}0.3$, top-3 ${\leq}50\,\%$)
%DIFDELCMD <   & Losses significant; demand uniform
%DIFDELCMD <   & Minutes \\[4pt]
%DIFDELCMD < L3
%DIFDELCMD <   & Heterogeneous demand
%DIFDELCMD <     (Gini ${>}0.3$ or top-3 ${>}50\,\%$) OR
%DIFDELCMD <     storage ${>}6$\,h
%DIFDELCMD <   & Spatial dispatch differences
%DIFDELCMD <   & Min- to hours \\[4pt]
%DIFDELCMD < $\Lplus$
%DIFDELCMD <   & Pipe ${>}\SI{5}{\kilo\metre}$ AND pumping ${>}1\,\%$
%DIFDELCMD <     of thermal losses
%DIFDELCMD <   & Extended hydraulic physics needed
%DIFDELCMD <   & Hours \\[4pt]
%DIFDELCMD < $\Lnl$
%DIFDELCMD <   & $\tau_{\max} > \SI{1}{\hour}$ (transport delay relevant);
%DIFDELCMD <     OR validation of $\Lplus$;
%DIFDELCMD <     OR research
%DIFDELCMD <   & Only level with active transport delay
%DIFDELCMD <   & Hours to days \\
%DIFDELCMD < \bottomrule
%DIFDELCMD < \end{tabularx}
%DIFDELCMD < \end{table*}
%DIFDELCMD < \subsection{\ce{CO2} policy implications}
%DIFDELCMD < \label{sec:disc_co2}
%DIFDELCMD < %%%
\DIFdelend \DIFaddbegin \DIFadd{Three of this paper's results are nulls: station-resolved hydraulics change dispatch cost
by under one percent, transport delay by exactly nothing, and the piecewise linearisation
by a few tenths of a percent. A null invites two readings. Either these phenomena are
genuinely inert in district-heating dispatch, or the assumptions under which they were tested
closed the channels through which they would have acted. The second reading is the
sceptical one; it is legitimate, and it deserves an answer more specific than a caveat.
}\DIFaddend 

\DIFdelbegin \DIFdel{Simplified models underestimate }%DIFDELCMD < \ce{CO2} %%%
\DIFdel{through compounding
channels: missed loss-compensation generation, missed pumping
electricity, and temporal correlation between additional demand
and grid emission intensity. In the primary case, L1-based
accounting underestimates emissions by $8.7\,\%$ (772
vs. \ 846\,t}\DIFdelend \DIFaddbegin \DIFadd{Four assumptions each close a specific channel. A }\emph{\DIFadd{prescribed heating curve}} \DIFadd{fixes the
supply and return temperatures, so the optimiser cannot trade thermal loss against pumping
energy, the trade in which hydraulics would acquire economic weight. }\emph{\DIFadd{Precomputed
heat-pump performance}} \DIFadd{makes the coefficient of performance independent of the network
state, so a heat pump cannot respond to a locally depressed supply temperature by shifting
its output in time or place. }\emph{\DIFadd{Fixed capacities}} \DIFadd{remove siting and sizing, so spatial
structure cannot express itself through where capacity is placed, only through where heat
is routed. And }\emph{\DIFadd{hourly resolution}} \DIFadd{discretises away every transient shorter than the
time step: With $\Delta t = 1$}\DIFaddend \,\DIFdelbegin %DIFDELCMD < \ce{CO2}%%%
\DIFdel{/yr) ; the synthetic parametric study shows this scaling
from ${<}5\,\%$ at 1}\DIFdelend \DIFaddbegin \DIFadd{h, the trunk transit time of about $35$}\DIFaddend \,\DIFdelbegin \DIFdel{km to ${>}30\,\%$ at 15}\DIFdelend \DIFaddbegin \DIFadd{min gives $k_p = 0$: sub-hourly
delays are not represented on the hourly grid, and the delay step therefore shifts nothing on
this network. This is the reason the measured delay effect in
Table~\ref{tab:linearisation} lies below the solver tolerance.
}

\DIFadd{The honest conclusion is therefore conditional, and stated that way throughout: within
the deterministic, hourly, fixed-capacity, fixed-temperature planning regime, extended
network physics adds little to a fixed schedule's forward-evaluated cost and reveals no
hydraulic violation. That regime is not a corner case. It is
where the overwhelming majority of published district-heating dispatch models operate, and
the practical guidance the nulls support (do not spend modelling effort on hydraulic
detail before loss representation) applies wherever such models are used.
}

\DIFadd{What raises this above a plausible story is that one of the four channels has been opened
and measured the result. Freeing the supply temperature, the lever most likely to break the
hydraulic null, does break it: pipe velocity becomes the binding constraint at a 20}\DIFaddend \,\DIFdelbegin \DIFdel{km.
}\DIFdelend \DIFaddbegin \DIFadd{K
reduction and pumping energy rises by a factor of four
(Section~\ref{subsec:tsup}). Hydraulics move from a negligible }\emph{\DIFadd{cost}} \DIFadd{to a binding
}\emph{\DIFadd{constraint}} \DIFadd{the moment the degree of freedom exists. That is direct evidence for the
mechanism rather than an appeal to it, and it sets the expectation for the three channels
left unopened: each should convert a null into a constraint rather than into a large cost
term, because the physics enters dispatch as a feasibility boundary and not as a driver of
the schedule.
}\DIFaddend 

\DIFdelbegin \DIFdel{These resultssuggest that }%DIFDELCMD < \ce{CO2} %%%
\DIFdel{reporting frameworks applied
to district heating may need to specify a minimum
network-resolution level. Based on the present evidence, L2
resolution
appears warranted whenever the pipe length exceeds
1--2}\DIFdelend \DIFaddbegin \DIFadd{One asymmetry is worth stating plainly because it determines how far the null results can be
trusted. Relaxing these assumptions would generally give the extended physics }\emph{\DIFadd{more}} \DIFadd{room
to matter, although not necessarily monotonically. The null results should therefore be
interpreted as conservative }\emph{\DIFadd{within the tested regime}}\DIFadd{; outside that regime, they are
indicative rather than guaranteed. The loss-dominance result is more robust: the decomposition
is an exact algebraic identity based on the four optimised costs, and the loss term dominates
empirically across all 135 synthetic networks, which span two orders of magnitude in loss
burden. Its magnitude remains conditional on the tested systems, but within that regime it
does not depend on any single one of the four assumptions above. This result, rather than the
null results, therefore motivates the title.
}}

\subsection{Implications for modelling practice}
\DIFadd{The results translate into concrete guidance for building district-heating dispatch models.
Zeroth, and most concretely: the decision can be made before any model is built. Memmingen has a loss number of $\lambda=0.12$ from its pipe geometry, insulation classes, temperatures and annual demand, which the design rule converts to a predicted copperplate error of 11}\DIFaddend \,\DIFdelbegin \DIFdel{km, though this recommendation rests on a single industrial
case and a synthetic parametric study covering radial-tree topologies with
fixed heating curves. Broader empirical validationacross diverse
network types is needed before this finding can support binding
regulatory specifications. }\DIFdelend \DIFaddbegin \DIFadd{\%, above the threshold at which a lumped representation is likely safe, a screening indication that node resolution is worth checking here, though a known transferable loss adder can still suffice. The realised error is 15\,\%. A planner reaches that indication from pipe lengths, insulation classes, temperatures and an annual demand figure, without solving anything. First, what a model must capture is the network loss, not the network graph. A per-node or
per-link loss term recovers almost all of the bias that separating a copperplate from a
node-resolved model reveals, whereas resolving the graph without losses recovers almost none
of it; a modeller on a budget should spend it on loss representation before spatial detail.
Second, the tempting shortcut of a single lumped loss adder is safe only where its coefficient
is already known for the network and operating regime at hand: calibrated on one configuration
it mis-estimates the loss on another by tens of percentage points of cost, so it is a
diagnostic, not a substitute for a model that computes the loss endogenously. Third, and most
practically, estimation accuracy and control adequacy are different requirements and should be
stated separately. A loss-ignoring model can be used for neither (it both misprices and
misschedules), but the general lesson is sharper: a model may be adequate for scheduling
(small regret, deliverable) while being unfit for cost forecasting, CO$_2$ accounting or tariff
setting (large bias), and a modeller should say which use a given fidelity level supports rather
than reporting a single ``accuracy''. Fourth, fidelity should be claimed only to the resolution
at which it can be validated: a model resolved below its metering cannot have its extra detail
verified, so for most operators the node or zone level is the finest }\emph{\DIFadd{defensible}} \DIFadd{level,
not merely the cheapest. Finally, the physics beyond losses (trunk pressure, station-resolved
hydraulics, and even a flexible supply temperature) changes dispatch cost by little and,
evaluated forward, discloses no hydraulic violation in the setting studied; its main role is as a feasibility
}\emph{\DIFadd{constraint}} \DIFadd{(velocity, differential pressure) rather than a driver of the schedule, so it
belongs in a model as bounds to be respected, not as detail to be optimised over.
}\begin{table}[t]
\centering
\caption{What each fidelity level supports, stated as \emph{observations within the tested
regime} (fixed capacities, radial, central generation, hourly resolution) rather than universal
guidelines. ``Scheduling'' = small regret and physically deliverable; ``Estimation'' = usable for
cost forecasting, CO$_2$ accounting and tariffs (small bias).}
\label{tab:criteria}
\begin{tabular}{lcc}
\hline
Model & Scheduling & Estimation \\
\hline
Copperplate ($\CP$) & no (under-delivers loss) & no (biased) \\
Copperplate + loss ($\CPL$) & yes, \emph{if} the loss adder is known & only where the adder transfers \\
Node-resolved + loss ($\Lone$) & yes & yes \\
$+$ pressure / station ($\Lthree$--$\Lfive$) & yes (no material gain) & yes (no material gain) \\
\hline
\end{tabular}
\end{table}
\DIFaddend 


\subsection{Limitations}
\DIFdelbegin %DIFDELCMD < \label{sec:discussion_limitations}
%DIFDELCMD < 

%DIFDELCMD < %%%
\paragraph{\DIFdel{Validation and generalisation.}}
%DIFAUXCMD
\addtocounter{paragraph}{-1}%DIFAUXCMD
\DIFdel{Direct validation covers the pre-upgrade network (Stage~1);
}\DIFdelend \DIFaddbegin \DIFadd{The negative results, namely that extended physics adds little to the fixed schedule's
forward-evaluated cost, are conditional on the regime studied. Fixed capacities, a fixed
heating curve, precomputed heat-pump performance, parameterised losses, and hourly resolution
together limit how much the extended physics }\emph{\DIFadd{can}} \DIFadd{move a dispatch decision. A setting with
optimised supply temperature, sub-hourly transients, demand-side flexibility, or joint sizing
could give hydraulics and temperature more room to matter, and only the first of these, a
flexible supply temperature, has been bounded through forward evaluation. However, this regime
is not a corner case but the standard planning setting addressed by the paper: fixed-asset
dispatch on a radial, centrally supplied network at hourly resolution. The conclusions
therefore apply where such models are used, while the boundaries above identify where a
modeller should re-examine them. A second, data-driven limitation concerns validation. The
network measurements predate the heat pump, electrode boiler, and storage. Although the
transport physics validated by these data is asset-independent, the dispatch of the new units
is modelled using standard component relations and checked for plausibility rather than
validated against }\DIFaddend post-upgrade \DIFdelbegin \DIFdel{HP/TES dispatch is verified through the
necessary
conditions only. The 13\,\% headline gap is from oneindustrial
network; the synthetic parametric study provides qualitative external validity,
but absolute figures should not be treated as universal constants. }\DIFdelend \DIFaddbegin \DIFadd{metering. Because the conclusions are comparative, considering
bias and regret across fidelity levels for a fixed portfolio and corroborating them through the
synthetic factorial, they do not depend on a measurement-accurate absolute schedule. A
post-upgrade measurement campaign that directly validates the dispatch interactions would be
the natural next step; no such validation is claimed here.
}{
\paragraph{\DIFadd{Uncertainty, reserves and sub-hourly operation.}}
\DIFadd{The dispatch is deterministic and hourly. Price and demand are perfect-foresight time
series; forecast error, reserve provision, intra-hour demand excursions and operational
disturbances such as unit outages are all outside the formulation. Each of these plausibly
raises the value of the physics this paper finds inert. Transport delay is a zero-timestep
shift on an hourly grid and would not be at a five-minute one; thermal inertia in the pipe
volume buffers exactly the sub-hourly excursions excluded here; and storage and heat-pump
flexibility are worth more when they can be sold into reserve markets than when they can
only arbitrage a known price curve. The results should therefore be read as establishing
that extended network physics is decision-irrelevant }\emph{\DIFadd{in the deterministic hourly
planning regime}}\DIFadd{, the regime in which the great majority of published district-heating
dispatch models operate, and not as a claim about real-time control. The direction of
the bias is worth stating plainly: relaxing these assumptions would generally give the
extended physics more room to matter, though not necessarily monotonically, so the nulls
reported here are best read as conservative within the tested regime; the loss-dominance
algebraic identity does not depend on them, though its empirical magnitude is conditional on
the tested systems.
}\DIFaddend 

\DIFdelbegin \paragraph{\DIFdel{Deterministic price and demand assumptions.}}
%DIFAUXCMD
\addtocounter{paragraph}{-1}%DIFAUXCMD
\DIFdel{Dispatch here is evaluated under deterministic price and demand
scenarios}\DIFdelend \DIFaddbegin \DIFadd{A related caution applies to the estimation-bias result in the opposite direction}\DIFaddend .
\citet{Koek2026} show \DIFdelbegin \DIFdel{, for stochastic }\DIFdelend \DIFaddbegin \DIFadd{for }\DIFaddend multistage investment planning \DIFdelbegin \DIFdel{of DH generation capacity, }\DIFdelend that ignoring price and
demand-trajectory uncertainty\DIFdelbegin \DIFdel{(the Value of the Stochastic
Solution) }\DIFdelend \DIFaddbegin \DIFadd{, the value of the stochastic solution, }\DIFaddend biases technology
selection by 27--39\,\%\DIFdelbegin \DIFdel{; while their
}\DIFdelend \DIFaddbegin \DIFadd{. Their }\DIFaddend setting is investment \DIFdelbegin \DIFdel{planning }\DIFdelend rather than dispatch \DIFdelbegin \DIFdel{, and network
topology is not a variable in their framework, the result }\DIFdelend \DIFaddbegin \DIFadd{and does not
vary network topology, but the comparison }\DIFaddend suggests that the \DIFdelbegin \DIFdel{topology }\DIFdelend \DIFaddbegin \DIFadd{resolution }\DIFaddend bias quantified
here and the uncertainty\DIFdelbegin \DIFdel{-driven}\DIFdelend  bias they quantify are \DIFdelbegin \DIFdel{likely }\DIFdelend additive rather than substitutable\DIFdelbegin \DIFdel{,
and jointly addressing both remains an openextension}\DIFdelend \DIFaddbegin \DIFadd{. A
model can be wrong about the network and wrong about the future independently, and
correcting one does not correct the other; addressing both within a single formulation
remains open}\DIFaddend .

\paragraph{\DIFdelbegin \DIFdel{Scope restrictions}\DIFdelend \DIFaddbegin \DIFadd{Demand-side flexibility}\DIFaddend .}
\DIFdelbegin \DIFdel{The framework is dispatch-only (fixed capacities); joint
dispatch-and-sizing may interact with topology choice. Supply
temperature follows a heating curve (not optimised) to preserve
MILP linearity; variable-temperature operation could amplify the L3$\to\Lplus$ effect . \mbox{%DIFAUXCMD
\citet{Wirtz2021_tempcontrol} }\hskip0pt%DIFAUXCMD
show, for a
5GDHC network , that co-optimizing network temperature within an
MILP via temperature-level discretisation is tractable and
cost-reducing relative to free-floating temperature control;
applying this discretisation approach within the present
topology-isolated framework is a natural extension}\DIFdelend \DIFaddbegin \DIFadd{All demands are exogenous time series and demand response is excluded. This is one of the
few exclusions that could plausibly }\emph{\DIFadd{increase}} \DIFadd{the topology effect rather than leave
it unchanged, because spatially distributed flexibility has location-dependent value: a
flexible load near the network far end can absorb local surplus without incurring the
trunk losses that serving it from the plant would. Whether that is enough to lift routing
above the decision threshold is untested here and is a well-posed question for a follow-up
study}\DIFaddend .

\paragraph{\DIFdelbegin \DIFdel{Demand concentration metric}\DIFdelend \DIFaddbegin \DIFadd{Network topology class and generation siting}\DIFaddend .}
\DIFdelbegin \DIFdel{The normalised variance index HI (defined in
}%DIFDELCMD < \Cref{app:hi_definition}%%%
\DIFdel{) is used as a parametric study variable but has a narrow dynamic range on realistic tree networks (HI\,=\,0.05 at
Gini\,=\,0.41 for the primary case). The selection criteria
(}%DIFDELCMD < \Cref{tab:criteria}%%%
\DIFdel{), therefore, use Gini and top-$k$ share.
An
empirical mapping from these metrics to the L2$\to$L3 gap is a natural follow-up requiring multiple real networks. }%DIFDELCMD < 

%DIFDELCMD < %%%
\paragraph{\DIFdel{$\Lplus\!\to\!\Lnl$ interpretation.}}
%DIFAUXCMD
\addtocounter{paragraph}{-1}%DIFAUXCMD
\DIFdel{The $+0.35$ to $+0.50\,\%$ gap across two independent winter
windows combines the linearisation error and transport-delay
activation; the two contributions cannot be separated without an
intermediate MIQCP run with delay disabled. The bound is, therefore,
an upper limit on }\DIFdelend \DIFaddbegin \DIFadd{Every case, real and synthetic, is a radial tree with unidirectional flow. Ring and meshed
topologies are not merely a larger instance of the same problem: with pipe-flow direction
fixed by the pipe list, }\DIFaddend the \DIFdelbegin \DIFdel{combined effect, not a precise
decomposition}\DIFdelend \DIFaddbegin \DIFadd{return-side pressure balance around a cycle forces the sum of
the loop's friction drops to be non-positive, which is infeasible for any non-zero flow.
Admitting loops requires signed-flow hydraulic variables, a per-pipe direction selection
that reorients each pressure drop and temperature mix, which is a distinct formulation
rather than a parameter change. The same boundary limits the generation-topology
moderator: the synthetic generator roots all injection at the primary producer, so the
distributed-generation case is posed as a controlled open question rather than answered,
and multi-source operation on a meshed network is left to the companion study. Because
generation is central in every case examined, the null reported for spatial routing is
established for centrally-supplied networks specifically}\DIFaddend .

\paragraph{\DIFdelbegin \DIFdel{$\Lnl$ tractability}\DIFdelend \DIFaddbegin \DIFadd{The nonlinear reference}\DIFaddend .}
The \DIFdelbegin \DIFdel{MIQCP }\DIFdelend \DIFaddbegin \DIFadd{nonlinear }\DIFaddend reference is solved on \DIFdelbegin \DIFdel{two winter windows (October to
February); a spring window (May) proved infeasible due to
near-zero-flow degeneracy in bilinear constraints. The bound is therefore validated for peak-flow winter conditions. Whether the linearisation error }\DIFdelend \DIFaddbegin \DIFadd{representative 72-hour windows rather than globally. The
full-year non-convex programme finds no incumbent within the solver budget, and a low-load
summer 72-hour window is infeasible under native physics with the linearised schedule fixed. Both
are reported as results in Section~\ref{subsec:linearisation} (the second is a solver outcome
whose physical-versus-numerical cause is not yet established; an IIS diagnosis is left to future
work), but they do bound the claim: the solved linearisation error is established for winter
and autumn operating points, and whether it }\DIFaddend remains sub-percent under shoulder-season
\DIFdelbegin \DIFdel{operating points }\DIFdelend \DIFaddbegin \DIFadd{operation }\DIFaddend with frequent on/off cycling is \DIFdelbegin \DIFdel{an openquestion
best addressed through warm-start heuristics or reformulation}\DIFdelend \DIFaddbegin \DIFadd{open. Warm-start heuristics or a reformulation
of the bilinear terms are the natural routes}\DIFaddend .

\paragraph{\DIFdelbegin \DIFdel{Demand-side flexibility}\DIFdelend \DIFaddbegin \DIFadd{One real network, and the synthetic generator's assumptions}\DIFaddend .}
\DIFdelbegin \DIFdel{All demands are exogenous time series; demand response is
excluded. If flexible loads are modelled, the topology effect may
increase because spatially distributed flexibility has
location-dependent value (a flexible load near j$_{12}$ can
absorb local HP surplus without trunk losses) . This interaction
between demand flexibility and topology is relevant for real-time
dispatch. }\DIFdelend \DIFaddbegin \DIFadd{The Memmingen figures come from a single industrial network; the synthetic factorial
supplies external validity for the }\emph{\DIFadd{structure}} \DIFadd{of the result (loss dominance,
adder non-transfer, the loss-number rule) but not for the absolute magnitudes, which
should not be treated as universal constants. The synthetic networks also share the
assumptions of the generator that produced them: radial trees, a common heating curve,
proportionally sized capacity under a stated convention. They are a controlled parametric
family, not a sample of real networks, and the balanced factorial buys unambiguous
attribution across the design at the cost of that representativeness.
}\DIFaddend 

\DIFdelbegin \paragraph{\DIFdel{Network topology class.}}
%DIFAUXCMD
\addtocounter{paragraph}{-1}%DIFAUXCMD
\DIFdel{All cases (primary and synthetic ) are radial trees. Ring topologies
with bidirectional flow may amplify hydraulic coupling effects and shift the L3$\to\Lplus$ threshold; this remains an open question.
%DIF < % ============================================================
%DIF < % 7. CONCLUSION
%DIF < % ============================================================
}%DIFDELCMD < \section{Conclusion}
%DIFDELCMD < \label{sec:conclusion}
%DIFDELCMD < %%%
\DIFdelend \DIFaddbegin \paragraph{\DIFadd{Demand-concentration metric.}}
\DIFadd{The normalised variance index used as a synthetic design factor has a narrow dynamic range
on realistic tree networks, and the model-selection observations in
Table~\ref{tab:criteria} therefore use the Gini coefficient and top-$k$ demand share
instead. An empirical mapping from either metric to the aggregation gap would require
several real networks and is a natural follow-up.
}}
\DIFaddend 

\DIFdelbegin \DIFdel{This paper presents a comparison that systematically
isolates topology abstraction, physics fidelity, and linearisation
error as independent experimental variables in
district heating
dispatch optimisation.
This study varied three dimensions, one at a time, within a single experimental framework }\DIFdelend \DIFaddbegin \section{Conclusions}\label{sec:conclusion}
{
\DIFadd{This paper set out to determine which modelling choices actually change a district-heating
dispatch decision, rather than which change a reported objective value. To do so it
separates two quantities that the fidelity literature has treated as one (estimation
bias, the difference in optimised cost between formulations, and decision regret, the cost
of executing a formulation's schedule under a common forward model) on a four-cell
decomposition whose controls each isolate one phenomenon, extended by a physics-evaluation
framework for the finer levels. A copperplate supplied with the
aggregate loss it cannot itself compute serves as a control condition, so that loss
representation and spatial resolution can be switched independently and their contributions
separated exactly rather than estimated. The design was applied to one industrial network
with measured data and to a complete 135-cell balanced factorial of synthetic networks}\DIFaddend .

\DIFdelbegin \paragraph{\DIFdel{Principal findings}}
%DIFAUXCMD
\addtocounter{paragraph}{-1}%DIFAUXCMD
%DIFDELCMD < \begin{enumerate}[nosep]
%DIFDELCMD <   \item %%%
\textbf{\DIFdel{Topology effect (RQ1):}} %DIFAUXCMD
\DIFdel{The
cumulative L1$\to$L3
        gap reaches 13.0}\DIFdelend \DIFaddbegin \DIFadd{What a simplified model must capture is loss visibility, not spatial topology. The
copperplate-to-baseline cost gap decomposes to machine precision into a loss main effect of
}\textbf{\DIFadd{95.8}}\DIFaddend \,\%\DIFdelbegin \DIFdel{in cost and 74\,t\,}%DIFDELCMD < \ce{CO2}%%%
\DIFdel{/year.
        Across 36~synthetic networks, it scales monotonically
        with pipe length (3.2}\DIFdelend \DIFaddbegin \DIFadd{, while resolving spatial topology alone, with physics held constant between the
copperplate and node-resolved configurations, adds a negligible }\textbf{\DIFadd{0.25}}\DIFaddend \,\%\DIFdelbegin \DIFdel{at 1\,km, 15.8}\DIFdelend \DIFaddbegin \DIFadd{; the same holds
across the synthetic factorial, whose variants share physics by construction. The distinction is
decision-relevant rather than merely book-keeping: a copperplate understates cost by
}\textbf{\DIFadd{15.1}}\DIFaddend \,\% \DIFdelbegin \DIFdel{at 5\,km,
        39.0}\DIFdelend \DIFaddbegin \DIFadd{yet carries a }\textbf{\DIFadd{46.1}}\DIFaddend \,\% \DIFdelbegin \DIFdel{at 15\,km; 36/36 positive).
  }%DIFDELCMD < \item %%%
\textbf{\DIFdel{Physical driver:}} %DIFAUXCMD
\DIFdel{Thermal losses account for
        nearly the entire effect (median $+20.1\,\%$); pure
        routing ($-0.2\,\%$) and pressure drop ($+0.02\,\%$)
        are negligible. The
key insight is that topology matters
        not because of routing constraints but because it is the
        only mechanism making pipe losses visible to the optimiser. }%DIFDELCMD < \item %%%
\textbf{\DIFdel{Physics fidelity (RQ2):}}
        %DIFAUXCMD
\DIFdel{Extended hydraulic and thermal physics do not materially
        change dispatch when topology is already resolved: L3$\to\Lplus$ adds $+0.11\,\%$, attributable entirely
        to pumping power.
}%DIFDELCMD < \item %%%
\textbf{\DIFdel{Linearisation and transport delay (RQ3):}}
        %DIFAUXCMD
\DIFdel{The combined $\Lplus\to\Lnl$ gap is bounded at
        ${\leq}\,0.5\,\%$ based on the February window
        ($+0.50\,\%$, MIP gap 0.46\,\%); the January window
        ($+0.35\,\%$) is consistent but falls below solver
        tolerance (MIP gap 1.99\,\%) and
is therefore not
        independently significant.
Both windows confirm that $\Lnl$ cannot materially outperform $\Lplus$, supporting
        MILP as adequate for planning-accuracy dispatch.
A precise
        separation of linearisation error from transport-delay
        activation remains open.
  }%DIFDELCMD < \item %%%
\textbf{\DIFdel{Generalisability (RQ4):}} %DIFAUXCMD
\DIFdel{Pipe length is }\DIFdelend \DIFaddbegin \DIFadd{forward-valued schedule gap, representing its
unprovisioned loss valued at }\DIFaddend the \DIFdelbegin \DIFdel{dominant predictor; other parameters show weak
        within-band correlations.
}%DIFDELCMD < \end{enumerate}
%DIFDELCMD < 

%DIFDELCMD < %%%
\paragraph{\DIFdel{Recommendations}}
%DIFAUXCMD
\addtocounter{paragraph}{-1}%DIFAUXCMD
%DIFDELCMD < \begin{itemize}[nosep]
%DIFDELCMD <   \item %%%
\DIFdel{L2 as minimum for pipe
${>}\SI{1}{\kilo\metre}$ or
        storage present. }%DIFDELCMD < \item %%%
\DIFdel{L3 for spatially heterogeneous networks (Gini ${>}0.3$)
.
  }%DIFDELCMD < \item %%%
\DIFdel{$\Lplus$ when pipe ${>}\SI{5}{\kilo\metre}$ or $\tau > \SI{1}{\hour}$. }%DIFDELCMD < \item %%%
\DIFdel{$\Lnl$ as a one-time validation benchmark.
}%DIFDELCMD < \end{itemize}
%DIFDELCMD < %%%
\DIFdelend \DIFaddbegin \DIFadd{marginal unit, with the two effects carrying opposite signs. The
schedules show why: blind to the network losses, it provisions the heat pump for a demand that is
too small, running it 246 fewer hours and at a ten-percentage-point lower share, and the shortfall
is then covered at the margin in winter hours it never planned for. A lumped per-network loss
correction captures most of the bias, but only where its coefficient is already known: frozen and
carried to another network, it mis-estimates the loss burden by tens of percentage points of cost,
and even where it reproduces the cost, it still runs a materially different generation mix.
Pushed to the opposite extreme, with the demand side resolved to every transmission station and
its service laterals using manufacturer component data and an independent nonlinear cross-check,
hydraulic detail changes the forward-valued cost of the fixed schedule by under one percent and reveals no evaluated hydraulic violation, suggesting limited re-optimisation value.
}\DIFaddend 

\DIFdelbegin \paragraph{\DIFdel{Future work}}
%DIFAUXCMD
\addtocounter{paragraph}{-1}%DIFAUXCMD
\DIFdel{The research identified several continues research topics: storage-dispatch sensitivity on networks with longer transport delays; ring and meshed topologies; variable transport delay; supply-temperature optimisation ; joint dispatch-sizing comparison; full MINLP with native exponentials; stochastic variants; sub-hourly resolution; and 5th-generation bidirectional networks. Independent separation of the linearisation error from transport-delay activation via intermediate MIQCP runs is
a
priority follow-up}\DIFdelend \DIFaddbegin \DIFadd{For model selection this yields a criterion that can be applied before any model is built.
The share of cost a copperplate misses is $b = \lambda/(1+\lambda)$, where the loss number
$\lambda$ is the ratio of annual network loss to annual demand and follows from the pipe
inventory alone. Across all 136 networks examined (135 synthetic plus the industrial case)
this parameter-free relation reproduces the computed burden with $R^2 = \textbf{0.87}$, so the
screening decision (copperplate, calibrated adder, or resolved nodes) is read off from
$\lambda$ alone. Where a model must be resolved, effort belongs in the loss representation
before the network graph and well before the hydraulics; and fidelity should be claimed only to
the resolution at which it can be verified, which for operators metering at billing grade is
the node or zone level rather than the substation}\DIFaddend .

\DIFdelbegin \paragraph{\DIFdel{Closing remark}}
%DIFAUXCMD
\addtocounter{paragraph}{-1}%DIFAUXCMD
\DIFdel{A single-node model underestimates annual operating costs by
13\,\% and }%DIFDELCMD < \ce{CO2} %%%
\DIFdel{emissions by 8.7\,\% on the primary
case, a gap that scales from 3\,\% to 45\,\% across the
36~synthetic configurations, driven almost entirely by thermal losses made
visible through spatial resolution. Within the tested scope
(radialtrees, fixed heating curve}\DIFdelend \DIFaddbegin \DIFadd{These conclusions are bounded by the regime that produced them, and the boundaries are
specific rather than decorative. They hold for radial}\DIFaddend , \DIFdelbegin \DIFdel{dispatch-only), extended
hydraulic and thermal physics add less than 1\,\% on top, }\DIFdelend \DIFaddbegin \DIFadd{centrally-supplied networks with
fixed capacities, a prescribed heating curve }\DIFaddend and \DIFdelbegin \DIFdel{MILP linearisation introduces a further sub-percent margin on the validation windows.
The practical implication for this regime is direct: invest modelling effort in spatial resolution rather than
physics refinement. Whether the same hierarchy holds for ring
topologies, variable supply temperature, or joint
dispatch--sizing problems }\DIFdelend \DIFaddbegin \DIFadd{hourly deterministic dispatch. Each of
those assumptions closes a channel through which extended physics could act, so the nulls
reported here are conservative }\emph{\DIFadd{within the tested regime}}\DIFadd{: relaxing these assumptions would
generally give the physics more room to matter, though not necessarily monotonically, as the
supply-temperature study demonstrates by opening one of
them and turning hydraulics from a negligible cost into the binding constraint. Ring and
bidirectional topologies require a signed-flow formulation absent here; distributed
generation remains an open question because every case studied is centrally supplied; and
post-upgrade measurement of the new assets' dispatch is the natural next validation step.
The sharpest boundary is to the companion study: whether network resolution changes not only
how a network is dispatched but where new capacity is sited, under joint
dispatch-and-sizing, }\DIFaddend is a question \DIFdelbegin \DIFdel{for the follow-up
studies identified above.
}\DIFdelend \DIFaddbegin \DIFadd{this paper deliberately does not settle.
}}
\DIFaddend 


%DIF < % ============================================================
%DIF < % ACKNOWLEDGMENTS
%DIF < % ============================================================
\section*{Acknowledgments}
\DIFaddbegin 

\DIFaddend This work was funded by the German Federal Ministry for Economic Affairs and Energy
(BMWE, reference IIC2) under grant number \DIFdelbegin \DIFdel{~}\DIFdelend \textbf{03EN6057B} (Verbundvorhaben eProNe;
project period 2026--2028). Project management was carried out by Forschungszentrum
J\"ulich GmbH (PT-J.ESN6). Open Access funding was enabled and organized by Projekt DEAL.
The e-con AG in Memmingen provided operational data and domain expertise. The authors
express their gratitude for the support provided.

%DIF < % ============================================================
%DIF < % DATA AVAILABILITY
%DIF < % ============================================================
\DIFaddbegin \section*{\DIFadd{Declaration of competing interest}}


\DIFadd{The authors declare that they have no known competing financial interests or personal
relationships that could have appeared to influence the work reported in this paper.
}

\DIFaddend \section*{Data availability}

The optimisation model source code (Python/Pyomo), all \DIFdelbegin \DIFdel{YAML
}\DIFdelend configuration files, \DIFdelbegin \DIFdel{and }\DIFdelend \DIFaddbegin \DIFadd{the forward
evaluator and all }\DIFaddend post-processing scripts are available on Zenodo
\DIFdelbegin \DIFdel{at \mbox{%DIFAUXCMD
\citet{ruess_2026_calion} }\hskip0pt%DIFAUXCMD
under "}\DIFdelend \DIFaddbegin \DIFadd{\mbox{%DIFAUXCMD
\citep{ruess_2026_calion} }\hskip0pt%DIFAUXCMD
under a }\DIFaddend Creative Commons Attribution 4.0 International \DIFdelbegin \DIFdel{".
Input }\DIFdelend \DIFaddbegin \DIFadd{licence.
The release additionally contains the decomposition and regret
analyses, the complete 135-network synthetic generator and configuration set, and the
scripts producing every table and figure in this paper. The Memmingen input }\DIFaddend time series
are subject to \DIFdelbegin \DIFdel{NDA}\DIFdelend \DIFaddbegin \DIFadd{a non-disclosure agreement and are not released}\DIFaddend ; anonymised summary
statistics \DIFdelbegin \DIFdel{are included in the repository}\DIFdelend \DIFaddbegin \DIFadd{sufficient to characterise the demand and price signals are included. The
synthetic networks carry no such restriction, so the entire synthetic factorial is
independently reproducible}\DIFaddend .

%DIF < % ============================================================
%DIF < % AI Use
%DIF < % ============================================================
\section*{Declaration of generative AI and AI-assisted technologies in the manuscript
preparation process}

During the preparation of this work, the \DIFdelbegin \DIFdel{author(s) used }\DIFdelend \DIFaddbegin \DIFadd{authors used }\DIFaddend \textbf{Claude Opus \DIFdelbegin \DIFdel{4.7}\DIFdelend \DIFaddbegin \DIFadd{5}\DIFaddend } (Anthropic)
for language improvement\DIFdelbegin \DIFdel{and
}\DIFdelend \DIFaddbegin \DIFadd{, }\DIFaddend manuscript structuring, \DIFdelbegin \DIFdel{and }\textbf{\DIFdel{Claude Sonnet 4.6}} %DIFAUXCMD
\DIFdel{(Anthropic)
for }\DIFdelend \DIFaddbegin \DIFadd{manuscript editing, }\DIFaddend code generation and debugging. All content was reviewed and edited by the authors, who take full responsibility for the published article.


%DIF < % ============================================================
%DIF < % APPENDICES
%DIF < % ============================================================
\DIFdelbegin %DIFDELCMD < 

%DIFDELCMD < %%%
\DIFdelend \appendix
\DIFdelbegin %DIFDELCMD < \newpage
%DIFDELCMD < %%%
\DIFdelend 

\DIFdelbegin %DIFDELCMD < \section{Nomenclature}
%DIFDELCMD < \label{app:nomenclature}
%DIFDELCMD < %%%
\DIFdelend \DIFaddbegin \section{Hydraulic cross-check against a nonlinear pipe-flow solver}\label{app:pandapipes}
\DIFadd{As an independent check on the linearised trunk hydraulics, the same network, flows and pipe
geometry were solved with }\texttt{\DIFadd{pandapipes}}\DIFadd{, a nonlinear steady-state pipe-flow solver, and
the resulting trunk pressure drops compared pipe by pipe. The two agree to better than
$0.007$\,bar across the trunk, confirming that the piecewise-linear pressure model used in the
optimisation is not the source of the low pumping-energy estimate; the low pumping energy
follows from the network's compactness and heavily oversized pumps (Section~\ref{sec:methodology}),
not from a hydraulic modelling error. The station-internal and service-lateral pressure terms,
which }\texttt{\DIFadd{pandapipes}} \DIFadd{does not resolve, are added from the manufacturer datasheets and the
DXF-reconstructed lateral geometry.
}\DIFaddend 

{
\DIFdelbegin %DIFDELCMD < \footnotesize
%DIFDELCMD < 

%DIFDELCMD < %%%
%DIF < %--- Sets and indices ---
%DIFDELCMD < \noindent%%%
\textbf{\DIFdel{Sets and indices}}%DIFAUXCMD
%DIFDELCMD < \par\nopagebreak\smallskip
%DIFDELCMD < \noindent
%DIFDELCMD < \begin{tabular*}{\columnwidth}{@{\extracolsep{\fill}} >{\raggedright}p{2.0cm} l @{}}
%DIFDELCMD < \toprule
%DIFDELCMD < Symbol & Description \\
%DIFDELCMD < \midrule
%DIFDELCMD < $t \in \mathcal{T}$                & Time steps ($|\mathcal{T}|=8760$) \\
%DIFDELCMD < $n \in \mathcal{N}$                & Network nodes \\
%DIFDELCMD < $p \in \mathcal{P}$                & Pipe segments \\
%DIFDELCMD < $g \in \mathcal{G}$                & Thermal generators \\
%DIFDELCMD < $h \in \mathcal{H}$                & Heat pumps \\
%DIFDELCMD < $s \in \mathcal{S}$                & Thermal storage units \\
%DIFDELCMD < $r \in \mathcal{R}$                & Electrode boilers (power-to-heat) \\
%DIFDELCMD < $k \in \{1,\ldots,K\}$            & PWL segment index \\
%DIFDELCMD < $\mathcal{G}^{\mathrm{CHP}}$      & Subset of CHP generators \\
%DIFDELCMD < $\mathcal{P}_n^{\mathrm{in/out}}$  & Pipes entering/leaving node~$n$ \\
%DIFDELCMD < $\mathcal{G}_n,\mathcal{S}_n$     & Generators/storage at node~$n$ \\
%DIFDELCMD < $\mathcal{T}_{\mathrm{val}}$      & Validation time steps \\
%DIFDELCMD < \bottomrule
%DIFDELCMD < \end{tabular*}
%DIFDELCMD < 

%DIFDELCMD < \bigskip
%DIFDELCMD < 

%DIFDELCMD < %%%
%DIF < %--- Key parameters ---
%DIFDELCMD < \noindent%%%
\textbf{\DIFdel{Key parameters}}%DIFAUXCMD
%DIFDELCMD < \par\nopagebreak\smallskip
%DIFDELCMD < \noindent
%DIFDELCMD < \begin{tabular*}{\columnwidth}{@{\extracolsep{\fill}} >{\raggedright}p{2.0cm} >{\raggedright}p{2.0cm} l @{}}
%DIFDELCMD < \toprule
%DIFDELCMD < Symbol & Unit & Description \\
%DIFDELCMD < \midrule
%DIFDELCMD < $U_p$                        & W\,m$^{-1}$\,K$^{-1}$ & Pipe heat transfer coeff.\ \\
%DIFDELCMD < $L_p$                        & m              & Pipe length \\
%DIFDELCMD < $D_p$                        & m              & Pipe inner diameter \\
%DIFDELCMD < $A_p$                        & m$^2$          & Pipe cross-sectional area \\
%DIFDELCMD < $R_p$                        & Pa\,MW$^{-2}$  & Pipe hydraulic resistance \\
%DIFDELCMD < $f_D$                        & --             & Darcy friction factor \\
%DIFDELCMD < $\rho$                       & kg\,m$^{-3}$   & Fluid density \\
%DIFDELCMD < $c_p$                        & J\,(kg\,K)$^{-1}$ & Specific heat capacity \\
%DIFDELCMD < $\tau_p$                     & h              & Transport delay \\
%DIFDELCMD < $k_p$                        & --             & Delay in time steps \\
%DIFDELCMD < $\phi_{p,t}$                 & --             & Temp.\ decay factor \\
%DIFDELCMD < $\phi_p^{\mathrm{nom}}$      & --             & Nominal decay factor \\
%DIFDELCMD < $T_{\mathrm{sup}}^{\mathrm{nom}}(t)$ & $^\circ$C & Nominal supply temp.\ \\
%DIFDELCMD < $T_{\mathrm{ret}}^{\mathrm{nom}}$    & $^\circ$C & Nominal return temp.\ \\
%DIFDELCMD < $T^{\mathrm{nom}}(t)$        & $^\circ$C      & Nominal operating temp.\ \\
%DIFDELCMD < $T_{\mathrm{ground}}(t)$     & $^\circ$C      & Ground temperature \\
%DIFDELCMD < $\bar{Q}_g,\bar{Q}_h,\bar{Q}_r$ & MW         & Installed thermal capacities \\
%DIFDELCMD < $\bar{Q}_p$                  & MW             & Rated pipe capacity \\
%DIFDELCMD < $\phi_h^{\mathrm{min}}$      & --             & HP minimum load fraction \\
%DIFDELCMD < $\eta_g$                     & --             & Generator thermal efficiency \\
%DIFDELCMD < $\eta_s^{\mathrm{ch/dis}}$   & --             & Storage charge/discharge eff.\ \\
%DIFDELCMD < $\eta_{\mathrm{pump}}$       & --             & Pump efficiency \\
%DIFDELCMD < $\eta_{\mathrm{Lorenz}}$     & --             & Lorenz cycle efficiency \\
%DIFDELCMD < $\alpha_s$                   & h$^{-1}$       & Storage self-discharge rate \\
%DIFDELCMD < $\mathrm{COP}_{h,t}^{\mathrm{wrg}}$ & --     & Waste-heat COP (time-varying) \\
%DIFDELCMD < $\mathrm{COP}_h^{\mathrm{def}}$     & --     & Default-source COP (constant) \\
%DIFDELCMD < $m_k^{(p)}$                  & --             & PWL segment slope \\
%DIFDELCMD < $\Delta T_{\mathrm{pp}}$     & K              & Heat pump pinch-point temp.\ \\
%DIFDELCMD < $\lambda_t^{\mathrm{buy/sell}}$ & \euro\,MWh$^{-1}$ & Electricity buy/sell price \\
%DIFDELCMD < $\lambda_g^{\mathrm{fuel}}$  & \euro\,MWh$^{-1}$ & Specific fuel cost \\
%DIFDELCMD < $\lambda^{\mathrm{dump}}$    & \euro\,MWh$^{-1}$ & Heat curtailment penalty \\
%DIFDELCMD < $\lambda^{\ce{CO2}}$         & \euro\,t$^{-1}$ & Carbon price \\
%DIFDELCMD < $\lambda^{\mathrm{dem}}$     & \euro\,MW$^{-1}$ & Demand charge rate \\
%DIFDELCMD < $e_t^{\mathrm{grid}}$        & kg\,MWh$^{-1}$ & Grid emission factor \\
%DIFDELCMD < $e_g^{\mathrm{fuel}}$        & kg\,MWh$^{-1}$ & Fuel emission intensity \\
%DIFDELCMD < $f^{\mathrm{yr}}$            & --             & Annual scaling factor \\
%DIFDELCMD < $M^{\mathrm{grid}}$          & MW             & Grid capacity (Big-M) \\
%DIFDELCMD < $\Delta t$                   & h              & Time step length ($= 1$\,h) \\
%DIFDELCMD < HI                           & --             & Demand heterog.\ index \\
%DIFDELCMD < $K$                          & --             & Number of PWL segments \\
%DIFDELCMD < \bottomrule
%DIFDELCMD < \end{tabular*}
%DIFDELCMD < 

%DIFDELCMD < \bigskip
%DIFDELCMD < 

%DIFDELCMD < %%%
%DIF < %--- Cost components ---
%DIFDELCMD < \noindent%%%
\textbf{\DIFdel{Cost components}}%DIFAUXCMD
%DIFDELCMD < \par\nopagebreak\smallskip
%DIFDELCMD < \noindent
%DIFDELCMD < \begin{tabular*}{\columnwidth}{@{\extracolsep{\fill}} >{\raggedright}p{2.0cm} l @{}}
%DIFDELCMD < \toprule
%DIFDELCMD < Symbol & Description \\
%DIFDELCMD < \midrule
%DIFDELCMD < $C^{\mathrm{en}}$    & Electricity procurement net cost \\
%DIFDELCMD < $C^{\mathrm{fuel}}$  & Fuel expenditure \\
%DIFDELCMD < $C^{\mathrm{dump}}$  & Heat curtailment penalty \\
%DIFDELCMD < $C^{\ce{CO2}}$       & Carbon cost \\
%DIFDELCMD < $C^{\mathrm{dem}}$   & Demand charge (peak import) \\
%DIFDELCMD < $C^{\mathrm{pump}}$  & Network pumping cost \\
%DIFDELCMD < $Z$                  & Total annual operational cost \\
%DIFDELCMD < \bottomrule
%DIFDELCMD < \end{tabular*}
%DIFDELCMD < 

%DIFDELCMD < \bigskip
%DIFDELCMD < 

%DIFDELCMD < %%%
%DIF < %--- Key decision variables ---
%DIFDELCMD < \noindent%%%
\textbf{\DIFdel{Key decision variables}}%DIFAUXCMD
%DIFDELCMD < \par\nopagebreak\smallskip
%DIFDELCMD < \noindent
%DIFDELCMD < \begin{tabular*}{\columnwidth}{@{\extracolsep{\fill}} >{\raggedright}p{2.0cm} >{\raggedright}p{2.0cm} l @{}}
%DIFDELCMD < \toprule
%DIFDELCMD < Symbol & Unit & Description \\
%DIFDELCMD < \midrule
%DIFDELCMD < $Q_{g,t}^{\mathrm{th}}$       & MW       & Generator thermal output \\
%DIFDELCMD < $Q_{h,t}$                     & MW       & Heat pump thermal output \\
%DIFDELCMD < $Q_{h,t}^{\mathrm{wrg/def}}$  & MW       & HP waste-heat/default output \\
%DIFDELCMD < $Q_{r,t}$                     & MW       & Electrode boiler output \\
%DIFDELCMD < $Q_{p,t}$                     & MW       & Pipe heat flow \\
%DIFDELCMD < $Q_{p,t}^{\mathrm{loss,pipe}}$& MW       & Pipe thermal loss \\
%DIFDELCMD < $Q_{s,t}^{\mathrm{ch/dis}}$   & MW       & Storage charge/discharge rate \\
%DIFDELCMD < $Q_{n,t}^{\mathrm{dem}}$      & MW       & Nodal heat demand \\
%DIFDELCMD < $Q_{n,t}^{\mathrm{dump}}$     & MW       & Curtailed heat at node~$n$ \\
%DIFDELCMD < $P_t^{\mathrm{buy/sell}}$     & MW       & Grid import/export \\
%DIFDELCMD < $P_{g,t}^{\mathrm{el}}$       & MW       & CHP electrical output \\
%DIFDELCMD < $P_{h,t}^{\mathrm{el}}$       & MW       & Heat pump electrical input \\
%DIFDELCMD < $P_{r,t}^{\mathrm{el}}$       & MW       & Electrode boiler electrical input \\
%DIFDELCMD < $P_t^{\mathrm{pump}}$         & MW       & Network pumping power \\
%DIFDELCMD < $P^{\mathrm{peak}}$           & MW       & Annual peak grid import \\
%DIFDELCMD < $\Delta p_{p,t}$              & Pa       & Pipe pressure drop \\
%DIFDELCMD < $\dot{m}_{p,t}$               & kg\,s$^{-1}$ & Pipe mass flow rate \\
%DIFDELCMD < $T_{n,t}^{\mathrm{sup}}$      & $^\circ$C & Node supply temperature \\
%DIFDELCMD < $T_{n_1,t}^{\mathrm{sup}}$    & $^\circ$C & Upstream node supply temp.\ \\
%DIFDELCMD < $T_{p,t}^{\mathrm{out}}$      & $^\circ$C & Pipe outlet temperature \\
%DIFDELCMD < $E_{s,t}$                     & MWh      & Storage state of charge \\
%DIFDELCMD < $F_{g,t}$                     & MW       & Fuel input (HHV basis) \\
%DIFDELCMD < $E_t^{\ce{CO2},\mathrm{grid}}$ & t       & Grid \ce{CO2} emissions \\
%DIFDELCMD < $E_{g,t}^{\ce{CO2},\mathrm{fuel}}$ & t   & Fuel \ce{CO2} emissions \\
%DIFDELCMD < $W_{p,t}$                     & MW$^2$   & Bilinear auxiliary ($= Q_{p,t}^2$) \\
%DIFDELCMD < $u_{h,t}$                     & $\{0,1\}$ & Heat pump on/off \\
%DIFDELCMD < $z_t$                         & $\{0,1\}$ & Grid import mode \\
%DIFDELCMD < $w_{p,t,k}$                   & $[0,1]$  & SOS2 weights \\
%DIFDELCMD < \bottomrule
%DIFDELCMD < \end{tabular*}
%DIFDELCMD < 

%DIFDELCMD < }%%%
%DIF <  end footnotesize
%DIFDELCMD < 

%DIFDELCMD < %%%
%DIF < % ============================================================
\DIFdelend \section{Model formulation details}
\label{app:model_details}

\DIFdelbegin %DIFDELCMD < \subsection{COP pre-computation}
%DIFDELCMD < %%%
\DIFdelend \DIFaddbegin \subsection{Heat-pump performance pre-computation}
\DIFaddend \label{app:cop}

The \DIFdelbegin \DIFdel{heat pump COP }\DIFdelend \DIFaddbegin \DIFadd{waste-heat coefficient of performance }\DIFaddend is pre-computed from an analytical Lorenz-cycle
model\DIFdelbegin \DIFdel{:
}\DIFdelend \DIFaddbegin \DIFadd{,
}\DIFaddend \begin{equation}
  \mathrm{COP}_{h,t}^{\mathrm{WHR}} = \eta_{\mathrm{Lorenz}}
  \cdot \frac{T_{\mathrm{sink,out}}}{T_{\mathrm{sink,out}} -
  T_{\mathrm{src},h}(t) + \Delta T_{\mathrm{pp}}}\DIFaddbegin \DIFadd{,
}\DIFaddend \end{equation}
with \DIFdelbegin \DIFdel{$\eta_{\mathrm{Lorenz}} = 0.6$ and $\Delta T_{\mathrm{pp}} = \SI{5}{\kelvin}$.  Combined error :
$\varepsilon_{\mathrm{COP}} \approx 3.6\,\%$}\DIFdelend \DIFaddbegin \DIFadd{all temperatures in kelvin, a Lorenz efficiency of $0.6$ and a pinch-point difference of $5$\,K, giving a combined
approximation error of about $3.6$\,\%. Pre-computing rather than co-optimising the
coefficient of performance keeps the formulation free of bilinear terms; it also decouples
heat-pump operation from the network state, which is one of the assumptions discussed in
Section~\ref{subsec:whynull} as limiting how far the extended physics can move a dispatch
decision}\DIFaddend .

\subsection{Linear component models}
\label{app:component_models}

{%
\DIFdelbegin %DIFDELCMD < \captionof{table}{Linear component models (all levels).}
%DIFDELCMD < %%%
\DIFdelend \DIFaddbegin \captionof{table}{Linear component models, applied identically at every fidelity level.}
\DIFaddend \label{tab:simple_components}
\centering
\footnotesize
\DIFdelbegin %DIFDELCMD < \begin{tabular*}{\columnwidth}{@{\extracolsep{\fill}} l l l @{}}
%DIFDELCMD < \toprule
%DIFDELCMD < Component & Equation & Bound \\
%DIFDELCMD < \midrule
%DIFDELCMD < Gas boiler    & $Q_{g,t}^{\mathrm{th}} = \eta_g F_{g,t}$ & $\leq \bar{Q}_g$ \\
%DIFDELCMD < Biomass       & $Q_{g,t}^{\mathrm{th}} = \eta_g F_{g,t}$ & $\leq \bar{Q}_g$ \\
%DIFDELCMD < Electrode     & $Q_{r,t} = \eta_r P_{r,t}^{\mathrm{el}}$ & $\leq \bar{Q}_r$ \\
%DIFDELCMD < CHP           & $Q^{\mathrm{th}} = \eta^{\mathrm{th}} F$  & $\sigma_g = \mathrm{const}$ \\
%DIFDELCMD < \bottomrule
%DIFDELCMD < \end{tabular*}
%DIFDELCMD < %%%
\DIFdelend \DIFaddbegin \begin{tabular*}{\columnwidth}{@{\extracolsep{\fill}} l l l @{}}
\toprule
Component & Equation & Bound \\
\midrule
Gas boiler       & $Q_{g,t}^{\mathrm{th}} = \eta_g F_{g,t}$ & $\leq \bar{Q}_g$ \\
Biomass boiler   & $Q_{g,t}^{\mathrm{th}} = \eta_g F_{g,t}$ & $\leq \bar{Q}_g$ \\
Electrode boiler & $Q_{r,t} = \eta_r P_{r,t}^{\mathrm{el}}$ & $\leq \bar{Q}_r$ \\
CHP              & $Q^{\mathrm{th}} = \eta^{\mathrm{th}} F$ & $\sigma_g$ constant \\
\bottomrule
\end{tabular*}
\DIFaddend }% end centering group

\bigskip

\DIFdelbegin %DIFDELCMD < \subsection{Extended physics: $\Lplus$ vs.\ $\Lnl$}
%DIFDELCMD < \label{app:extended_formulations}
%DIFDELCMD < 

%DIFDELCMD < {%%%
%DIF < 
%DIFDELCMD < \captionof{table}{Mathematical treatment of extended physics.}
%DIFDELCMD < \label{tab:lplus_vs_lnl}
%DIFDELCMD < \centering
%DIFDELCMD < \footnotesize
%DIFDELCMD < \begin{threeparttable}
%DIFDELCMD < \begin{tabular*}{\columnwidth}{@{\extracolsep{\fill}} l l l @{}}
%DIFDELCMD < \toprule
%DIFDELCMD < Phenomenon & $\Lplus$ & $\Lnl$ \\
%DIFDELCMD < \midrule
%DIFDELCMD < Pressure drop   & PWL ($K{=}3$)     & Quad.\ $RQ^2$ \\
%DIFDELCMD < Pumping power   & PWL of cubic      & Bilin.\ $Q{\cdot}W$ \\
%DIFDELCMD < Temp.\ decay    & PWL exp.\tnote{*} & PWL exp.\tnote{*} \\
%DIFDELCMD < Temp.\ mixing   & Taylor lin.       & Bilin.\ $T{\cdot}\phi$ \\
%DIFDELCMD < Transport delay & Bypassed          & Active ($k_p$) \\
%DIFDELCMD < Optimality      & Global (LP)       & Global (branch.) \\
%DIFDELCMD < \bottomrule
%DIFDELCMD < \end{tabular*}
%DIFDELCMD < \begin{tablenotes}\scriptsize
%DIFDELCMD <   \item[*] %%%
\DIFdelFL{Shared; excluded from linearisation error.
}%DIFDELCMD < \end{tablenotes}
%DIFDELCMD < \end{threeparttable}
%DIFDELCMD < }%%%
%DIF <  end centering group
%DIFDELCMD < 

%DIFDELCMD < \bigskip
%DIFDELCMD < %%%
\DIFdelend \DIFaddbegin \DIFadd{Because these models are identical at every level, they cannot contribute to any
level-to-level difference. This is what allows the contrasts of
Table~\ref{tab:contrasts} to be attributed to the network representation alone.
}\DIFaddend 

\DIFdelbegin %DIFDELCMD < \subsection{Linearisation error bounds}
%DIFDELCMD < %%%
\DIFdelend \DIFaddbegin \subsection{Piecewise-linearisation error bound}
\DIFaddend \label{app:pwl}

For \DIFaddbegin \DIFadd{the quadratic pressure relation }\DIFaddend $\Delta p = R Q^2$ \DIFdelbegin \DIFdel{with }\DIFdelend \DIFaddbegin \DIFadd{approximated over }\DIFaddend $K$ equal
\DIFdelbegin \DIFdel{PWL segments:
}\DIFdelend \DIFaddbegin \DIFadd{segments, the maximum approximation error is
}\DIFaddend \begin{equation}
  \varepsilon_{\Delta p}^{\max} = \frac{R \bar{Q}^2}{4K^2}\DIFaddbegin \DIFadd{,
}\DIFaddend \end{equation}
\DIFdelbegin \DIFdel{With $K{=}3$: 2.8}\DIFdelend \DIFaddbegin \DIFadd{which at $K = 3$ is $2.8$}\DIFaddend \,\% of \DIFaddbegin \DIFadd{the }\DIFaddend peak pressure drop. The \DIFdelbegin \DIFdel{cost gap
$\Delta Z$ is well below $\varepsilon^{\max}$ }\DIFdelend \DIFaddbegin \DIFadd{resulting cost deviation is far
smaller than this bound, }\DIFaddend because the optimiser compensates by adjusting \DIFdelbegin \DIFdel{dispatch }\DIFdelend \DIFaddbegin \DIFadd{the schedule }\DIFaddend rather
than by violating feasibility\DIFdelbegin \DIFdel{; the PWL error thus }\DIFdelend \DIFaddbegin \DIFadd{: the approximation error }\DIFaddend manifests as a minor scheduling
shift, not as a proportional cost deviation. \DIFaddbegin \DIFadd{The cost deviation itself is not estimated from
this bound but measured directly, by re-solving the continuous problem under native
non-linearised physics with the mixed-integer schedule fixed
(Section~\ref{subsec:linearisation}).
}\DIFaddend 

\medskip
{%
\DIFdelbegin %DIFDELCMD < \captionof{table}{PWL approximation error and problem size
%DIFDELCMD < vs.\ number of segments $K$ (primary case).}
%DIFDELCMD < %%%
\DIFdelend \DIFaddbegin \captionof{table}{Piecewise-linear approximation error and problem size against the number
of segments $K$. The cost effect is not inferred from this table; it is measured by the
solved nonlinear reference in Section~\ref{subsec:linearisation}.}
\DIFaddend \label{tab:pwl_sensitivity}
\centering
\footnotesize
\begin{threeparttable}
\DIFdelbegin %DIFDELCMD < \begin{tabular*}{\columnwidth}{@{\extracolsep{\fill}} c c r l @{}}
%DIFDELCMD < \toprule
%DIFDELCMD < $K$ & $\varepsilon^{\max}$ & Vars\tnote{a} & $\Delta Z$ \\
%DIFDELCMD < \midrule
%DIFDELCMD < 3 & 2.8\,\% & 753k  & $+0.35\,\%$ \\
%DIFDELCMD < 5 & 1.0\,\% & 998k\tnote{b}  & -- \\
%DIFDELCMD < 8 & 0.4\,\% & 1366k\tnote{b} & -- \\
%DIFDELCMD < \bottomrule
%DIFDELCMD < \end{tabular*}
%DIFDELCMD < %%%
\DIFdelend \DIFaddbegin \begin{tabular*}{\columnwidth}{@{\extracolsep{\fill}} c c r @{}}
\toprule
$K$ & $\varepsilon^{\max}$ & Additional variables\tnote{a} \\
\midrule
3 & 2.8\,\% & 753k \\
5 & 1.0\,\% & 998k\tnote{b} \\
8 & 0.4\,\% & 1366k\tnote{b} \\
\bottomrule
\end{tabular*}
\DIFaddend \begin{tablenotes}\scriptsize
  \item[a] \DIFdelbeginFL \DIFdelFL{Total additional variables relative to L3
    (binary + continuous )}\DIFdelendFL \DIFaddbeginFL \DIFaddFL{Relative to the node-resolved baseline, binary and continuous combined}\DIFaddendFL .
  \item[b] Extrapolated \DIFdelbeginFL \DIFdelFL{from the per-segment increment
    (${\approx}\,122$k }\DIFdelendFL \DIFaddbeginFL \DIFaddFL{at approximately 122k }\DIFaddendFL variables per additional segment\DIFdelbeginFL \DIFdelFL{)}\DIFdelendFL ; not
    simulated.
\end{tablenotes}
\end{threeparttable}
}% end centering group

\DIFdelbegin %DIFDELCMD < \subsection{Taylor linearisation}
%DIFDELCMD < %%%
\DIFdelend \DIFaddbegin \subsection{Taylor linearisation of the temperature product}
\DIFaddend \label{app:taylor_derivation}

The bilinear product \DIFdelbegin \DIFdel{$T_{n_1} \cdot \phi_p$ is expanded around
}\DIFdelend \DIFaddbegin \DIFadd{of upstream temperature and decay factor is expanded about the
nominal operating point }\DIFaddend $(T_{n_1}^0, \phi_p^0)$:
\begin{equation}
  T_{n_1} \DIFdelbegin \DIFdel{\cdot }\DIFdelend \phi_p \approx T_{n_1}^0 \phi_p^0
  + \phi_p^0 (T_{n_1} - T_{n_1}^0)
  + T_{n_1}^0 (\phi_p - \phi_p^0)\DIFaddbegin \DIFadd{.
}\DIFaddend \end{equation}
\DIFdelbegin \DIFdel{Neglected }\DIFdelend \DIFaddbegin \DIFadd{The neglected second-order }\DIFaddend cross-term \DIFdelbegin \DIFdel{:
$|\varepsilon| \leq |\Delta T_{\max}| \cdot |\Delta\phi_{\max}| < 2.5\,$K .
}%DIFDELCMD < 

%DIFDELCMD < \subsection{Electricity selling price}
%DIFDELCMD < %%%
\DIFdelend \DIFaddbegin \DIFadd{is bounded by
$|\varepsilon| \leq |\Delta T_{\max}| |\Delta\phi_{\max}| < 2.5$\,K over the operating
envelope. The exponential decay factor is approximated piecewise-linearly only in the
PWL-linearised temperature-propagating level ($\Ltwo$); the nonlinear reference treats it
natively (Section~\ref{subsec:evaluator}). The measured linearisation gap of
$-0.15$/$-0.33$\,\% (Section~\ref{subsec:linearisation}) therefore already contains this
exponential-approximation component (native versus piecewise-linear) rather than cancelling
it, so the reported figure is the combined PWL-versus-native gap, not an isolated one.
}}
{
\section{Electricity selling price}
\DIFaddend \label{app:selling_price}

\DIFaddbegin \DIFadd{The price received for exported electricity applies the spot price less a marketing spread,
floored, net of a self-consumption share and a fixed fee, and cannot go negative:
}\DIFaddend \begin{equation}
  \lambda_t^{\mathrm{sell}} = \max\!\Bigl(
    \max\bigl(\lambda_t^{\mathrm{spot}} - \lambda^{\mathrm{spread}},\;
      \lambda^{\mathrm{floor}}\bigr)
    (1 - h^{\mathrm{sell}})
    - \lambda^{\mathrm{sell,fee}},\; 0\Bigr)\DIFaddbegin \DIFadd{.
}\DIFaddend \end{equation}

\DIFdelbegin %DIFDELCMD < \subsection{Per-node MAE breakdown}
%DIFDELCMD < %%%
\DIFdelend \DIFaddbegin \section{Per-node temperature error breakdown}
\DIFaddend \label{app:per_node_mae}

\DIFdelbegin %DIFDELCMD < \Cref{tab:per_node_mae} %%%
\DIFdelend \DIFaddbegin \DIFadd{Table~\ref{tab:per_node_mae} }\DIFaddend reports supply-temperature \DIFdelbegin \DIFdel{MAE and bias
for all 14~}\DIFdelend \DIFaddbegin \DIFadd{error at all fourteen }\DIFaddend network nodes
over \DIFdelbegin \DIFdel{the 744\,h winter validation window, }\DIFdelend \DIFaddbegin \DIFadd{a 744-hour winter window, evaluated on the temperature-propagating formulation and
}\DIFaddend sorted by path length from the source. Three \DIFdelbegin \DIFdel{node }\DIFdelend groups are distinguished\DIFdelbegin \DIFdel{:
%DIF < 
}%DIFDELCMD < \begin{itemize}[nosep]
%DIFDELCMD <   \item %%%
\textbf{\DIFdel{Validation nodes}} %DIFAUXCMD
\DIFdel{(val., 6~nodes) : consumer sensors with }\DIFdelend \DIFaddbegin \DIFadd{, and the distinction
matters for how the aggregate figures in Section~\ref{subsec:validation} should be read.
}

\emph{\DIFadd{Validation nodes}} \DIFadd{(six) are consumer sensors whose }\DIFaddend mixing-valve offset \DIFdelbegin \DIFdel{${<}\SI{2}{\celsius}$,
        independent of the calibrationprocedure.  Mean MAE
        }%DIFDELCMD < \SI{1.32}{\celsius}%%%
\DIFdel{; these are the nodes that enter the
        gate evaluation in }%DIFDELCMD < \Cref{tab:validation_stage1}%%%
\DIFdel{.
}%DIFDELCMD < \item %%%
\textbf{\DIFdel{Calibration nodes}} %DIFAUXCMD
\DIFdel{(cal., 4~nodes) : }\DIFdelend \DIFaddbegin \DIFadd{is below
$2$\,°C and which took no part in the calibration; their mean error is $1.32$\,°C.
}\emph{\DIFadd{Calibration nodes}} \DIFadd{(four) were }\DIFaddend used in the branch-sequential \DIFdelbegin \DIFdel{$U_p$-tuning of
        \mbox{%DIFAUXCMD
\citet{Maldonado2024}}\hskip0pt%DIFAUXCMD
; }\DIFdelend \DIFaddbegin \DIFadd{coefficient tuning, so
}\DIFaddend their residuals are in-sample by construction\DIFdelbegin \DIFdel{.  Mean MAE }%DIFDELCMD < \SI{2.32}{\celsius}%%%
\DIFdel{.  Node j$_7$ at
}%DIFDELCMD < \SI{3.07}{\celsius} %%%
\DIFdel{would fail the worst-node gate of
        }%DIFDELCMD < \SI{2.5}{\celsius} %%%
\DIFdel{if held out for validation. This
        }\DIFdelend \DIFaddbegin \DIFadd{; their mean is $2.32$\,°C, and one node at
$3.07$\,°C would fail a $2.5$\,°C held-out criterion. That }\DIFaddend in-sample deterioration is
expected\DIFdelbegin \DIFdel{and motivates the }\DIFdelend \DIFaddbegin \DIFadd{, and it is the reason a }\DIFaddend strict calibration/validation split \DIFdelbegin \DIFdel{.
  }%DIFDELCMD < \item %%%
\textbf{\DIFdel{Excluded nodes}} %DIFAUXCMD
\DIFdel{(excl., 4~nodes) : excluded }\DIFdelend \DIFaddbegin \DIFadd{matters.
}\emph{\DIFadd{Excluded nodes}} \DIFadd{(four) are set aside }\DIFaddend for measurement reasons unrelated to model
quality\DIFdelbegin \DIFdel{.  Nodes
        j$_3$ (}%DIFDELCMD < \SI{0.58}{\celsius}%%%
\DIFdel{) and j$_4$
        (}%DIFDELCMD < \SI{0.37}{\celsius}%%%
\DIFdel{) are the two best-performing nodes
        in the entire network; j$_4$ has }\DIFdelend \DIFaddbegin \DIFadd{: two exhibit mixing-valve absorption above the $2$\,°C threshold and cannot be read
as junction temperatures, one is a boundary-condition reference, and one serves }\DIFaddend four
consumers with a \DIFdelbegin %DIFDELCMD < \SI{20.9}{\celsius} %%%
\DIFdelend \DIFaddbegin \DIFadd{$20.9$\,°C }\DIFaddend intra-node spread, so no single \DIFdelbegin \DIFdel{junction-reference sensor exists.  j$_3$ is a
        boundary-condition reference.  j$_2$ and j$_5$ exhibit
        mixing-valve absorption above the }%DIFDELCMD < \SI{2}{\celsius}
%DIFDELCMD <         %%%
\DIFdel{threshold and are therefore not interpretable as junction temperatures.
}%DIFDELCMD < \end{itemize}
%DIFDELCMD < %%%
%DIF < 
\DIFdel{The all-14-node mean MAE is }%DIFDELCMD < \SI{1.68}{\celsius}%%%
\DIFdel{, which sits
}\DIFdelend \DIFaddbegin \DIFadd{junction reference exists.
}

\DIFadd{The mean across all fourteen nodes is $1.68$\,°C, }\DIFaddend between the validation-node mean \DIFdelbegin \DIFdel{(}%DIFDELCMD < \SI{1.32}{\celsius}%%%
\DIFdel{) }\DIFdelend \DIFaddbegin \DIFadd{of
$1.32$\,°C }\DIFaddend and the calibration-node mean \DIFdelbegin \DIFdel{(}%DIFDELCMD < \SI{2.32}{\celsius}%%%
\DIFdel{) and confirms that the
model performance is not artificially inflated by selective node
choice: the }\DIFdelend \DIFaddbegin \DIFadd{of $2.32$\,°C. The }\DIFaddend validation nodes are \DIFaddbegin \DIFadd{therefore
}\DIFaddend not systematically the best-performing ones \DIFaddbegin \DIFadd{(two of the excluded nodes are in fact the
best in the network), so the group means are not an artefact of favourable node selection}\DIFaddend .

\medskip
{%
\DIFdelbegin %DIFDELCMD < \captionof{table}{Per-node supply-temperature MAE and bias
%DIFDELCMD < (744\,h winter window, $\Lnl$).}
%DIFDELCMD < %%%
\DIFdelend \DIFaddbegin \captionof{table}{Per-node supply-temperature error over a 744\,h winter window, evaluated
on the temperature-propagating formulation.}
\DIFaddend \label{tab:per_node_mae}
\centering
\footnotesize
\begin{threeparttable}
\DIFdelbegin %DIFDELCMD < \begin{tabular*}{\columnwidth}{@{\extracolsep{\fill}} l c c c c c @{}}
%DIFDELCMD < \toprule
%DIFDELCMD < Node & Path [m] & Role & MAE [$^\circ$C] & Bias [$^\circ$C] & In gate? \\
%DIFDELCMD < \midrule
%DIFDELCMD < j$_2$    &  350 & excl. & 2.99 & $+2.98$ & No  \\
%DIFDELCMD < j$_3$    &  650 & excl. & 0.58 & $+0.07$ & No  \\
%DIFDELCMD < j$_4$    &  910 & excl. & 0.37 & $-0.18$ & No  \\
%DIFDELCMD < j$_9$    & 1100 & val.  & 0.87 & $+0.41$ & Yes \\
%DIFDELCMD < j$_5$    & 1150 & excl. & 2.35 & $+1.51$ & No  \\
%DIFDELCMD < j$_6$    & 1300 & cal.  & 1.62 & $-1.62$ & No  \\
%DIFDELCMD < j$_{10}$ & 1300 & val.  & 2.27 & $+1.78$ & Yes \\
%DIFDELCMD < j$_7$    & 1370 & cal.  & 3.07\tnote{a} & $+2.53$ & No  \\
%DIFDELCMD < j$_8$    & 1450 & cal.  & 2.29 & $+1.00$ & No  \\
%DIFDELCMD < j$_{11}$ & 1530 & val.  & 1.41 & $+0.71$ & Yes \\
%DIFDELCMD < j$_{12}$ & 1780 & val.  & 0.69 & $-0.29$ & Yes \\
%DIFDELCMD < j$_{13}$ & 2000 & val.  & 1.09 & $-1.02$ & Yes \\
%DIFDELCMD < j$_{15}$ & 2125 & val.  & 1.58 & $-0.24$ & Yes \\
%DIFDELCMD < j$_{14}$ & 2180 & cal.  & 2.31 & $+1.02$ & No  \\
%DIFDELCMD < \midrule
%DIFDELCMD < \multicolumn{3}{@{}l}{Group mean, val.\ (6~nodes)}  & 1.32 & & \\
%DIFDELCMD < \multicolumn{3}{@{}l}{Group mean, cal.\ (4~nodes)}  & 2.32 & & \\
%DIFDELCMD < \multicolumn{3}{@{}l}{Group mean, excl.\ (4~nodes)} & 1.57 & & \\
%DIFDELCMD < \multicolumn{3}{@{}l}{All 14~nodes}                  & 1.68 & & \\
%DIFDELCMD < \bottomrule
%DIFDELCMD < \end{tabular*}
%DIFDELCMD < %%%
\DIFdelend \DIFaddbegin \begin{tabular*}{\columnwidth}{@{\extracolsep{\fill}} l c c c c c @{}}
\toprule
Node & Path [m] & Role & MAE [$^\circ$C] & Bias [$^\circ$C] & In gate? \\
\midrule
j$_2$    &  350 & excl. & 2.99 & $+2.98$ & No  \\
j$_3$    &  650 & excl. & 0.58 & $+0.07$ & No  \\
j$_4$    &  910 & excl. & 0.37 & $-0.18$ & No  \\
j$_9$    & 1100 & val.  & 0.87 & $+0.41$ & Yes \\
j$_5$    & 1150 & excl. & 2.35 & $+1.51$ & No  \\
j$_6$    & 1300 & cal.  & 1.62 & $-1.62$ & No  \\
j$_{10}$ & 1300 & val.  & 2.27 & $+1.78$ & Yes \\
j$_7$    & 1370 & cal.  & 3.07\tnote{a} & $+2.53$ & No  \\
j$_8$    & 1450 & cal.  & 2.29 & $+1.00$ & No  \\
j$_{11}$ & 1530 & val.  & 1.41 & $+0.71$ & Yes \\
j$_{12}$ & 1780 & val.  & 0.69 & $-0.29$ & Yes \\
j$_{13}$ & 2000 & val.  & 1.09 & $-1.02$ & Yes \\
j$_{15}$ & 2125 & val.  & 1.58 & $-0.24$ & Yes \\
j$_{14}$ & 2180 & cal.  & 2.31 & $+1.02$ & No  \\
\midrule
\multicolumn{3}{@{}l}{Group mean, validation (6 nodes)}  & 1.32 & & \\
\multicolumn{3}{@{}l}{Group mean, calibration (4 nodes)} & 2.32 & & \\
\multicolumn{3}{@{}l}{Group mean, excluded (4 nodes)}    & 1.57 & & \\
\multicolumn{3}{@{}l}{All 14 nodes}                       & 1.68 & & \\
\bottomrule
\end{tabular*}
\DIFaddend \begin{tablenotes}\scriptsize
  \item[a] In-sample residual \DIFdelbeginFL \DIFdelFL{(calibration node)}\DIFdelendFL \DIFaddbeginFL \DIFaddFL{on a calibration node}\DIFaddendFL ; would exceed \DIFdelbeginFL \DIFdelFL{the worst-node validation gate of }%DIFDELCMD < \SI{2.5}{\celsius}
%DIFDELCMD <     %%%
\DIFdelFL{if held out for validation}\DIFdelendFL \DIFaddbeginFL \DIFaddFL{a $2.5$\,°C held-out
    criterion}\DIFaddendFL .
\end{tablenotes}
\end{threeparttable}
}% end centering group

\DIFdelbegin %DIFDELCMD < \bigskip
%DIFDELCMD < 

%DIFDELCMD < %%%
\DIFdelend \section{Demand heterogeneity index}
\label{app:hi_definition}

\DIFdelbegin \DIFdel{The }\DIFdelend \DIFaddbegin \DIFadd{With $d_n$ the share of total annual demand at node $n$, the }\DIFaddend normalised variance index \DIFdelbegin \DIFdel{HI is
defined as:
}\DIFdelend \DIFaddbegin \DIFadd{is
}\DIFaddend \begin{equation}
  d_n = \frac{\sum_t Q_{n,t}^{\mathrm{dem}}}{\sum_n \sum_t Q_{n,t}^{\mathrm{dem}}},
  \DIFdelbegin \DIFdel{\quad
  }\DIFdelend \DIFaddbegin \DIFadd{\qquad
  }\DIFaddend \mathrm{HI} = \DIFdelbegin \DIFdel{\frac{N}{N{-}1} }\DIFdelend \DIFaddbegin \DIFadd{\frac{N}{N-1} }\DIFaddend \sum_{n=1}^{N}
    \left(d_n - \frac{1}{N}\right)\DIFdelbegin \DIFdel{^{\!2}
  }\DIFdelend \DIFaddbegin \DIFadd{^{2}.
  }\DIFaddend \label{eq:hi}
\end{equation}
\DIFdelbegin \DIFdel{HI }\DIFdelend \DIFaddbegin \DIFadd{The index }\DIFaddend maps monotonically to \DIFdelbegin \DIFdel{Gini for }\DIFdelend \DIFaddbegin \DIFadd{the Gini coefficient on }\DIFaddend tree topologies but compresses the
realistic range\DIFdelbegin \DIFdel{(HI\,$= 0.05$ for the primary case despite
Gini \,$= 0.41$). The
synthetic-study values (HI\,$= 0.1, 0.4, 0.8$) }\DIFdelend \DIFaddbegin \DIFadd{: the industrial case has $\mathrm{HI} = 0.05$ against a Gini of $0.41$. The
synthetic values $0.1$, $0.4$ and $0.8$ }\DIFaddend correspond approximately to Gini ranges of
\DIFdelbegin \DIFdel{0.15--0.25, 0.40--0.55, and 0.70--0.85.
}%DIFDELCMD < 

%DIFDELCMD < \section{BCM calibration details}
%DIFDELCMD < \label{app:bcm_calibration}
%DIFDELCMD < 

%DIFDELCMD < %%%
\DIFdel{The BCM forward reconstruction distributes heat loss uniformly
across the full }%DIFDELCMD < \SI{2125}{\metre} %%%
\DIFdel{trunk corridor via a global
multiplier (${\times}1.330$ on nominal EN~253 $U$-values for all
8~trunk pipes). The multiplier was determined by golden-section
search, minimising MAE on the Oct--Feb winter subset.  This
approach differs from the
branch-sequential optimisation -model
calibration (which concentrates losses in the terminal pipe at
$4.7{\times}$~nominal); both methods are valid for their
respective purposes (BCM: forward physics; L3/$\Lnl$: dispatch
optimisation ).  RMSE: }%DIFDELCMD < \SI{2.15}{\celsius}%%%
\DIFdel{; bias:
${+}\SI{0.72}{\celsius}$; 60\,\% of hourly steps within
${\pm}\SI{1.5}{\celsius}$.
}%DIFDELCMD < 

%DIFDELCMD < %%%
%DIF < % ============================================================
%DIFDELCMD < \section{Synthetic parametric study: detailed parameter effects}
%DIFDELCMD < \label{app:synth_detail}
%DIFDELCMD < 

%DIFDELCMD < %%%
\DIFdel{The following tables provide the quantitative evidence underlying
the parameter-ranking statements in }%DIFDELCMD < \Cref{sec:results_generalizability}%%%
\DIFdel{.
}%DIFDELCMD < 

%DIFDELCMD < %%%
\paragraph{\DIFdel{Implementation note.}}
%DIFAUXCMD
\addtocounter{paragraph}{-1}%DIFAUXCMD
\DIFdel{In the synthetic parametric study, the L3 configuration includes
steady-state pressure-drop calculation (Darcy--Weisbach) as part
of the pipe model, unlike the primary-case L3, which excludes
pressure drop (}%DIFDELCMD < \Cref{tab:model_scope}%%%
\DIFdel{). This reassignment is
required by the additive decomposition
(}%DIFDELCMD < \Cref{eq:cost_decomposition,tab:synth_level_mapping}%%%
\DIFdel{).
Consequently, the L3$\to\Lplus$ increment in the synthetic study
isolates }\emph{\DIFdel{only}} %DIFAUXCMD
\DIFdel{transport-delay activation, which produces
zero cost difference at hourly resolution with
$\tau_{\max} < 1$\,h. The primary-case L3$\to\Lplus$ gap
($+0.11\,\%$, }%DIFDELCMD < \Cref{tab:cost_extended}%%%
\DIFdel{) remains the valid
reference for the isolated pressure-drop effect.
}%DIFDELCMD < 

%DIFDELCMD < %%%
%DIF < %--- Tab A: HI effect within pipe-length bands ---
%DIFDELCMD < \medskip
%DIFDELCMD < {%%%
%DIF < 
%DIFDELCMD < \captionof{table}{L1$\to$L3 cost gap [\%] by pipe length and
%DIFDELCMD < demand heterogeneity index (median over node count and storage
%DIFDELCMD < capacity).}
%DIFDELCMD < \label{tab:synth_hi_effect}
%DIFDELCMD < \centering
%DIFDELCMD < \begin{threeparttable}
%DIFDELCMD < \footnotesize
%DIFDELCMD < \begin{tabular*}{\columnwidth}{@{\extracolsep{\fill}} l c c c c @{}}
%DIFDELCMD < \toprule
%DIFDELCMD < Pipe [km] & HI\,=\,0.1 & HI\,=\,0.4\tnote{a} & HI\,=\,0.8
%DIFDELCMD <   & $\Delta$(High--Low) \\
%DIFDELCMD < \midrule
%DIFDELCMD < 1   & $+3.23$ & -- & $+3.19$ & $-0.04$\,pp \\
%DIFDELCMD < 5   & $+15.88$ & -- & $+14.55$ & $-1.33$\,pp \\
%DIFDELCMD < 15  & $+38.46$ & -- & $+38.63$ & $+0.17$\,pp \\
%DIFDELCMD < \bottomrule
%DIFDELCMD < \end{tabular*}
%DIFDELCMD < \begin{tablenotes}\scriptsize
%DIFDELCMD <   \item[a] %%%
\DIFdelFL{Median not reported where $n < 3$ feasible
    configurations per cell.
}%DIFDELCMD < \end{tablenotes}
%DIFDELCMD < \end{threeparttable}
%DIFDELCMD < }%%%
%DIF <  end group
%DIFDELCMD < 

%DIFDELCMD < \bigskip
%DIFDELCMD < 

%DIFDELCMD < %%%
%DIF < %--- Tab B: Decomposition extremes ---
%DIFDELCMD < {%%%
%DIF < 
%DIFDELCMD < \captionof{table}{Additive gap decomposition for the three
%DIFDELCMD < largest-gap and three smallest-gap configurations.  Heat losses
%DIFDELCMD < account for 100--108\,\% of the net gap; routing is slightly
%DIFDELCMD < negative in all cases.}
%DIFDELCMD < \label{tab:synth_decomp_extremes}
%DIFDELCMD < \centering
%DIFDELCMD < \footnotesize
%DIFDELCMD < \setlength{\tabcolsep}{2pt}
%DIFDELCMD < \begin{tabular*}{\columnwidth}{@{\extracolsep{\fill}}
%DIFDELCMD <   c c c c c c c c @{}}
%DIFDELCMD < \toprule
%DIFDELCMD < Rank & Pipe & HI & Nodes & Stor.
%DIFDELCMD <   & Routing & Loss & $\Delta P$ \\
%DIFDELCMD <   & [km] & & & [h] & [\%] & [\%] & [\%] \\
%DIFDELCMD < \midrule
%DIFDELCMD < \multicolumn{8}{@{}l}{\textit{Top-3 (largest total gap)}} \\
%DIFDELCMD < 1 & 15 & 0.1 & 30 & 12 & $-$3.5 & $+$47.2 & $+$1.1 \\
%DIFDELCMD < 2 & 15 & 0.8 & 30 & 12 & $-$4.1 & $+$46.8 & $+$0.9 \\
%DIFDELCMD < 3 & 15 & 0.1 & 15 & 12 & $-$2.8 & $+$44.1 & $+$1.5 \\
%DIFDELCMD < \midrule
%DIFDELCMD < \multicolumn{8}{@{}l}{\textit{Bottom-3 (smallest total gap)}} \\
%DIFDELCMD < 34 & 1 & 0.8 & 5  & 6  & $-$0.2 & $+$3.3 & $+$0.0 \\
%DIFDELCMD < 35 & 1 & 0.1 & 5  & 2  & $-$0.1 & $+$3.1 & $+$0.0 \\
%DIFDELCMD < 36 & 1 & 0.1 & 5  & 6  & $-$0.1 & $+$3.0 & $+$0.0 \\
%DIFDELCMD < \bottomrule
%DIFDELCMD < \end{tabular*}
%DIFDELCMD < }%%%
%DIF <  end group
%DIFDELCMD < 

%DIFDELCMD < \bigskip
%DIFDELCMD < 

%DIFDELCMD < \bigskip
%DIFDELCMD < 

%DIFDELCMD < %%%
%DIF < %--- Tab D: Parameter ranking ---
%DIFDELCMD < {%%%
%DIF < 
%DIFDELCMD < \captionof{table}{Marginal effect of each design parameter on
%DIFDELCMD < the L1$\to$L3 cost gap (median gap at each parameter level,
%DIFDELCMD < averaged over remaining parameters; sorted by effect size).}
%DIFDELCMD < \label{tab:synth_marginal}
%DIFDELCMD < \centering
%DIFDELCMD < \begin{threeparttable}
%DIFDELCMD < \footnotesize
%DIFDELCMD < \setlength{\tabcolsep}{3pt}
%DIFDELCMD < \begin{tabular*}{\columnwidth}{@{\extracolsep{\fill}}
%DIFDELCMD <   c
%DIFDELCMD <   >{\raggedright\arraybackslash}p{0.20\columnwidth}
%DIFDELCMD <   >{\centering\arraybackslash}p{0.11\columnwidth}
%DIFDELCMD <   >{\centering\arraybackslash}p{0.11\columnwidth}
%DIFDELCMD <   >{\centering\arraybackslash}p{0.11\columnwidth}
%DIFDELCMD <   >{\centering\arraybackslash}p{0.14\columnwidth} @{}}
%DIFDELCMD < \toprule
%DIFDELCMD < Rank & Parameter & Low & Med & High
%DIFDELCMD <   & $\Delta$\,(H--L) \\
%DIFDELCMD < \midrule
%DIFDELCMD < 1 & Pipe length [km]
%DIFDELCMD <   & $+3.2\,\%$ & $+15.8\,\%$ & $+39.0\,\%$
%DIFDELCMD <   & $+35.8$\,pp \\
%DIFDELCMD < 2 & Storage cap.\ [h]\tnote{a}
%DIFDELCMD <   & $+9.4\,\%$ & $+17.2\,\%$ & $+39.0\,\%$
%DIFDELCMD <   & $+29.6$\,pp \\
%DIFDELCMD < 3 & Node count
%DIFDELCMD <   & $+12.9\,\%$ & $+16.8\,\%$ & $+20.0\,\%$
%DIFDELCMD <   & $+7.1$\,pp \\
%DIFDELCMD < 4 & Demand HI
%DIFDELCMD <   & $+20.8\,\%$ & -- & $+16.0\,\%$
%DIFDELCMD <   & $-4.8$\,pp \\
%DIFDELCMD < \bottomrule
%DIFDELCMD < \end{tabular*}
%DIFDELCMD < \begin{tablenotes}\scriptsize
%DIFDELCMD <   \item[a] %%%
\DIFdelFL{Storage-capacity effect partially confounded with
    pipe length: 2\,h configurations are unevenly distributed
    across pipe-length bands (4/6 occur at $L \geq 5$\,km).
    The isolated storage effect within fixed pipe-length bands
    is smaller (${\approx}5$}\DIFdelendFL \DIFaddbeginFL \DIFaddFL{$0.15$}\DIFaddendFL --\DIFdelbeginFL \DIFdelFL{$8$\,pp).
}%DIFDELCMD < \end{tablenotes}
%DIFDELCMD < \end{threeparttable}
%DIFDELCMD < }%%%
%DIF <  end group
%DIF < % ============================================================
%DIFDELCMD < \section{Case study supplementary data}
%DIFDELCMD < \label{app:case_data}
%DIFDELCMD < 

%DIFDELCMD < \subsection{Heating curve}
%DIFDELCMD < \label{app:heating_curve}
%DIFDELCMD < 

%DIFDELCMD < {%%%
%DIF < 
%DIFDELCMD < \captionof{table}{PWL temperature profile (heating curve).}
%DIFDELCMD < \label{tab:pwl}
%DIFDELCMD < \centering
%DIFDELCMD < \footnotesize
%DIFDELCMD < \begin{tabular*}{\columnwidth}{@{\extracolsep{\fill}} l c c c c @{}}
%DIFDELCMD < \toprule
%DIFDELCMD < Load fraction & 0.0 & 0.3 & 0.6 & 1.0 \\
%DIFDELCMD < \midrule
%DIFDELCMD < $T_{\mathrm{sup}}$ [$^\circ$C] & 70 & 80 & 90 & 100 \\
%DIFDELCMD < $T_{\mathrm{ret}}$ [$^\circ$C] & 55 & 50 & 45 & 40 \\
%DIFDELCMD < $\Delta T$ [$^\circ$C]         & 15 & 30 & 45 & 60 \\
%DIFDELCMD < \bottomrule
%DIFDELCMD < \end{tabular*}
%DIFDELCMD < }%%%
%DIF <  end centering group
%DIFDELCMD < 

%DIFDELCMD < \bigskip
%DIFDELCMD < 

%DIFDELCMD < \subsection{L2 zone mapping}
%DIFDELCMD < \label{app:l2_zone_mapping}
%DIFDELCMD < 

%DIFDELCMD < {%%%
%DIF < 
%DIFDELCMD < \captionof{table}{L2 zone mapping (15~junctions $\to$ 7~zones).}
%DIFDELCMD < \label{tab:l2_zones}
%DIFDELCMD < \centering
%DIFDELCMD < \footnotesize
%DIFDELCMD < \begin{tabular*}{\columnwidth}{@{\extracolsep{\fill}} l l l l l @{}}
%DIFDELCMD < \toprule
%DIFDELCMD < Zone & Nodes & Cons. & Assets & Pipe \\
%DIFDELCMD < \midrule
%DIFDELCMD < A & j$_1$              & V$_1$              & CHP,TES,Gas,Bio,HP,EB & -- \\
%DIFDELCMD < B & j$_2$--j$_3$       & V$_2$--V$_3$       & --                    & 650\,m \\
%DIFDELCMD < C & j$_4$--j$_5$       & V$_4$--V$_9$       & --                    & 500\,m \\
%DIFDELCMD < D & j$_6$--j$_8$       & V$_{10}$--V$_{14}$ & --                    & 450\,m \\
%DIFDELCMD < E & j$_9$--j$_{11}$    & V$_{15}$--V$_{18}$ & --                    & 880\,m \\
%DIFDELCMD < F & j$_{12}$--j$_{13}$ & V$_{19}$--V$_{24}$ & --                    & 470\,m \\
%DIFDELCMD < G & j$_{14}$--j$_{15}$ & V$_{25}$--V$_{27}$ & --                    & 305\,m \\
%DIFDELCMD < \bottomrule
%DIFDELCMD < \end{tabular*}
%DIFDELCMD < }%%%
%DIF <  end centering group
%DIFDELCMD < 

%DIFDELCMD < \bigskip
%DIFDELCMD < 

%DIFDELCMD < \subsection{Post-processing validation}
%DIFDELCMD < \label{app:state_validation}
%DIFDELCMD < 

%DIFDELCMD < \noindent%%%
\textit{\DIFdel{Pipe velocity:}}
%DIFAUXCMD
\DIFdel{$v_{p,t} \leq \SI{2.5}{\metre\per\second}$; no violations.
}\textit{\DIFdel{Pumping (L1/L2/L3):}}
%DIFAUXCMD
\DIFdel{${<}5\,$kW peak, ${<}1\,\%$ of thermal losses.
}\textit{\DIFdel{$\Lnl$ convergence:}}
%DIFAUXCMD
\DIFdel{target ${\leq}\,0.5\,\%$ gap; achieved 1.99\,\% at 4\,h time
limit for 744\,h window}\DIFdelend \DIFaddbegin \DIFadd{$0.25$, $0.40$--$0.55$ and $0.70$--$0.85$. Because of that compression, the
model-selection observations in Table~\ref{tab:criteria} are expressed in terms of the Gini
coefficient and the top-$k$ demand share rather than this index}\DIFaddend .

\DIFdelbegin %DIFDELCMD < \subsection{Gap stability across scenarios}
%DIFDELCMD < \label{app:gap_stability}
%DIFDELCMD < 

%DIFDELCMD < {%%%
%DIF < 
%DIFDELCMD < \captionof{table}{Gap stability across 11~scenarios (\%).}
%DIFDELCMD < \label{tab:gap_stability}
%DIFDELCMD < \centering
%DIFDELCMD < \footnotesize
%DIFDELCMD < \begin{threeparttable}
%DIFDELCMD < \begin{tabular*}{\columnwidth}{@{\extracolsep{\fill}} l c c c @{}}
%DIFDELCMD < \toprule
%DIFDELCMD < Scenario & L1--L3 & L2--L3 & L3--$\Lplus$ \\
%DIFDELCMD < \midrule
%DIFDELCMD < Baseline     & +13.0 & +2.6 & +0.11\tnote{b} \\
%DIFDELCMD < Gas high     & +12.4 & +2.4 & +0.18 \\
%DIFDELCMD < Gas low      & +13.5 & +2.8 & +0.05 \\
%DIFDELCMD < Elec.\ high  & +12.5 & +2.2 & +0.00 \\
%DIFDELCMD < Elec.\ low   & +12.5 & +1.9 & +0.10 \\
%DIFDELCMD < CO$_2$ high  & +12.1 & +2.2 & +0.08 \\
%DIFDELCMD < CO$_2$ low   & +14.3 & +3.0 & +0.16 \\
%DIFDELCMD < Cold year    & --\tnote{a} & +0.6 & +0.11 \\
%DIFDELCMD < Warm year    & --\tnote{a} & $-$2.2 & +0.00 \\
%DIFDELCMD < HP COP low   & +13.0 & +2.5 & +0.01 \\
%DIFDELCMD < Biomass exp. & +13.0 & +2.5 & +0.20 \\
%DIFDELCMD < \bottomrule
%DIFDELCMD < \end{tabular*}
%DIFDELCMD < \begin{tablenotes}\scriptsize
%DIFDELCMD <   \item[a] %%%
\DIFdelFL{Uniform scaling; L1--L3 not comparable.
  }%DIFDELCMD < \item[b] %%%
\DIFdelFL{Value from primary-case reference run
    (}%DIFDELCMD < \Cref{tab:cost_extended}%%%
\DIFdelFL{); sensitivity scenarios use
    independent solver instances.
}%DIFDELCMD < \end{tablenotes}
%DIFDELCMD < \end{threeparttable}
%DIFDELCMD < %%%
\DIFdelend }
%DIF <  end centering group

%DIF < % ============================================================
\DIFdelbegin \section*{\DIFdel{CRediT authorship contribution statement}}
%DIFAUXCMD
\textbf{\DIFdel{Lukas Ruess:}} %DIFAUXCMD
\DIFdel{Conceptualisation, Methodology, Software,
Validation, Formal analysis, Investigation, Data curation,
Writing -- original draft, Visualisation, Funding acquisition.
}\textbf{\DIFdel{Alexander Sauer:}} %DIFAUXCMD
\DIFdel{Supervision, Writing -- review \& editing.
}%DIFDELCMD < 

%DIFDELCMD < %%%
%DIF < % ============================================================
\DIFdelend \DIFaddbegin \printcredits
\DIFaddend \bibliographystyle{elsarticle-num-names}
\DIFdelbegin %DIFDELCMD < \bibliography{paper1_dh_fidelity/Paper20_Literatur}
%DIFDELCMD <  %%%
\DIFdelend \DIFaddbegin \bibliography{Paper20_Literatur}
 \DIFaddend\end{document}
